%% For double-blind review submission, w/o CCS and ACM Reference (max submission space)
%\documentclass[sigplan,10pt,review,anonymous]{acmart}
%\settopmatter{printfolios=false,printccs=false,printacmref=false}
%% For double-blind review submission, w/ CCS and ACM Reference
%\documentclass[sigplan,review,anonymous]{acmart}\settopmatter{printfolios=true}
%% For single-blind review submission, w/o CCS and ACM Reference (max submission space)
%\documentclass[sigplan,review]{acmart}\settopmatter{printfolios=true,printccs=false,printacmref=false}
%% For single-blind review submission, w/ CCS and ACM Reference
%\documentclass[sigplan,review]{acmart}\settopmatter{printfolios=true}
%% For final camera-ready submission, w/ required CCS and ACM Reference
%\documentclass[sigplan,nonacm]{acmart}
\documentclass[sigplan,review,anonymous,acmsmall]{acmart}\settopmatter{printfolios=false,printccs=false,printacmref=false}


%% Conference information
%% Supplied to authors by publisher for camera-ready submission;
%% use defaults for review submission.
\acmConference[SPLASH'24]{ACM SIGPLAN conference on Systems, Programming, Languages, and Applications: Software for Humanity}{October 22-27, 2024}{Pasadena, California, United States}
%\acmYear{2018}
%\acmISBN{} % \acmISBN{978-x-xxxx-xxxx-x/YY/MM}
%\acmDOI{} % \acmDOI{10.1145/nnnnnnn.nnnnnnn}
%\startPage{1}

%% Copyright information
%% Supplied to authors (based on authors' rights management selection;
%% see authors.acm.org) by publisher for camera-ready submission;
%% use 'none' for review submission.
\setcopyright{none}
%\setcopyright{acmcopyright}
%\setcopyright{acmlicensed}
%\setcopyright{rightsretained}
%\copyrightyear{2018}           %% If different from \acmYear

%% Bibliography style
\bibliographystyle{acmart}

\usepackage{dsfont}
\usepackage{stmaryrd}
\usepackage{colortbl}
\usepackage{hyperref}

\usepackage{amsmath}
\DeclareMathOperator*{\argmax}{argmax}
\DeclareMathOperator*{\argmin}{argmin}
\usepackage{amssymb}

\usepackage[dvipsnames, table]{xcolor}
\usepackage{textcomp}

% Packages
\usepackage[pdf]{graphviz}
\usepackage{mathrsfs}

\newcommand*\circled[1]{\tikz[baseline=-0.1cm]{
  \node[shape=circle,draw,inner sep=0.48pt] (char) {\fontsize{7}{12}\textsf{#1}};}}

\DeclareMathAlphabet{\mathcal}{OMS}{cmsy}{m}{n}
\usepackage{cancel}
\newcommand\ccancel[2][red]{\renewcommand\CancelColor{\color{#1}}\cancel{#2}}
\newcommand{\nDownarrow}{\ensuremath{\text{ }\cancel{\Downarrow}\text{ }}}
\usepackage{centernot}

\usepackage{pgfplots, pgfplotstable}
\pgfplotsset{compat=1.7}
\usepgfplotslibrary{fillbetween}
\usetikzlibrary{patterns}
\pgfmathdeclarefunction{gauss}{2}{\pgfmathparse{1/(#2*sqrt(2*pi))*exp(-((x-#1)^2)/(2*#2^2))}}
\pgfmathdeclarefunction{nil}{1}{\pgfmathparse{0.001}}

\usepackage{arydshln}
\usepackage{adjustbox}
\usepackage{enumerate}
\usepackage{enumitem}
\usepackage{tikz-cd}
\usetikzlibrary{calc}
\usepackage{amsfonts}
%\usepackage{prooftrees}
\usepackage{bussproofs}
\renewcommand{\sectionautorefname}{\S}
\renewcommand{\subsectionautorefname}{\S}
\usepackage{float}

\usepackage{tikz-3dplot}
\usetikzlibrary{3d}
\usetikzlibrary{calligraphy}
\newif\ifshowcellnumber
\showcellnumbertrue

\usepackage{algorithm}
\usepackage{algpseudocode}
\usepackage{algorithmicx}
\usepackage{sourcecodepro}
\usepackage{tikz-qtree}
\usepackage{amsthm}
\usepackage{bm}
\usetikzlibrary{bayesnet}
\usetikzlibrary{arrows}
\usepackage{subcaption}
\usetikzlibrary{backgrounds}
\usetikzlibrary{tikzmark}
\usetikzlibrary{hobby}

\usepackage{mwe}

\newcommand{\E}{\mathbb{E}}
\newcommand{\Var}{\mathrm{Var}}
\newcommand{\Cov}{\mathrm{Cov}}

\newcommand{\CompOrder}{\mathcal{O}}
\def\graphspace{\mathbf{G}}
\def\Uniform{\mbox{\rm Uniform}}
\def\Gaussian{\mbox{\rm Gaussian}}
\def\Bernoulli{\mbox{\rm Bernoulli}}
\def\Dirichlet{\mbox{\rm Dirichlet}}

\usepackage{mathtools}% superior to amsmath
\usepackage{tikz}
% Packages
\usepackage{listings}
\DeclareRobustCommand{\hlred}[1]{{\sethlcolor{pink}\hl{#1}}}
\usepackage{fontspec}

\setmonofont[Scale=0.8]{JetBrainsMono}[
  Contextuals={Alternate},
  Path=./font/,
  Extension = .ttf,
  UprightFont=*-Regular,
  BoldFont=*-Bold,
  ItalicFont=*-Italic,
  BoldItalicFont=*-BoldItalic
]

\usepackage[skins,breakable,listings]{tcolorbox}

\lstdefinelanguage{kotlin}{
  comment=[l]{//},
  commentstyle={\color{gray}\ttfamily},
  emph={delegate, filter, firstOrNull, forEach, it, lazy, mapNotNull, println, repeat, assert, with, head, tail, len, return@},
  numberstyle=\noncopyable,
  identifierstyle=\color{black},
  keywords={abstract, actual, as, as?, break, by, class, companion, continue, data, do, dynamic, else, enum, expect, false, final, for, fun, get, if, import, in, infix, interface, internal, is, null, object, open, operator, override, package, private, public, return, sealed, set, super, suspend, this, throw, true, try, catch, typealias, val, var, vararg, when, where, while, tailrec, reified},
  keywordstyle={\bfseries},
  morecomment=[s]{/*}{*/},
  morestring=[b]",
  morestring=[s]{"""*}{*"""},
  ndkeywords={@Deprecated, @JvmField, @JvmName, @JvmOverloads, @JvmStatic, @JvmSynthetic, Array, Byte, Double, Float, Boolean, Int, Integer, Iterable, Long, Runnable, Short, String, int},
  ndkeywordstyle={\bfseries},
  sensitive=true,
  stringstyle={\ttfamily},
  literate={`}{{\char0}}1,
  escapeinside={(*@}{@*)}
}
\lstdefinelanguage{tidy}{
  comment=[l]{//},
  commentstyle={\color{gray}\ttfamily},
  emph={|, ->, ---},
  emphstyle={\color{red}},
  identifierstyle=\color{black},
  keywords={\|, ->, ---},
  otherkeywords={|,->},
  morekeywords={|,->},
  keywordstyle={\color{blue}\bfseries},
  morecomment=[s]{/*}{*/},
  morestring=[b]",
  morestring=[s]{"""*}{*"""},
  ndkeywords={@Deprecated, @JvmField, @JvmName, @JvmOverloads, @JvmStatic, @JvmSynthetic, Array, Byte, Double, Float, Int, Integer, Iterable, Long, Runnable, Short, String},
  ndkeywordstyle={\color{orange}\bfseries},
  sensitive=true,
  stringstyle={\color{green}\ttfamily},
  literate={`}{{\char0}}1
}

%%%%%%%%%%%%%%%%%%%%%%%%%%%%%%%%%%%%%%%%%%%
%
% Color boxes
%
%%%%%%%%%%%%%%%%%%%%%%%%%%%%%%%%%%%%%%%%%%%

\tcbset{
  enhanced jigsaw,
  breakable,
  listing only,
%  boxsep=-1pt,
%  top=-1pt,
  bottom=0.1cm,
  right=0.5cm,
  overlay first={
    \node[black!50] (S) at (frame.south) {\Large\ding{34}};
    \draw[dashed,black!50] (frame.south west) -- (S) -- (frame.south east);
  },
  overlay middle={
    \node[black!50] (S) at (frame.south) {\Large\ding{34}};
    \draw[dashed,black!50] (frame.south west) -- (S) -- (frame.south east);
    \node[black!50] (S) at (frame.north) {\Large\ding{34}};
    \draw[dashed,black!50] (frame.north west) -- (S) -- (frame.north east);
  },
  overlay last={
    \node[black!50] (S) at (frame.north) {\Large\ding{34}};
    \draw[dashed,black!50] (frame.north west) -- (S) -- (frame.north east);
  },
  before={\par\vspace{5pt}},
  after={\par\vspace{\parskip}\noindent}
}

\definecolor{slightgray}{rgb}{0.90, 0.90, 0.90}

\usepackage{soul}
\makeatletter
\def\SOUL@hlpreamble{%
  \setul{}{3.0ex}%
  \let\SOUL@stcolor\SOUL@hlcolor%
  \SOUL@stpreamble%
}
\makeatother

\newcommand{\inline}[1]{%
  \begingroup%
  \sethlcolor{slightgray}%
  \hl{\ttfamily\footnotesize #1}%
  \endgroup
}

\newcommand{\tinline}[1]{%
  \begingroup%
  \sethlcolor{slightgray}%
  \hl{\ttfamily\tiny #1}%
  \endgroup
}

\newtcblisting{halftidyinput}[1][]{%
  left skip=0.7cm,
  left=0.35cm,
  width=6cm,
%  left=-0.01cm,
  top=-0.1cm,
  bottom=-0.35cm,
  listing options={
    language=tidy,
    basicstyle=\ttfamily\small,
%numberstyle=\footnotesize,
    showstringspaces=false,
    tabsize=2,
    breaklines=true,
    numbers=none,
    inputencoding=utf8,
    escapeinside={(*@}{@*)},
    #1
  },
  underlay unbroken and first={%
    \path[draw=none] (interior.north west) rectangle node[white]{\includegraphics[width=4mm]{../figures/tidyparse_logo.png}} ([xshift=-10mm,yshift=-7mm]interior.north west);
  }
}

\newtcblisting{wholetidyinput}[1][]{%
  left skip=0.7cm,
  left=0.35cm,
  top=0.1cm,
  middle=0mm,
  boxsep=0mm,
  listing options={
    language=tidy,
    basicstyle=\ttfamily\small,
%numberstyle=\footnotesize,
    showstringspaces=false,
    tabsize=2,
    breaklines=true,
    numbers=none,
    inputencoding=utf8,
    escapeinside={(*@}{@*)},
    #1
  },
  underlay unbroken and first={%
      \path[draw=none] (interior.north west) rectangle node[white]{\includegraphics[width=4mm]{../figures/tidyparse_logo.png}} ([xshift=-10mm,yshift=-9mm]interior.north west);
  }
}

\definecolor{A}{RGB}{6,150,104}
\definecolor{B}{RGB}{196,74,137}
\definecolor{C}{RGB}{117,237,133}
\definecolor{D}{RGB}{246,46,243}
\definecolor{E}{RGB}{89,162,12}
\definecolor{F}{RGB}{113,12,158}
\definecolor{G}{RGB}{191,205,142}
\definecolor{H}{RGB}{51,58,158}
\definecolor{I}{RGB}{244,212,3}
\definecolor{J}{RGB}{37,36,249}
\definecolor{K}{RGB}{253,165,71}
\definecolor{L}{RGB}{27,81,29}
\colorlet{LA}{A!30}
\colorlet{LB}{B!30}
\colorlet{LC}{C!30}
\colorlet{LD}{D!30}
\colorlet{LE}{E!30}
\colorlet{LF}{F!30}
\colorlet{LG}{G!30}
\colorlet{LH}{H!30}
\colorlet{LI}{I!30}
\colorlet{LJ}{J!30}
\colorlet{LK}{K!30}
\colorlet{LL}{L!30}
\newcommand{\hiliA}[1]{%
  \colorbox{LA}{$#1$}}
\newcommand{\hiliB}[1]{%
  \colorbox{LB}{$#1$}}
\newcommand{\hiliC}[1]{%
  \colorbox{LC}{$#1$}}
\newcommand{\hiliD}[1]{%
  \colorbox{LD}{$#1$}}
\newcommand{\hiliE}[1]{%
  \colorbox{LE}{$#1$}}
\newcommand{\hiliF}[1]{%
  \colorbox{LF}{$#1$}}
\newcommand{\hiliG}[1]{%
  \colorbox{LG}{$#1$}}
\newcommand{\hiliH}[1]{%
  \colorbox{LH}{$#1$}}
\newcommand{\hiliI}[1]{%
  \colorbox{LI}{$#1$}}
\newcommand{\hiliJ}[1]{%
  \colorbox{LJ}{$#1$}}
\newcommand{\hiliK}[1]{%
  \colorbox{LK}{$#1$}}
\newcommand{\hiliL}[1]{%
  \colorbox{LL}{$#1$}}
\newcommand{\highlight}[1]{%
  \colorbox{lgray}{$#1$}}
\colorlet{lred}{red!30}
\colorlet{lorange}{orange!30}
\colorlet{lgreen}{green!30}
\colorlet{lgray}{black!15}
\colorlet{dgray}{black!75}
\DeclareRobustCommand{\hlred}[1]{{\sethlcolor{lred}\hl{#1}}}
\DeclareRobustCommand{\hlorange}[1]{{\sethlcolor{lorange}\hl{#1}}}
\DeclareRobustCommand{\hlgreen}[1]{{\sethlcolor{lgreen}\hl{#1}}}
\DeclareRobustCommand{\hlgray}[1]{{\sethlcolor{lgray}\hl{#1}}}
\DeclareRobustCommand{\caret}[1]{{\sethlcolor{dgray}\textcolor{white}{\hl{#1}}}}

\usepackage{url}
\usepackage{qtree}

\usepackage{filecontents}
\usepackage{pstricks-add}
\usepackage{emoji}
\usepackage{alltt}
\usepackage{nicematrix}
\usepackage{graphicx}
\usepackage{ulem}
\usepackage{upquote}
\tikzstyle{every picture}+=[remember picture]
\usepackage{menukeys}
\pgfplotstableread[col sep=comma,]{timings_loc.csv}\loctimings
\pgfplotstableread[col sep=comma,]{timings_unloc.csv}\unloctimings

\makeatletter
\DeclareRobustCommand{\cev}[1]{%
  {\mathpalette\do@cev{#1}}%
}
\newcommand{\do@cev}[2]{%
  \vbox{\offinterlineskip
  \sbox\z@{$\m@th#1 x$}%
  \ialign{##\cr
  \hidewidth\reflectbox{$\m@th#1\vec{}\mkern4mu$}\hidewidth\cr
  \noalign{\kern-\ht\z@}
    $\m@th#1#2$\cr
  }%
  }%
}
\makeatother

\makeatletter
\DeclareRobustCommand{\pder}[1]{%
  \@ifnextchar\bgroup{\@pder{#1}}{\@pder{}{#1}}}
\newcommand{\@pder}[2]{\frac{\partial#1}{\partial#2}}
\makeatother

\newcommand{\shup}{\shortuparrow}
\newcommand{\shri}{\shortrightarrow}
\newcommand{\shur}{\shup\hspace{-5pt}\shri}

\makeatletter
\def\squigglyred{\bgroup \markoverwith{\textcolor{red}{\lower3\p@\hbox{\sixly \char58}}}\ULon}
\makeatother

\makeatletter
\def\squigglyblu{\bgroup \markoverwith{\textcolor{blue}{\lower3\p@\hbox{\sixly \char58}}}\ULon}
\makeatother

\makeatletter
\def\squigglyora{\bgroup \markoverwith{\textcolor{orange}{\lower3\p@\hbox{\sixly \char58}}}\ULon}
\makeatother

\newcommand{\err}[1]{\smash{\squigglyred{#1}{}}}
\newcommand{\erb}[1]{\smash{\squigglyblu{#1}{}}}
\newcommand{\ero}[1]{\smash{\squigglyora{#1}{}}}
\newcommand{\stirlingii}{\genfrac{\{}{\}}{0pt}{}}

%======== Arrows =========
\newcommand{\knightarrow}{
  \tikz{
    \fill (0pt,0pt) circle [radius = 1pt];
    \fill (0pt,6pt) circle [radius = 1pt];
    \fill (6pt,0pt) circle [radius = 1pt];
    \fill (6pt,6pt) circle [radius = 1pt];
    \fill (12pt,0pt) circle [radius = 1pt];
    \fill (12pt,6pt) circle [radius = 1pt];
    \fill (6pt,0pt) circle [radius = 1pt];
    \fill (12pt,0pt) circle [radius = 1pt];
    \draw [-to] (0pt,0pt) -- (12pt,6pt);
  }
}

\newcommand{\kingarrow}{
  \tikz{
    \fill (0pt,0pt) circle [radius = 1pt];
    \fill (6pt,0pt) circle [radius = 1pt];
    \fill (0pt,6pt) circle [radius = 1pt];
    \fill (6pt,6pt) circle [radius = 1pt];
    \draw [-to] (0pt,0pt) -- (6pt,6pt);
    \draw [-to] (0pt,0pt) -- (0pt,6pt);
    \draw [-to] (0pt,0pt) -- (6pt,0pt);
  }
}

\newcommand{\duparrow}{
  \tikz{
    \fill (0pt,0pt) circle [radius = 1pt];
    \fill (0pt,6pt) circle [radius = 1pt];
    \draw [-to] (0pt,0pt) -- (0pt,6pt);
  }
}

\newcommand{\drightarrow}{
  \tikz{
    \fill (0pt,0pt) circle [radius = 1pt];
    \fill (6pt,0pt) circle [radius = 1pt];
    \draw [-to] (0pt,0pt) -- (6pt,0pt);
  }
}

\newcommand{\ddiagarrow}{
  \tikz{
    \fill (0pt,0pt) circle [radius = 1pt];
    \fill (6pt,0pt) circle [radius = 1pt];
    \fill (0pt,6pt) circle [radius = 1pt];
    \fill (6pt,6pt) circle [radius = 1pt];
    \draw [-to] (0pt,0pt) -- (6pt,6pt);
  }
}

\newcommand{\knightkingarrow}{
  \tikz{
    \fill (0pt,0pt) circle [radius = 1pt];
    \fill (0pt,6pt) circle [radius = 1pt];
    \fill (6pt,0pt) circle [radius = 1pt];
    \fill (6pt,6pt) circle [radius = 1pt];
    \fill (12pt,0pt) circle [radius = 1pt];
    \fill (12pt,6pt) circle [radius = 1pt];
    \draw [-to] (0pt,0pt) -- (6pt,6pt);
    \draw [-to] (0pt,0pt) -- (0pt,6pt);
    \draw [-to] (0pt,0pt) -- (6pt,0pt);
    \draw [-to] (0pt,0pt) -- (12pt,6pt);
  }
}

%======== Arrows =========

\usetikzlibrary{decorations.pathreplacing,automata,calc,positioning,matrix,fit,decorations.pathmorphing}

\usepackage{wrapfig}

\newcommand{\mkTrellis}[1]{
  \begin{tikzpicture}
    \def\dx{20pt}
    \def\dy{30pt}
    \newcounter{i}
    \stepcounter{i}
    \node[circle, draw, fill=black!30] (\arabic{i}) at (0,0){};
    \foreach [count=\i] \x in {2,...,#1}{
      \pgfmathsetmacro{\lox}{\x-1}%
      \pgfmathsetmacro{\loxt}{\x-3}%
      \foreach [count=\j] \xx in {-\lox,-\loxt,...,\lox}{
        \pgfmathsetmacro{\jj}{\j-1}%
        \stepcounter{i}
        \pgfmathsetmacro{\kk}{\xx-2}%
        \pgfmathsetmacro{\lbl}{\lox!/(\jj!*(\lox-\jj)!)}
        \ifnum\x<\kk
        \pgfmath\node[circle, draw]  (\arabic{i}) at (\xx*\dx, -\lox*\dy) {};
        \else
        \pgfmath\node[circle, draw, fill=black!30]  (\arabic{i}) at (\xx*\dx, -\lox*\dy) {};
        \fi
      }
    }
    \newcounter{z}
    \newcounter{xn}
    \newcounter{xnn}
    \pgfmathsetmacro{\maxx}{#1 - 1}
    \foreach \x in {1,...,\maxx}{
      \foreach \xx in {1,...,\x}{
        \stepcounter{z}
        \setcounter{xn}{\arabic{z}}
        \addtocounter{xn}{\x}
        \setcounter{xnn}{\arabic{xn}}
        \stepcounter{xnn}
        \draw [<-] (\arabic{z}) -- (\arabic{xn});
        \draw [<-] (\arabic{z}) -- (\arabic{xnn});
      }
    }
  \end{tikzpicture}
}

\newcommand{\dx}{20pt}
\newcommand{\dy}{30pt}
\newcounter{i}
\newcounter{z}
\newcounter{xn}
\newcounter{xnn}
\newcommand{\mkTrellisAppend}[1]{
  \begin{tikzpicture}
    \setcounter{i}{0}
    \setcounter{z}{0}
    \setcounter{xn}{0}
    \setcounter{xnn}{0}
    \stepcounter{i}
    \node[circle, draw] (\arabic{i}) at (0,0){};
    \foreach [count=\i] \x in {2,...,#1}{
      \pgfmathsetmacro{\lox}{\x-1}%
      \pgfmathsetmacro{\loxt}{\x-3}%
      \foreach [count=\j] \xx in {-\lox,-\loxt,...,\lox}{
        \pgfmathsetmacro{\jj}{\j-1}%
        \stepcounter{i}
        \pgfmathsetmacro{\kk}{\xx+2}%
        \pgfmathsetmacro{\lbl}{\lox!/(\jj!*(\lox-\jj)!)}
        \ifnum\x>\kk
        \pgfmath\node[circle, draw, fill=black!30]  (\arabic{i}) at (\xx*\dx, -\lox*\dy) {};
        \else
        \pgfmath\node[circle, draw]  (\arabic{i}) at (\xx*\dx, -\lox*\dy) {};
        \fi
      }
    }
    \pgfmathsetmacro{\maxx}{#1 - 1}
    \foreach \x in {1,...,\maxx}{
      \foreach \xx in {1,...,\x}{
        \stepcounter{z}
        \setcounter{xn}{\arabic{z}}
        \addtocounter{xn}{\x}
        \setcounter{xnn}{\arabic{xn}}
        \stepcounter{xnn}
        \draw [<-] (\arabic{z}) -- (\arabic{xn});
        \draw [<-] (\arabic{z}) -- (\arabic{xnn});
      }
    }
  \end{tikzpicture}
}

\newcommand{\mkTrellisInsert}[1]{
  \begin{tikzpicture}
    \setcounter{i}{0}
    \setcounter{z}{0}
    \setcounter{xn}{0}
    \setcounter{xnn}{0}
    \stepcounter{i}
    \node[circle, draw] (\arabic{i}) at (0,0){};
    \foreach [count=\i] \x in {2,...,#1}{
      \pgfmathsetmacro{\lox}{\x-1}%
      \pgfmathsetmacro{\loxt}{\x-3}%
      \foreach [count=\j] \xx in {-\lox,-\loxt,...,\lox}{
        \pgfmathsetmacro{\jj}{\j-1}%
        \stepcounter{i}
        \pgfmathsetmacro{\mp}{\xx+#1}%
        \pgfmathsetmacro{\mq}{\xx+\x}%
        \pgfmathsetmacro{\lbl}{\lox!/(\jj!*(\lox-\jj)!)}
        \ifnum\x>\mp
        \pgfmath\node[circle, draw, fill=black!30]  (\arabic{i}) at (\xx*\dx, -\lox*\dy) {};
        \else
        \ifnum#1<\mq
        \pgfmath\node[circle, draw, fill=black!30]  (\arabic{i}) at (\xx*\dx, -\lox*\dy) {};
        \else
        \pgfmath\node[circle, draw]  (\arabic{i}) at (\xx*\dx, -\lox*\dy) {};
        \fi
        \fi

      }
    }
    \pgfmathsetmacro{\maxx}{#1 - 1}
    \foreach \x in {1,...,\maxx}{
      \foreach \xx in {1,...,\x}{
        \stepcounter{z}
        \setcounter{xn}{\arabic{z}}
        \addtocounter{xn}{\x}
        \setcounter{xnn}{\arabic{xn}}
        \stepcounter{xnn}
        \draw [<-] (\arabic{z}) -- (\arabic{xn});
        \draw [<-] (\arabic{z}) -- (\arabic{xnn});
      }
    }
  \end{tikzpicture}
}

\usetikzlibrary{automata, positioning, arrows}

\newcommand{\nobarfrac}{\genfrac{}{}{0pt}{}}
\pgfplotstableread[col sep=comma,]{timings_loc.csv}\loctimings
\pgfplotstableread[col sep=comma,]{timings_unloc.csv}\unloctimings
\pgfplotstableread[col sep=comma,]{natural_errors.csv}\naturalerrors
\pgfplotstableread[col sep=comma,]{synthetic_errors.csv}\syntheticerrors

\usepackage[all,pdf]{xy}

\newcommand{\bs}{\blacksquare}
\newcommand{\ws}{\square}
\begin{document}
%
  \title{Syntax Repair as Language Intersection}
  %
  \begin{abstract}
    We introduce a new technique for correcting syntax errors in arbitrary context-free languages. Our work comes from the observation that syntax errors with a small repair typically have very few unique small repairs, which can usually be enumerated within a small edit distance then ranked within a short amount of time. Furthermore, we place a heavy emphasis on precision: the enumerated set must contain every possible repair within a few edits and no invalid repairs. To do so, we reduce CFL recognition onto Boolean tensor completion, then model error correction as a language intersection problem between a Levenshtein automaton and a context-free grammar. To decode the solutions, we then sample trees without replacement from the intersection grammar, which yields valid repairs within a certain Levenshtein distance. Finally, we rank all repairs discovered within 60 seconds by a Markov chain.
    \keywords{Error correction \and CFL reachability \and Langauge games.}
  \end{abstract}
%
%\titlerunning{Abbreviated paper title}
% If the paper title is too long for the running head, you can set
% an abbreviated paper title here
%
%  \author{Breandan Considine\inst{1} \and
%  Jin Guo\inst{1}\and
%  Xujie Si\inst{2}}
%
%  \authorrunning{Considine et al.}
% First names are abbreviated in the running head.
% If there are more than two authors, 'et al.' is used.
%
%  \institute{McGill University, Montr\'eal, QC H2R 2Z4, Canada\\
%  \email{\{breandan.considine@mail, jguo@cs\}.mcgill.ca}\and
%  University of Toronto, Toronto, ON, M5S 1A1 Canada\\
%  \email{six@utoronto.ca}}
%
  \maketitle              % typeset the header of the contribution


  \section{Introduction}

  Syntax repair is the problem of modifying an invalid sentence so it conforms to some grammar. Prior work has been devoted to fixing syntax errors using handcrafted heuristics. This work features a variety of approaches including rule-based systems and statistical language models. However, these techniques are often brittle, and are susceptible to misgeneralization. In a prior paper published in SPLASH, the authors sample production rules from an error correcting grammar. While theoretically sound, this technique is incomplete, i.e., not guaranteed to sample all edits within a certain Levenshtein distance, and no more. In this paper, we demonsrate it is possible to attain a significant advantage by synthesizing and scoring all repairs within a certain Levenshtein distance. Not only does this technique guarantee perfect generalization, but also helps with precision.

  We take a first-principles approach making no assumptions about the sentence or grammar and focuses on correctness and end-to-end latency. Our technique is simple:

  \begin{enumerate}
    \item We first reduce the problem of CFL recognition to Boolean tensor completion, then use that to compute the Parikh image of the CFL. This follows from a straightforward extension of the Chomsky-Sch\"utzenberger enumeration theorem.
    \item We then model syntax correction as a language intersection problem between a Levenshtein automaton and a context-free grammar, which we explicitly materialize using a specialized version of the Bar-Hillel construction to Levenshtein intersections. This greatly reduces the number of generated productions.
    \item To decode the members from the intersection grammar, we sample trees without replacement by constructing a bijection between syntax trees and the integers, then sampling integers uniformly without replacement from a finite range. This yields concrete repairs within a certain Levenshtein distance.
    \item Finally, we rerank all repairs found before a fixed timeout by n-gram perplexity.
  \end{enumerate}

  Though simple, this technique outperforms SoTA syntax repair techniques. Its efficacy owes to the fact it does not sample edits or nor productions, but unique, fully formed repairs within a certain Levenshtein distance. It is sound and complete up to a Levenshtein bound - i.e., it will find all repairs within an arbitrary Levenshtein distance, and no more. Often the langauge of small repairs is surpisingly small compared with the langauge of all possible edits, enabling us to efficiently synthesize and score every possible solution. This offers a significant advantage over memoryless sampling techniques.

  \pagebreak

  \section{Example}

Consider the following Python snippet, which contains a small syntax error:\\

\noindent\texttt{def prepend(i, k, L=[]) n and [prepend(i - 1, k, [b] + L) for b in range(k)]}\\

We can fix it by inserting a colon after the function definition, yielding:\\

\noindent\texttt{def prepend(i, k, L=[])\hlgreen{:} n and [prepend(i - 1, k, [b] + L) for b in range(k)]}\\

A careful observer will note that there is only one way to repair this Python snippet by making a single edit. In fact, many programming languages share this curious property: syntax errors with a small repair have few uniquely small repairs. Valid sentences corrupted by a few small errors rarely have many small corrections. We call such sentences \textit{metastable}, since they are relatively stable to small perturbations, as likely to be incurred by a careless typist or novice programmer.
%  Consider the following Kotlin snippet:\\
%
%  \texttt{fun main() = try \{ fetch() \} except(e: Exception) \{ handle(e) \}}\\
%
%  \noindent Again, there are thousands of possible single-token edits, only one of which is a valid repair:\\
%
%  \texttt{fun main() = try \{ fetch() \} \hlorange{catch}(e: Exception) \{ handle(e) \}}\\

Let us consider a slightly more ambiguous error: \texttt{v = df.iloc(5:, 2:)}. Assuming an alphabet of just one hundred lexical tokens, this tiny statement has millions of possible two-token edits, yet only six of those possibilities are accepted by the Python parser:

%\setlength{\columnsep}{-10pt}
%\setlength{\columnseprule}{-10pt}
%\noindent\begin{multicols}{3}
%  \begin{enumerate}
%    \item\texttt{v = df.iloc(5\hlred{:}, 2\hlorange{,})}\\
%    \item\texttt{v = df.iloc(5\hlgreen{[}:, 2:\hlgreen{]})}\\
%    \item\texttt{v = df.iloc\hlorange{[}5:, 2:\hlorange{]}}\\
%    \item\texttt{v = df.iloc(5\hlorange{)}, 2\hlorange{(})}\\
%    \item\texttt{v = df.iloc(5\hlred{:}, 2\hlred{:})}\\
%    \item\texttt{v = df.iloc(5\hlgreen{[}:, 2\hlorange{]})}
%  \end{enumerate}
%\end{multicols}
\begin{figure}[h!]
  \noindent\begin{tabular}{@{}l@{\hspace{10pt}}l@{\hspace{10pt}}l@{}}
  (1) \texttt{v = df.iloc(5\hlred{:}, 2\hlorange{,})} & (3) \texttt{v = df.iloc(5\hlgreen{[}:, 2:\hlgreen{]})} & (5) \texttt{v = df.iloc\hlorange{[}5:, 2:\hlorange{]}} \\\\
  (2) \texttt{v = df.iloc(5\hlorange{)}, 2\hlorange{(})} & (4) \texttt{v = df.iloc(5\hlred{:}, 2\hlred{:})} & (6) \texttt{v = df.iloc(5\hlgreen{[}:, 2\hlorange{]})} \\
  \end{tabular}
\end{figure}

With some typing information, we could easily narrow the results, but even in the absence of semantic constraints, one can probably rule out (2, 3, 6) given that \texttt{5[} and \texttt{2(} are rare bigrams in Python, and knowing \texttt{df.iloc} is often followed by \texttt{[}, determine (4) is most natural. This is the key insight behind our approach: we can usually locate the intended fix by exhaustively searching small repairs. As the set of small repairs is itself often small, if only we had some procedure to distinguish valid from invalid patches, the resulting solutions could be simply ranked by naturalness.

The trouble is that any such procedure must be highly sample-efficient. We cannot afford to sample the universe of possible d-token edits, then reject invalid ones -- assuming it takes just 10ms to generate and check each sample, (1-6) could take 24+ hours to find. The hardness of brute-force search grows superpolynomially with edit distance, sentence length and alphabet size. We need a more efficient procedure for sampling all and only small valid repairs.

  \section{Problem}

  We can model syntax repair as a language intersection problem between a context-free language (CFL) and a regular language.

  \begin{definition}[Bounded Levenshtein-CFL reachability]
    Given a CFL $\ell$ and an invalid string $\err{\sigma}: \ell^\complement$, the BCFLR problem is to find every valid string reachable within $d$ edits of $\err{\sigma}$, i.e., letting $\Delta$ be the Levenshtein metric and $L(\err\sigma, d) \coloneqq \{\sigma \mid \Delta(\err{\sigma}, \sigma) \leq d\}$, we seek to find $A = L(\err\sigma, d) \cap \ell$.
  \end{definition}

%  To solve this problem, it is convenient to first consider intersections with a finite-length string with holes, then turn our attention back to BCFLR.
%

  As the admissible set $A$ is typically under-determined, we want a procedure which prioritizes natural and valid repairs over unnatural but valid repairs:

  \begin{definition}[Ranked repair]\label{def:ranked-repair}
    Given a finite language $A = L(\err\sigma, d) \cap \ell$ and a probabilistic language model $\text{P}_\theta: \Sigma^* \rightarrow [0, 1] \subset \mathbb{R}$, the ranked repair problem is to find the top-$k$ maximum likelihood repairs under the language model. That is,
    \begin{equation}
      R(A, P_\theta) \coloneqq \argmax_{\{\bm{\sigma} \mid \bm{\sigma} \subseteq A, |\bm{\sigma}| \leq k\}} \sum_{\sigma \in \bm{\sigma}}\text{P}_\theta(\sigma\mid\err\sigma)
    \end{equation}
    % On average, across all $G, \sigma$ $\hat{R}$ should approximate $R$.
%    We want a procedure $\hat{R}$, minimizing $\mathbb{E}_{G, \sigma}\big[D_{\text{KL}}(\hat{R} \parallel R)\big]$ and wallclock runtime.
  \end{definition}

  The problem of ranked repair is typically solved by learning a distribution over strings, however this approach can be sample-inefficient and generalize poorly to new langauges. Representing the problem as a distribution over $\Sigma^*$ forces the language model to jointly learn syntax and stylometry. As our work demonstrates, this problem can be factorized into a bilevel objective: first retrieval, then ranking. By ensuring retrieval is sufficiently precise and exhaustive, maximizing likliehood over the admissible set can be achieved with a much simpler, syntax-oblivious language model.

  Even with an extremely efficient approximate sampler for $\sigma \sim \ell_\cap$, due to the size of $\ell$ and $L(\err\sigma, d)$, it would be intractable to sample either $\ell$ or $L(\err\sigma, d)$, then reject invalid ($\sigma \notin \ell$) or unreachable ($\sigma \notin L(\err\sigma, d)$) edits, and completely out of the question to sample $\sigma \sim \Sigma^*$ as do many large language models. Instead, we will explicitly construct a grammar generating $\ell \cap L(\err\sigma, d)$, sample from it, then rerank all repairs after a fixed timeout. So long as $|\ell_\cap|$ is sufficiently small and recognizes all and only small repairs, our sampler is sure to retrieve the most natural repair and terminate quickly.

  \section{Method}

  \begin{wrapfigure}{r}{0.45\textwidth}
%\begin{figure}[h!]
    \vspace{-1.2cm}
    \begin{center}
      \resizebox{0.45\textwidth}{!}{
        \begin{tikzpicture}[node distance=2cm]
          \node (start) [startstop] {Start};
          \node (pro1) [process, below of=start, yshift=-0.5cm] {$G_\cap \leftarrow G\cap\Delta(\err\sigma, d)$};
          \node (pcfg) [io, left of=pro1, xshift=-3cm] {[P]CFG};
          \node (lnfa) [io, right of=pro1, xshift=3cm] {L-WFA};

          \node (code) [io, right of=start,xshift=3cm] {Code};
          \node (synt) [io, left of=start,xshift=-3cm] {Syntax};

          \node (dec1) [decision, below of=pro1, yshift=-0.5cm] {$[G_{\cap} = \varnothing]$};

          \node (pro2b) [process, right of=dec1, xshift=3cm] {Increase radius, $d$};

          \node (dec2) [decision, below of=dec1, yshift=-1cm] {$|\mathcal{L}(G_\cap)|$};

          \node (samp1) [process, left of=dec2, xshift=-3cm] {Enumerate $\sigma' \in \mathcal{L}(G_\cap)$};
          \node (samp2) [process, right of=dec2, xshift=3cm] {Sample $\sigma' \sim P(G_\cap)$};

          \node (rank) [process, below of=dec2, yshift=-2.5cm] {Rerank by $L_\theta(\sigma')$};
          \node (vlmc) [io, above of=rank, yshift=-0.1cm] {VLMC};

%  \node (out1) [io, below of=pro2a] {Output};
          \node (stop) [startstop, below of=rank] {Done};

%  \draw [arrow] (dec0) -- node[anchor=east] {no} (pro1);

          \draw [->,thick] (0, 1.3) -- (start);

          \draw [arrow] (start) -- (code);
          \draw [arrow] (start) -- (synt);
          \draw [arrow] (code) -- (lnfa);
          \draw [arrow] (synt) -- (pcfg);
          \draw [arrow] (lnfa) -- (pro1);
          \draw [arrow] (pcfg) -- (pro1);

%  \draw [arrow] (in1) -- (pro1);
          \draw [arrow] (pro1) -- (dec1);
          \draw [arrow] (dec1) -- node[anchor=south] {yes} (pro2b);
          \draw [arrow] (dec1) -- node[anchor=east] {no} (dec2);
          \draw [arrow] (pro2b) -- (lnfa);
          \draw [arrow] (dec2) -- node[anchor=south] {small} (samp1);
          \draw [arrow] (dec2) -- node[anchor=south] {large} (samp2);

          \draw [arrow] (vlmc) -- (rank);
          \draw [arrow] (samp1) |- (rank);
          \draw [arrow] (samp2) |- (rank);
%  \draw [arrow] (pro2a) -- (out1);
          \draw [arrow] (rank) -- (stop);
%          \draw [arrow] (dec2) -- node[anchor=east] {1} (stop);

        \end{tikzpicture}
      }
    \end{center}
    \caption{Flowchart of our proposed method.}\label{fig:flowchart}
    \vspace{-1cm}
  \end{wrapfigure}

  The syntax of most programming languages is context-free. Our proposed method is simple. We construct a context-free grammar representing the intersection between the langauge syntax and an automaton recognizing the Levenshtein ball of a given radius. Since CFLs are closed under intersection with regular languages, this is admissible. Three outcomes are possible:

  \begin{enumerate}
    \item $G_\cap$ is empty, in which case there is no repair within the given radius. In this case, we simply increase the radius and try again.
    \item $\mathcal{L}(G_\cap)$ is small, in which case we enumerate all possible repairs. Enumeration is tractable for $\sim 80\%$ of the dataset in $\leq 90$s.
    \item $\mathcal{L}(G_\cap)$ is too large to enumerate, so we sample from the intersection grammar $G_\cap$. Sampling is necessary for $\sim20\%$ of the dataset.
  \end{enumerate}

  When ambiguous, we use an n-gram model to rank and return the top-k results by likelihood. This procedure is depicted in Fig.~\ref{fig:flowchart}. In the following section, we will describe how we construct the intersecection grammar (\ref{sec:lev_nfa}, \ref{sec:lev_bh}), then, provide a method for extracting explicit repairs from it.

  \subsection{Preliminaries}

  Recall that a CFG, $\mathcal{G} = \langle \Sigma, V, P, S\rangle$, is a quadruple consisting of terminals $(\Sigma)$, nonterminals $(V)$, productions $(P\colon V \rightarrow (V \mid \Sigma)^*)$, and a start symbol, $(S)$. Every CFG is reducible to \textit{Chomsky Normal Form}, $P'\colon V \rightarrow (V^2 \mid \Sigma)$, in which every production takes one of two forms, either $w \rightarrow xz$, or $w \rightarrow t$, where $w, x, z: V$ and $t: \Sigma$. For example:\vspace{-3pt}

  \begin{table}[H]
    \begin{tabular}{llll}
      $G\coloneqq\big\{\;S \rightarrow S\:S \mid (\:S\:) \mid (\:)\;\big\} \Longrightarrow \big\{\;S\rightarrow Q\:R \mid S\:S \mid L\:R,$ & $R \rightarrow\:),$ & $L \rightarrow (,$ & $Q\rightarrow L\:S\;\big\}$
    \end{tabular}
  \end{table}\vspace{-8pt}

  Likewise, a finite state automaton is a quintuple $\mathcal{A} = \langle Q, \Sigma, \delta, I, F\rangle$, where $Q$ is a finite set of states, $\Sigma$ is a finite alphabet, $\delta \subseteq Q \times \Sigma \times Q$ is the transition function, and $I, F \subseteq Q$ are the set of initial and final states, respectively. We will use this notation in the following sections.

  \pagebreak\subsection{The Nominal Levenshtein Automaton}\label{sec:lev_nfa}

  \begin{wrapfigure}{r}{0.5\textwidth}
    \vspace{-0.3cm}
    \begin{center}
        \resizebox{\textwidth}{!}{
  \begin{tikzpicture}[
%->, % makes the edges directed
  >=stealth',
  node distance=2.5cm, % specifies the minimum distance between two nodes. Change if necessary.
%  every state/.style={thick, fill=gray!10}, % sets the properties for each ’state’ node
  initial text=$ $, % sets the text that appears on the start arrow
  ]
  \node[state, initial]                (00) {$q_{0,0}$};
  \node[state, right of=00]            (10) {$q_{1,0}$};
  \node[accepting, state, right of=10] (20) {$q_{2,0}$};
  \node[accepting, state, right of=20] (30) {$q_{3,0}$};
  \node[accepting, state, right of=30] (40) {$q_{4,0}$};
  \node[accepting, state, right of=40] (50) {$q_{5,0}$};

  \node[state, above of=00, shift={(-1.5cm,-0.0cm)}] (01) {$q_{0,1}$};
  \node[state, right of=01]                          (11) {$q_{1,1}$};
  \node[state, right of=11]                          (21) {$q_{2,1}$};
  \node[accepting, state, right of=21]               (31) {$q_{3,1}$};
  \node[accepting, state, right of=31]               (41) {$q_{4,1}$};
  \node[accepting, state, right of=41]               (51) {$q_{5,1}$};

\node[state, above of=01, shift={(-1.5cm,-0.0cm)}] (0j) {$q_{0,2}$};
\node[state, right of=0j]                          (1j) {$q_{1,2}$};
\node[state, right of=1j]                          (2j) {$q_{2,2}$};
\node[state, right of=2j]                          (3j) {$q_{3,2}$};
\node[accepting, state, right of=3j]               (4j) {$q_{4,2}$};
\node[accepting, state, right of=4j]               (5j) {$q_{5,2}$};

\node[state, above of=0j, shift={(-1.5cm,0.0cm)}] (0k) {$q_{0,3}$};
\node[state, right of=0k]                         (1k) {$q_{1,3}$};
\node[state, right of=1k]                         (2k) {$q_{2,3}$};
\node[state, right of=2k]                         (3k) {$q_{3,3}$};
\node[state, right of=3k]                         (4k) {$q_{4,3}$};
\node[accepting, state, right of=4k]              (5k) {$q_{5,3}$};

\draw [->] (00) edge[below] node{$\sigma_1$} (10);
\draw [->] (10) edge[below] node{$\sigma_2$} (20);
\draw [->] (20) edge[below] node{$\sigma_3$} (30);
\draw [->] (30) edge[below] node{$\sigma_4$} (40);
\draw [->] (40) edge[below] node{$\sigma_5$} (50);

\draw [->] (01) edge[below] node{$\sigma_1$} (11);
\draw [->] (11) edge[below] node{$\sigma_2$} (21);
\draw [->] (21) edge[below] node{$\sigma_3$} (31);
\draw [->] (31) edge[below] node{$\sigma_4$} (41);
\draw [->] (41) edge[below] node{$\sigma_5$} (51);

\draw [->] (0j) edge[below] node{$\sigma_1$} (1j);
\draw [->] (1j) edge[below] node{$\sigma_2$} (2j);
\draw [->] (2j) edge[below] node{$\sigma_3$} (3j);
\draw [->] (3j) edge[below] node{$\sigma_4$} (4j);
\draw [->] (4j) edge[below] node{$\sigma_5$} (5j);

\draw [->] (0k) edge[below] node{$\sigma_1$} (1k);
\draw [->] (1k) edge[below] node{$\sigma_2$} (2k);
\draw [->] (2k) edge[below] node{$\sigma_3$} (3k);
\draw [->] (3k) edge[below] node{$\sigma_4$} (4k);
\draw [->] (4k) edge[below] node{$\sigma_5$} (5k);

\draw [->] (00) edge[left] node{$\phantom{\cdot}$} (11);
\draw [->] (10) edge[left] node{$\phantom{\cdot}$} (21);
\draw [->] (20) edge[left] node{$\phantom{\cdot}$} (31);
\draw [->] (30) edge[left] node{$\phantom{\cdot}$} (41);
\draw [->] (40) edge[left] node{$\phantom{\cdot}$} (51);

% Super-knight arcs
\draw [->, red] (00) edge[bend left=3] node[east, shift={(-0.3cm,-0.65cm)}]{$\color{red}\sigma_3$}         (3j);
\draw [->, red] (10) edge[bend left=3] node[east, shift={(-0.3cm,-0.65cm)}]{$\color{red}\sigma_4$}         (4j);
\draw [->, red] (20) edge[bend left=3] node[east, shift={(-0.3cm,-0.65cm)}]{$\color{red}\sigma_5$}         (5j);
\draw [->, red] (01) edge[bend left=3] node[east, shift={(-0.3cm,-0.65cm)}]{$\color{red}\sigma_3$}         (3k);
\draw [->, red] (11) edge[bend left=3] node[east, shift={(-0.3cm,-0.65cm)}]{$\color{red}\sigma_4$}         (4k);
\draw [->, red] (21) edge[bend left=3] node[east, shift={(-0.3cm,-0.65cm)}]{$\color{red}\sigma_5$}         (5k);

\draw [->, violet] (00) edge[bend right=1] node[east, shift={(0cm,0.35cm)}]{$\color{violet}\sigma_4$}  (4k);
\draw [->, violet] (10) edge[bend right=1] node[east, shift={(0cm,0.35cm)}]{$\color{violet}\sigma_5$}  (5k);

\draw [->] (01) edge[left] node{$\phantom{\cdot}$} (1j);
\draw [->] (11) edge[left] node{$\phantom{\cdot}$} (2j);
\draw [->] (21) edge[left] node{$\phantom{\cdot}$} (3j);
\draw [->] (31) edge[left] node{$\phantom{\cdot}$} (4j);
\draw [->] (41) edge[left] node{$\phantom{\cdot}$} (5j);

\draw [->] (0j) edge[left] node{$\phantom{\cdot}$} (1k);
\draw [->] (1j) edge[left] node{$\phantom{\cdot}$} (2k);
\draw [->] (2j) edge[left] node{$\phantom{\cdot}$} (3k);
\draw [->] (3j) edge[left] node{$\phantom{\cdot}$} (4k);
\draw [->] (4j) edge[left] node{$\phantom{\cdot}$} (5k);

\draw [->] (00) edge[bend left=10, left] node{$\phantom{\cdot}$} (01);
\draw [->] (10) edge[bend left=10, left] node{$\phantom{\cdot}$} (11);
\draw [->] (20) edge[bend left=10, left] node{$\phantom{\cdot}$} (21);
\draw [->] (30) edge[bend left=10, left] node{$\phantom{\cdot}$} (31);
\draw [->] (40) edge[bend left=10, left] node{$\phantom{\cdot}$} (41);
\draw [->] (50) edge[bend left=10, left] node{$\phantom{\cdot}$} (51);

\draw [->] (01) edge[bend left=10, left] node{$\phantom{\cdot}$} (0j);
\draw [->] (11) edge[bend left=10, left] node{$\phantom{\cdot}$} (1j);
\draw [->] (21) edge[bend left=10, left] node{$\phantom{\cdot}$} (2j);
\draw [->] (31) edge[bend left=10, left] node{$\phantom{\cdot}$} (3j);
\draw [->] (41) edge[bend left=10, left] node{$\phantom{\cdot}$} (4j);
\draw [->] (51) edge[bend left=10, left] node{$\phantom{\cdot}$} (5j);

\draw [->] (0j) edge[bend left=10, left] node{$\phantom{\cdot}$} (0k);
\draw [->] (1j) edge[bend left=10, left] node{$\phantom{\cdot}$} (1k);
\draw [->] (2j) edge[bend left=10, left] node{$\phantom{\cdot}$} (2k);
\draw [->] (3j) edge[bend left=10, left] node{$\phantom{\cdot}$} (3k);
\draw [->] (4j) edge[bend left=10, left] node{$\phantom{\cdot}$} (4k);
\draw [->] (5j) edge[bend left=10, left] node{$\phantom{\cdot}$} (5k);

\draw [->, blue] (00) edge[bend left=2,below] node[shift={(-0.7cm,-0.5cm)}]{$\color{blue}\sigma_2$}    (21);
\draw [->, blue] (10) edge[bend left=2,below] node[shift={(-0.7cm,-0.5cm)}]{$\color{blue}\sigma_3$}    (31);
\draw [->, blue] (20) edge[bend left=2,below] node[shift={(-0.7cm,-0.5cm)}]{$\color{blue}\sigma_4$}    (41);
\draw [->, blue] (30) edge[bend left=2,below] node[shift={(-0.7cm,-0.5cm)}]{$\color{blue}\sigma_5$}    (51);

\draw [->, blue] (01) edge[bend left=2,below] node[shift={(-0.7cm,-0.5cm)}]{$\color{blue}\sigma_2$}    (2j);
\draw [->, blue] (11) edge[bend left=2,below] node[shift={(-0.7cm,-0.5cm)}]{$\color{blue}\sigma_3$}    (3j);
\draw [->, blue] (21) edge[bend left=2,below] node[shift={(-0.7cm,-0.5cm)}]{$\color{blue}\sigma_4$}    (4j);
\draw [->, blue] (31) edge[bend left=2,below] node[shift={(-0.7cm,-0.5cm)}]{$\color{blue}\sigma_4$}    (5j);

\draw [->, blue] (0j) edge[bend left=2,below] node[shift={(-0.7cm,-0.5cm)}]{$\color{blue}\sigma_2$}    (2k);
\draw [->, blue] (1j) edge[bend left=2,below] node[shift={(-0.7cm,-0.5cm)}]{$\color{blue}\sigma_3$}    (3k);
\draw [->, blue] (2j) edge[bend left=2,below] node[shift={(-0.7cm,-0.5cm)}]{$\color{blue}\sigma_4$}    (4k);
\draw [->, blue] (3j) edge[bend left=2,below] node[shift={(-0.7cm,-0.5cm)}]{$\color{blue}\sigma_5$}    (5k);

%https://tex.stackexchange.com/a/20986/139648
\draw [decorate,decoration={brace,amplitude=10pt,raise=10pt,mirror}] (00.south west) -- (50.south east) node[midway,yshift=-3em]{\textbf{String length}};
\draw [decorate,decoration={brace,amplitude=10pt,raise=20pt}] (00.south west) -- (0k.north west) node[midway,xshift=-1.2cm,yshift=-0.6cm,rotate=-60]{\textbf{Edit distance}};
\end{tikzpicture}
}
    \end{center}
    \caption{NFA recognizing Levenshtein $\Delta(\sigma: \Sigma^5, 3)$.}\label{fig:lev_nfa}
    \vspace{-0.5cm}
  \end{wrapfigure}

  Levenshtein edits are recognized by an automaton known as the Levenshtein automaton. As the original construction defined by Schultz and Mihov~\cite{schulz2002fast} contains cycles and $\varepsilon$-transitions, we propose a variant which is $\varepsilon$-free and acyclic. Furthermore, we adpot a nominal form which supports infinite alphabets and considerably simplifies the langauge intersection. Illustrated in Fig.~\ref{fig:lev_nfa} is an example of a small Levenshtein automaton recognizing $\Delta(\sigma: \Sigma^5, 3)$. Unlabeled arcs accept any terminal from the alphabet, $\Sigma$.

  \noindent Alternatively, this transition system can be viewed as a kind of proof system in an unlabeled lattice. This is equivalent to Schultz and Mihov's automaton, but more amenable to our purposes as it does not any contain $\varepsilon$-arcs, and uses skip connections to represent consecutive deletions.

  \begin{prooftree}
    \AxiomC{$s\in\Sigma \phantom{\land} i \in [0, n] \phantom{\land} j \in [1, k]$}
    \RightLabel{$\duparrow$}
    \UnaryInfC{$(q_{i, j-1} \overset{s}{\rightarrow} q_{i,j}) \in \delta$}
    \DisplayProof
    \hskip 1.5em
    \AxiomC{$s\in\Sigma \phantom{\land} i \in [1, n] \phantom{\land} j \in [1, k]$}
    \RightLabel{$\ddiagarrow$}
    \UnaryInfC{$(q_{i-1, j-1} \overset{s}{\rightarrow} q_{i,j}) \in \delta$}
  \end{prooftree}
  \begin{prooftree}
    \AxiomC{$i \in [1, n] \phantom{\land} j \in [0, k]$}
    \RightLabel{$\drightarrow$}
    \UnaryInfC{$(q_{i-1, j} \overset{\sigma_i}{\rightarrow} q_{i,j}) \in \delta$}
    \DisplayProof
    \hskip 1.5em
    \AxiomC{$d \in [1, d_{\max}] \phantom{\land} i \in [d + 1, n] \phantom{\land} j \in [d, k]$}
    \RightLabel{$\knightarrow$}
    \UnaryInfC{$(q_{i-d-1, j-d} \overset{\sigma_i}{\rightarrow} q_{i,j}) \in \delta$}
  \end{prooftree}
  \begin{prooftree}
    \AxiomC{$\vphantom{|}$}
    \RightLabel{$\textsc{Init}$}
    \UnaryInfC{$q_{0,0} \in I$}
    \DisplayProof
    \hskip 1.5em
    \AxiomC{$q_{i, j}$}
    \AxiomC{$|n-i+j| \leq k$}
    \RightLabel{$\textsc{Done}$}
    \BinaryInfC{$q_{i, j}\in F$}
  \end{prooftree}

  \newcommand{\substitutionExample}{
    \tikz{
      \foreach \x in {0,8,16,24,32,40}{
        \fill (\x pt,0pt) circle [radius = 1pt];
        \fill (\x pt,8pt) circle [radius = 1pt];
      }
      \phantom{\fill (0pt,-8pt) circle [radius = 1pt];}
      \draw [-to] (0pt,0pt) -- (8pt,0pt);
      \draw [-to] (8pt,0pt) -- (16pt,0pt);
      \draw [-to] (16pt,0pt) -- (24pt,8pt);
      \draw [-to] (24pt,8pt) -- (32pt,8pt);
      \draw [-to] (32pt,8pt) -- (40pt,8pt);
    }
  }

  \newcommand{\insertionExample}{
    \tikz{
      \foreach \x in {0,8,16,24,32,40}{
        \fill (\x pt,0pt) circle [radius = 1pt];
        \fill (\x pt,8pt) circle [radius = 1pt];
      }
      \phantom{\fill (0pt,-8pt) circle [radius = 1pt];}
      \fill[white] (16pt,0pt) circle [radius = 1.2pt];
      \fill[white] (24pt,8pt) circle [radius = 1.2pt];
      \draw [-to] (0pt,0pt) -- (8pt,0pt);
      \draw [-to] (8pt,0pt) -- (24pt,0pt);
      \draw [-to] (24pt,0pt) -- (16pt,8pt);
      \draw [-to] (16pt,8pt) -- (32pt,8pt);
      \draw [-to] (32pt,8pt) -- (40pt,8pt);
    }
  }

  \newcommand{\deletionExample}{
    \tikz{
      \foreach \x in {0,8,16,24,32,40}{
        \fill (\x pt,0pt) circle [radius = 1pt];
        \fill (\x pt,8pt) circle [radius = 1pt];
      }
      \phantom{\fill (0pt,-8pt) circle [radius = 1pt];}
      \draw [-to] (0pt,0pt) -- (8pt,0pt);
      \draw [-to] (8pt,0pt) -- (16pt,0pt);
      \draw [-to] (16pt,0pt) -- (24pt,0pt);
      \draw [-to] (24pt,0pt) -- (40pt,8pt);
    }
  }

  \newcommand{\doubleDeletionExample}{
    \tikz{
      \foreach \x in {0,8,16,24,32,40}{
        \fill (\x pt,0pt) circle [radius = 1pt];
        \fill (\x pt,8pt) circle [radius = 1pt];
        \fill (\x pt,16pt) circle [radius = 1pt];
      }
      \draw [-to] (0pt,0pt) -- (24pt,16pt);
      \draw [-to] (24pt,16pt) -- (32pt,16pt);
      \draw [-to] (32pt,16pt) -- (40pt,16pt);
    }
  }

  \newcommand{\subDelExample}{
    \tikz{
      \foreach \x in {0,8,16,24,32,40}{
        \fill (\x pt,0pt) circle [radius = 1pt];
        \fill (\x pt,8pt) circle [radius = 1pt];
        \fill (\x pt,16pt) circle [radius = 1pt];
      }
      \draw [-to] (0pt,0pt) -- (8pt,0pt);
      \draw [-to] (8pt,0pt) -- (16pt,8pt);
      \draw [-to] (16pt,8pt) -- (32pt,16pt);
      \draw [-to] (32pt,16pt) -- (40pt,16pt);
    }
  }

  \newcommand{\subSubExample}{
    \tikz{
      \foreach \x in {0,8,16,24,32,40}{
        \fill (\x pt,0pt) circle [radius = 1pt];
        \fill (\x pt,8pt) circle [radius = 1pt];
        \fill (\x pt,16pt) circle [radius = 1pt];
      }
      \draw [-to] (0pt,0pt) -- (8pt,0pt);
      \draw [-to] (8pt,0pt) -- (16pt,8pt);
      \draw [-to] (16pt,8pt) -- (24pt,16pt);
      \draw [-to] (24pt,16pt) -- (32pt,16pt);
      \draw [-to] (32pt,16pt) -- (40pt,16pt);
    }
  }

  \newcommand{\insertDeleteExample}{
    \tikz{
      \foreach \x in {0,8,16,24,32,40,48}{
        \fill (\x pt,0pt) circle [radius = 1pt];
        \fill (\x pt,8pt) circle [radius = 1pt];
        \fill (\x pt,16pt) circle [radius = 1pt];
      }
      \fill[white] (16pt,16pt) circle [radius = 1.2pt];
      \fill[white] (8pt,0pt) circle [radius = 1.2pt];
      \fill[white] (16pt,8pt) circle [radius = 1.2pt];
      \draw [-to] (0pt,0pt) -- (16pt,0pt);
      \draw [-to] (16pt,0pt) -- (8pt,8pt);
      \draw [-to] (8pt,8pt) -- (24pt,8pt);
      \draw [-to] (24pt,8pt) -- (40pt,16pt);
      \draw [-to] (40pt,16pt) -- (48pt,16pt);
    }
  }

  Each arc plays a specific role. $\duparrow$ handles insertions, $\ddiagarrow$ handles substitutions, $\duparrow$ handles insertions and $\knightarrow$ handles [consecutive] deletions of various lengths. Let us consider some illustrative cases.

  \begin{table}[h!]
    \begin{tabular}{ccccccc}

      \texttt{f\hspace{3pt}.\hspace{3pt}\hlorange{[}\hspace{3pt}x\hspace{3pt})} &
      \texttt{f\hspace{3pt}.\hspace{3pt}\phantom{(}\hspace{3pt}x\hspace{3pt})} &
      \texttt{f\hspace{3pt}.\hspace{3pt}(\hspace{3pt}\hlred{x}\hspace{3pt})} &
      \texttt{\hlred{.}\hspace{3pt}\hlred{+}\hspace{3pt}(\hspace{3pt}x\hspace{3pt})} &
      \texttt{f\hspace{3pt}\hlorange{.}\hspace{3pt}\hlred{(}\hspace{3pt}x\hspace{3pt};} &
      \texttt{[\hspace{3pt}\hlorange{,}\hspace{3pt}\hlorange{x}\hspace{3pt}y\hspace{3pt}]} &
      \texttt{[\hspace{3pt}\phantom{,}\hspace{3pt},\hspace{3pt}\hlred{x}\hspace{3pt}y\hspace{3pt}]} \\

      \texttt{f\hspace{3pt}.\hspace{3pt}\hlorange{(}\hspace{3pt}x\hspace{3pt})} &
      \texttt{f\hspace{3pt}.\hspace{3pt}\hlgreen{(}\hspace{3pt}x\hspace{3pt})} &
      \texttt{f\hspace{3pt}.\hspace{3pt}(\hspace{3pt}\phantom{x}\hspace{3pt})} &
      \texttt{\phantom{f}\hspace{3pt}\phantom{.}\hspace{3pt}(\hspace{3pt}x\hspace{3pt})} &
      \texttt{f\hspace{3pt}\hlorange{*}\hspace{3pt}\phantom{(}\hspace{3pt}x\hspace{3pt};} &
      \texttt{[\hspace{3pt}\hlorange{x}\hspace{3pt}\hlorange{,}\hspace{3pt}y\hspace{3pt}]} &
      \texttt{[\hspace{3pt}\hlgreen{x}\hspace{3pt},\hspace{3pt}\phantom{x}\hspace{3pt}y\hspace{3pt}]} \\

      \substitutionExample & \insertionExample & \deletionExample & \doubleDeletionExample & \subDelExample & \subSubExample & \insertDeleteExample
    \end{tabular}
  \end{table}

  Note that the same patch can have multiple Levenshtein alignments. $\textsc{Done}$ constructs the final states, which are all states accepting strings $\sigma'$ such that Levenshtein distance of $\Delta(\sigma, \sigma') \leq d_\max$.

  To avoid creating a parallel bundle of arcs for each insertion and substutition point, we instead decorate each arc with a nominal predicate, accepting or rejecting $\sigma_i$. To distinguish this nominal variant from the original construction, we highlight the modified rules in orange below.

  \begin{prooftree}
    \AxiomC{$i \in [0, n] \phantom{\land} j \in [1, k]$}
    \RightLabel{$\duparrow$}
    \UnaryInfC{$(q_{i, j-1} \overset{{\color{orange}[\neq \sigma_i]}}{\rightarrow} q_{i,j}) \in \delta$}
    \DisplayProof
    \hskip 1.5em
    \AxiomC{$i \in [1, n] \phantom{\land} j \in [1, k]$}
    \RightLabel{$\ddiagarrow$}
    \UnaryInfC{$(q_{i-1, j-1} \overset{{\color{orange}[\neq \sigma_i]}}{\rightarrow} q_{i,j}) \in \delta$}
  \end{prooftree}
  \begin{prooftree}
    \AxiomC{$i \in [1, n] \phantom{\land} j \in [0, k]$}
    \RightLabel{$\drightarrow$}
    \UnaryInfC{$(q_{i-1, j} \overset{{\color{orange}[=\sigma_i]}}{\rightarrow} q_{i,j}) \in \delta$}
    \DisplayProof
    \hskip 1.5em
    \AxiomC{$d \in [1, d_{\max}] \phantom{\land} i \in [d + 1, n] \phantom{\land} j \in [d, k]$}
    \RightLabel{$\knightarrow$}
    \UnaryInfC{$(q_{i-d-1, j-d} \overset{{\color{orange}[=\sigma_i]}}{\rightarrow} q_{i,j}) \in \delta$}
  \end{prooftree}

  Nominalizing the NFA eliminates the creation of $e=|\Sigma - 1|\cdot|\sigma|\cdot2d_\max$ unnecessary arcs over the entire Levenshtein automaton and drastically reduces the size of the construction to follow, but does not affect the underlying semantics. Thus, it is essential to first nominalize the automaton before proceeding to avoid a large blowup in the intermediate grammar.

  \subsection{Levenshtein-Bar-Hillel Construction}\label{sec:lev_bh}

  We now describe the Bar-Hillel construction, which generates a grammar recognizing the intersection between a regular and a context-free language, then specialize it to Levenshtein intersections.

  \begin{lemma}\label{lemma:bar-hillel}
  For any context-free language $\ell$ and finite state automaton $\alpha$, there exists a context-free grammar $G_\cap$ such that $\mathcal{L}(G_\cap) = \ell \cap \mathcal{L}(\alpha)$. See Bar-Hillel~\cite{bar1961formal}.
  \end{lemma}

  \noindent Although Bar-Hillel~\cite{bar1961formal} lacks an explicit construction, Beigel and Gasarch~\cite{beigelproof} construct $G_\cap$ like so:

  \begin{prooftree}
    \AxiomC{$q \in I \phantom{\land} r \in F\vphantom{\overset{a}{\rightarrow}}$}
%  \RightLabel{$\varepsilon^+\textsc{-int}$}
    \UnaryInfC{$\big(S\rightarrow q S r\big) \in P_\cap$}
    \DisplayProof
    \hskip 1em
    \AxiomC{$(A \rightarrow a) \in P$}
    \AxiomC{$(q\overset{a}{\rightarrow}r) \in \delta$}
%  \RightLabel{$\varepsilon^+\textsc{-int}$}
    \BinaryInfC{$\big(qAr\rightarrow a\big)\in P_\cap$}
    \DisplayProof
    \hskip 1em
%\end{prooftree}
%\begin{prooftree}
    \AxiomC{$(w \rightarrow xz) \in P\vphantom{\overset{a}{\rightarrow}}$}
    \AxiomC{$p,q,r \in Q$}
    \RightLabel{\Join}
    \BinaryInfC{$\big(pwr\rightarrow (pxq)(qzr)\big) \in P_\cap$}
  \end{prooftree}

  This, now standard, Bar-Hillel construction applies to any CFL and REG language intersection, but generates a grammar whose cardinality is approximately $|P_\cap|=|I|\cdot|F| + |P|\cdot|\Sigma|\cdot|\sigma|\cdot2d_{\max} + |P|\cdot|Q|^3$. Applying the BH construction directly to practical langauges and code snippets can generate hundreds of trillions of productions for even moderately sized grammars and Levenshtein automata. Instead, we will describe a kind of reachability analysis that elides many unreachable productions in the case of Levenshtein intersection.

  Consider $\Join$, the most expensive rule. What $\Join$ tells us is each nonterminal in the intersection grammar, $P_\cap$, matches a substring simultaneously recognized by (1) a pair of states in the original NFA and (2) a nonterminal in the original CFG. A key observation is that $\Join$ considers the intersection between every such triple, but this is a gross overapproximation for most NFAs and CFGs, as the vast majority of all state pairs and nonterminals recognize no strings in common.

  To identify these triples, we define an interval domain that soundly overapproximates the Parikh image, encoding the minimum and maximum number of terminals each nonterminal can generate. Since some intervals may be right-unbounded, we write $\mathbb{N}^*=\mathbb{N} \cup \{\infty\}$ to denote the upper bound, and $\Pi = \{[a, b] \in \mathbb{N} \times \mathbb{N}^* \mid a \leq b\}^{|\Sigma|}$ to denote the Parikh map of all terminals.

  \begin{definition}[Parikh mapping of a nonterminal]
    Let $p: \Sigma^*\rightarrow\mathbb{N}^{|\Sigma|}$ be the Parikh operator~\cite{parikh1966context}, which counts the frequency of terminals in a string. Let $\pi: V \rightarrow \Pi$ be a function returning the smallest interval such that $\forall \sigma: \Sigma^*, \forall v: V$, $v \Rightarrow^* \sigma \vdash p(s) \in \pi(v)$.
  \end{definition}

  In other words, the Parikh mapping computes the greatest lower and least upper bound of the Parikh image over all strings in the language of a nonterminal. The infimum of a nonterminal's Parikh interval tells us how many of each terminal a nonterminal \textit{must} generate, and the supremum tells us how many it \textit{can} generate. Likewise, we define a similar relation over NFA state pairs:

  \begin{definition}[Parikh mapping of NFA states]
    We define $\pi: Q\times Q \rightarrow \Pi$ as returning the smallest interval such that $\forall \sigma: \Sigma^*, \forall q, q': Q$, $q \overset{\sigma}{\Longrightarrow} q' \vdash p(\sigma) \in \pi(q, q')$.
  \end{definition}

  Next, we will define a measure on Parikh intervals, which will represent the minimum total edits required to transform a string in one Parikh interval to a string in another.

  \begin{definition}[Parikh divergence]
    Given two Parikh intervals $\pi, \pi': \Pi$, we define the divergence between them as $\pi \parallel \pi' = \sum_{n=1}^{|\Sigma|} \min_{(i, i') \in \pi[n]\times \pi'[n]} |i - i'|$.
  \end{definition}

  Now, we know that if the Parikh divergence between two intervals exceeds the Levenshtein margin between two states in a Lev-NFA, those intervals must be incompatible as no two strings, one from each Parikh interval, can be transformed into the other with fewer than $\pi \parallel \pi'$ edits.

  \begin{definition}[Levenshtein-Parikh compatibility]
    Let $q = q_{h,i}, q'=q_{j,k}$ be two states in a Lev-NFA and V be a CFG nonterminal. We say that $(q, v, q'): Q\times V\times Q$ are compatible iff the Parikh divergence is bounded by the Levenshtein margin $k-i$, i.e., $v \lhd qq' \iff (\pi(v) \parallel \pi(q, q')) \leq k-i$.
  \end{definition}

  We define the modified Bar-Hillel construction for Levenshtein intersections as follows:

\begin{prooftree}
  \def\defaultHypSeparation{\hskip 0.15cm}
  \AxiomC{$(A \rightarrow a) \in P$}
  \AxiomC{$\color{orange}S.a$}
  \AxiomC{$(q\overset{{\color{orange}S}}{\rightarrow}r) \in \delta$}
%  \RightLabel{$\varepsilon^+\textsc{-int}$}
  \TrinaryInfC{$\big(qAr\rightarrow a\big)\in P_\cap$}
  \DisplayProof
  \AxiomC{$\vphantom{\overset{S}{\rightarrow}}\color{orange} w \lhd pr \phantom{\land} x \lhd pq \phantom{\land} z \lhd qr$}
  \AxiomC{$(w \rightarrow xz) \in P\vphantom{\overset{a}{\rightarrow}}$}
  \AxiomC{$p,q,r \in Q$}
  \RightLabel{\Join}
  \TrinaryInfC{$\big(pwr\rightarrow (pxq)(qzr)\big) \in P_\cap$}
\end{prooftree}

  Finally, once $G_\cap$ is constructed, it is renormalized by removing all unreachable and non-generating productions following~\cite{firsov2015certified} to obtain $G_\cap^*$, which is usually several orders of magnitude smaller. When this grammar is sufficiently small, we can extract all relevant repairs from it, otherwise we sample from it to obtain a subset of candidate repairs.

%Specifically, we compute Parikh intervals generated by every path though the Levenshtein automaton, then intersect the Parikh intervals for the candidate nonterminals in question. For example, suppose we have a $p, q, r: Q$ and $w \rightarrow x z$. Let's check...

%To generate edits from it, we can use the same procedure as before, but instead of interleaving $\err\sigma$ with $\varepsilon$ and introducing holes, we simply use $A\big((\_)^{|\err{\sigma}| + d}\big, G_\cap)$.


  \subsection{Repair as idempotent matrix completion}

  Now that we have a language that recognizes local repairs, we need a method to extract the repairs themselves. We impose specific criteria on such a procedure: it must generate (1) only valid repairs and (2) all repairs in the langauge. We also require the sampler to have (3) parallel decomposability and generate samples (4) uniformly, both (5) with and (6) without replacement in (7) constant time and space. These requirements motivate the need for a specialized generating function.

  In this section, we will introduce the porous completion problem and show how it can be translated to a kind of idempotent matrix completion, whose roots are valid strings in a context-free language. This technique is convenient for its geometric interpretability, amenability to parallelization, and generalizability to any CFG, regardless of finitude or ambiguity. We will see how, by redefining the algebraic operations $\oplus, \otimes$ over different carrier sets, one can obtain a recognizer, parser, generator, Parikh map and other convenient structures for working with CFLs.

  Given a CFG, $G' : \mathcal{G}$ in Chomsky Normal Form (CNF), we can construct a recognizer $R: \mathcal{G} \rightarrow \Sigma^n \rightarrow \mathbb{B}$ for strings $\sigma: \Sigma^n$ as follows. Let $2^V$ be our domain, $0$ be $\varnothing$, $\oplus$ be $\cup$, and $\otimes$ be defined as:\vspace{-10pt}

  \begin{align}
    X \otimes Z \coloneqq \big\{\;w \mid \langle x, z\rangle \in X \times Z, (w\rightarrow xz) \in P\;\big\}
  \end{align}

  \noindent If we define $\hat\sigma_r \coloneqq \{w \mid (w \rightarrow \sigma_r) \in P\}$, then construct a matrix with nonterminals on the superdiagonal representing each token, $M_0[r+1=c](G', \sigma) \coloneqq \;\hat\sigma_r$, the fixpoint $M_{i+1} = M_i + M_i^2$ is uniquely determined by the superdiagonal entries, which \vspace{-10pt}

  \begin{align*}
    M_0\coloneqq
    \begin{pNiceMatrix}[nullify-dots,xdots/line-style=loosely dotted]
      \varnothing & \hat\sigma_1 & \varnothing & \Cdots & \varnothing \\
      \Vdots      & \Ddots   & \Ddots      & \Ddots & \Vdots\\
      &          &             &        & \varnothing\\
      &          &             &        & \hat\sigma_n \\
      \varnothing & \Cdots   &             &        & \varnothing
    \end{pNiceMatrix} &\Rightarrow
    \begin{pNiceMatrix}[nullify-dots,xdots/line-style=loosely dotted]
      \varnothing & \hat\sigma_1 & \Lambda & \Cdots & \varnothing \\
      \Vdots      & \Ddots   & \Ddots  & \Ddots & \Vdots\\
      &          &         &        & \Lambda\\
      &          &         &        & \hat\sigma_n \\
      \varnothing & \Cdots   &         &        & \varnothing
    \end{pNiceMatrix} &\Rightarrow \ldots \Rightarrow M_\infty =
    \begin{pNiceMatrix}[nullify-dots,xdots/line-style=loosely dotted]
      \varnothing & \hat\sigma_1 & \Lambda & \Cdots & \Lambda^*_\sigma\\
      \Vdots      & \Ddots   & \Ddots  & \Ddots & \Vdots\\
      &          &         &        & \Lambda\\
      &          &         &        & \hat\sigma_n \\
      \varnothing & \Cdots   &         &        & \varnothing
    \end{pNiceMatrix}
  \end{align*}

  \noindent The northeasternmost entry $\Lambda^*_\sigma$ gives us the recognizer $R(G', \sigma) \coloneqq [S \in \Lambda^*_\sigma] \Leftrightarrow [\sigma \in \mathcal{L}(G)]$~\footnote{Hereinafter, we use Iverson brackets to denote the indicator function of a predicate with free variables, i.e., $[P] \Leftrightarrow \mathds{1}(P)$.}.

  This procedure can be lifted to the domain of strings containing free variables, which we call the \textit{porous completion problem}, whose solutions are the set of all syntactically valid strings consistent with the input string.

  \begin{definition}[Porous completion]
    Let $\underline\Sigma \coloneqq \Sigma \cup \{\_\}$, where $\_$ denotes a hole. We denote $\sqsubseteq: \Sigma^n \times \underline\Sigma^n$ as the relation $\{\langle\sigma', \sigma\rangle \mid \sigma_i \in \Sigma \implies \sigma_i' = \sigma_i\}$ and the set of all inhabitants $\{\sigma': \Sigma^+ \mid \sigma' \sqsubseteq \sigma\}$ as $\text{H}(\sigma)$. Given a \textit{porous string}, $\sigma: \underline\Sigma^*$ we seek all syntactically valid inhabitants, i.e., $A(\sigma)\coloneqq\text{H}(\sigma)\cap\ell$.
  \end{definition}

  Let us consider an example with two holes, $\sigma = 1$ \_ \_, and the grammar being $G\coloneqq\{S\rightarrow N O N, O \rightarrow + \mid \times, N \rightarrow 0 \mid 1\}$. This can be rewritten into CNF as $G'\coloneqq \{S \rightarrow N L, N \rightarrow 0 \mid 1, O \rightarrow \times \mid +, L \rightarrow O N\}$. Using the algebra where $\oplus=\cup$, $X \otimes Z = \big\{\;w \mid \langle x, z\rangle \in X \times Z, (w\rightarrow xz) \in P\;\big\}$, the fixpoint $M' = M + M^2$ can be computed as follows, shown in the leftmost column:\\

  \begin{small}
  {\renewcommand{\arraystretch}{1.2}
  \noindent\phantom{...}\begin{tabular}{|c|c|c|c|}
                          \hline
                          & $2^V$ & $\mathbb{Z}_2^{|V|}$ & $\mathbb{Z}_2^{|V|}\rightarrow\mathbb{Z}_2^{|V|}$\\\hline
                          $M_0$ & \begin{pmatrix}
                                    \phantom{V} & \tiny{\{N\}} &         &             \\
                                    &              & \{N,O\} &             \\
                                    &              &         & \{N,O\} \\
                                    &              &         &
                          \end{pmatrix} & \begin{pmatrix}
                                            \phantom{V} & \ws\bs\ws\ws &              &              \\
                                            &              & \ws\bs\bs\ws &              \\
                                            &              &              & \ws\bs\bs\ws \\
                                            &              &              &
                          \end{pmatrix} & \begin{pmatrix}
                                            \phantom{V} & V_{0, 1} &          &          \\
                                            &          & V_{1, 2} &          \\
                                            &          &          & V_{2, 3} \\
                                            &          &          &
                          \end{pmatrix} \\\hline
                          $M_1$ & \begin{pmatrix}
                                    \phantom{V} & \tiny{\{N\}} & \varnothing &         \\
                                    &              & \{N,O\}     & \{L\}   \\
                                    &              &             & \{N,O\} \\
                                    &              &             &
                          \end{pmatrix} & \begin{pmatrix}
                                            \phantom{V} & \ws\bs\ws\ws & \ws\ws\ws\ws &              \\
                                            &              & \ws\bs\bs\ws & \bs\ws\ws\ws \\
                                            &              &              & \ws\bs\bs\ws \\
                                            &              &              &
                          \end{pmatrix} & \begin{pmatrix}
                                            \phantom{V} & V_{0, 1} & V_{0, 2} &          \\
                                            &          & V_{1, 2} & V_{1, 3} \\
                                            &          &          & V_{2, 3} \\
                                            &          &          &
                          \end{pmatrix} \\\hline
                          $M_\infty$ & \begin{pmatrix}
                                         \phantom{V} & \tiny{\{N\}} & \varnothing & \{S\}   \\
                                         &              & \{N,O\}     & \{L\}   \\
                                         &              &             & \{N,O\} \\
                                         &              &             &
                          \end{pmatrix} & \begin{pmatrix}
                                            \phantom{V} & \ws\bs\ws\ws & \ws\ws\ws\ws & \ws\ws\ws\bs \\
                                            &              & \ws\bs\bs\ws & \bs\ws\ws\ws \\
                                            &              &              & \ws\bs\bs\ws \\
                                            &              &              &
                          \end{pmatrix} & \begin{pmatrix}
                                            \phantom{V} & V_{0, 1} & V_{0, 2} & V_{0, 3} \\
                                            &          & V_{1, 2} & V_{1, 3} \\
                                            &          &          & V_{2, 3} \\
                                            &          &          &
                          \end{pmatrix}\\\hline
  \end{tabular}\\
  }
  \end{small}

  The same procedure can be translated, without loss of generality, into the bit domain ($\mathbb{Z}_2^{|V|}$) using a lexicographic ordering, however $M_\infty$ in both $2^V$ and $\mathbb{Z}_2^{|V|}$ represents a decision procedure, i.e., $[S\in V_{0, 3}]\Leftrightarrow [V_{0, 3, 3}=\bs] \Leftrightarrow [A(\sigma) \neq \varnothing]$. Since $V_{0, 3} = \{S\}$, we know there exists at least one $\sigma' \in A$, but $M_\infty$ does not reveal its identity.

%$\{\text{xor}, \land, \top\}$ is a functionally complete set is equivalent to $\mathbb{Z}_2$ $\top := 1, \land := \times, \text{xor} := +$. We can define $=$ as $(a = b) \Leftrightarrow (a \text{ xor } b) \text{ xor } \top \Leftrightarrow (a + b) + \top$.

  In order to extract the inhabitants, we can translate the bitwise procedure into an equation with free variables. Here, we can encode the idempotency constraint directly as $M = M^2$. We first define $X \boxtimes Z = [X_2 \land Z_1, \bot, \bot, X_1 \land Z_0]$ and $X \boxplus Z = [X_i \lor Z_i]_{i \in [0, |V|)}$. Since the unit nonterminals $O, N$ can only occur on the superdiagonal, they may be safely ignored by $\boxtimes$. To solve for $M_\infty$, we proceed by first computing $V_{0, 2}, V_{1, 3}$ as follows:

  \begin{small}
  \begin{align*}
    V_{0, 2} &= V_{0, j} \cdot V_{j, 2} = V_{0, 1} \boxtimes V_{1, 2}                         &  V_{1, 3} &= V_{1, j} \cdot V_{j, 3} = V_{1, 2} \boxtimes V_{2, 3}\\
    &= [L \in V_{0, 2}, \bot, \bot, S \in V_{0, 2}]                                           &  &= [L \in V_{1, 3}, \bot, \bot, S \in V_{1, 3}]\\
    &= [O \in V_{0, 1} \land N \in V_{1, 2}, \bot, \bot, N \in V_{0, 1} \land L \in V_{1, 2}] &  &= [O \in V_{1, 2} \land N \in V_{2, 3}, \bot, \bot, N \in V_{1, 2} \land L \in V_{2, 3}]\\
    &= [V_{0, 1, 2} \land V_{1, 2, 1}, \bot, \bot, V_{0, 1, 1} \land V_{1, 2, 0}]             &  &= [V_{1, 2, 2} \land V_{2, 3, 1}, \bot, \bot, V_{1, 2, 1} \land V_{2, 3, 0}]
  \end{align*}
  \end{small}

  Now we solve for the corner entry $V_{0, 3}$ by taking the bitwise dot product between the first row and last column, yielding:

  \begin{align*}
    V_{0, 3} &= V_{0, j} \cdot V_{j, 3} = V_{0, 1} \boxtimes V_{1, 3} \boxplus V_{0, 2} \boxtimes V_{2, 3}\\
%  &= [V_{0, 1, 2} \land V_{1, 3, 1}, \bot, \bot, V_{0, 1, 1} \land V_{1, 3, 0}] + [V_{0, 2, 2} \land V_{2, 3, 1}, \bot, \bot, V_{0, 2, 1} \land V_{2, 3, 0}]\\
    &= [V_{0, 1, 2} \land V_{1, 3, 1} \lor V_{0, 2, 2} \land V_{2, 3, 1}, \bot, \bot, V_{0, 1, 1} \land V_{1, 3, 0} \lor V_{0, 2, 1} \land V_{2, 3, 0}]
  \end{align*}

  \noindent Since we only care about $V_{0, 3, 3} \Leftrightarrow [S \in V_{0, 3}]$, so we can ignore the first three entries and solve for:

  \begin{align*}
    V_{0, 3, 3} &= V_{0, 1, 1} \land V_{1, 3, 0} \lor V_{0, 2, 1} \land V_{2, 3, 0}\\
    &= V_{0, 1, 1} \land (V_{1, 2, 2} \land V_{2, 3, 1}) \lor V_{0, 2, 1} \land \bot\\
    &= V_{0, 1, 1} \land V_{1, 2, 2} \land V_{2, 3, 1}\\
    &= [N \in V_{0, 1}] \land [O \in V_{1, 2}] \land [N \in V_{2, 3}]
  \end{align*}

  Now we know that $\sigma =$ 1 \underline{O} \underline{N} is a valid solution, and we can take the product $\{1\}\times \hat\sigma_2^{-1}(O) \times \hat\sigma_3^{-1}(N)$ to recover the admissible set, yielding $A=\{1+0, 1+1, 1\times 0, 1\times 1\}$. In this case, since $G$ is unambiguous, there is only one parse tree satisfying $V_{0, |\sigma|, 3}$, but in general, there can be multiple valid parse trees.

  This procedure generates solutions satisfying the matrix fixpoint, but forgets their provenance. The question naturally arises, is there a way to solve for the parse trees directly? This would allow us to handle ambiguous grammars, whilst preserving the natural structure of the parsing domain.

  \begin{wrapfigure}{r}{0.45\textwidth}
    \vspace{-5pt}
    \resizebox{0.45\textwidth}{!}{
      \begin{tikzpicture}
      [
        grow                    = right,
        sibling distance        = 3em,
        level distance          = 5em,
        edge from parent/.style = {draw, -latex},
        every node/.style       = {font=\footnotesize},
        sloped,
        treenode/.style = {shape=rectangle, rounded corners,
        draw, align=center,
        top color=white, bottom color=blue!20},
        root/.style     = {treenode, font=\tiny, bottom color=red!30},
        env/.style      = {treenode, font=\tiny},
        dummy/.style    = {circle,draw}
      ]
        \node [root] {S}
        child { node [env] {BC}
        child { node [root] {B}
        child { node [env] {RD}
        child { node [root] {R} edge from parent node [above] }
        child { node [root] {D} edge from parent node [above] }
        edge from parent node [above] }
        edge from parent node [below] }
        child { node [root] {C}
        child { node [env] {\ldots\vphantom{BB}} edge from parent node [below] }
%  child { edge from parent node [above] {\ldots} }
        edge from parent node [below] }
        edge from parent node [above] }
        child { node [env] {\ldots\vphantom{BB}} edge from parent node [below] }
        child { node [env] {AB}
        child { node [root] {A}
        child {
          node [env] {QC}
          child { node [root] {Q} edge from parent node [above] }
          child { node [root] {C} edge from parent node [above] }
          edge from parent node [above]
        }
%    child { node [env] {ZQ} edge from parent node [above] }
        child { node [env] {\ldots\vphantom{BB}} edge from parent node [below] }
        edge from parent node [below] }
        child { node [root] {B}
        child { node [env] {RD}
        child { node [root] {R} edge from parent node [above] }
        child { node [root] {D} edge from parent node [above] }
        edge from parent node [above] }
        edge from parent node [below] }
        edge from parent node [above] };
      \end{tikzpicture}
    }
    \caption{A partial $\mathbb{T}_2$ for the grammar $P=\{S \rightarrow BC \mid \ldots \mid AB, B\rightarrow RD \mid \ldots, A\rightarrow QC \mid \ldots\}$.}
    \label{fig:ptree}
    \vspace{-10pt}
  \end{wrapfigure}

  We will now describe a datatype for compactly representing candidate parse forests, then redefine the semiring algebra over indexed parse forests. This construction is particularly convenient for tracking provenance under ambiguity, Parikh mapping, counting the size of a finite CFL, and sampling trees with replacement using a PCFG, or without replacement using an integer pairing function.

  We first define a datatype $\mathbb{T}_3 = (V \cup \Sigma) \rightharpoonup \mathbb{T}_2$ where $\mathbb{T}_2 = (V \cup \Sigma) \times (\mathbb{N} \rightharpoonup \mathbb{T}_2\times\mathbb{T}_2)$\footnote{Given a $T:\mathbb{T}_2$, we may also refer to $\pi_1(T), \pi_2(T)$ as $\texttt{root}(T)$ and $\texttt{children}(T)$ respectively.}. Morally, we can think of $\mathbb{T}_2$ as an implicit set of possible trees that can be generated by a CFG in CNF, consistent with a finite-length porous string. We may also interpret $\mathbb{T}_2$ as as an algebraic data type corresponding to the fixpoints of the following recurrence, which tells us that each $\mathbb{T}_2$ can be a terminal, nonterminal, or a nonterminal and a sequence consisting of nonterminal pairs and their two children:\vspace{-10pt}

  \begin{equation}
    L(p) = 1 + p L(p) \phantom{addspace} P(a) = \Sigma + V + V L\big(V^2P(a)^2\big)
  \end{equation}

  Given a $\sigma: \underline\Sigma$, we construct $\mathbb{T}_2$ from the bottom-up, and sample from the top-down. Depicted in Fig.~\ref{fig:ptree} is a partial $\mathbb{T}_2$, where red nodes are \texttt{root}s and blue nodes are \texttt{children}.

  We construct the first upper diagonal $\hat\sigma_r = \Lambda(\sigma_r)$ as follows:

\vspace{-10pt}\begin{equation}
  \begin{footnotesize}
\Lambda(s: \underline\Sigma) \mapsto \begin{cases}
\bigoplus_{s'\in \Sigma} \Lambda(s') & \text{if $s$ is a hole,} \vspace{5pt}\\
\big\{\mathbb{T}_2\big(w, \big[\langle\mathbb{T}_2(s), \mathbb{T}_2(\varepsilon)\rangle\big]\big) \mid (w \rightarrow s)\in P\big\} & \text{otherwise.}
\end{cases}
  \end{footnotesize}
\end{equation}

\noindent This initializes the superdiagonal entries, enabling us to compute the fixpoint $M_\infty$ by redefining $\oplus, \otimes: \mathbb{T}_3 \times \mathbb{T}_3 \rightarrow \mathbb{T}_3$ as:

\begin{align}
  X \oplus Z &\mapsto \bigcup_{\mathclap{k\in \pi_1(X \cup Z)}}\Big\{k \Rightarrow \mathbb{T}_2(k, x \cup z) \mid x \in \pi_2(X\circ k), z \in \pi_2(Z\circ k)\Big\}\\
  X \otimes Z &\mapsto \bigoplus_{\mathclap{(w\rightarrow xz) \in P}}\Big\{\mathbb{T}_2\Big(w, \big[\langle X\circ x, Z\circ z\rangle\big]\Big) \mid x \in \pi_1(X), z \in \pi_1(Z)\Big\}
\end{align}

These operators group subtrees by their root nonterminal, then aggregate their children. Instead of tracking sets, each $\Lambda$ now becomes a dictionary of $\mathbb{T}_2$, indexed by their root nonterminals.


  $\mathbb{T}_2$ is a convenient datatype for many operations. We can use it to approximate the Parikh image, compute the size of a finite CFG, and sample parse trees with or without replacement. For example, to obtain the Parikh map, we may use the following recurrence,

\begin{equation}
  \pi(T: \mathbb{T}_2) \mapsto \begin{cases}
%    \big|\{s \mid \big(\texttt{root}(T) \rightarrow s\big) \in P^\cap\}\big| & \text{if $T$ is a leaf,} \\
  \big[[1, 1] \text{ if } \texttt{root}(T) = s \text{ else } [0, 0]\big]_{s\in \Sigma}  & \text{if $T$ is a leaf,} \\
  \bigoplus_{\langle T_1, T_2\rangle \in \texttt{children}(T)} \pi(T_1) \otimes \pi(T_2) & \text{otherwise.}
  \end{cases}
\end{equation}

  %infix fun IntRange.merge(other: IntRange) =
  %  minOf(start, other.first)..maxOf(last, other.last)
  %
  %operator fun ParikhBounds.plus(other: ParikhBounds) =
  %  (keys + other.keys).associateWith {
  %    (get(it) ?: 0..0) merge (other[it] ?: 0..0)
  %  }
  %
  %operator fun ParikhBounds.times(other: ParikhBounds) =
  %  (keys + other.keys).associateWith {
  %    (get(it) ?: 0..0) join (other[it] ?: 0..0)
  %  }
  %
  %infix fun IntRange.join(other: IntRange) =
  %  (first + other.first)..(last + other.last)

  \noindent where $\oplus, \otimes: \Pi \times \Pi \rightarrow \Pi$ are defined as follows:

  \begin{align}
      X \oplus Z &\mapsto \big[[\min(X_s \cup Z_s), \max(X_s \cup Z_s)]\big]_{s \in \Sigma}\\ X \otimes Z &\mapsto \big[[\min(X_s) + \min(Z_s), \max(X_s) + \max(Z_s)]\big]_{s \in \Sigma}
  \end{align}

  To obtain the Parikh interval for a CFG up to finite length bounds, we abstractly parse the string $(\_)^n$ for all $n \in [1, l_{\max}]$ and take the union of the intervals across $[|\err\sigma| - d_{\max}, |\err\sigma| + d_{\max}]$, representing every repair length in the Levenshtein ball.

  $\mathbb{T}_2$ also allows us to sample whole parse trees by obtaining $(\Lambda_\sigma^* \circ S): \mathbb{T}_2$. Given a probabilistic CFG whose productions indexed by each nonterminal are decorated with a probability vector $\mathbf{p}$ (this may be uniform in the non-probabilistic case), we define a tree sampler $\Gamma: (\mathbb{T}_2 \mid \mathbb{T}_2^2) \rightsquigarrow \mathbb{T}$ which recursively samples children according to a Multinoulli distribution:

\begin{equation}
  \Gamma(T) \mapsto \begin{cases}
        \texttt{BTree}\Big(\texttt{root}(T), \Gamma\big(\text{Multi}(\texttt{children}(T), \mathbf{p})\big)\Big) & \text{ if $T: \mathbb{T}_2$ } \\
        \big\langle \Gamma\big(\pi_1(T)\big), \Gamma\big(\pi_2(T)\big) \big\rangle & \text{ if $T: \mathbb{T}_2\times\mathbb{T}_2$ }
  \end{cases}
\end{equation}

This is closely related to the generating function for the ordinary Boltzmann sampler from analytic combinatorics,

\begin{equation}
  \Gamma C(x) \mapsto \begin{cases}
  \text{Bern} \left(\frac{A(x)}{A(x) + B(x)}\right) \rightarrow \Gamma A(x) \mid \Gamma B(x) & \text{ if } \mathcal{C}=\mathcal{A}+\mathcal{B} \\
  \big\langle \Gamma A(x), \Gamma B(x)\big\rangle & \text{ if } \mathcal{C}=\mathcal{A} \times \mathcal{B}
  \end{cases}
\end{equation}

\noindent however unlike Duchon et al.~\cite{duchon2004boltzmann}, our work does not depend on rejection to guarantee exact-size sampling, as all trees contained in $\mathbb{T}_2$ will necessarily be the same width.


The number of binary trees inhabiting a single instance of $\mathbb{T}_2$ is sensititive to the number of nonterminals and rule expansions in the grammar. To obtain the total number of trees with breadth $n$, we abstractly parse the porous string $\sigma = (\_)^n$, letting $T=\Lambda_{\sigma}^* \circ S$, then use the recurrence below to compute the total number of trees:

\begin{equation}
  |T: \mathbb{T}_2| \mapsto \begin{cases}
%    \big|\{s \mid \big(\texttt{root}(T) \rightarrow s\big) \in P^\cap\}\big| & \text{if $T$ is a leaf,} \\
    1 & \text{if $T$ is a leaf,} \\
    \sum_{\langle T_1, T_2\rangle \in \texttt{children}(T)} |T_1| \cdot |T_2| & \text{otherwise.}
  \end{cases}
\end{equation}

To sample all trees in a given $T: \mathbb{T}_2$ uniformly without replacement, we then construct a modular pairing function $\varphi: \mathbb{T}_2 \rightarrow \mathbb{Z}_{|T|} \rightarrow \texttt{BTree}$, that we define as follows:

\begin{equation}
  \varphi(T: \mathbb{T}_2, i: \mathbb{Z}_{|T|}) \mapsto \begin{cases}
  \Big\langle\texttt{BTree}\big(\texttt{root}(T)\big), i\Big\rangle & \text{if $T$ is a leaf,} \vspace{5pt}\\
  \text{Let } b = |\texttt{children}(T)|,\\
  \phantom{\text{Let }} q_1, r=\big\langle\lfloor\frac{i}{b}\rfloor, i \pmod{b}\big\rangle,\\
  \phantom{\text{Let }} lb, rb = \texttt{children}[r],\\
  \phantom{\text{Let }} T_1, q_2 = \varphi(lb, q_1),\\
  \phantom{\text{Let }} T_2, q_3 = \varphi(rb, q_2) \text{ in } \\
  \Big\langle\texttt{BTree}\big(\texttt{root}(T), T_1, T_2\big), q_3\Big\rangle & \text{otherwise.} \\
  \end{cases}
\end{equation}

If the language is suffiently small, instead of sampling trees, we can sample integers uniformly without replacement from $\mathbb{Z}_{|T|}$ then decode them into trees using $\varphi$. This procedure is the basis for our sampling algorithm and the method we use to decode repairs from the intersection grammar.

  \subsection{Ranked Repair}

  Returning to the ranked repair problem (Def.~\ref{def:ranked-repair}), the above procedure returns a set of strings, and we need an ordering over them. We note that any metric is sufficient, such as the perplexity of the string under a large language model, or the probability of the string under a PCFG. We the simplest possible solution: the negative log likelihood of the string computed by a variable-order Markov chain. This is a simple, fast, and effective metric for ranking strings, and as we will show, already yields competitive results in practice.

  \section{Evaluation}

  We consider the following research questions for our evaluation:

  \begin{itemize}
    \item \textbf{RQ 1}: What properties do natural repairs exhibit? (e.g., error frequency, edit distance)
    \item \textbf{RQ 2}: How precise is our technique at fixing syntax errors? (Precision@k vs. Seq2Parse)
    \item \textbf{RQ 3}: Which design choices are most significant? (e.g., search vs. sampling, n-gram order)
  \end{itemize}

  For our evaluation, we use the StackOverflow dataset from~\cite{hindle2012naturalness}. We preprocess the dataset to lexicalize both the broken and fixed code snippets, then filter the dataset by length and edit distance, in which all Python snippets whose broken form is fewer than 80 lexical tokens and whose human fix is under four Levenshtein edits is retained.

  \subsection{Dataset}

  In the following experiments, we use two datasets of Python snippets. The first is a set of 5,600 pairwise-aligned (broken, fixed) Python code snippets from Wong et al.'s StackOverflow dataset~\cite{wong2019syntax} shorter than 40 lexical tokens, whose patch sizes are less than five lexical tokens ($|\Sigma| = 50, |\err{\sigma}| \leq 40, \Delta(\err{\sigma}, \ell) \leq 4$).

  The StackOverflow dataset is comprised of 500k Python code snippets, each of which has been annotated with a human repair. We depict the normalized edit loations relative to the snippet length below.

  \begin{figure}[h!]
    \begin{tikzpicture}
      \begin{axis}[
      ybar,
      bar width=5pt,
      xlabel={Length},
      ylabel={Frequency},
      title={Snippet Length},
      axis x line*=bottom,
      axis y line*=left,
      ymin=0,
      ymax=65,
      xtick=data,
      xticklabels={,\leq 20,,\leq 40,,\leq 60,,\leq 80,,\leq 100},
      ymajorgrids=true,
      grid style=dashed,
      width=0.45\textwidth,
      height=0.3\textwidth
      ]

      \addplot table {
        X Y
        1 7.60
        2 14.52
        3 22.01
        4 30.54
        5 37.82
        6 44.30
        7 49.68
        8 55.21
        9 59.75
        10 63.59
      };
      \end{axis}
    \end{tikzpicture}
    \begin{tikzpicture}
      \begin{axis}[
        ybar,
        bar width=5pt,
        title={Human Repair Size},
        xlabel={$\Delta(\err\sigma, \sigma)$},
        ylabel={Frequency},
        axis x line*=bottom,
        axis y line*=left,
        xtick=data,
        ymajorgrids=true,
        grid style=dashed,
        xticklabels={,\leq 2,,\leq 4,,\leq 6,,\leq 8,,\leq 10},
        width=0.45\textwidth,
        height=0.3\textwidth
      ]
        \addplot table {
          X Y
          1  31.48
          2  47.52
          3  54.89
          4  60.44
          5  63.88
          6  66.38
          7  68.02
          8  70.04
          9  71.49
          10 72.22
        };
      \end{axis}
    \end{tikzpicture}
    \begin{tikzpicture}
      \begin{axis}[
      ybar,
      bar width=5pt,
      xlabel={Beginning $\longleftrightarrow$ End},
      ylabel={Frequency},
      title={Normalized edit locations},
      axis x line*=bottom,
      axis y line*=left,
      ymin=0,
      ymax=35,
      xtick=data,
      xticklabels={,20\%,,40\%,,60\%,,80\%,,100\%},
      ymajorgrids=true,
      grid style=dashed,
      width=0.45\textwidth,
      height=0.3\textwidth
      ]

      \addplot table {
        X Y
        10 11.6539
        20 5.7252
        30 6.2087
        40 5.9542
        50 5.5980
        60 7.9389
        70 7.0738
        80 6.9466
        90 12.4173
        100 30.4835
      };
      \end{axis}
    \end{tikzpicture}
    \begin{tikzpicture}
      \begin{axis}[
        ybar,
        bar width=5pt,
        title={Intra-patch edit distance},
        xlabel={Caret distance},
        ylabel={Frequency},
        axis x line*=bottom,
        axis y line*=left,
        xtick=data,
        ymajorgrids=true,
        grid style=dashed,
        xticklabels={1,2,3,4,5,6,7,8,9,10+},
        width=0.45\textwidth,
        height=0.3\textwidth
      ]

        \addplot table {
          X Y
          1 40.66
          2 15.00
          3 5.80
          4 4.86
          5 4.26
          6 2.98
          7 2.05
          8 2.73
          9 1.62
          10 13.64
        };
      \end{axis}
    \end{tikzpicture}
    \caption{Statistics on properties of each repair, across the StackOverflow pairwise dataset.}\label{fig:patch_stats}
  \end{figure}

  In the first set of experiments, we measure the precision@k of our repair procedure at recovering human repairs of varying edit distances and latency cutoffs, where both the broken and fixed code are written by a human.

  In the second set of experiments, we use a dataset of valid Python snippets from BIFI~\cite{yasunaga2021break}, then synthetically corrupt them by introducing syntax errors sampled from a distribution trained from the StackOverflow dataset. We then measure the precision@k of our repair procedure at recovering the original, uncorrupted snippet.

  To train the scoring function, we use a length-5 variable-order Markov (VOM) chain implemented using a count-min sketch based on Apache Datasketches~\cite{apache2022datasketches}. Training on 55 million StackOverflow tokens took roughly 10 minutes, after which calculating perplexity is nearly instantaneous. Sequences are scored using negative log likelihood with Laplace smoothing and our evaluation measures the precision@\{1, 5, 10, All\} for samples at varying latency cutoffs.

  Both sets of experiments were conducted on 40 Intel Skylake cores running at 2.4 GHz, with 16 GB of RAM, running bytecode compiled for JVM 17.0.2. We measure the precision using abstract lexical matching, following the Seq2Parse~\cite{sakkas2022seq2parse} evaluation, and give a baseline for their approach on the same dataset.

  \subsection{StackOverflow Evaluation}

  For our first experiment, we run the sampler until the human repair is detected, then measure the number of samples required to draw the exact human repair across varying Levenshtein radii.

  \begin{figure}[h!]
%  % This file was created with tikzplotlib v0.10.1.
\begin{tikzpicture}

\definecolor{darkgray176}{RGB}{176,176,176}
\definecolor{darkorange25512714}{RGB}{255,127,14}
\definecolor{forestgreen4416044}{RGB}{44,160,44}
\definecolor{steelblue31119180}{RGB}{31,119,180}

\begin{axis}[
width=\textwidth,
height=.4\textwidth,
log basis x={10},
tick align=outside,
tick pos=left,
axis lines*=left,
title={Sample Efficiency},
legend style={fill opacity=0.8, draw opacity=1, text opacity=1, legend columns=1, legend pos=north west, at={(0.055, 0.96)}},
x grid style={darkgray176},
xlabel={Samples drawn (log scale)},
xmin=0.34892211545542, xmax=1001,
xmode=log,
%xtick style={color=black},
xtick={0.01,0.1,1,10,100,1000,10000},
xticklabels={
  \(\displaystyle 10^-2\),
  \(\displaystyle 10^-1\),
  \(\displaystyle 10^2\),
  \(\displaystyle 10^3\),
  \(\displaystyle 10^4\),
  \(\displaystyle 10^5\),
  \(\displaystyle 10^6\)
},
y grid style={darkgray176},
ylabel={Precision@All},
ymin=0, ymax=101,
%ytick style={color=black}
ytick={0,20,40,60,80,100}
]
\draw[draw=none,fill=steelblue31119180] (axis cs:-0.5,0) rectangle (axis cs:0.5,100);
\addlegendimage{ybar,ybar legend,draw=none,fill=steelblue31119180}
\addlegendentry{$\Delta_L=1$}

\draw[draw=none,fill=steelblue31119180] (axis cs:0.5,0) rectangle (axis cs:1.5,100);
\draw[draw=none,fill=steelblue31119180] (axis cs:1.5,0) rectangle (axis cs:2.5,100);
\draw[draw=none,fill=steelblue31119180] (axis cs:2.5,0) rectangle (axis cs:3.5,100);
\draw[draw=none,fill=steelblue31119180] (axis cs:3.5,0) rectangle (axis cs:4.5,100);
\draw[draw=none,fill=steelblue31119180] (axis cs:4.5,0) rectangle (axis cs:5.5,100);
\draw[draw=none,fill=steelblue31119180] (axis cs:5.5,0) rectangle (axis cs:6.5,100);
\draw[draw=none,fill=steelblue31119180] (axis cs:6.5,0) rectangle (axis cs:7.5,100);
\draw[draw=none,fill=steelblue31119180] (axis cs:7.5,0) rectangle (axis cs:8.5,100);
\draw[draw=none,fill=steelblue31119180] (axis cs:8.5,0) rectangle (axis cs:9.5,100);
\draw[draw=none,fill=steelblue31119180] (axis cs:9.5,0) rectangle (axis cs:10.5,100);
\draw[draw=none,fill=steelblue31119180] (axis cs:10.5,0) rectangle (axis cs:11.5,100);
\draw[draw=none,fill=steelblue31119180] (axis cs:11.5,0) rectangle (axis cs:12.5,100);
\draw[draw=none,fill=steelblue31119180] (axis cs:12.5,0) rectangle (axis cs:13.5,100);
\draw[draw=none,fill=steelblue31119180] (axis cs:13.5,0) rectangle (axis cs:14.5,100);
\draw[draw=none,fill=steelblue31119180] (axis cs:14.5,0) rectangle (axis cs:15.5,100);
\draw[draw=none,fill=steelblue31119180] (axis cs:15.5,0) rectangle (axis cs:16.5,100);
\draw[draw=none,fill=steelblue31119180] (axis cs:16.5,0) rectangle (axis cs:17.5,100);
\draw[draw=none,fill=steelblue31119180] (axis cs:17.5,0) rectangle (axis cs:18.5,100);
\draw[draw=none,fill=steelblue31119180] (axis cs:18.5,0) rectangle (axis cs:19.5,100);
\draw[draw=none,fill=steelblue31119180] (axis cs:19.5,0) rectangle (axis cs:20.5,100);
\draw[draw=none,fill=steelblue31119180] (axis cs:20.5,0) rectangle (axis cs:21.5,100);
\draw[draw=none,fill=steelblue31119180] (axis cs:21.5,0) rectangle (axis cs:22.5,100);
\draw[draw=none,fill=steelblue31119180] (axis cs:22.5,0) rectangle (axis cs:23.5,100);
\draw[draw=none,fill=steelblue31119180] (axis cs:23.5,0) rectangle (axis cs:24.5,100);
\draw[draw=none,fill=steelblue31119180] (axis cs:24.5,0) rectangle (axis cs:25.5,100);
\draw[draw=none,fill=steelblue31119180] (axis cs:25.5,0) rectangle (axis cs:26.5,100);
\draw[draw=none,fill=steelblue31119180] (axis cs:26.5,0) rectangle (axis cs:27.5,100);
\draw[draw=none,fill=steelblue31119180] (axis cs:27.5,0) rectangle (axis cs:28.5,100);
\draw[draw=none,fill=steelblue31119180] (axis cs:28.5,0) rectangle (axis cs:29.5,100);
\draw[draw=none,fill=steelblue31119180] (axis cs:29.5,0) rectangle (axis cs:30.5,100);
\draw[draw=none,fill=steelblue31119180] (axis cs:30.5,0) rectangle (axis cs:31.5,100);
\draw[draw=none,fill=steelblue31119180] (axis cs:31.5,0) rectangle (axis cs:32.5,100);
\draw[draw=none,fill=steelblue31119180] (axis cs:32.5,0) rectangle (axis cs:33.5,100);
\draw[draw=none,fill=steelblue31119180] (axis cs:33.5,0) rectangle (axis cs:34.5,100);
\draw[draw=none,fill=steelblue31119180] (axis cs:34.5,0) rectangle (axis cs:35.5,100);
\draw[draw=none,fill=steelblue31119180] (axis cs:35.5,0) rectangle (axis cs:36.5,100);
\draw[draw=none,fill=steelblue31119180] (axis cs:36.5,0) rectangle (axis cs:37.5,100);
\draw[draw=none,fill=steelblue31119180] (axis cs:37.5,0) rectangle (axis cs:38.5,100);
\draw[draw=none,fill=steelblue31119180] (axis cs:38.5,0) rectangle (axis cs:39.5,100);
\draw[draw=none,fill=steelblue31119180] (axis cs:39.5,0) rectangle (axis cs:40.5,100);
\draw[draw=none,fill=steelblue31119180] (axis cs:40.5,0) rectangle (axis cs:41.5,100);
\draw[draw=none,fill=steelblue31119180] (axis cs:41.5,0) rectangle (axis cs:42.5,100);
\draw[draw=none,fill=steelblue31119180] (axis cs:42.5,0) rectangle (axis cs:43.5,100);
\draw[draw=none,fill=steelblue31119180] (axis cs:43.5,0) rectangle (axis cs:44.5,100);
\draw[draw=none,fill=steelblue31119180] (axis cs:44.5,0) rectangle (axis cs:45.5,100);
\draw[draw=none,fill=steelblue31119180] (axis cs:45.5,0) rectangle (axis cs:46.5,100);
\draw[draw=none,fill=steelblue31119180] (axis cs:46.5,0) rectangle (axis cs:47.5,100);
\draw[draw=none,fill=steelblue31119180] (axis cs:47.5,0) rectangle (axis cs:48.5,100);
\draw[draw=none,fill=steelblue31119180] (axis cs:48.5,0) rectangle (axis cs:49.5,100);
\draw[draw=none,fill=steelblue31119180] (axis cs:49.5,0) rectangle (axis cs:50.5,100);
\draw[draw=none,fill=steelblue31119180] (axis cs:50.5,0) rectangle (axis cs:51.5,100);
\draw[draw=none,fill=steelblue31119180] (axis cs:51.5,0) rectangle (axis cs:52.5,100);
\draw[draw=none,fill=steelblue31119180] (axis cs:52.5,0) rectangle (axis cs:53.5,100);
\draw[draw=none,fill=steelblue31119180] (axis cs:53.5,0) rectangle (axis cs:54.5,100);
\draw[draw=none,fill=steelblue31119180] (axis cs:54.5,0) rectangle (axis cs:55.5,100);
\draw[draw=none,fill=steelblue31119180] (axis cs:55.5,0) rectangle (axis cs:56.5,100);
\draw[draw=none,fill=steelblue31119180] (axis cs:56.5,0) rectangle (axis cs:57.5,100);
\draw[draw=none,fill=steelblue31119180] (axis cs:57.5,0) rectangle (axis cs:58.5,100);
\draw[draw=none,fill=steelblue31119180] (axis cs:58.5,0) rectangle (axis cs:59.5,100);
\draw[draw=none,fill=steelblue31119180] (axis cs:59.5,0) rectangle (axis cs:60.5,100);
\draw[draw=none,fill=steelblue31119180] (axis cs:60.5,0) rectangle (axis cs:61.5,100);
\draw[draw=none,fill=steelblue31119180] (axis cs:61.5,0) rectangle (axis cs:62.5,100);
\draw[draw=none,fill=steelblue31119180] (axis cs:62.5,0) rectangle (axis cs:63.5,100);
\draw[draw=none,fill=steelblue31119180] (axis cs:63.5,0) rectangle (axis cs:64.5,100);
\draw[draw=none,fill=steelblue31119180] (axis cs:64.5,0) rectangle (axis cs:65.5,100);
\draw[draw=none,fill=steelblue31119180] (axis cs:65.5,0) rectangle (axis cs:66.5,100);
\draw[draw=none,fill=steelblue31119180] (axis cs:66.5,0) rectangle (axis cs:67.5,100);
\draw[draw=none,fill=steelblue31119180] (axis cs:67.5,0) rectangle (axis cs:68.5,100);
\draw[draw=none,fill=steelblue31119180] (axis cs:68.5,0) rectangle (axis cs:69.5,100);
\draw[draw=none,fill=steelblue31119180] (axis cs:69.5,0) rectangle (axis cs:70.5,100);
\draw[draw=none,fill=steelblue31119180] (axis cs:70.5,0) rectangle (axis cs:71.5,100);
\draw[draw=none,fill=steelblue31119180] (axis cs:71.5,0) rectangle (axis cs:72.5,100);
\draw[draw=none,fill=steelblue31119180] (axis cs:72.5,0) rectangle (axis cs:73.5,100);
\draw[draw=none,fill=steelblue31119180] (axis cs:73.5,0) rectangle (axis cs:74.5,100);
\draw[draw=none,fill=steelblue31119180] (axis cs:74.5,0) rectangle (axis cs:75.5,100);
\draw[draw=none,fill=steelblue31119180] (axis cs:75.5,0) rectangle (axis cs:76.5,100);
\draw[draw=none,fill=steelblue31119180] (axis cs:76.5,0) rectangle (axis cs:77.5,100);
\draw[draw=none,fill=steelblue31119180] (axis cs:77.5,0) rectangle (axis cs:78.5,100);
\draw[draw=none,fill=steelblue31119180] (axis cs:78.5,0) rectangle (axis cs:79.5,100);
\draw[draw=none,fill=steelblue31119180] (axis cs:79.5,0) rectangle (axis cs:80.5,100);
\draw[draw=none,fill=steelblue31119180] (axis cs:80.5,0) rectangle (axis cs:81.5,100);
\draw[draw=none,fill=steelblue31119180] (axis cs:81.5,0) rectangle (axis cs:82.5,100);
\draw[draw=none,fill=steelblue31119180] (axis cs:82.5,0) rectangle (axis cs:83.5,100);
\draw[draw=none,fill=steelblue31119180] (axis cs:83.5,0) rectangle (axis cs:84.5,100);
\draw[draw=none,fill=steelblue31119180] (axis cs:84.5,0) rectangle (axis cs:85.5,100);
\draw[draw=none,fill=steelblue31119180] (axis cs:85.5,0) rectangle (axis cs:86.5,100);
\draw[draw=none,fill=steelblue31119180] (axis cs:86.5,0) rectangle (axis cs:87.5,100);
\draw[draw=none,fill=steelblue31119180] (axis cs:87.5,0) rectangle (axis cs:88.5,100);
\draw[draw=none,fill=steelblue31119180] (axis cs:88.5,0) rectangle (axis cs:89.5,100);
\draw[draw=none,fill=steelblue31119180] (axis cs:89.5,0) rectangle (axis cs:90.5,100);
\draw[draw=none,fill=steelblue31119180] (axis cs:90.5,0) rectangle (axis cs:91.5,100);
\draw[draw=none,fill=steelblue31119180] (axis cs:91.5,0) rectangle (axis cs:92.5,100);
\draw[draw=none,fill=steelblue31119180] (axis cs:92.5,0) rectangle (axis cs:93.5,100);
\draw[draw=none,fill=steelblue31119180] (axis cs:93.5,0) rectangle (axis cs:94.5,100);
\draw[draw=none,fill=steelblue31119180] (axis cs:94.5,0) rectangle (axis cs:95.5,100);
\draw[draw=none,fill=steelblue31119180] (axis cs:95.5,0) rectangle (axis cs:96.5,100);
\draw[draw=none,fill=steelblue31119180] (axis cs:96.5,0) rectangle (axis cs:97.5,100);
\draw[draw=none,fill=steelblue31119180] (axis cs:97.5,0) rectangle (axis cs:98.5,100);
\draw[draw=none,fill=steelblue31119180] (axis cs:98.5,0) rectangle (axis cs:99.5,100);
\draw[draw=none,fill=steelblue31119180] (axis cs:99.5,0) rectangle (axis cs:100.5,100);
\draw[draw=none,fill=steelblue31119180] (axis cs:100.5,0) rectangle (axis cs:101.5,100);
\draw[draw=none,fill=steelblue31119180] (axis cs:101.5,0) rectangle (axis cs:102.5,100);
\draw[draw=none,fill=steelblue31119180] (axis cs:102.5,0) rectangle (axis cs:103.5,100);
\draw[draw=none,fill=steelblue31119180] (axis cs:103.5,0) rectangle (axis cs:104.5,100);
\draw[draw=none,fill=steelblue31119180] (axis cs:104.5,0) rectangle (axis cs:105.5,100);
\draw[draw=none,fill=steelblue31119180] (axis cs:105.5,0) rectangle (axis cs:106.5,100);
\draw[draw=none,fill=steelblue31119180] (axis cs:106.5,0) rectangle (axis cs:107.5,100);
\draw[draw=none,fill=steelblue31119180] (axis cs:107.5,0) rectangle (axis cs:108.5,100);
\draw[draw=none,fill=steelblue31119180] (axis cs:108.5,0) rectangle (axis cs:109.5,100);
\draw[draw=none,fill=steelblue31119180] (axis cs:109.5,0) rectangle (axis cs:110.5,100);
\draw[draw=none,fill=steelblue31119180] (axis cs:110.5,0) rectangle (axis cs:111.5,100);
\draw[draw=none,fill=steelblue31119180] (axis cs:111.5,0) rectangle (axis cs:112.5,100);
\draw[draw=none,fill=steelblue31119180] (axis cs:112.5,0) rectangle (axis cs:113.5,100);
\draw[draw=none,fill=steelblue31119180] (axis cs:113.5,0) rectangle (axis cs:114.5,100);
\draw[draw=none,fill=steelblue31119180] (axis cs:114.5,0) rectangle (axis cs:115.5,100);
\draw[draw=none,fill=steelblue31119180] (axis cs:115.5,0) rectangle (axis cs:116.5,100);
\draw[draw=none,fill=steelblue31119180] (axis cs:116.5,0) rectangle (axis cs:117.5,100);
\draw[draw=none,fill=steelblue31119180] (axis cs:117.5,0) rectangle (axis cs:118.5,100);
\draw[draw=none,fill=steelblue31119180] (axis cs:118.5,0) rectangle (axis cs:119.5,100);
\draw[draw=none,fill=steelblue31119180] (axis cs:119.5,0) rectangle (axis cs:120.5,100);
\draw[draw=none,fill=steelblue31119180] (axis cs:120.5,0) rectangle (axis cs:121.5,100);
\draw[draw=none,fill=steelblue31119180] (axis cs:121.5,0) rectangle (axis cs:122.5,100);
\draw[draw=none,fill=steelblue31119180] (axis cs:122.5,0) rectangle (axis cs:123.5,100);
\draw[draw=none,fill=steelblue31119180] (axis cs:123.5,0) rectangle (axis cs:124.5,100);
\draw[draw=none,fill=steelblue31119180] (axis cs:124.5,0) rectangle (axis cs:125.5,100);
\draw[draw=none,fill=steelblue31119180] (axis cs:125.5,0) rectangle (axis cs:126.5,100);
\draw[draw=none,fill=steelblue31119180] (axis cs:126.5,0) rectangle (axis cs:127.5,100);
\draw[draw=none,fill=steelblue31119180] (axis cs:127.5,0) rectangle (axis cs:128.5,100);
\draw[draw=none,fill=steelblue31119180] (axis cs:128.5,0) rectangle (axis cs:129.5,100);
\draw[draw=none,fill=steelblue31119180] (axis cs:129.5,0) rectangle (axis cs:130.5,100);
\draw[draw=none,fill=steelblue31119180] (axis cs:130.5,0) rectangle (axis cs:131.5,100);
\draw[draw=none,fill=steelblue31119180] (axis cs:131.5,0) rectangle (axis cs:132.5,100);
\draw[draw=none,fill=steelblue31119180] (axis cs:132.5,0) rectangle (axis cs:133.5,100);
\draw[draw=none,fill=steelblue31119180] (axis cs:133.5,0) rectangle (axis cs:134.5,100);
\draw[draw=none,fill=steelblue31119180] (axis cs:134.5,0) rectangle (axis cs:135.5,100);
\draw[draw=none,fill=steelblue31119180] (axis cs:135.5,0) rectangle (axis cs:136.5,100);
\draw[draw=none,fill=steelblue31119180] (axis cs:136.5,0) rectangle (axis cs:137.5,100);
\draw[draw=none,fill=steelblue31119180] (axis cs:137.5,0) rectangle (axis cs:138.5,100);
\draw[draw=none,fill=steelblue31119180] (axis cs:138.5,0) rectangle (axis cs:139.5,100);
\draw[draw=none,fill=steelblue31119180] (axis cs:139.5,0) rectangle (axis cs:140.5,100);
\draw[draw=none,fill=steelblue31119180] (axis cs:140.5,0) rectangle (axis cs:141.5,100);
\draw[draw=none,fill=steelblue31119180] (axis cs:141.5,0) rectangle (axis cs:142.5,100);
\draw[draw=none,fill=steelblue31119180] (axis cs:142.5,0) rectangle (axis cs:143.5,100);
\draw[draw=none,fill=steelblue31119180] (axis cs:143.5,0) rectangle (axis cs:144.5,100);
\draw[draw=none,fill=steelblue31119180] (axis cs:144.5,0) rectangle (axis cs:145.5,100);
\draw[draw=none,fill=steelblue31119180] (axis cs:145.5,0) rectangle (axis cs:146.5,100);
\draw[draw=none,fill=steelblue31119180] (axis cs:146.5,0) rectangle (axis cs:147.5,100);
\draw[draw=none,fill=steelblue31119180] (axis cs:147.5,0) rectangle (axis cs:148.5,100);
\draw[draw=none,fill=steelblue31119180] (axis cs:148.5,0) rectangle (axis cs:149.5,100);
\draw[draw=none,fill=steelblue31119180] (axis cs:149.5,0) rectangle (axis cs:150.5,100);
\draw[draw=none,fill=steelblue31119180] (axis cs:150.5,0) rectangle (axis cs:151.5,100);
\draw[draw=none,fill=steelblue31119180] (axis cs:151.5,0) rectangle (axis cs:152.5,100);
\draw[draw=none,fill=steelblue31119180] (axis cs:152.5,0) rectangle (axis cs:153.5,100);
\draw[draw=none,fill=steelblue31119180] (axis cs:153.5,0) rectangle (axis cs:154.5,100);
\draw[draw=none,fill=steelblue31119180] (axis cs:154.5,0) rectangle (axis cs:155.5,100);
\draw[draw=none,fill=steelblue31119180] (axis cs:155.5,0) rectangle (axis cs:156.5,100);
\draw[draw=none,fill=steelblue31119180] (axis cs:156.5,0) rectangle (axis cs:157.5,100);
\draw[draw=none,fill=steelblue31119180] (axis cs:157.5,0) rectangle (axis cs:158.5,100);
\draw[draw=none,fill=steelblue31119180] (axis cs:158.5,0) rectangle (axis cs:159.5,100);
\draw[draw=none,fill=steelblue31119180] (axis cs:159.5,0) rectangle (axis cs:160.5,100);
\draw[draw=none,fill=steelblue31119180] (axis cs:160.5,0) rectangle (axis cs:161.5,100);
\draw[draw=none,fill=steelblue31119180] (axis cs:161.5,0) rectangle (axis cs:162.5,100);
\draw[draw=none,fill=steelblue31119180] (axis cs:162.5,0) rectangle (axis cs:163.5,100);
\draw[draw=none,fill=steelblue31119180] (axis cs:163.5,0) rectangle (axis cs:164.5,100);
\draw[draw=none,fill=steelblue31119180] (axis cs:164.5,0) rectangle (axis cs:165.5,100);
\draw[draw=none,fill=steelblue31119180] (axis cs:165.5,0) rectangle (axis cs:166.5,100);
\draw[draw=none,fill=steelblue31119180] (axis cs:166.5,0) rectangle (axis cs:167.5,100);
\draw[draw=none,fill=steelblue31119180] (axis cs:167.5,0) rectangle (axis cs:168.5,100);
\draw[draw=none,fill=steelblue31119180] (axis cs:168.5,0) rectangle (axis cs:169.5,100);
\draw[draw=none,fill=steelblue31119180] (axis cs:169.5,0) rectangle (axis cs:170.5,100);
\draw[draw=none,fill=steelblue31119180] (axis cs:170.5,0) rectangle (axis cs:171.5,100);
\draw[draw=none,fill=steelblue31119180] (axis cs:171.5,0) rectangle (axis cs:172.5,100);
\draw[draw=none,fill=steelblue31119180] (axis cs:172.5,0) rectangle (axis cs:173.5,100);
\draw[draw=none,fill=steelblue31119180] (axis cs:173.5,0) rectangle (axis cs:174.5,100);
\draw[draw=none,fill=steelblue31119180] (axis cs:174.5,0) rectangle (axis cs:175.5,100);
\draw[draw=none,fill=steelblue31119180] (axis cs:175.5,0) rectangle (axis cs:176.5,100);
\draw[draw=none,fill=steelblue31119180] (axis cs:176.5,0) rectangle (axis cs:177.5,100);
\draw[draw=none,fill=steelblue31119180] (axis cs:177.5,0) rectangle (axis cs:178.5,100);
\draw[draw=none,fill=steelblue31119180] (axis cs:178.5,0) rectangle (axis cs:179.5,100);
\draw[draw=none,fill=steelblue31119180] (axis cs:179.5,0) rectangle (axis cs:180.5,100);
\draw[draw=none,fill=steelblue31119180] (axis cs:180.5,0) rectangle (axis cs:181.5,100);
\draw[draw=none,fill=steelblue31119180] (axis cs:181.5,0) rectangle (axis cs:182.5,100);
\draw[draw=none,fill=steelblue31119180] (axis cs:182.5,0) rectangle (axis cs:183.5,100);
\draw[draw=none,fill=steelblue31119180] (axis cs:183.5,0) rectangle (axis cs:184.5,100);
\draw[draw=none,fill=steelblue31119180] (axis cs:184.5,0) rectangle (axis cs:185.5,100);
\draw[draw=none,fill=steelblue31119180] (axis cs:185.5,0) rectangle (axis cs:186.5,100);
\draw[draw=none,fill=steelblue31119180] (axis cs:186.5,0) rectangle (axis cs:187.5,100);
\draw[draw=none,fill=steelblue31119180] (axis cs:187.5,0) rectangle (axis cs:188.5,100);
\draw[draw=none,fill=steelblue31119180] (axis cs:188.5,0) rectangle (axis cs:189.5,100);
\draw[draw=none,fill=steelblue31119180] (axis cs:189.5,0) rectangle (axis cs:190.5,100);
\draw[draw=none,fill=steelblue31119180] (axis cs:190.5,0) rectangle (axis cs:191.5,100);
\draw[draw=none,fill=steelblue31119180] (axis cs:191.5,0) rectangle (axis cs:192.5,100);
\draw[draw=none,fill=steelblue31119180] (axis cs:192.5,0) rectangle (axis cs:193.5,100);
\draw[draw=none,fill=steelblue31119180] (axis cs:193.5,0) rectangle (axis cs:194.5,100);
\draw[draw=none,fill=steelblue31119180] (axis cs:194.5,0) rectangle (axis cs:195.5,100);
\draw[draw=none,fill=steelblue31119180] (axis cs:195.5,0) rectangle (axis cs:196.5,100);
\draw[draw=none,fill=steelblue31119180] (axis cs:196.5,0) rectangle (axis cs:197.5,100);
\draw[draw=none,fill=steelblue31119180] (axis cs:197.5,0) rectangle (axis cs:198.5,100);
\draw[draw=none,fill=steelblue31119180] (axis cs:198.5,0) rectangle (axis cs:199.5,100);
\draw[draw=none,fill=steelblue31119180] (axis cs:199.5,0) rectangle (axis cs:200.5,100);
\draw[draw=none,fill=steelblue31119180] (axis cs:200.5,0) rectangle (axis cs:201.5,100);
\draw[draw=none,fill=steelblue31119180] (axis cs:201.5,0) rectangle (axis cs:202.5,100);
\draw[draw=none,fill=steelblue31119180] (axis cs:202.5,0) rectangle (axis cs:203.5,100);
\draw[draw=none,fill=steelblue31119180] (axis cs:203.5,0) rectangle (axis cs:204.5,100);
\draw[draw=none,fill=steelblue31119180] (axis cs:204.5,0) rectangle (axis cs:205.5,100);
\draw[draw=none,fill=steelblue31119180] (axis cs:205.5,0) rectangle (axis cs:206.5,100);
\draw[draw=none,fill=steelblue31119180] (axis cs:206.5,0) rectangle (axis cs:207.5,100);
\draw[draw=none,fill=steelblue31119180] (axis cs:207.5,0) rectangle (axis cs:208.5,100);
\draw[draw=none,fill=steelblue31119180] (axis cs:208.5,0) rectangle (axis cs:209.5,100);
\draw[draw=none,fill=steelblue31119180] (axis cs:209.5,0) rectangle (axis cs:210.5,100);
\draw[draw=none,fill=steelblue31119180] (axis cs:210.5,0) rectangle (axis cs:211.5,100);
\draw[draw=none,fill=steelblue31119180] (axis cs:211.5,0) rectangle (axis cs:212.5,100);
\draw[draw=none,fill=steelblue31119180] (axis cs:212.5,0) rectangle (axis cs:213.5,100);
\draw[draw=none,fill=steelblue31119180] (axis cs:213.5,0) rectangle (axis cs:214.5,100);
\draw[draw=none,fill=steelblue31119180] (axis cs:214.5,0) rectangle (axis cs:215.5,100);
\draw[draw=none,fill=steelblue31119180] (axis cs:215.5,0) rectangle (axis cs:216.5,100);
\draw[draw=none,fill=steelblue31119180] (axis cs:216.5,0) rectangle (axis cs:217.5,100);
\draw[draw=none,fill=steelblue31119180] (axis cs:217.5,0) rectangle (axis cs:218.5,100);
\draw[draw=none,fill=steelblue31119180] (axis cs:218.5,0) rectangle (axis cs:219.5,100);
\draw[draw=none,fill=steelblue31119180] (axis cs:219.5,0) rectangle (axis cs:220.5,100);
\draw[draw=none,fill=steelblue31119180] (axis cs:220.5,0) rectangle (axis cs:221.5,100);
\draw[draw=none,fill=steelblue31119180] (axis cs:221.5,0) rectangle (axis cs:222.5,100);
\draw[draw=none,fill=steelblue31119180] (axis cs:222.5,0) rectangle (axis cs:223.5,100);
\draw[draw=none,fill=steelblue31119180] (axis cs:223.5,0) rectangle (axis cs:224.5,100);
\draw[draw=none,fill=steelblue31119180] (axis cs:224.5,0) rectangle (axis cs:225.5,100);
\draw[draw=none,fill=steelblue31119180] (axis cs:225.5,0) rectangle (axis cs:226.5,100);
\draw[draw=none,fill=steelblue31119180] (axis cs:226.5,0) rectangle (axis cs:227.5,100);
\draw[draw=none,fill=steelblue31119180] (axis cs:227.5,0) rectangle (axis cs:228.5,100);
\draw[draw=none,fill=steelblue31119180] (axis cs:228.5,0) rectangle (axis cs:229.5,100);
\draw[draw=none,fill=steelblue31119180] (axis cs:229.5,0) rectangle (axis cs:230.5,100);
\draw[draw=none,fill=steelblue31119180] (axis cs:230.5,0) rectangle (axis cs:231.5,100);
\draw[draw=none,fill=steelblue31119180] (axis cs:231.5,0) rectangle (axis cs:232.5,100);
\draw[draw=none,fill=steelblue31119180] (axis cs:232.5,0) rectangle (axis cs:233.5,100);
\draw[draw=none,fill=steelblue31119180] (axis cs:233.5,0) rectangle (axis cs:234.5,100);
\draw[draw=none,fill=steelblue31119180] (axis cs:234.5,0) rectangle (axis cs:235.5,100);
\draw[draw=none,fill=steelblue31119180] (axis cs:235.5,0) rectangle (axis cs:236.5,100);
\draw[draw=none,fill=steelblue31119180] (axis cs:236.5,0) rectangle (axis cs:237.5,100);
\draw[draw=none,fill=steelblue31119180] (axis cs:237.5,0) rectangle (axis cs:238.5,100);
\draw[draw=none,fill=steelblue31119180] (axis cs:238.5,0) rectangle (axis cs:239.5,100);
\draw[draw=none,fill=steelblue31119180] (axis cs:239.5,0) rectangle (axis cs:240.5,100);
\draw[draw=none,fill=steelblue31119180] (axis cs:240.5,0) rectangle (axis cs:241.5,100);
\draw[draw=none,fill=steelblue31119180] (axis cs:241.5,0) rectangle (axis cs:242.5,100);
\draw[draw=none,fill=steelblue31119180] (axis cs:242.5,0) rectangle (axis cs:243.5,100);
\draw[draw=none,fill=steelblue31119180] (axis cs:243.5,0) rectangle (axis cs:244.5,100);
\draw[draw=none,fill=steelblue31119180] (axis cs:244.5,0) rectangle (axis cs:245.5,100);
\draw[draw=none,fill=steelblue31119180] (axis cs:245.5,0) rectangle (axis cs:246.5,100);
\draw[draw=none,fill=steelblue31119180] (axis cs:246.5,0) rectangle (axis cs:247.5,100);
\draw[draw=none,fill=steelblue31119180] (axis cs:247.5,0) rectangle (axis cs:248.5,100);
\draw[draw=none,fill=steelblue31119180] (axis cs:248.5,0) rectangle (axis cs:249.5,100);
\draw[draw=none,fill=steelblue31119180] (axis cs:249.5,0) rectangle (axis cs:250.5,100);
\draw[draw=none,fill=steelblue31119180] (axis cs:250.5,0) rectangle (axis cs:251.5,100);
\draw[draw=none,fill=steelblue31119180] (axis cs:251.5,0) rectangle (axis cs:252.5,100);
\draw[draw=none,fill=steelblue31119180] (axis cs:252.5,0) rectangle (axis cs:253.5,100);
\draw[draw=none,fill=steelblue31119180] (axis cs:253.5,0) rectangle (axis cs:254.5,100);
\draw[draw=none,fill=steelblue31119180] (axis cs:254.5,0) rectangle (axis cs:255.5,100);
\draw[draw=none,fill=steelblue31119180] (axis cs:255.5,0) rectangle (axis cs:256.5,100);
\draw[draw=none,fill=steelblue31119180] (axis cs:256.5,0) rectangle (axis cs:257.5,100);
\draw[draw=none,fill=steelblue31119180] (axis cs:257.5,0) rectangle (axis cs:258.5,100);
\draw[draw=none,fill=steelblue31119180] (axis cs:258.5,0) rectangle (axis cs:259.5,100);
\draw[draw=none,fill=steelblue31119180] (axis cs:259.5,0) rectangle (axis cs:260.5,100);
\draw[draw=none,fill=steelblue31119180] (axis cs:260.5,0) rectangle (axis cs:261.5,100);
\draw[draw=none,fill=steelblue31119180] (axis cs:261.5,0) rectangle (axis cs:262.5,100);
\draw[draw=none,fill=steelblue31119180] (axis cs:262.5,0) rectangle (axis cs:263.5,100);
\draw[draw=none,fill=steelblue31119180] (axis cs:263.5,0) rectangle (axis cs:264.5,100);
\draw[draw=none,fill=steelblue31119180] (axis cs:264.5,0) rectangle (axis cs:265.5,100);
\draw[draw=none,fill=steelblue31119180] (axis cs:265.5,0) rectangle (axis cs:266.5,100);
\draw[draw=none,fill=steelblue31119180] (axis cs:266.5,0) rectangle (axis cs:267.5,100);
\draw[draw=none,fill=steelblue31119180] (axis cs:267.5,0) rectangle (axis cs:268.5,100);
\draw[draw=none,fill=steelblue31119180] (axis cs:268.5,0) rectangle (axis cs:269.5,100);
\draw[draw=none,fill=steelblue31119180] (axis cs:269.5,0) rectangle (axis cs:270.5,100);
\draw[draw=none,fill=steelblue31119180] (axis cs:270.5,0) rectangle (axis cs:271.5,100);
\draw[draw=none,fill=steelblue31119180] (axis cs:271.5,0) rectangle (axis cs:272.5,100);
\draw[draw=none,fill=steelblue31119180] (axis cs:272.5,0) rectangle (axis cs:273.5,100);
\draw[draw=none,fill=steelblue31119180] (axis cs:273.5,0) rectangle (axis cs:274.5,100);
\draw[draw=none,fill=steelblue31119180] (axis cs:274.5,0) rectangle (axis cs:275.5,100);
\draw[draw=none,fill=steelblue31119180] (axis cs:275.5,0) rectangle (axis cs:276.5,100);
\draw[draw=none,fill=steelblue31119180] (axis cs:276.5,0) rectangle (axis cs:277.5,100);
\draw[draw=none,fill=steelblue31119180] (axis cs:277.5,0) rectangle (axis cs:278.5,100);
\draw[draw=none,fill=steelblue31119180] (axis cs:278.5,0) rectangle (axis cs:279.5,100);
\draw[draw=none,fill=steelblue31119180] (axis cs:279.5,0) rectangle (axis cs:280.5,100);
\draw[draw=none,fill=steelblue31119180] (axis cs:280.5,0) rectangle (axis cs:281.5,100);
\draw[draw=none,fill=steelblue31119180] (axis cs:281.5,0) rectangle (axis cs:282.5,100);
\draw[draw=none,fill=steelblue31119180] (axis cs:282.5,0) rectangle (axis cs:283.5,100);
\draw[draw=none,fill=steelblue31119180] (axis cs:283.5,0) rectangle (axis cs:284.5,100);
\draw[draw=none,fill=steelblue31119180] (axis cs:284.5,0) rectangle (axis cs:285.5,100);
\draw[draw=none,fill=steelblue31119180] (axis cs:285.5,0) rectangle (axis cs:286.5,100);
\draw[draw=none,fill=steelblue31119180] (axis cs:286.5,0) rectangle (axis cs:287.5,100);
\draw[draw=none,fill=steelblue31119180] (axis cs:287.5,0) rectangle (axis cs:288.5,100);
\draw[draw=none,fill=steelblue31119180] (axis cs:288.5,0) rectangle (axis cs:289.5,100);
\draw[draw=none,fill=steelblue31119180] (axis cs:289.5,0) rectangle (axis cs:290.5,100);
\draw[draw=none,fill=steelblue31119180] (axis cs:290.5,0) rectangle (axis cs:291.5,100);
\draw[draw=none,fill=steelblue31119180] (axis cs:291.5,0) rectangle (axis cs:292.5,100);
\draw[draw=none,fill=steelblue31119180] (axis cs:292.5,0) rectangle (axis cs:293.5,100);
\draw[draw=none,fill=steelblue31119180] (axis cs:293.5,0) rectangle (axis cs:294.5,100);
\draw[draw=none,fill=steelblue31119180] (axis cs:294.5,0) rectangle (axis cs:295.5,100);
\draw[draw=none,fill=steelblue31119180] (axis cs:295.5,0) rectangle (axis cs:296.5,100);
\draw[draw=none,fill=steelblue31119180] (axis cs:296.5,0) rectangle (axis cs:297.5,100);
\draw[draw=none,fill=steelblue31119180] (axis cs:297.5,0) rectangle (axis cs:298.5,100);
\draw[draw=none,fill=steelblue31119180] (axis cs:298.5,0) rectangle (axis cs:299.5,100);
\draw[draw=none,fill=steelblue31119180] (axis cs:299.5,0) rectangle (axis cs:300.5,100);
\draw[draw=none,fill=steelblue31119180] (axis cs:300.5,0) rectangle (axis cs:301.5,100);
\draw[draw=none,fill=steelblue31119180] (axis cs:301.5,0) rectangle (axis cs:302.5,100);
\draw[draw=none,fill=steelblue31119180] (axis cs:302.5,0) rectangle (axis cs:303.5,100);
\draw[draw=none,fill=steelblue31119180] (axis cs:303.5,0) rectangle (axis cs:304.5,100);
\draw[draw=none,fill=steelblue31119180] (axis cs:304.5,0) rectangle (axis cs:305.5,100);
\draw[draw=none,fill=steelblue31119180] (axis cs:305.5,0) rectangle (axis cs:306.5,100);
\draw[draw=none,fill=steelblue31119180] (axis cs:306.5,0) rectangle (axis cs:307.5,100);
\draw[draw=none,fill=steelblue31119180] (axis cs:307.5,0) rectangle (axis cs:308.5,100);
\draw[draw=none,fill=steelblue31119180] (axis cs:308.5,0) rectangle (axis cs:309.5,100);
\draw[draw=none,fill=steelblue31119180] (axis cs:309.5,0) rectangle (axis cs:310.5,100);
\draw[draw=none,fill=steelblue31119180] (axis cs:310.5,0) rectangle (axis cs:311.5,100);
\draw[draw=none,fill=steelblue31119180] (axis cs:311.5,0) rectangle (axis cs:312.5,100);
\draw[draw=none,fill=steelblue31119180] (axis cs:312.5,0) rectangle (axis cs:313.5,100);
\draw[draw=none,fill=steelblue31119180] (axis cs:313.5,0) rectangle (axis cs:314.5,100);
\draw[draw=none,fill=steelblue31119180] (axis cs:314.5,0) rectangle (axis cs:315.5,100);
\draw[draw=none,fill=steelblue31119180] (axis cs:315.5,0) rectangle (axis cs:316.5,100);
\draw[draw=none,fill=steelblue31119180] (axis cs:316.5,0) rectangle (axis cs:317.5,100);
\draw[draw=none,fill=steelblue31119180] (axis cs:317.5,0) rectangle (axis cs:318.5,100);
\draw[draw=none,fill=steelblue31119180] (axis cs:318.5,0) rectangle (axis cs:319.5,100);
\draw[draw=none,fill=steelblue31119180] (axis cs:319.5,0) rectangle (axis cs:320.5,100);
\draw[draw=none,fill=steelblue31119180] (axis cs:320.5,0) rectangle (axis cs:321.5,100);
\draw[draw=none,fill=steelblue31119180] (axis cs:321.5,0) rectangle (axis cs:322.5,100);
\draw[draw=none,fill=steelblue31119180] (axis cs:322.5,0) rectangle (axis cs:323.5,100);
\draw[draw=none,fill=steelblue31119180] (axis cs:323.5,0) rectangle (axis cs:324.5,100);
\draw[draw=none,fill=steelblue31119180] (axis cs:324.5,0) rectangle (axis cs:325.5,100);
\draw[draw=none,fill=steelblue31119180] (axis cs:325.5,0) rectangle (axis cs:326.5,100);
\draw[draw=none,fill=steelblue31119180] (axis cs:326.5,0) rectangle (axis cs:327.5,100);
\draw[draw=none,fill=steelblue31119180] (axis cs:327.5,0) rectangle (axis cs:328.5,100);
\draw[draw=none,fill=steelblue31119180] (axis cs:328.5,0) rectangle (axis cs:329.5,100);
\draw[draw=none,fill=steelblue31119180] (axis cs:329.5,0) rectangle (axis cs:330.5,100);
\draw[draw=none,fill=steelblue31119180] (axis cs:330.5,0) rectangle (axis cs:331.5,100);
\draw[draw=none,fill=steelblue31119180] (axis cs:331.5,0) rectangle (axis cs:332.5,100);
\draw[draw=none,fill=steelblue31119180] (axis cs:332.5,0) rectangle (axis cs:333.5,100);
\draw[draw=none,fill=steelblue31119180] (axis cs:333.5,0) rectangle (axis cs:334.5,100);
\draw[draw=none,fill=steelblue31119180] (axis cs:334.5,0) rectangle (axis cs:335.5,100);
\draw[draw=none,fill=steelblue31119180] (axis cs:335.5,0) rectangle (axis cs:336.5,100);
\draw[draw=none,fill=steelblue31119180] (axis cs:336.5,0) rectangle (axis cs:337.5,100);
\draw[draw=none,fill=steelblue31119180] (axis cs:337.5,0) rectangle (axis cs:338.5,100);
\draw[draw=none,fill=steelblue31119180] (axis cs:338.5,0) rectangle (axis cs:339.5,100);
\draw[draw=none,fill=steelblue31119180] (axis cs:339.5,0) rectangle (axis cs:340.5,100);
\draw[draw=none,fill=steelblue31119180] (axis cs:340.5,0) rectangle (axis cs:341.5,100);
\draw[draw=none,fill=steelblue31119180] (axis cs:341.5,0) rectangle (axis cs:342.5,100);
\draw[draw=none,fill=steelblue31119180] (axis cs:342.5,0) rectangle (axis cs:343.5,100);
\draw[draw=none,fill=steelblue31119180] (axis cs:343.5,0) rectangle (axis cs:344.5,100);
\draw[draw=none,fill=steelblue31119180] (axis cs:344.5,0) rectangle (axis cs:345.5,100);
\draw[draw=none,fill=steelblue31119180] (axis cs:345.5,0) rectangle (axis cs:346.5,100);
\draw[draw=none,fill=steelblue31119180] (axis cs:346.5,0) rectangle (axis cs:347.5,100);
\draw[draw=none,fill=steelblue31119180] (axis cs:347.5,0) rectangle (axis cs:348.5,100);
\draw[draw=none,fill=steelblue31119180] (axis cs:348.5,0) rectangle (axis cs:349.5,100);
\draw[draw=none,fill=steelblue31119180] (axis cs:349.5,0) rectangle (axis cs:350.5,100);
\draw[draw=none,fill=steelblue31119180] (axis cs:350.5,0) rectangle (axis cs:351.5,100);
\draw[draw=none,fill=steelblue31119180] (axis cs:351.5,0) rectangle (axis cs:352.5,100);
\draw[draw=none,fill=steelblue31119180] (axis cs:352.5,0) rectangle (axis cs:353.5,100);
\draw[draw=none,fill=steelblue31119180] (axis cs:353.5,0) rectangle (axis cs:354.5,100);
\draw[draw=none,fill=steelblue31119180] (axis cs:354.5,0) rectangle (axis cs:355.5,100);
\draw[draw=none,fill=steelblue31119180] (axis cs:355.5,0) rectangle (axis cs:356.5,100);
\draw[draw=none,fill=steelblue31119180] (axis cs:356.5,0) rectangle (axis cs:357.5,100);
\draw[draw=none,fill=steelblue31119180] (axis cs:357.5,0) rectangle (axis cs:358.5,100);
\draw[draw=none,fill=steelblue31119180] (axis cs:358.5,0) rectangle (axis cs:359.5,100);
\draw[draw=none,fill=steelblue31119180] (axis cs:359.5,0) rectangle (axis cs:360.5,100);
\draw[draw=none,fill=steelblue31119180] (axis cs:360.5,0) rectangle (axis cs:361.5,100);
\draw[draw=none,fill=steelblue31119180] (axis cs:361.5,0) rectangle (axis cs:362.5,100);
\draw[draw=none,fill=steelblue31119180] (axis cs:362.5,0) rectangle (axis cs:363.5,100);
\draw[draw=none,fill=steelblue31119180] (axis cs:363.5,0) rectangle (axis cs:364.5,100);
\draw[draw=none,fill=steelblue31119180] (axis cs:364.5,0) rectangle (axis cs:365.5,100);
\draw[draw=none,fill=steelblue31119180] (axis cs:365.5,0) rectangle (axis cs:366.5,100);
\draw[draw=none,fill=steelblue31119180] (axis cs:366.5,0) rectangle (axis cs:367.5,100);
\draw[draw=none,fill=steelblue31119180] (axis cs:367.5,0) rectangle (axis cs:368.5,100);
\draw[draw=none,fill=steelblue31119180] (axis cs:368.5,0) rectangle (axis cs:369.5,100);
\draw[draw=none,fill=steelblue31119180] (axis cs:369.5,0) rectangle (axis cs:370.5,100);
\draw[draw=none,fill=steelblue31119180] (axis cs:370.5,0) rectangle (axis cs:371.5,100);
\draw[draw=none,fill=steelblue31119180] (axis cs:371.5,0) rectangle (axis cs:372.5,100);
\draw[draw=none,fill=steelblue31119180] (axis cs:372.5,0) rectangle (axis cs:373.5,100);
\draw[draw=none,fill=steelblue31119180] (axis cs:373.5,0) rectangle (axis cs:374.5,100);
\draw[draw=none,fill=steelblue31119180] (axis cs:374.5,0) rectangle (axis cs:375.5,100);
\draw[draw=none,fill=steelblue31119180] (axis cs:375.5,0) rectangle (axis cs:376.5,100);
\draw[draw=none,fill=steelblue31119180] (axis cs:376.5,0) rectangle (axis cs:377.5,100);
\draw[draw=none,fill=steelblue31119180] (axis cs:377.5,0) rectangle (axis cs:378.5,100);
\draw[draw=none,fill=steelblue31119180] (axis cs:378.5,0) rectangle (axis cs:379.5,100);
\draw[draw=none,fill=steelblue31119180] (axis cs:379.5,0) rectangle (axis cs:380.5,100);
\draw[draw=none,fill=steelblue31119180] (axis cs:380.5,0) rectangle (axis cs:381.5,100);
\draw[draw=none,fill=steelblue31119180] (axis cs:381.5,0) rectangle (axis cs:382.5,100);
\draw[draw=none,fill=steelblue31119180] (axis cs:382.5,0) rectangle (axis cs:383.5,100);
\draw[draw=none,fill=steelblue31119180] (axis cs:383.5,0) rectangle (axis cs:384.5,100);
\draw[draw=none,fill=steelblue31119180] (axis cs:384.5,0) rectangle (axis cs:385.5,100);
\draw[draw=none,fill=steelblue31119180] (axis cs:385.5,0) rectangle (axis cs:386.5,100);
\draw[draw=none,fill=steelblue31119180] (axis cs:386.5,0) rectangle (axis cs:387.5,100);
\draw[draw=none,fill=steelblue31119180] (axis cs:387.5,0) rectangle (axis cs:388.5,100);
\draw[draw=none,fill=steelblue31119180] (axis cs:388.5,0) rectangle (axis cs:389.5,100);
\draw[draw=none,fill=steelblue31119180] (axis cs:389.5,0) rectangle (axis cs:390.5,100);
\draw[draw=none,fill=steelblue31119180] (axis cs:390.5,0) rectangle (axis cs:391.5,100);
\draw[draw=none,fill=steelblue31119180] (axis cs:391.5,0) rectangle (axis cs:392.5,100);
\draw[draw=none,fill=steelblue31119180] (axis cs:392.5,0) rectangle (axis cs:393.5,100);
\draw[draw=none,fill=steelblue31119180] (axis cs:393.5,0) rectangle (axis cs:394.5,100);
\draw[draw=none,fill=steelblue31119180] (axis cs:394.5,0) rectangle (axis cs:395.5,100);
\draw[draw=none,fill=steelblue31119180] (axis cs:395.5,0) rectangle (axis cs:396.5,100);
\draw[draw=none,fill=steelblue31119180] (axis cs:396.5,0) rectangle (axis cs:397.5,100);
\draw[draw=none,fill=steelblue31119180] (axis cs:397.5,0) rectangle (axis cs:398.5,100);
\draw[draw=none,fill=steelblue31119180] (axis cs:398.5,0) rectangle (axis cs:399.5,100);
\draw[draw=none,fill=steelblue31119180] (axis cs:399.5,0) rectangle (axis cs:400.5,100);
\draw[draw=none,fill=steelblue31119180] (axis cs:400.5,0) rectangle (axis cs:401.5,100);
\draw[draw=none,fill=steelblue31119180] (axis cs:401.5,0) rectangle (axis cs:402.5,100);
\draw[draw=none,fill=steelblue31119180] (axis cs:402.5,0) rectangle (axis cs:403.5,100);
\draw[draw=none,fill=steelblue31119180] (axis cs:403.5,0) rectangle (axis cs:404.5,100);
\draw[draw=none,fill=steelblue31119180] (axis cs:404.5,0) rectangle (axis cs:405.5,100);
\draw[draw=none,fill=steelblue31119180] (axis cs:405.5,0) rectangle (axis cs:406.5,100);
\draw[draw=none,fill=steelblue31119180] (axis cs:406.5,0) rectangle (axis cs:407.5,100);
\draw[draw=none,fill=steelblue31119180] (axis cs:407.5,0) rectangle (axis cs:408.5,100);
\draw[draw=none,fill=steelblue31119180] (axis cs:408.5,0) rectangle (axis cs:409.5,100);
\draw[draw=none,fill=steelblue31119180] (axis cs:409.5,0) rectangle (axis cs:410.5,100);
\draw[draw=none,fill=steelblue31119180] (axis cs:410.5,0) rectangle (axis cs:411.5,100);
\draw[draw=none,fill=steelblue31119180] (axis cs:411.5,0) rectangle (axis cs:412.5,100);
\draw[draw=none,fill=steelblue31119180] (axis cs:412.5,0) rectangle (axis cs:413.5,100);
\draw[draw=none,fill=steelblue31119180] (axis cs:413.5,0) rectangle (axis cs:414.5,100);
\draw[draw=none,fill=steelblue31119180] (axis cs:414.5,0) rectangle (axis cs:415.5,100);
\draw[draw=none,fill=steelblue31119180] (axis cs:415.5,0) rectangle (axis cs:416.5,100);
\draw[draw=none,fill=steelblue31119180] (axis cs:416.5,0) rectangle (axis cs:417.5,100);
\draw[draw=none,fill=steelblue31119180] (axis cs:417.5,0) rectangle (axis cs:418.5,100);
\draw[draw=none,fill=steelblue31119180] (axis cs:418.5,0) rectangle (axis cs:419.5,100);
\draw[draw=none,fill=steelblue31119180] (axis cs:419.5,0) rectangle (axis cs:420.5,100);
\draw[draw=none,fill=steelblue31119180] (axis cs:420.5,0) rectangle (axis cs:421.5,100);
\draw[draw=none,fill=steelblue31119180] (axis cs:421.5,0) rectangle (axis cs:422.5,100);
\draw[draw=none,fill=steelblue31119180] (axis cs:422.5,0) rectangle (axis cs:423.5,100);
\draw[draw=none,fill=steelblue31119180] (axis cs:423.5,0) rectangle (axis cs:424.5,100);
\draw[draw=none,fill=steelblue31119180] (axis cs:424.5,0) rectangle (axis cs:425.5,100);
\draw[draw=none,fill=steelblue31119180] (axis cs:425.5,0) rectangle (axis cs:426.5,100);
\draw[draw=none,fill=steelblue31119180] (axis cs:426.5,0) rectangle (axis cs:427.5,100);
\draw[draw=none,fill=steelblue31119180] (axis cs:427.5,0) rectangle (axis cs:428.5,100);
\draw[draw=none,fill=steelblue31119180] (axis cs:428.5,0) rectangle (axis cs:429.5,100);
\draw[draw=none,fill=steelblue31119180] (axis cs:429.5,0) rectangle (axis cs:430.5,100);
\draw[draw=none,fill=steelblue31119180] (axis cs:430.5,0) rectangle (axis cs:431.5,100);
\draw[draw=none,fill=steelblue31119180] (axis cs:431.5,0) rectangle (axis cs:432.5,100);
\draw[draw=none,fill=steelblue31119180] (axis cs:432.5,0) rectangle (axis cs:433.5,100);
\draw[draw=none,fill=steelblue31119180] (axis cs:433.5,0) rectangle (axis cs:434.5,100);
\draw[draw=none,fill=steelblue31119180] (axis cs:434.5,0) rectangle (axis cs:435.5,100);
\draw[draw=none,fill=steelblue31119180] (axis cs:435.5,0) rectangle (axis cs:436.5,100);
\draw[draw=none,fill=steelblue31119180] (axis cs:436.5,0) rectangle (axis cs:437.5,100);
\draw[draw=none,fill=steelblue31119180] (axis cs:437.5,0) rectangle (axis cs:438.5,100);
\draw[draw=none,fill=steelblue31119180] (axis cs:438.5,0) rectangle (axis cs:439.5,100);
\draw[draw=none,fill=steelblue31119180] (axis cs:439.5,0) rectangle (axis cs:440.5,100);
\draw[draw=none,fill=steelblue31119180] (axis cs:440.5,0) rectangle (axis cs:441.5,100);
\draw[draw=none,fill=steelblue31119180] (axis cs:441.5,0) rectangle (axis cs:442.5,100);
\draw[draw=none,fill=steelblue31119180] (axis cs:442.5,0) rectangle (axis cs:443.5,100);
\draw[draw=none,fill=steelblue31119180] (axis cs:443.5,0) rectangle (axis cs:444.5,100);
\draw[draw=none,fill=steelblue31119180] (axis cs:444.5,0) rectangle (axis cs:445.5,100);
\draw[draw=none,fill=steelblue31119180] (axis cs:445.5,0) rectangle (axis cs:446.5,100);
\draw[draw=none,fill=steelblue31119180] (axis cs:446.5,0) rectangle (axis cs:447.5,100);
\draw[draw=none,fill=steelblue31119180] (axis cs:447.5,0) rectangle (axis cs:448.5,100);
\draw[draw=none,fill=steelblue31119180] (axis cs:448.5,0) rectangle (axis cs:449.5,100);
\draw[draw=none,fill=steelblue31119180] (axis cs:449.5,0) rectangle (axis cs:450.5,100);
\draw[draw=none,fill=steelblue31119180] (axis cs:450.5,0) rectangle (axis cs:451.5,100);
\draw[draw=none,fill=steelblue31119180] (axis cs:451.5,0) rectangle (axis cs:452.5,100);
\draw[draw=none,fill=steelblue31119180] (axis cs:452.5,0) rectangle (axis cs:453.5,100);
\draw[draw=none,fill=steelblue31119180] (axis cs:453.5,0) rectangle (axis cs:454.5,100);
\draw[draw=none,fill=steelblue31119180] (axis cs:454.5,0) rectangle (axis cs:455.5,100);
\draw[draw=none,fill=steelblue31119180] (axis cs:455.5,0) rectangle (axis cs:456.5,100);
\draw[draw=none,fill=steelblue31119180] (axis cs:456.5,0) rectangle (axis cs:457.5,100);
\draw[draw=none,fill=steelblue31119180] (axis cs:457.5,0) rectangle (axis cs:458.5,100);
\draw[draw=none,fill=steelblue31119180] (axis cs:458.5,0) rectangle (axis cs:459.5,100);
\draw[draw=none,fill=steelblue31119180] (axis cs:459.5,0) rectangle (axis cs:460.5,100);
\draw[draw=none,fill=steelblue31119180] (axis cs:460.5,0) rectangle (axis cs:461.5,100);
\draw[draw=none,fill=steelblue31119180] (axis cs:461.5,0) rectangle (axis cs:462.5,100);
\draw[draw=none,fill=steelblue31119180] (axis cs:462.5,0) rectangle (axis cs:463.5,100);
\draw[draw=none,fill=steelblue31119180] (axis cs:463.5,0) rectangle (axis cs:464.5,100);
\draw[draw=none,fill=steelblue31119180] (axis cs:464.5,0) rectangle (axis cs:465.5,100);
\draw[draw=none,fill=steelblue31119180] (axis cs:465.5,0) rectangle (axis cs:466.5,100);
\draw[draw=none,fill=steelblue31119180] (axis cs:466.5,0) rectangle (axis cs:467.5,100);
\draw[draw=none,fill=steelblue31119180] (axis cs:467.5,0) rectangle (axis cs:468.5,100);
\draw[draw=none,fill=steelblue31119180] (axis cs:468.5,0) rectangle (axis cs:469.5,100);
\draw[draw=none,fill=steelblue31119180] (axis cs:469.5,0) rectangle (axis cs:470.5,100);
\draw[draw=none,fill=steelblue31119180] (axis cs:470.5,0) rectangle (axis cs:471.5,100);
\draw[draw=none,fill=steelblue31119180] (axis cs:471.5,0) rectangle (axis cs:472.5,100);
\draw[draw=none,fill=steelblue31119180] (axis cs:472.5,0) rectangle (axis cs:473.5,100);
\draw[draw=none,fill=steelblue31119180] (axis cs:473.5,0) rectangle (axis cs:474.5,100);
\draw[draw=none,fill=steelblue31119180] (axis cs:474.5,0) rectangle (axis cs:475.5,100);
\draw[draw=none,fill=steelblue31119180] (axis cs:475.5,0) rectangle (axis cs:476.5,100);
\draw[draw=none,fill=steelblue31119180] (axis cs:476.5,0) rectangle (axis cs:477.5,100);
\draw[draw=none,fill=steelblue31119180] (axis cs:477.5,0) rectangle (axis cs:478.5,100);
\draw[draw=none,fill=steelblue31119180] (axis cs:478.5,0) rectangle (axis cs:479.5,100);
\draw[draw=none,fill=steelblue31119180] (axis cs:479.5,0) rectangle (axis cs:480.5,100);
\draw[draw=none,fill=steelblue31119180] (axis cs:480.5,0) rectangle (axis cs:481.5,100);
\draw[draw=none,fill=steelblue31119180] (axis cs:481.5,0) rectangle (axis cs:482.5,100);
\draw[draw=none,fill=steelblue31119180] (axis cs:482.5,0) rectangle (axis cs:483.5,100);
\draw[draw=none,fill=steelblue31119180] (axis cs:483.5,0) rectangle (axis cs:484.5,100);
\draw[draw=none,fill=steelblue31119180] (axis cs:484.5,0) rectangle (axis cs:485.5,100);
\draw[draw=none,fill=steelblue31119180] (axis cs:485.5,0) rectangle (axis cs:486.5,100);
\draw[draw=none,fill=steelblue31119180] (axis cs:486.5,0) rectangle (axis cs:487.5,100);
\draw[draw=none,fill=steelblue31119180] (axis cs:487.5,0) rectangle (axis cs:488.5,100);
\draw[draw=none,fill=steelblue31119180] (axis cs:488.5,0) rectangle (axis cs:489.5,100);
\draw[draw=none,fill=steelblue31119180] (axis cs:489.5,0) rectangle (axis cs:490.5,100);
\draw[draw=none,fill=steelblue31119180] (axis cs:490.5,0) rectangle (axis cs:491.5,100);
\draw[draw=none,fill=steelblue31119180] (axis cs:491.5,0) rectangle (axis cs:492.5,100);
\draw[draw=none,fill=steelblue31119180] (axis cs:492.5,0) rectangle (axis cs:493.5,100);
\draw[draw=none,fill=steelblue31119180] (axis cs:493.5,0) rectangle (axis cs:494.5,100);
\draw[draw=none,fill=steelblue31119180] (axis cs:494.5,0) rectangle (axis cs:495.5,100);
\draw[draw=none,fill=steelblue31119180] (axis cs:495.5,0) rectangle (axis cs:496.5,100);
\draw[draw=none,fill=steelblue31119180] (axis cs:496.5,0) rectangle (axis cs:497.5,100);
\draw[draw=none,fill=steelblue31119180] (axis cs:497.5,0) rectangle (axis cs:498.5,100);
\draw[draw=none,fill=steelblue31119180] (axis cs:498.5,0) rectangle (axis cs:499.5,100);
\draw[draw=none,fill=steelblue31119180] (axis cs:499.5,0) rectangle (axis cs:500.5,100);
\draw[draw=none,fill=steelblue31119180] (axis cs:500.5,0) rectangle (axis cs:501.5,100);
\draw[draw=none,fill=steelblue31119180] (axis cs:501.5,0) rectangle (axis cs:502.5,100);
\draw[draw=none,fill=steelblue31119180] (axis cs:502.5,0) rectangle (axis cs:503.5,100);
\draw[draw=none,fill=steelblue31119180] (axis cs:503.5,0) rectangle (axis cs:504.5,100);
\draw[draw=none,fill=steelblue31119180] (axis cs:504.5,0) rectangle (axis cs:505.5,100);
\draw[draw=none,fill=steelblue31119180] (axis cs:505.5,0) rectangle (axis cs:506.5,100);
\draw[draw=none,fill=steelblue31119180] (axis cs:506.5,0) rectangle (axis cs:507.5,100);
\draw[draw=none,fill=steelblue31119180] (axis cs:507.5,0) rectangle (axis cs:508.5,100);
\draw[draw=none,fill=steelblue31119180] (axis cs:508.5,0) rectangle (axis cs:509.5,100);
\draw[draw=none,fill=steelblue31119180] (axis cs:509.5,0) rectangle (axis cs:510.5,100);
\draw[draw=none,fill=steelblue31119180] (axis cs:510.5,0) rectangle (axis cs:511.5,100);
\draw[draw=none,fill=steelblue31119180] (axis cs:511.5,0) rectangle (axis cs:512.5,100);
\draw[draw=none,fill=steelblue31119180] (axis cs:512.5,0) rectangle (axis cs:513.5,100);
\draw[draw=none,fill=steelblue31119180] (axis cs:513.5,0) rectangle (axis cs:514.5,100);
\draw[draw=none,fill=steelblue31119180] (axis cs:514.5,0) rectangle (axis cs:515.5,100);
\draw[draw=none,fill=steelblue31119180] (axis cs:515.5,0) rectangle (axis cs:516.5,100);
\draw[draw=none,fill=steelblue31119180] (axis cs:516.5,0) rectangle (axis cs:517.5,100);
\draw[draw=none,fill=steelblue31119180] (axis cs:517.5,0) rectangle (axis cs:518.5,100);
\draw[draw=none,fill=steelblue31119180] (axis cs:518.5,0) rectangle (axis cs:519.5,100);
\draw[draw=none,fill=steelblue31119180] (axis cs:519.5,0) rectangle (axis cs:520.5,100);
\draw[draw=none,fill=steelblue31119180] (axis cs:520.5,0) rectangle (axis cs:521.5,100);
\draw[draw=none,fill=steelblue31119180] (axis cs:521.5,0) rectangle (axis cs:522.5,100);
\draw[draw=none,fill=steelblue31119180] (axis cs:522.5,0) rectangle (axis cs:523.5,100);
\draw[draw=none,fill=steelblue31119180] (axis cs:523.5,0) rectangle (axis cs:524.5,100);
\draw[draw=none,fill=steelblue31119180] (axis cs:524.5,0) rectangle (axis cs:525.5,100);
\draw[draw=none,fill=steelblue31119180] (axis cs:525.5,0) rectangle (axis cs:526.5,100);
\draw[draw=none,fill=steelblue31119180] (axis cs:526.5,0) rectangle (axis cs:527.5,100);
\draw[draw=none,fill=steelblue31119180] (axis cs:527.5,0) rectangle (axis cs:528.5,100);
\draw[draw=none,fill=steelblue31119180] (axis cs:528.5,0) rectangle (axis cs:529.5,100);
\draw[draw=none,fill=steelblue31119180] (axis cs:529.5,0) rectangle (axis cs:530.5,100);
\draw[draw=none,fill=steelblue31119180] (axis cs:530.5,0) rectangle (axis cs:531.5,100);
\draw[draw=none,fill=steelblue31119180] (axis cs:531.5,0) rectangle (axis cs:532.5,100);
\draw[draw=none,fill=steelblue31119180] (axis cs:532.5,0) rectangle (axis cs:533.5,100);
\draw[draw=none,fill=steelblue31119180] (axis cs:533.5,0) rectangle (axis cs:534.5,100);
\draw[draw=none,fill=steelblue31119180] (axis cs:534.5,0) rectangle (axis cs:535.5,100);
\draw[draw=none,fill=steelblue31119180] (axis cs:535.5,0) rectangle (axis cs:536.5,100);
\draw[draw=none,fill=steelblue31119180] (axis cs:536.5,0) rectangle (axis cs:537.5,100);
\draw[draw=none,fill=steelblue31119180] (axis cs:537.5,0) rectangle (axis cs:538.5,100);
\draw[draw=none,fill=steelblue31119180] (axis cs:538.5,0) rectangle (axis cs:539.5,100);
\draw[draw=none,fill=steelblue31119180] (axis cs:539.5,0) rectangle (axis cs:540.5,100);
\draw[draw=none,fill=steelblue31119180] (axis cs:540.5,0) rectangle (axis cs:541.5,100);
\draw[draw=none,fill=steelblue31119180] (axis cs:541.5,0) rectangle (axis cs:542.5,100);
\draw[draw=none,fill=steelblue31119180] (axis cs:542.5,0) rectangle (axis cs:543.5,100);
\draw[draw=none,fill=steelblue31119180] (axis cs:543.5,0) rectangle (axis cs:544.5,100);
\draw[draw=none,fill=steelblue31119180] (axis cs:544.5,0) rectangle (axis cs:545.5,100);
\draw[draw=none,fill=steelblue31119180] (axis cs:545.5,0) rectangle (axis cs:546.5,100);
\draw[draw=none,fill=steelblue31119180] (axis cs:546.5,0) rectangle (axis cs:547.5,100);
\draw[draw=none,fill=steelblue31119180] (axis cs:547.5,0) rectangle (axis cs:548.5,100);
\draw[draw=none,fill=steelblue31119180] (axis cs:548.5,0) rectangle (axis cs:549.5,100);
\draw[draw=none,fill=steelblue31119180] (axis cs:549.5,0) rectangle (axis cs:550.5,100);
\draw[draw=none,fill=steelblue31119180] (axis cs:550.5,0) rectangle (axis cs:551.5,100);
\draw[draw=none,fill=steelblue31119180] (axis cs:551.5,0) rectangle (axis cs:552.5,100);
\draw[draw=none,fill=steelblue31119180] (axis cs:552.5,0) rectangle (axis cs:553.5,100);
\draw[draw=none,fill=steelblue31119180] (axis cs:553.5,0) rectangle (axis cs:554.5,100);
\draw[draw=none,fill=steelblue31119180] (axis cs:554.5,0) rectangle (axis cs:555.5,100);
\draw[draw=none,fill=steelblue31119180] (axis cs:555.5,0) rectangle (axis cs:556.5,100);
\draw[draw=none,fill=steelblue31119180] (axis cs:556.5,0) rectangle (axis cs:557.5,100);
\draw[draw=none,fill=steelblue31119180] (axis cs:557.5,0) rectangle (axis cs:558.5,100);
\draw[draw=none,fill=steelblue31119180] (axis cs:558.5,0) rectangle (axis cs:559.5,100);
\draw[draw=none,fill=steelblue31119180] (axis cs:559.5,0) rectangle (axis cs:560.5,100);
\draw[draw=none,fill=steelblue31119180] (axis cs:560.5,0) rectangle (axis cs:561.5,100);
\draw[draw=none,fill=steelblue31119180] (axis cs:561.5,0) rectangle (axis cs:562.5,100);
\draw[draw=none,fill=steelblue31119180] (axis cs:562.5,0) rectangle (axis cs:563.5,100);
\draw[draw=none,fill=steelblue31119180] (axis cs:563.5,0) rectangle (axis cs:564.5,100);
\draw[draw=none,fill=steelblue31119180] (axis cs:564.5,0) rectangle (axis cs:565.5,100);
\draw[draw=none,fill=steelblue31119180] (axis cs:565.5,0) rectangle (axis cs:566.5,100);
\draw[draw=none,fill=steelblue31119180] (axis cs:566.5,0) rectangle (axis cs:567.5,100);
\draw[draw=none,fill=steelblue31119180] (axis cs:567.5,0) rectangle (axis cs:568.5,100);
\draw[draw=none,fill=steelblue31119180] (axis cs:568.5,0) rectangle (axis cs:569.5,100);
\draw[draw=none,fill=steelblue31119180] (axis cs:569.5,0) rectangle (axis cs:570.5,100);
\draw[draw=none,fill=steelblue31119180] (axis cs:570.5,0) rectangle (axis cs:571.5,100);
\draw[draw=none,fill=steelblue31119180] (axis cs:571.5,0) rectangle (axis cs:572.5,100);
\draw[draw=none,fill=steelblue31119180] (axis cs:572.5,0) rectangle (axis cs:573.5,100);
\draw[draw=none,fill=steelblue31119180] (axis cs:573.5,0) rectangle (axis cs:574.5,100);
\draw[draw=none,fill=steelblue31119180] (axis cs:574.5,0) rectangle (axis cs:575.5,100);
\draw[draw=none,fill=steelblue31119180] (axis cs:575.5,0) rectangle (axis cs:576.5,100);
\draw[draw=none,fill=steelblue31119180] (axis cs:576.5,0) rectangle (axis cs:577.5,100);
\draw[draw=none,fill=steelblue31119180] (axis cs:577.5,0) rectangle (axis cs:578.5,100);
\draw[draw=none,fill=steelblue31119180] (axis cs:578.5,0) rectangle (axis cs:579.5,100);
\draw[draw=none,fill=steelblue31119180] (axis cs:579.5,0) rectangle (axis cs:580.5,100);
\draw[draw=none,fill=steelblue31119180] (axis cs:580.5,0) rectangle (axis cs:581.5,100);
\draw[draw=none,fill=steelblue31119180] (axis cs:581.5,0) rectangle (axis cs:582.5,100);
\draw[draw=none,fill=steelblue31119180] (axis cs:582.5,0) rectangle (axis cs:583.5,100);
\draw[draw=none,fill=steelblue31119180] (axis cs:583.5,0) rectangle (axis cs:584.5,100);
\draw[draw=none,fill=steelblue31119180] (axis cs:584.5,0) rectangle (axis cs:585.5,100);
\draw[draw=none,fill=steelblue31119180] (axis cs:585.5,0) rectangle (axis cs:586.5,100);
\draw[draw=none,fill=steelblue31119180] (axis cs:586.5,0) rectangle (axis cs:587.5,100);
\draw[draw=none,fill=steelblue31119180] (axis cs:587.5,0) rectangle (axis cs:588.5,100);
\draw[draw=none,fill=steelblue31119180] (axis cs:588.5,0) rectangle (axis cs:589.5,100);
\draw[draw=none,fill=steelblue31119180] (axis cs:589.5,0) rectangle (axis cs:590.5,100);
\draw[draw=none,fill=steelblue31119180] (axis cs:590.5,0) rectangle (axis cs:591.5,100);
\draw[draw=none,fill=steelblue31119180] (axis cs:591.5,0) rectangle (axis cs:592.5,100);
\draw[draw=none,fill=steelblue31119180] (axis cs:592.5,0) rectangle (axis cs:593.5,100);
\draw[draw=none,fill=steelblue31119180] (axis cs:593.5,0) rectangle (axis cs:594.5,100);
\draw[draw=none,fill=steelblue31119180] (axis cs:594.5,0) rectangle (axis cs:595.5,100);
\draw[draw=none,fill=steelblue31119180] (axis cs:595.5,0) rectangle (axis cs:596.5,100);
\draw[draw=none,fill=steelblue31119180] (axis cs:596.5,0) rectangle (axis cs:597.5,100);
\draw[draw=none,fill=steelblue31119180] (axis cs:597.5,0) rectangle (axis cs:598.5,100);
\draw[draw=none,fill=steelblue31119180] (axis cs:598.5,0) rectangle (axis cs:599.5,100);
\draw[draw=none,fill=steelblue31119180] (axis cs:599.5,0) rectangle (axis cs:600.5,100);
\draw[draw=none,fill=steelblue31119180] (axis cs:600.5,0) rectangle (axis cs:601.5,100);
\draw[draw=none,fill=steelblue31119180] (axis cs:601.5,0) rectangle (axis cs:602.5,100);
\draw[draw=none,fill=steelblue31119180] (axis cs:602.5,0) rectangle (axis cs:603.5,100);
\draw[draw=none,fill=steelblue31119180] (axis cs:603.5,0) rectangle (axis cs:604.5,100);
\draw[draw=none,fill=steelblue31119180] (axis cs:604.5,0) rectangle (axis cs:605.5,100);
\draw[draw=none,fill=steelblue31119180] (axis cs:605.5,0) rectangle (axis cs:606.5,100);
\draw[draw=none,fill=steelblue31119180] (axis cs:606.5,0) rectangle (axis cs:607.5,100);
\draw[draw=none,fill=steelblue31119180] (axis cs:607.5,0) rectangle (axis cs:608.5,100);
\draw[draw=none,fill=steelblue31119180] (axis cs:608.5,0) rectangle (axis cs:609.5,100);
\draw[draw=none,fill=steelblue31119180] (axis cs:609.5,0) rectangle (axis cs:610.5,100);
\draw[draw=none,fill=steelblue31119180] (axis cs:610.5,0) rectangle (axis cs:611.5,100);
\draw[draw=none,fill=steelblue31119180] (axis cs:611.5,0) rectangle (axis cs:612.5,100);
\draw[draw=none,fill=steelblue31119180] (axis cs:612.5,0) rectangle (axis cs:613.5,100);
\draw[draw=none,fill=steelblue31119180] (axis cs:613.5,0) rectangle (axis cs:614.5,100);
\draw[draw=none,fill=steelblue31119180] (axis cs:614.5,0) rectangle (axis cs:615.5,100);
\draw[draw=none,fill=steelblue31119180] (axis cs:615.5,0) rectangle (axis cs:616.5,100);
\draw[draw=none,fill=steelblue31119180] (axis cs:616.5,0) rectangle (axis cs:617.5,100);
\draw[draw=none,fill=steelblue31119180] (axis cs:617.5,0) rectangle (axis cs:618.5,100);
\draw[draw=none,fill=steelblue31119180] (axis cs:618.5,0) rectangle (axis cs:619.5,100);
\draw[draw=none,fill=steelblue31119180] (axis cs:619.5,0) rectangle (axis cs:620.5,100);
\draw[draw=none,fill=steelblue31119180] (axis cs:620.5,0) rectangle (axis cs:621.5,100);
\draw[draw=none,fill=steelblue31119180] (axis cs:621.5,0) rectangle (axis cs:622.5,100);
\draw[draw=none,fill=steelblue31119180] (axis cs:622.5,0) rectangle (axis cs:623.5,100);
\draw[draw=none,fill=steelblue31119180] (axis cs:623.5,0) rectangle (axis cs:624.5,100);
\draw[draw=none,fill=steelblue31119180] (axis cs:624.5,0) rectangle (axis cs:625.5,100);
\draw[draw=none,fill=steelblue31119180] (axis cs:625.5,0) rectangle (axis cs:626.5,100);
\draw[draw=none,fill=steelblue31119180] (axis cs:626.5,0) rectangle (axis cs:627.5,100);
\draw[draw=none,fill=steelblue31119180] (axis cs:627.5,0) rectangle (axis cs:628.5,100);
\draw[draw=none,fill=steelblue31119180] (axis cs:628.5,0) rectangle (axis cs:629.5,100);
\draw[draw=none,fill=steelblue31119180] (axis cs:629.5,0) rectangle (axis cs:630.5,100);
\draw[draw=none,fill=steelblue31119180] (axis cs:630.5,0) rectangle (axis cs:631.5,100);
\draw[draw=none,fill=steelblue31119180] (axis cs:631.5,0) rectangle (axis cs:632.5,100);
\draw[draw=none,fill=steelblue31119180] (axis cs:632.5,0) rectangle (axis cs:633.5,100);
\draw[draw=none,fill=steelblue31119180] (axis cs:633.5,0) rectangle (axis cs:634.5,100);
\draw[draw=none,fill=steelblue31119180] (axis cs:634.5,0) rectangle (axis cs:635.5,100);
\draw[draw=none,fill=steelblue31119180] (axis cs:635.5,0) rectangle (axis cs:636.5,100);
\draw[draw=none,fill=steelblue31119180] (axis cs:636.5,0) rectangle (axis cs:637.5,100);
\draw[draw=none,fill=steelblue31119180] (axis cs:637.5,0) rectangle (axis cs:638.5,100);
\draw[draw=none,fill=steelblue31119180] (axis cs:638.5,0) rectangle (axis cs:639.5,100);
\draw[draw=none,fill=steelblue31119180] (axis cs:639.5,0) rectangle (axis cs:640.5,100);
\draw[draw=none,fill=steelblue31119180] (axis cs:640.5,0) rectangle (axis cs:641.5,100);
\draw[draw=none,fill=steelblue31119180] (axis cs:641.5,0) rectangle (axis cs:642.5,100);
\draw[draw=none,fill=steelblue31119180] (axis cs:642.5,0) rectangle (axis cs:643.5,100);
\draw[draw=none,fill=steelblue31119180] (axis cs:643.5,0) rectangle (axis cs:644.5,100);
\draw[draw=none,fill=steelblue31119180] (axis cs:644.5,0) rectangle (axis cs:645.5,100);
\draw[draw=none,fill=steelblue31119180] (axis cs:645.5,0) rectangle (axis cs:646.5,100);
\draw[draw=none,fill=steelblue31119180] (axis cs:646.5,0) rectangle (axis cs:647.5,100);
\draw[draw=none,fill=steelblue31119180] (axis cs:647.5,0) rectangle (axis cs:648.5,100);
\draw[draw=none,fill=steelblue31119180] (axis cs:648.5,0) rectangle (axis cs:649.5,100);
\draw[draw=none,fill=steelblue31119180] (axis cs:649.5,0) rectangle (axis cs:650.5,100);
\draw[draw=none,fill=steelblue31119180] (axis cs:650.5,0) rectangle (axis cs:651.5,100);
\draw[draw=none,fill=steelblue31119180] (axis cs:651.5,0) rectangle (axis cs:652.5,100);
\draw[draw=none,fill=steelblue31119180] (axis cs:652.5,0) rectangle (axis cs:653.5,100);
\draw[draw=none,fill=steelblue31119180] (axis cs:653.5,0) rectangle (axis cs:654.5,100);
\draw[draw=none,fill=steelblue31119180] (axis cs:654.5,0) rectangle (axis cs:655.5,100);
\draw[draw=none,fill=steelblue31119180] (axis cs:655.5,0) rectangle (axis cs:656.5,100);
\draw[draw=none,fill=steelblue31119180] (axis cs:656.5,0) rectangle (axis cs:657.5,100);
\draw[draw=none,fill=steelblue31119180] (axis cs:657.5,0) rectangle (axis cs:658.5,100);
\draw[draw=none,fill=steelblue31119180] (axis cs:658.5,0) rectangle (axis cs:659.5,100);
\draw[draw=none,fill=steelblue31119180] (axis cs:659.5,0) rectangle (axis cs:660.5,100);
\draw[draw=none,fill=steelblue31119180] (axis cs:660.5,0) rectangle (axis cs:661.5,100);
\draw[draw=none,fill=steelblue31119180] (axis cs:661.5,0) rectangle (axis cs:662.5,100);
\draw[draw=none,fill=steelblue31119180] (axis cs:662.5,0) rectangle (axis cs:663.5,100);
\draw[draw=none,fill=steelblue31119180] (axis cs:663.5,0) rectangle (axis cs:664.5,100);
\draw[draw=none,fill=steelblue31119180] (axis cs:664.5,0) rectangle (axis cs:665.5,100);
\draw[draw=none,fill=steelblue31119180] (axis cs:665.5,0) rectangle (axis cs:666.5,100);
\draw[draw=none,fill=steelblue31119180] (axis cs:666.5,0) rectangle (axis cs:667.5,100);
\draw[draw=none,fill=steelblue31119180] (axis cs:667.5,0) rectangle (axis cs:668.5,100);
\draw[draw=none,fill=steelblue31119180] (axis cs:668.5,0) rectangle (axis cs:669.5,100);
\draw[draw=none,fill=steelblue31119180] (axis cs:669.5,0) rectangle (axis cs:670.5,100);
\draw[draw=none,fill=steelblue31119180] (axis cs:670.5,0) rectangle (axis cs:671.5,100);
\draw[draw=none,fill=steelblue31119180] (axis cs:671.5,0) rectangle (axis cs:672.5,100);
\draw[draw=none,fill=steelblue31119180] (axis cs:672.5,0) rectangle (axis cs:673.5,100);
\draw[draw=none,fill=steelblue31119180] (axis cs:673.5,0) rectangle (axis cs:674.5,100);
\draw[draw=none,fill=steelblue31119180] (axis cs:674.5,0) rectangle (axis cs:675.5,100);
\draw[draw=none,fill=steelblue31119180] (axis cs:675.5,0) rectangle (axis cs:676.5,100);
\draw[draw=none,fill=steelblue31119180] (axis cs:676.5,0) rectangle (axis cs:677.5,100);
\draw[draw=none,fill=steelblue31119180] (axis cs:677.5,0) rectangle (axis cs:678.5,100);
\draw[draw=none,fill=steelblue31119180] (axis cs:678.5,0) rectangle (axis cs:679.5,100);
\draw[draw=none,fill=steelblue31119180] (axis cs:679.5,0) rectangle (axis cs:680.5,100);
\draw[draw=none,fill=steelblue31119180] (axis cs:680.5,0) rectangle (axis cs:681.5,100);
\draw[draw=none,fill=steelblue31119180] (axis cs:681.5,0) rectangle (axis cs:682.5,100);
\draw[draw=none,fill=steelblue31119180] (axis cs:682.5,0) rectangle (axis cs:683.5,100);
\draw[draw=none,fill=steelblue31119180] (axis cs:683.5,0) rectangle (axis cs:684.5,100);
\draw[draw=none,fill=steelblue31119180] (axis cs:684.5,0) rectangle (axis cs:685.5,100);
\draw[draw=none,fill=steelblue31119180] (axis cs:685.5,0) rectangle (axis cs:686.5,100);
\draw[draw=none,fill=steelblue31119180] (axis cs:686.5,0) rectangle (axis cs:687.5,100);
\draw[draw=none,fill=steelblue31119180] (axis cs:687.5,0) rectangle (axis cs:688.5,100);
\draw[draw=none,fill=steelblue31119180] (axis cs:688.5,0) rectangle (axis cs:689.5,100);
\draw[draw=none,fill=steelblue31119180] (axis cs:689.5,0) rectangle (axis cs:690.5,100);
\draw[draw=none,fill=steelblue31119180] (axis cs:690.5,0) rectangle (axis cs:691.5,100);
\draw[draw=none,fill=steelblue31119180] (axis cs:691.5,0) rectangle (axis cs:692.5,100);
\draw[draw=none,fill=steelblue31119180] (axis cs:692.5,0) rectangle (axis cs:693.5,100);
\draw[draw=none,fill=steelblue31119180] (axis cs:693.5,0) rectangle (axis cs:694.5,100);
\draw[draw=none,fill=steelblue31119180] (axis cs:694.5,0) rectangle (axis cs:695.5,100);
\draw[draw=none,fill=steelblue31119180] (axis cs:695.5,0) rectangle (axis cs:696.5,100);
\draw[draw=none,fill=steelblue31119180] (axis cs:696.5,0) rectangle (axis cs:697.5,100);
\draw[draw=none,fill=steelblue31119180] (axis cs:697.5,0) rectangle (axis cs:698.5,100);
\draw[draw=none,fill=steelblue31119180] (axis cs:698.5,0) rectangle (axis cs:699.5,100);
\draw[draw=none,fill=steelblue31119180] (axis cs:699.5,0) rectangle (axis cs:700.5,100);
\draw[draw=none,fill=steelblue31119180] (axis cs:700.5,0) rectangle (axis cs:701.5,100);
\draw[draw=none,fill=steelblue31119180] (axis cs:701.5,0) rectangle (axis cs:702.5,100);
\draw[draw=none,fill=steelblue31119180] (axis cs:702.5,0) rectangle (axis cs:703.5,100);
\draw[draw=none,fill=steelblue31119180] (axis cs:703.5,0) rectangle (axis cs:704.5,100);
\draw[draw=none,fill=steelblue31119180] (axis cs:704.5,0) rectangle (axis cs:705.5,100);
\draw[draw=none,fill=steelblue31119180] (axis cs:705.5,0) rectangle (axis cs:706.5,100);
\draw[draw=none,fill=steelblue31119180] (axis cs:706.5,0) rectangle (axis cs:707.5,100);
\draw[draw=none,fill=steelblue31119180] (axis cs:707.5,0) rectangle (axis cs:708.5,100);
\draw[draw=none,fill=steelblue31119180] (axis cs:708.5,0) rectangle (axis cs:709.5,100);
\draw[draw=none,fill=steelblue31119180] (axis cs:709.5,0) rectangle (axis cs:710.5,100);
\draw[draw=none,fill=steelblue31119180] (axis cs:710.5,0) rectangle (axis cs:711.5,100);
\draw[draw=none,fill=steelblue31119180] (axis cs:711.5,0) rectangle (axis cs:712.5,100);
\draw[draw=none,fill=steelblue31119180] (axis cs:712.5,0) rectangle (axis cs:713.5,100);
\draw[draw=none,fill=steelblue31119180] (axis cs:713.5,0) rectangle (axis cs:714.5,100);
\draw[draw=none,fill=steelblue31119180] (axis cs:714.5,0) rectangle (axis cs:715.5,100);
\draw[draw=none,fill=steelblue31119180] (axis cs:715.5,0) rectangle (axis cs:716.5,100);
\draw[draw=none,fill=steelblue31119180] (axis cs:716.5,0) rectangle (axis cs:717.5,100);
\draw[draw=none,fill=steelblue31119180] (axis cs:717.5,0) rectangle (axis cs:718.5,100);
\draw[draw=none,fill=steelblue31119180] (axis cs:718.5,0) rectangle (axis cs:719.5,100);
\draw[draw=none,fill=steelblue31119180] (axis cs:719.5,0) rectangle (axis cs:720.5,100);
\draw[draw=none,fill=steelblue31119180] (axis cs:720.5,0) rectangle (axis cs:721.5,100);
\draw[draw=none,fill=steelblue31119180] (axis cs:721.5,0) rectangle (axis cs:722.5,100);
\draw[draw=none,fill=steelblue31119180] (axis cs:722.5,0) rectangle (axis cs:723.5,100);
\draw[draw=none,fill=steelblue31119180] (axis cs:723.5,0) rectangle (axis cs:724.5,100);
\draw[draw=none,fill=steelblue31119180] (axis cs:724.5,0) rectangle (axis cs:725.5,100);
\draw[draw=none,fill=steelblue31119180] (axis cs:725.5,0) rectangle (axis cs:726.5,100);
\draw[draw=none,fill=steelblue31119180] (axis cs:726.5,0) rectangle (axis cs:727.5,100);
\draw[draw=none,fill=steelblue31119180] (axis cs:727.5,0) rectangle (axis cs:728.5,100);
\draw[draw=none,fill=steelblue31119180] (axis cs:728.5,0) rectangle (axis cs:729.5,100);
\draw[draw=none,fill=steelblue31119180] (axis cs:729.5,0) rectangle (axis cs:730.5,100);
\draw[draw=none,fill=steelblue31119180] (axis cs:730.5,0) rectangle (axis cs:731.5,100);
\draw[draw=none,fill=steelblue31119180] (axis cs:731.5,0) rectangle (axis cs:732.5,100);
\draw[draw=none,fill=steelblue31119180] (axis cs:732.5,0) rectangle (axis cs:733.5,100);
\draw[draw=none,fill=steelblue31119180] (axis cs:733.5,0) rectangle (axis cs:734.5,100);
\draw[draw=none,fill=steelblue31119180] (axis cs:734.5,0) rectangle (axis cs:735.5,100);
\draw[draw=none,fill=steelblue31119180] (axis cs:735.5,0) rectangle (axis cs:736.5,100);
\draw[draw=none,fill=steelblue31119180] (axis cs:736.5,0) rectangle (axis cs:737.5,100);
\draw[draw=none,fill=steelblue31119180] (axis cs:737.5,0) rectangle (axis cs:738.5,100);
\draw[draw=none,fill=steelblue31119180] (axis cs:738.5,0) rectangle (axis cs:739.5,100);
\draw[draw=none,fill=steelblue31119180] (axis cs:739.5,0) rectangle (axis cs:740.5,100);
\draw[draw=none,fill=steelblue31119180] (axis cs:740.5,0) rectangle (axis cs:741.5,100);
\draw[draw=none,fill=steelblue31119180] (axis cs:741.5,0) rectangle (axis cs:742.5,100);
\draw[draw=none,fill=steelblue31119180] (axis cs:742.5,0) rectangle (axis cs:743.5,100);
\draw[draw=none,fill=steelblue31119180] (axis cs:743.5,0) rectangle (axis cs:744.5,100);
\draw[draw=none,fill=steelblue31119180] (axis cs:744.5,0) rectangle (axis cs:745.5,100);
\draw[draw=none,fill=steelblue31119180] (axis cs:745.5,0) rectangle (axis cs:746.5,100);
\draw[draw=none,fill=steelblue31119180] (axis cs:746.5,0) rectangle (axis cs:747.5,100);
\draw[draw=none,fill=steelblue31119180] (axis cs:747.5,0) rectangle (axis cs:748.5,100);
\draw[draw=none,fill=steelblue31119180] (axis cs:748.5,0) rectangle (axis cs:749.5,100);
\draw[draw=none,fill=steelblue31119180] (axis cs:749.5,0) rectangle (axis cs:750.5,100);
\draw[draw=none,fill=steelblue31119180] (axis cs:750.5,0) rectangle (axis cs:751.5,100);
\draw[draw=none,fill=steelblue31119180] (axis cs:751.5,0) rectangle (axis cs:752.5,100);
\draw[draw=none,fill=steelblue31119180] (axis cs:752.5,0) rectangle (axis cs:753.5,100);
\draw[draw=none,fill=steelblue31119180] (axis cs:753.5,0) rectangle (axis cs:754.5,100);
\draw[draw=none,fill=steelblue31119180] (axis cs:754.5,0) rectangle (axis cs:755.5,100);
\draw[draw=none,fill=steelblue31119180] (axis cs:755.5,0) rectangle (axis cs:756.5,100);
\draw[draw=none,fill=steelblue31119180] (axis cs:756.5,0) rectangle (axis cs:757.5,100);
\draw[draw=none,fill=steelblue31119180] (axis cs:757.5,0) rectangle (axis cs:758.5,100);
\draw[draw=none,fill=steelblue31119180] (axis cs:758.5,0) rectangle (axis cs:759.5,100);
\draw[draw=none,fill=steelblue31119180] (axis cs:759.5,0) rectangle (axis cs:760.5,100);
\draw[draw=none,fill=steelblue31119180] (axis cs:760.5,0) rectangle (axis cs:761.5,100);
\draw[draw=none,fill=steelblue31119180] (axis cs:761.5,0) rectangle (axis cs:762.5,100);
\draw[draw=none,fill=steelblue31119180] (axis cs:762.5,0) rectangle (axis cs:763.5,100);
\draw[draw=none,fill=steelblue31119180] (axis cs:763.5,0) rectangle (axis cs:764.5,100);
\draw[draw=none,fill=steelblue31119180] (axis cs:764.5,0) rectangle (axis cs:765.5,100);
\draw[draw=none,fill=steelblue31119180] (axis cs:765.5,0) rectangle (axis cs:766.5,100);
\draw[draw=none,fill=steelblue31119180] (axis cs:766.5,0) rectangle (axis cs:767.5,100);
\draw[draw=none,fill=steelblue31119180] (axis cs:767.5,0) rectangle (axis cs:768.5,100);
\draw[draw=none,fill=steelblue31119180] (axis cs:768.5,0) rectangle (axis cs:769.5,100);
\draw[draw=none,fill=steelblue31119180] (axis cs:769.5,0) rectangle (axis cs:770.5,100);
\draw[draw=none,fill=steelblue31119180] (axis cs:770.5,0) rectangle (axis cs:771.5,100);
\draw[draw=none,fill=steelblue31119180] (axis cs:771.5,0) rectangle (axis cs:772.5,100);
\draw[draw=none,fill=steelblue31119180] (axis cs:772.5,0) rectangle (axis cs:773.5,100);
\draw[draw=none,fill=steelblue31119180] (axis cs:773.5,0) rectangle (axis cs:774.5,100);
\draw[draw=none,fill=steelblue31119180] (axis cs:774.5,0) rectangle (axis cs:775.5,100);
\draw[draw=none,fill=steelblue31119180] (axis cs:775.5,0) rectangle (axis cs:776.5,100);
\draw[draw=none,fill=steelblue31119180] (axis cs:776.5,0) rectangle (axis cs:777.5,100);
\draw[draw=none,fill=steelblue31119180] (axis cs:777.5,0) rectangle (axis cs:778.5,100);
\draw[draw=none,fill=steelblue31119180] (axis cs:778.5,0) rectangle (axis cs:779.5,100);
\draw[draw=none,fill=steelblue31119180] (axis cs:779.5,0) rectangle (axis cs:780.5,100);
\draw[draw=none,fill=steelblue31119180] (axis cs:780.5,0) rectangle (axis cs:781.5,100);
\draw[draw=none,fill=steelblue31119180] (axis cs:781.5,0) rectangle (axis cs:782.5,100);
\draw[draw=none,fill=steelblue31119180] (axis cs:782.5,0) rectangle (axis cs:783.5,100);
\draw[draw=none,fill=steelblue31119180] (axis cs:783.5,0) rectangle (axis cs:784.5,100);
\draw[draw=none,fill=steelblue31119180] (axis cs:784.5,0) rectangle (axis cs:785.5,100);
\draw[draw=none,fill=steelblue31119180] (axis cs:785.5,0) rectangle (axis cs:786.5,100);
\draw[draw=none,fill=darkorange25512714] (axis cs:-0.5,0) rectangle (axis cs:0.5,9.87261146496815);
\addlegendimage{ybar,ybar legend,draw=none,fill=darkorange25512714}
\addlegendentry{$\Delta_L=2$}

\draw[draw=none,fill=darkorange25512714] (axis cs:0.5,0) rectangle (axis cs:1.5,19.7452229299363);
\draw[draw=none,fill=darkorange25512714] (axis cs:1.5,0) rectangle (axis cs:2.5,26.7515923566879);
\draw[draw=none,fill=darkorange25512714] (axis cs:2.5,0) rectangle (axis cs:3.5,31.5286624203822);
\draw[draw=none,fill=darkorange25512714] (axis cs:3.5,0) rectangle (axis cs:4.5,35.3503184713376);
\draw[draw=none,fill=darkorange25512714] (axis cs:4.5,0) rectangle (axis cs:5.5,38.2165605095541);
\draw[draw=none,fill=darkorange25512714] (axis cs:5.5,0) rectangle (axis cs:6.5,40.4458598726115);
\draw[draw=none,fill=darkorange25512714] (axis cs:6.5,0) rectangle (axis cs:7.5,44.9044585987261);
\draw[draw=none,fill=darkorange25512714] (axis cs:7.5,0) rectangle (axis cs:8.5,48.0891719745223);
\draw[draw=none,fill=darkorange25512714] (axis cs:8.5,0) rectangle (axis cs:9.5,51.5923566878981);
\draw[draw=none,fill=darkorange25512714] (axis cs:9.5,0) rectangle (axis cs:10.5,54.7770700636943);
\draw[draw=none,fill=darkorange25512714] (axis cs:10.5,0) rectangle (axis cs:11.5,57.6433121019108);
\draw[draw=none,fill=darkorange25512714] (axis cs:11.5,0) rectangle (axis cs:12.5,58.5987261146497);
\draw[draw=none,fill=darkorange25512714] (axis cs:12.5,0) rectangle (axis cs:13.5,61.1464968152866);
\draw[draw=none,fill=darkorange25512714] (axis cs:13.5,0) rectangle (axis cs:14.5,63.3757961783439);
\draw[draw=none,fill=darkorange25512714] (axis cs:14.5,0) rectangle (axis cs:15.5,66.8789808917197);
\draw[draw=none,fill=darkorange25512714] (axis cs:15.5,0) rectangle (axis cs:16.5,69.7452229299363);
\draw[draw=none,fill=darkorange25512714] (axis cs:16.5,0) rectangle (axis cs:17.5,70.3821656050955);
\draw[draw=none,fill=darkorange25512714] (axis cs:17.5,0) rectangle (axis cs:18.5,72.9299363057325);
\draw[draw=none,fill=darkorange25512714] (axis cs:18.5,0) rectangle (axis cs:19.5,75.1592356687898);
\draw[draw=none,fill=darkorange25512714] (axis cs:19.5,0) rectangle (axis cs:20.5,76.1146496815287);
\draw[draw=none,fill=darkorange25512714] (axis cs:20.5,0) rectangle (axis cs:21.5,78.343949044586);
\draw[draw=none,fill=darkorange25512714] (axis cs:21.5,0) rectangle (axis cs:22.5,79.2993630573248);
\draw[draw=none,fill=darkorange25512714] (axis cs:22.5,0) rectangle (axis cs:23.5,80.2547770700637);
\draw[draw=none,fill=darkorange25512714] (axis cs:23.5,0) rectangle (axis cs:24.5,80.5732484076433);
\draw[draw=none,fill=darkorange25512714] (axis cs:24.5,0) rectangle (axis cs:25.5,82.1656050955414);
\draw[draw=none,fill=darkorange25512714] (axis cs:25.5,0) rectangle (axis cs:26.5,82.484076433121);
\draw[draw=none,fill=darkorange25512714] (axis cs:26.5,0) rectangle (axis cs:27.5,82.8025477707006);
\draw[draw=none,fill=darkorange25512714] (axis cs:27.5,0) rectangle (axis cs:28.5,84.0764331210191);
\draw[draw=none,fill=darkorange25512714] (axis cs:28.5,0) rectangle (axis cs:29.5,85.6687898089172);
\draw[draw=none,fill=darkorange25512714] (axis cs:29.5,0) rectangle (axis cs:30.5,86.3057324840764);
\draw[draw=none,fill=darkorange25512714] (axis cs:30.5,0) rectangle (axis cs:31.5,86.3057324840764);
\draw[draw=none,fill=darkorange25512714] (axis cs:31.5,0) rectangle (axis cs:32.5,87.5796178343949);
\draw[draw=none,fill=darkorange25512714] (axis cs:32.5,0) rectangle (axis cs:33.5,89.171974522293);
\draw[draw=none,fill=darkorange25512714] (axis cs:33.5,0) rectangle (axis cs:34.5,90.1273885350319);
\draw[draw=none,fill=darkorange25512714] (axis cs:34.5,0) rectangle (axis cs:35.5,90.4458598726115);
\draw[draw=none,fill=darkorange25512714] (axis cs:35.5,0) rectangle (axis cs:36.5,90.7643312101911);
\draw[draw=none,fill=darkorange25512714] (axis cs:36.5,0) rectangle (axis cs:37.5,91.7197452229299);
\draw[draw=none,fill=darkorange25512714] (axis cs:37.5,0) rectangle (axis cs:38.5,91.7197452229299);
\draw[draw=none,fill=darkorange25512714] (axis cs:38.5,0) rectangle (axis cs:39.5,92.0382165605096);
\draw[draw=none,fill=darkorange25512714] (axis cs:39.5,0) rectangle (axis cs:40.5,92.9936305732484);
\draw[draw=none,fill=darkorange25512714] (axis cs:40.5,0) rectangle (axis cs:41.5,93.312101910828);
\draw[draw=none,fill=darkorange25512714] (axis cs:41.5,0) rectangle (axis cs:42.5,93.312101910828);
\draw[draw=none,fill=darkorange25512714] (axis cs:42.5,0) rectangle (axis cs:43.5,93.312101910828);
\draw[draw=none,fill=darkorange25512714] (axis cs:43.5,0) rectangle (axis cs:44.5,93.6305732484076);
\draw[draw=none,fill=darkorange25512714] (axis cs:44.5,0) rectangle (axis cs:45.5,95.2229299363057);
\draw[draw=none,fill=darkorange25512714] (axis cs:45.5,0) rectangle (axis cs:46.5,95.2229299363057);
\draw[draw=none,fill=darkorange25512714] (axis cs:46.5,0) rectangle (axis cs:47.5,95.2229299363057);
\draw[draw=none,fill=darkorange25512714] (axis cs:47.5,0) rectangle (axis cs:48.5,95.2229299363057);
\draw[draw=none,fill=darkorange25512714] (axis cs:48.5,0) rectangle (axis cs:49.5,95.2229299363057);
\draw[draw=none,fill=darkorange25512714] (axis cs:49.5,0) rectangle (axis cs:50.5,95.5414012738854);
\draw[draw=none,fill=darkorange25512714] (axis cs:50.5,0) rectangle (axis cs:51.5,95.5414012738854);
\draw[draw=none,fill=darkorange25512714] (axis cs:51.5,0) rectangle (axis cs:52.5,96.1783439490446);
\draw[draw=none,fill=darkorange25512714] (axis cs:52.5,0) rectangle (axis cs:53.5,96.1783439490446);
\draw[draw=none,fill=darkorange25512714] (axis cs:53.5,0) rectangle (axis cs:54.5,96.1783439490446);
\draw[draw=none,fill=darkorange25512714] (axis cs:54.5,0) rectangle (axis cs:55.5,96.1783439490446);
\draw[draw=none,fill=darkorange25512714] (axis cs:55.5,0) rectangle (axis cs:56.5,96.1783439490446);
\draw[draw=none,fill=darkorange25512714] (axis cs:56.5,0) rectangle (axis cs:57.5,96.8152866242038);
\draw[draw=none,fill=darkorange25512714] (axis cs:57.5,0) rectangle (axis cs:58.5,97.1337579617834);
\draw[draw=none,fill=darkorange25512714] (axis cs:58.5,0) rectangle (axis cs:59.5,97.1337579617834);
\draw[draw=none,fill=darkorange25512714] (axis cs:59.5,0) rectangle (axis cs:60.5,97.1337579617834);
\draw[draw=none,fill=darkorange25512714] (axis cs:60.5,0) rectangle (axis cs:61.5,97.1337579617834);
\draw[draw=none,fill=darkorange25512714] (axis cs:61.5,0) rectangle (axis cs:62.5,97.4522292993631);
\draw[draw=none,fill=darkorange25512714] (axis cs:62.5,0) rectangle (axis cs:63.5,98.0891719745223);
\draw[draw=none,fill=darkorange25512714] (axis cs:63.5,0) rectangle (axis cs:64.5,98.0891719745223);
\draw[draw=none,fill=darkorange25512714] (axis cs:64.5,0) rectangle (axis cs:65.5,98.4076433121019);
\draw[draw=none,fill=darkorange25512714] (axis cs:65.5,0) rectangle (axis cs:66.5,98.7261146496815);
\draw[draw=none,fill=darkorange25512714] (axis cs:66.5,0) rectangle (axis cs:67.5,98.7261146496815);
\draw[draw=none,fill=darkorange25512714] (axis cs:67.5,0) rectangle (axis cs:68.5,98.7261146496815);
\draw[draw=none,fill=darkorange25512714] (axis cs:68.5,0) rectangle (axis cs:69.5,99.0445859872611);
\draw[draw=none,fill=darkorange25512714] (axis cs:69.5,0) rectangle (axis cs:70.5,99.0445859872611);
\draw[draw=none,fill=darkorange25512714] (axis cs:70.5,0) rectangle (axis cs:71.5,99.0445859872611);
\draw[draw=none,fill=darkorange25512714] (axis cs:71.5,0) rectangle (axis cs:72.5,99.0445859872611);
\draw[draw=none,fill=darkorange25512714] (axis cs:72.5,0) rectangle (axis cs:73.5,99.3630573248408);
\draw[draw=none,fill=darkorange25512714] (axis cs:73.5,0) rectangle (axis cs:74.5,99.3630573248408);
\draw[draw=none,fill=darkorange25512714] (axis cs:74.5,0) rectangle (axis cs:75.5,99.3630573248408);
\draw[draw=none,fill=darkorange25512714] (axis cs:75.5,0) rectangle (axis cs:76.5,99.3630573248408);
\draw[draw=none,fill=darkorange25512714] (axis cs:76.5,0) rectangle (axis cs:77.5,99.3630573248408);
\draw[draw=none,fill=darkorange25512714] (axis cs:77.5,0) rectangle (axis cs:78.5,99.6815286624204);
\draw[draw=none,fill=darkorange25512714] (axis cs:78.5,0) rectangle (axis cs:79.5,100);
\draw[draw=none,fill=darkorange25512714] (axis cs:79.5,0) rectangle (axis cs:80.5,100);
\draw[draw=none,fill=darkorange25512714] (axis cs:80.5,0) rectangle (axis cs:81.5,100);
\draw[draw=none,fill=darkorange25512714] (axis cs:81.5,0) rectangle (axis cs:82.5,100);
\draw[draw=none,fill=darkorange25512714] (axis cs:82.5,0) rectangle (axis cs:83.5,100);
\draw[draw=none,fill=darkorange25512714] (axis cs:83.5,0) rectangle (axis cs:84.5,100);
\draw[draw=none,fill=darkorange25512714] (axis cs:84.5,0) rectangle (axis cs:85.5,100);
\draw[draw=none,fill=darkorange25512714] (axis cs:85.5,0) rectangle (axis cs:86.5,100);
\draw[draw=none,fill=darkorange25512714] (axis cs:86.5,0) rectangle (axis cs:87.5,100);
\draw[draw=none,fill=darkorange25512714] (axis cs:87.5,0) rectangle (axis cs:88.5,100);
\draw[draw=none,fill=darkorange25512714] (axis cs:88.5,0) rectangle (axis cs:89.5,100);
\draw[draw=none,fill=darkorange25512714] (axis cs:89.5,0) rectangle (axis cs:90.5,100);
\draw[draw=none,fill=darkorange25512714] (axis cs:90.5,0) rectangle (axis cs:91.5,100);
\draw[draw=none,fill=darkorange25512714] (axis cs:91.5,0) rectangle (axis cs:92.5,100);
\draw[draw=none,fill=darkorange25512714] (axis cs:92.5,0) rectangle (axis cs:93.5,100);
\draw[draw=none,fill=darkorange25512714] (axis cs:93.5,0) rectangle (axis cs:94.5,100);
\draw[draw=none,fill=darkorange25512714] (axis cs:94.5,0) rectangle (axis cs:95.5,100);
\draw[draw=none,fill=darkorange25512714] (axis cs:95.5,0) rectangle (axis cs:96.5,100);
\draw[draw=none,fill=darkorange25512714] (axis cs:96.5,0) rectangle (axis cs:97.5,100);
\draw[draw=none,fill=darkorange25512714] (axis cs:97.5,0) rectangle (axis cs:98.5,100);
\draw[draw=none,fill=darkorange25512714] (axis cs:98.5,0) rectangle (axis cs:99.5,100);
\draw[draw=none,fill=darkorange25512714] (axis cs:99.5,0) rectangle (axis cs:100.5,100);
\draw[draw=none,fill=darkorange25512714] (axis cs:100.5,0) rectangle (axis cs:101.5,100);
\draw[draw=none,fill=darkorange25512714] (axis cs:101.5,0) rectangle (axis cs:102.5,100);
\draw[draw=none,fill=darkorange25512714] (axis cs:102.5,0) rectangle (axis cs:103.5,100);
\draw[draw=none,fill=darkorange25512714] (axis cs:103.5,0) rectangle (axis cs:104.5,100);
\draw[draw=none,fill=darkorange25512714] (axis cs:104.5,0) rectangle (axis cs:105.5,100);
\draw[draw=none,fill=darkorange25512714] (axis cs:105.5,0) rectangle (axis cs:106.5,100);
\draw[draw=none,fill=darkorange25512714] (axis cs:106.5,0) rectangle (axis cs:107.5,100);
\draw[draw=none,fill=darkorange25512714] (axis cs:107.5,0) rectangle (axis cs:108.5,100);
\draw[draw=none,fill=darkorange25512714] (axis cs:108.5,0) rectangle (axis cs:109.5,100);
\draw[draw=none,fill=darkorange25512714] (axis cs:109.5,0) rectangle (axis cs:110.5,100);
\draw[draw=none,fill=darkorange25512714] (axis cs:110.5,0) rectangle (axis cs:111.5,100);
\draw[draw=none,fill=darkorange25512714] (axis cs:111.5,0) rectangle (axis cs:112.5,100);
\draw[draw=none,fill=darkorange25512714] (axis cs:112.5,0) rectangle (axis cs:113.5,100);
\draw[draw=none,fill=darkorange25512714] (axis cs:113.5,0) rectangle (axis cs:114.5,100);
\draw[draw=none,fill=darkorange25512714] (axis cs:114.5,0) rectangle (axis cs:115.5,100);
\draw[draw=none,fill=darkorange25512714] (axis cs:115.5,0) rectangle (axis cs:116.5,100);
\draw[draw=none,fill=darkorange25512714] (axis cs:116.5,0) rectangle (axis cs:117.5,100);
\draw[draw=none,fill=darkorange25512714] (axis cs:117.5,0) rectangle (axis cs:118.5,100);
\draw[draw=none,fill=darkorange25512714] (axis cs:118.5,0) rectangle (axis cs:119.5,100);
\draw[draw=none,fill=darkorange25512714] (axis cs:119.5,0) rectangle (axis cs:120.5,100);
\draw[draw=none,fill=darkorange25512714] (axis cs:120.5,0) rectangle (axis cs:121.5,100);
\draw[draw=none,fill=darkorange25512714] (axis cs:121.5,0) rectangle (axis cs:122.5,100);
\draw[draw=none,fill=darkorange25512714] (axis cs:122.5,0) rectangle (axis cs:123.5,100);
\draw[draw=none,fill=darkorange25512714] (axis cs:123.5,0) rectangle (axis cs:124.5,100);
\draw[draw=none,fill=darkorange25512714] (axis cs:124.5,0) rectangle (axis cs:125.5,100);
\draw[draw=none,fill=darkorange25512714] (axis cs:125.5,0) rectangle (axis cs:126.5,100);
\draw[draw=none,fill=darkorange25512714] (axis cs:126.5,0) rectangle (axis cs:127.5,100);
\draw[draw=none,fill=darkorange25512714] (axis cs:127.5,0) rectangle (axis cs:128.5,100);
\draw[draw=none,fill=darkorange25512714] (axis cs:128.5,0) rectangle (axis cs:129.5,100);
\draw[draw=none,fill=darkorange25512714] (axis cs:129.5,0) rectangle (axis cs:130.5,100);
\draw[draw=none,fill=darkorange25512714] (axis cs:130.5,0) rectangle (axis cs:131.5,100);
\draw[draw=none,fill=darkorange25512714] (axis cs:131.5,0) rectangle (axis cs:132.5,100);
\draw[draw=none,fill=darkorange25512714] (axis cs:132.5,0) rectangle (axis cs:133.5,100);
\draw[draw=none,fill=darkorange25512714] (axis cs:133.5,0) rectangle (axis cs:134.5,100);
\draw[draw=none,fill=darkorange25512714] (axis cs:134.5,0) rectangle (axis cs:135.5,100);
\draw[draw=none,fill=darkorange25512714] (axis cs:135.5,0) rectangle (axis cs:136.5,100);
\draw[draw=none,fill=darkorange25512714] (axis cs:136.5,0) rectangle (axis cs:137.5,100);
\draw[draw=none,fill=darkorange25512714] (axis cs:137.5,0) rectangle (axis cs:138.5,100);
\draw[draw=none,fill=darkorange25512714] (axis cs:138.5,0) rectangle (axis cs:139.5,100);
\draw[draw=none,fill=darkorange25512714] (axis cs:139.5,0) rectangle (axis cs:140.5,100);
\draw[draw=none,fill=darkorange25512714] (axis cs:140.5,0) rectangle (axis cs:141.5,100);
\draw[draw=none,fill=darkorange25512714] (axis cs:141.5,0) rectangle (axis cs:142.5,100);
\draw[draw=none,fill=darkorange25512714] (axis cs:142.5,0) rectangle (axis cs:143.5,100);
\draw[draw=none,fill=darkorange25512714] (axis cs:143.5,0) rectangle (axis cs:144.5,100);
\draw[draw=none,fill=darkorange25512714] (axis cs:144.5,0) rectangle (axis cs:145.5,100);
\draw[draw=none,fill=darkorange25512714] (axis cs:145.5,0) rectangle (axis cs:146.5,100);
\draw[draw=none,fill=darkorange25512714] (axis cs:146.5,0) rectangle (axis cs:147.5,100);
\draw[draw=none,fill=darkorange25512714] (axis cs:147.5,0) rectangle (axis cs:148.5,100);
\draw[draw=none,fill=darkorange25512714] (axis cs:148.5,0) rectangle (axis cs:149.5,100);
\draw[draw=none,fill=darkorange25512714] (axis cs:149.5,0) rectangle (axis cs:150.5,100);
\draw[draw=none,fill=darkorange25512714] (axis cs:150.5,0) rectangle (axis cs:151.5,100);
\draw[draw=none,fill=darkorange25512714] (axis cs:151.5,0) rectangle (axis cs:152.5,100);
\draw[draw=none,fill=darkorange25512714] (axis cs:152.5,0) rectangle (axis cs:153.5,100);
\draw[draw=none,fill=darkorange25512714] (axis cs:153.5,0) rectangle (axis cs:154.5,100);
\draw[draw=none,fill=darkorange25512714] (axis cs:154.5,0) rectangle (axis cs:155.5,100);
\draw[draw=none,fill=darkorange25512714] (axis cs:155.5,0) rectangle (axis cs:156.5,100);
\draw[draw=none,fill=darkorange25512714] (axis cs:156.5,0) rectangle (axis cs:157.5,100);
\draw[draw=none,fill=darkorange25512714] (axis cs:157.5,0) rectangle (axis cs:158.5,100);
\draw[draw=none,fill=darkorange25512714] (axis cs:158.5,0) rectangle (axis cs:159.5,100);
\draw[draw=none,fill=darkorange25512714] (axis cs:159.5,0) rectangle (axis cs:160.5,100);
\draw[draw=none,fill=darkorange25512714] (axis cs:160.5,0) rectangle (axis cs:161.5,100);
\draw[draw=none,fill=darkorange25512714] (axis cs:161.5,0) rectangle (axis cs:162.5,100);
\draw[draw=none,fill=darkorange25512714] (axis cs:162.5,0) rectangle (axis cs:163.5,100);
\draw[draw=none,fill=darkorange25512714] (axis cs:163.5,0) rectangle (axis cs:164.5,100);
\draw[draw=none,fill=darkorange25512714] (axis cs:164.5,0) rectangle (axis cs:165.5,100);
\draw[draw=none,fill=darkorange25512714] (axis cs:165.5,0) rectangle (axis cs:166.5,100);
\draw[draw=none,fill=darkorange25512714] (axis cs:166.5,0) rectangle (axis cs:167.5,100);
\draw[draw=none,fill=darkorange25512714] (axis cs:167.5,0) rectangle (axis cs:168.5,100);
\draw[draw=none,fill=darkorange25512714] (axis cs:168.5,0) rectangle (axis cs:169.5,100);
\draw[draw=none,fill=darkorange25512714] (axis cs:169.5,0) rectangle (axis cs:170.5,100);
\draw[draw=none,fill=darkorange25512714] (axis cs:170.5,0) rectangle (axis cs:171.5,100);
\draw[draw=none,fill=darkorange25512714] (axis cs:171.5,0) rectangle (axis cs:172.5,100);
\draw[draw=none,fill=darkorange25512714] (axis cs:172.5,0) rectangle (axis cs:173.5,100);
\draw[draw=none,fill=darkorange25512714] (axis cs:173.5,0) rectangle (axis cs:174.5,100);
\draw[draw=none,fill=darkorange25512714] (axis cs:174.5,0) rectangle (axis cs:175.5,100);
\draw[draw=none,fill=darkorange25512714] (axis cs:175.5,0) rectangle (axis cs:176.5,100);
\draw[draw=none,fill=darkorange25512714] (axis cs:176.5,0) rectangle (axis cs:177.5,100);
\draw[draw=none,fill=darkorange25512714] (axis cs:177.5,0) rectangle (axis cs:178.5,100);
\draw[draw=none,fill=darkorange25512714] (axis cs:178.5,0) rectangle (axis cs:179.5,100);
\draw[draw=none,fill=darkorange25512714] (axis cs:179.5,0) rectangle (axis cs:180.5,100);
\draw[draw=none,fill=darkorange25512714] (axis cs:180.5,0) rectangle (axis cs:181.5,100);
\draw[draw=none,fill=darkorange25512714] (axis cs:181.5,0) rectangle (axis cs:182.5,100);
\draw[draw=none,fill=darkorange25512714] (axis cs:182.5,0) rectangle (axis cs:183.5,100);
\draw[draw=none,fill=darkorange25512714] (axis cs:183.5,0) rectangle (axis cs:184.5,100);
\draw[draw=none,fill=darkorange25512714] (axis cs:184.5,0) rectangle (axis cs:185.5,100);
\draw[draw=none,fill=darkorange25512714] (axis cs:185.5,0) rectangle (axis cs:186.5,100);
\draw[draw=none,fill=darkorange25512714] (axis cs:186.5,0) rectangle (axis cs:187.5,100);
\draw[draw=none,fill=darkorange25512714] (axis cs:187.5,0) rectangle (axis cs:188.5,100);
\draw[draw=none,fill=darkorange25512714] (axis cs:188.5,0) rectangle (axis cs:189.5,100);
\draw[draw=none,fill=darkorange25512714] (axis cs:189.5,0) rectangle (axis cs:190.5,100);
\draw[draw=none,fill=darkorange25512714] (axis cs:190.5,0) rectangle (axis cs:191.5,100);
\draw[draw=none,fill=darkorange25512714] (axis cs:191.5,0) rectangle (axis cs:192.5,100);
\draw[draw=none,fill=darkorange25512714] (axis cs:192.5,0) rectangle (axis cs:193.5,100);
\draw[draw=none,fill=darkorange25512714] (axis cs:193.5,0) rectangle (axis cs:194.5,100);
\draw[draw=none,fill=darkorange25512714] (axis cs:194.5,0) rectangle (axis cs:195.5,100);
\draw[draw=none,fill=darkorange25512714] (axis cs:195.5,0) rectangle (axis cs:196.5,100);
\draw[draw=none,fill=darkorange25512714] (axis cs:196.5,0) rectangle (axis cs:197.5,100);
\draw[draw=none,fill=darkorange25512714] (axis cs:197.5,0) rectangle (axis cs:198.5,100);
\draw[draw=none,fill=darkorange25512714] (axis cs:198.5,0) rectangle (axis cs:199.5,100);
\draw[draw=none,fill=darkorange25512714] (axis cs:199.5,0) rectangle (axis cs:200.5,100);
\draw[draw=none,fill=darkorange25512714] (axis cs:200.5,0) rectangle (axis cs:201.5,100);
\draw[draw=none,fill=darkorange25512714] (axis cs:201.5,0) rectangle (axis cs:202.5,100);
\draw[draw=none,fill=darkorange25512714] (axis cs:202.5,0) rectangle (axis cs:203.5,100);
\draw[draw=none,fill=darkorange25512714] (axis cs:203.5,0) rectangle (axis cs:204.5,100);
\draw[draw=none,fill=darkorange25512714] (axis cs:204.5,0) rectangle (axis cs:205.5,100);
\draw[draw=none,fill=darkorange25512714] (axis cs:205.5,0) rectangle (axis cs:206.5,100);
\draw[draw=none,fill=darkorange25512714] (axis cs:206.5,0) rectangle (axis cs:207.5,100);
\draw[draw=none,fill=darkorange25512714] (axis cs:207.5,0) rectangle (axis cs:208.5,100);
\draw[draw=none,fill=darkorange25512714] (axis cs:208.5,0) rectangle (axis cs:209.5,100);
\draw[draw=none,fill=darkorange25512714] (axis cs:209.5,0) rectangle (axis cs:210.5,100);
\draw[draw=none,fill=darkorange25512714] (axis cs:210.5,0) rectangle (axis cs:211.5,100);
\draw[draw=none,fill=darkorange25512714] (axis cs:211.5,0) rectangle (axis cs:212.5,100);
\draw[draw=none,fill=darkorange25512714] (axis cs:212.5,0) rectangle (axis cs:213.5,100);
\draw[draw=none,fill=darkorange25512714] (axis cs:213.5,0) rectangle (axis cs:214.5,100);
\draw[draw=none,fill=darkorange25512714] (axis cs:214.5,0) rectangle (axis cs:215.5,100);
\draw[draw=none,fill=darkorange25512714] (axis cs:215.5,0) rectangle (axis cs:216.5,100);
\draw[draw=none,fill=darkorange25512714] (axis cs:216.5,0) rectangle (axis cs:217.5,100);
\draw[draw=none,fill=darkorange25512714] (axis cs:217.5,0) rectangle (axis cs:218.5,100);
\draw[draw=none,fill=darkorange25512714] (axis cs:218.5,0) rectangle (axis cs:219.5,100);
\draw[draw=none,fill=darkorange25512714] (axis cs:219.5,0) rectangle (axis cs:220.5,100);
\draw[draw=none,fill=darkorange25512714] (axis cs:220.5,0) rectangle (axis cs:221.5,100);
\draw[draw=none,fill=darkorange25512714] (axis cs:221.5,0) rectangle (axis cs:222.5,100);
\draw[draw=none,fill=darkorange25512714] (axis cs:222.5,0) rectangle (axis cs:223.5,100);
\draw[draw=none,fill=darkorange25512714] (axis cs:223.5,0) rectangle (axis cs:224.5,100);
\draw[draw=none,fill=darkorange25512714] (axis cs:224.5,0) rectangle (axis cs:225.5,100);
\draw[draw=none,fill=darkorange25512714] (axis cs:225.5,0) rectangle (axis cs:226.5,100);
\draw[draw=none,fill=darkorange25512714] (axis cs:226.5,0) rectangle (axis cs:227.5,100);
\draw[draw=none,fill=darkorange25512714] (axis cs:227.5,0) rectangle (axis cs:228.5,100);
\draw[draw=none,fill=darkorange25512714] (axis cs:228.5,0) rectangle (axis cs:229.5,100);
\draw[draw=none,fill=darkorange25512714] (axis cs:229.5,0) rectangle (axis cs:230.5,100);
\draw[draw=none,fill=darkorange25512714] (axis cs:230.5,0) rectangle (axis cs:231.5,100);
\draw[draw=none,fill=darkorange25512714] (axis cs:231.5,0) rectangle (axis cs:232.5,100);
\draw[draw=none,fill=darkorange25512714] (axis cs:232.5,0) rectangle (axis cs:233.5,100);
\draw[draw=none,fill=darkorange25512714] (axis cs:233.5,0) rectangle (axis cs:234.5,100);
\draw[draw=none,fill=darkorange25512714] (axis cs:234.5,0) rectangle (axis cs:235.5,100);
\draw[draw=none,fill=darkorange25512714] (axis cs:235.5,0) rectangle (axis cs:236.5,100);
\draw[draw=none,fill=darkorange25512714] (axis cs:236.5,0) rectangle (axis cs:237.5,100);
\draw[draw=none,fill=darkorange25512714] (axis cs:237.5,0) rectangle (axis cs:238.5,100);
\draw[draw=none,fill=darkorange25512714] (axis cs:238.5,0) rectangle (axis cs:239.5,100);
\draw[draw=none,fill=darkorange25512714] (axis cs:239.5,0) rectangle (axis cs:240.5,100);
\draw[draw=none,fill=darkorange25512714] (axis cs:240.5,0) rectangle (axis cs:241.5,100);
\draw[draw=none,fill=darkorange25512714] (axis cs:241.5,0) rectangle (axis cs:242.5,100);
\draw[draw=none,fill=darkorange25512714] (axis cs:242.5,0) rectangle (axis cs:243.5,100);
\draw[draw=none,fill=darkorange25512714] (axis cs:243.5,0) rectangle (axis cs:244.5,100);
\draw[draw=none,fill=darkorange25512714] (axis cs:244.5,0) rectangle (axis cs:245.5,100);
\draw[draw=none,fill=darkorange25512714] (axis cs:245.5,0) rectangle (axis cs:246.5,100);
\draw[draw=none,fill=darkorange25512714] (axis cs:246.5,0) rectangle (axis cs:247.5,100);
\draw[draw=none,fill=darkorange25512714] (axis cs:247.5,0) rectangle (axis cs:248.5,100);
\draw[draw=none,fill=darkorange25512714] (axis cs:248.5,0) rectangle (axis cs:249.5,100);
\draw[draw=none,fill=darkorange25512714] (axis cs:249.5,0) rectangle (axis cs:250.5,100);
\draw[draw=none,fill=darkorange25512714] (axis cs:250.5,0) rectangle (axis cs:251.5,100);
\draw[draw=none,fill=darkorange25512714] (axis cs:251.5,0) rectangle (axis cs:252.5,100);
\draw[draw=none,fill=darkorange25512714] (axis cs:252.5,0) rectangle (axis cs:253.5,100);
\draw[draw=none,fill=darkorange25512714] (axis cs:253.5,0) rectangle (axis cs:254.5,100);
\draw[draw=none,fill=darkorange25512714] (axis cs:254.5,0) rectangle (axis cs:255.5,100);
\draw[draw=none,fill=darkorange25512714] (axis cs:255.5,0) rectangle (axis cs:256.5,100);
\draw[draw=none,fill=darkorange25512714] (axis cs:256.5,0) rectangle (axis cs:257.5,100);
\draw[draw=none,fill=darkorange25512714] (axis cs:257.5,0) rectangle (axis cs:258.5,100);
\draw[draw=none,fill=darkorange25512714] (axis cs:258.5,0) rectangle (axis cs:259.5,100);
\draw[draw=none,fill=darkorange25512714] (axis cs:259.5,0) rectangle (axis cs:260.5,100);
\draw[draw=none,fill=darkorange25512714] (axis cs:260.5,0) rectangle (axis cs:261.5,100);
\draw[draw=none,fill=darkorange25512714] (axis cs:261.5,0) rectangle (axis cs:262.5,100);
\draw[draw=none,fill=darkorange25512714] (axis cs:262.5,0) rectangle (axis cs:263.5,100);
\draw[draw=none,fill=darkorange25512714] (axis cs:263.5,0) rectangle (axis cs:264.5,100);
\draw[draw=none,fill=darkorange25512714] (axis cs:264.5,0) rectangle (axis cs:265.5,100);
\draw[draw=none,fill=darkorange25512714] (axis cs:265.5,0) rectangle (axis cs:266.5,100);
\draw[draw=none,fill=darkorange25512714] (axis cs:266.5,0) rectangle (axis cs:267.5,100);
\draw[draw=none,fill=darkorange25512714] (axis cs:267.5,0) rectangle (axis cs:268.5,100);
\draw[draw=none,fill=darkorange25512714] (axis cs:268.5,0) rectangle (axis cs:269.5,100);
\draw[draw=none,fill=darkorange25512714] (axis cs:269.5,0) rectangle (axis cs:270.5,100);
\draw[draw=none,fill=darkorange25512714] (axis cs:270.5,0) rectangle (axis cs:271.5,100);
\draw[draw=none,fill=darkorange25512714] (axis cs:271.5,0) rectangle (axis cs:272.5,100);
\draw[draw=none,fill=darkorange25512714] (axis cs:272.5,0) rectangle (axis cs:273.5,100);
\draw[draw=none,fill=darkorange25512714] (axis cs:273.5,0) rectangle (axis cs:274.5,100);
\draw[draw=none,fill=darkorange25512714] (axis cs:274.5,0) rectangle (axis cs:275.5,100);
\draw[draw=none,fill=darkorange25512714] (axis cs:275.5,0) rectangle (axis cs:276.5,100);
\draw[draw=none,fill=darkorange25512714] (axis cs:276.5,0) rectangle (axis cs:277.5,100);
\draw[draw=none,fill=darkorange25512714] (axis cs:277.5,0) rectangle (axis cs:278.5,100);
\draw[draw=none,fill=darkorange25512714] (axis cs:278.5,0) rectangle (axis cs:279.5,100);
\draw[draw=none,fill=darkorange25512714] (axis cs:279.5,0) rectangle (axis cs:280.5,100);
\draw[draw=none,fill=darkorange25512714] (axis cs:280.5,0) rectangle (axis cs:281.5,100);
\draw[draw=none,fill=darkorange25512714] (axis cs:281.5,0) rectangle (axis cs:282.5,100);
\draw[draw=none,fill=darkorange25512714] (axis cs:282.5,0) rectangle (axis cs:283.5,100);
\draw[draw=none,fill=darkorange25512714] (axis cs:283.5,0) rectangle (axis cs:284.5,100);
\draw[draw=none,fill=darkorange25512714] (axis cs:284.5,0) rectangle (axis cs:285.5,100);
\draw[draw=none,fill=darkorange25512714] (axis cs:285.5,0) rectangle (axis cs:286.5,100);
\draw[draw=none,fill=darkorange25512714] (axis cs:286.5,0) rectangle (axis cs:287.5,100);
\draw[draw=none,fill=darkorange25512714] (axis cs:287.5,0) rectangle (axis cs:288.5,100);
\draw[draw=none,fill=darkorange25512714] (axis cs:288.5,0) rectangle (axis cs:289.5,100);
\draw[draw=none,fill=darkorange25512714] (axis cs:289.5,0) rectangle (axis cs:290.5,100);
\draw[draw=none,fill=darkorange25512714] (axis cs:290.5,0) rectangle (axis cs:291.5,100);
\draw[draw=none,fill=darkorange25512714] (axis cs:291.5,0) rectangle (axis cs:292.5,100);
\draw[draw=none,fill=darkorange25512714] (axis cs:292.5,0) rectangle (axis cs:293.5,100);
\draw[draw=none,fill=darkorange25512714] (axis cs:293.5,0) rectangle (axis cs:294.5,100);
\draw[draw=none,fill=darkorange25512714] (axis cs:294.5,0) rectangle (axis cs:295.5,100);
\draw[draw=none,fill=darkorange25512714] (axis cs:295.5,0) rectangle (axis cs:296.5,100);
\draw[draw=none,fill=darkorange25512714] (axis cs:296.5,0) rectangle (axis cs:297.5,100);
\draw[draw=none,fill=darkorange25512714] (axis cs:297.5,0) rectangle (axis cs:298.5,100);
\draw[draw=none,fill=darkorange25512714] (axis cs:298.5,0) rectangle (axis cs:299.5,100);
\draw[draw=none,fill=darkorange25512714] (axis cs:299.5,0) rectangle (axis cs:300.5,100);
\draw[draw=none,fill=darkorange25512714] (axis cs:300.5,0) rectangle (axis cs:301.5,100);
\draw[draw=none,fill=darkorange25512714] (axis cs:301.5,0) rectangle (axis cs:302.5,100);
\draw[draw=none,fill=darkorange25512714] (axis cs:302.5,0) rectangle (axis cs:303.5,100);
\draw[draw=none,fill=darkorange25512714] (axis cs:303.5,0) rectangle (axis cs:304.5,100);
\draw[draw=none,fill=darkorange25512714] (axis cs:304.5,0) rectangle (axis cs:305.5,100);
\draw[draw=none,fill=darkorange25512714] (axis cs:305.5,0) rectangle (axis cs:306.5,100);
\draw[draw=none,fill=darkorange25512714] (axis cs:306.5,0) rectangle (axis cs:307.5,100);
\draw[draw=none,fill=darkorange25512714] (axis cs:307.5,0) rectangle (axis cs:308.5,100);
\draw[draw=none,fill=darkorange25512714] (axis cs:308.5,0) rectangle (axis cs:309.5,100);
\draw[draw=none,fill=darkorange25512714] (axis cs:309.5,0) rectangle (axis cs:310.5,100);
\draw[draw=none,fill=darkorange25512714] (axis cs:310.5,0) rectangle (axis cs:311.5,100);
\draw[draw=none,fill=darkorange25512714] (axis cs:311.5,0) rectangle (axis cs:312.5,100);
\draw[draw=none,fill=darkorange25512714] (axis cs:312.5,0) rectangle (axis cs:313.5,100);
\draw[draw=none,fill=darkorange25512714] (axis cs:313.5,0) rectangle (axis cs:314.5,100);
\draw[draw=none,fill=darkorange25512714] (axis cs:314.5,0) rectangle (axis cs:315.5,100);
\draw[draw=none,fill=darkorange25512714] (axis cs:315.5,0) rectangle (axis cs:316.5,100);
\draw[draw=none,fill=darkorange25512714] (axis cs:316.5,0) rectangle (axis cs:317.5,100);
\draw[draw=none,fill=darkorange25512714] (axis cs:317.5,0) rectangle (axis cs:318.5,100);
\draw[draw=none,fill=darkorange25512714] (axis cs:318.5,0) rectangle (axis cs:319.5,100);
\draw[draw=none,fill=darkorange25512714] (axis cs:319.5,0) rectangle (axis cs:320.5,100);
\draw[draw=none,fill=darkorange25512714] (axis cs:320.5,0) rectangle (axis cs:321.5,100);
\draw[draw=none,fill=darkorange25512714] (axis cs:321.5,0) rectangle (axis cs:322.5,100);
\draw[draw=none,fill=darkorange25512714] (axis cs:322.5,0) rectangle (axis cs:323.5,100);
\draw[draw=none,fill=darkorange25512714] (axis cs:323.5,0) rectangle (axis cs:324.5,100);
\draw[draw=none,fill=darkorange25512714] (axis cs:324.5,0) rectangle (axis cs:325.5,100);
\draw[draw=none,fill=darkorange25512714] (axis cs:325.5,0) rectangle (axis cs:326.5,100);
\draw[draw=none,fill=darkorange25512714] (axis cs:326.5,0) rectangle (axis cs:327.5,100);
\draw[draw=none,fill=darkorange25512714] (axis cs:327.5,0) rectangle (axis cs:328.5,100);
\draw[draw=none,fill=darkorange25512714] (axis cs:328.5,0) rectangle (axis cs:329.5,100);
\draw[draw=none,fill=darkorange25512714] (axis cs:329.5,0) rectangle (axis cs:330.5,100);
\draw[draw=none,fill=darkorange25512714] (axis cs:330.5,0) rectangle (axis cs:331.5,100);
\draw[draw=none,fill=darkorange25512714] (axis cs:331.5,0) rectangle (axis cs:332.5,100);
\draw[draw=none,fill=darkorange25512714] (axis cs:332.5,0) rectangle (axis cs:333.5,100);
\draw[draw=none,fill=darkorange25512714] (axis cs:333.5,0) rectangle (axis cs:334.5,100);
\draw[draw=none,fill=darkorange25512714] (axis cs:334.5,0) rectangle (axis cs:335.5,100);
\draw[draw=none,fill=darkorange25512714] (axis cs:335.5,0) rectangle (axis cs:336.5,100);
\draw[draw=none,fill=darkorange25512714] (axis cs:336.5,0) rectangle (axis cs:337.5,100);
\draw[draw=none,fill=darkorange25512714] (axis cs:337.5,0) rectangle (axis cs:338.5,100);
\draw[draw=none,fill=darkorange25512714] (axis cs:338.5,0) rectangle (axis cs:339.5,100);
\draw[draw=none,fill=darkorange25512714] (axis cs:339.5,0) rectangle (axis cs:340.5,100);
\draw[draw=none,fill=darkorange25512714] (axis cs:340.5,0) rectangle (axis cs:341.5,100);
\draw[draw=none,fill=darkorange25512714] (axis cs:341.5,0) rectangle (axis cs:342.5,100);
\draw[draw=none,fill=darkorange25512714] (axis cs:342.5,0) rectangle (axis cs:343.5,100);
\draw[draw=none,fill=darkorange25512714] (axis cs:343.5,0) rectangle (axis cs:344.5,100);
\draw[draw=none,fill=darkorange25512714] (axis cs:344.5,0) rectangle (axis cs:345.5,100);
\draw[draw=none,fill=darkorange25512714] (axis cs:345.5,0) rectangle (axis cs:346.5,100);
\draw[draw=none,fill=darkorange25512714] (axis cs:346.5,0) rectangle (axis cs:347.5,100);
\draw[draw=none,fill=darkorange25512714] (axis cs:347.5,0) rectangle (axis cs:348.5,100);
\draw[draw=none,fill=darkorange25512714] (axis cs:348.5,0) rectangle (axis cs:349.5,100);
\draw[draw=none,fill=darkorange25512714] (axis cs:349.5,0) rectangle (axis cs:350.5,100);
\draw[draw=none,fill=darkorange25512714] (axis cs:350.5,0) rectangle (axis cs:351.5,100);
\draw[draw=none,fill=darkorange25512714] (axis cs:351.5,0) rectangle (axis cs:352.5,100);
\draw[draw=none,fill=darkorange25512714] (axis cs:352.5,0) rectangle (axis cs:353.5,100);
\draw[draw=none,fill=darkorange25512714] (axis cs:353.5,0) rectangle (axis cs:354.5,100);
\draw[draw=none,fill=darkorange25512714] (axis cs:354.5,0) rectangle (axis cs:355.5,100);
\draw[draw=none,fill=darkorange25512714] (axis cs:355.5,0) rectangle (axis cs:356.5,100);
\draw[draw=none,fill=darkorange25512714] (axis cs:356.5,0) rectangle (axis cs:357.5,100);
\draw[draw=none,fill=darkorange25512714] (axis cs:357.5,0) rectangle (axis cs:358.5,100);
\draw[draw=none,fill=darkorange25512714] (axis cs:358.5,0) rectangle (axis cs:359.5,100);
\draw[draw=none,fill=darkorange25512714] (axis cs:359.5,0) rectangle (axis cs:360.5,100);
\draw[draw=none,fill=darkorange25512714] (axis cs:360.5,0) rectangle (axis cs:361.5,100);
\draw[draw=none,fill=darkorange25512714] (axis cs:361.5,0) rectangle (axis cs:362.5,100);
\draw[draw=none,fill=darkorange25512714] (axis cs:362.5,0) rectangle (axis cs:363.5,100);
\draw[draw=none,fill=darkorange25512714] (axis cs:363.5,0) rectangle (axis cs:364.5,100);
\draw[draw=none,fill=darkorange25512714] (axis cs:364.5,0) rectangle (axis cs:365.5,100);
\draw[draw=none,fill=darkorange25512714] (axis cs:365.5,0) rectangle (axis cs:366.5,100);
\draw[draw=none,fill=darkorange25512714] (axis cs:366.5,0) rectangle (axis cs:367.5,100);
\draw[draw=none,fill=darkorange25512714] (axis cs:367.5,0) rectangle (axis cs:368.5,100);
\draw[draw=none,fill=darkorange25512714] (axis cs:368.5,0) rectangle (axis cs:369.5,100);
\draw[draw=none,fill=darkorange25512714] (axis cs:369.5,0) rectangle (axis cs:370.5,100);
\draw[draw=none,fill=darkorange25512714] (axis cs:370.5,0) rectangle (axis cs:371.5,100);
\draw[draw=none,fill=darkorange25512714] (axis cs:371.5,0) rectangle (axis cs:372.5,100);
\draw[draw=none,fill=darkorange25512714] (axis cs:372.5,0) rectangle (axis cs:373.5,100);
\draw[draw=none,fill=darkorange25512714] (axis cs:373.5,0) rectangle (axis cs:374.5,100);
\draw[draw=none,fill=darkorange25512714] (axis cs:374.5,0) rectangle (axis cs:375.5,100);
\draw[draw=none,fill=darkorange25512714] (axis cs:375.5,0) rectangle (axis cs:376.5,100);
\draw[draw=none,fill=darkorange25512714] (axis cs:376.5,0) rectangle (axis cs:377.5,100);
\draw[draw=none,fill=darkorange25512714] (axis cs:377.5,0) rectangle (axis cs:378.5,100);
\draw[draw=none,fill=darkorange25512714] (axis cs:378.5,0) rectangle (axis cs:379.5,100);
\draw[draw=none,fill=darkorange25512714] (axis cs:379.5,0) rectangle (axis cs:380.5,100);
\draw[draw=none,fill=darkorange25512714] (axis cs:380.5,0) rectangle (axis cs:381.5,100);
\draw[draw=none,fill=darkorange25512714] (axis cs:381.5,0) rectangle (axis cs:382.5,100);
\draw[draw=none,fill=darkorange25512714] (axis cs:382.5,0) rectangle (axis cs:383.5,100);
\draw[draw=none,fill=darkorange25512714] (axis cs:383.5,0) rectangle (axis cs:384.5,100);
\draw[draw=none,fill=darkorange25512714] (axis cs:384.5,0) rectangle (axis cs:385.5,100);
\draw[draw=none,fill=darkorange25512714] (axis cs:385.5,0) rectangle (axis cs:386.5,100);
\draw[draw=none,fill=darkorange25512714] (axis cs:386.5,0) rectangle (axis cs:387.5,100);
\draw[draw=none,fill=darkorange25512714] (axis cs:387.5,0) rectangle (axis cs:388.5,100);
\draw[draw=none,fill=darkorange25512714] (axis cs:388.5,0) rectangle (axis cs:389.5,100);
\draw[draw=none,fill=darkorange25512714] (axis cs:389.5,0) rectangle (axis cs:390.5,100);
\draw[draw=none,fill=darkorange25512714] (axis cs:390.5,0) rectangle (axis cs:391.5,100);
\draw[draw=none,fill=darkorange25512714] (axis cs:391.5,0) rectangle (axis cs:392.5,100);
\draw[draw=none,fill=darkorange25512714] (axis cs:392.5,0) rectangle (axis cs:393.5,100);
\draw[draw=none,fill=darkorange25512714] (axis cs:393.5,0) rectangle (axis cs:394.5,100);
\draw[draw=none,fill=darkorange25512714] (axis cs:394.5,0) rectangle (axis cs:395.5,100);
\draw[draw=none,fill=darkorange25512714] (axis cs:395.5,0) rectangle (axis cs:396.5,100);
\draw[draw=none,fill=darkorange25512714] (axis cs:396.5,0) rectangle (axis cs:397.5,100);
\draw[draw=none,fill=darkorange25512714] (axis cs:397.5,0) rectangle (axis cs:398.5,100);
\draw[draw=none,fill=darkorange25512714] (axis cs:398.5,0) rectangle (axis cs:399.5,100);
\draw[draw=none,fill=darkorange25512714] (axis cs:399.5,0) rectangle (axis cs:400.5,100);
\draw[draw=none,fill=darkorange25512714] (axis cs:400.5,0) rectangle (axis cs:401.5,100);
\draw[draw=none,fill=darkorange25512714] (axis cs:401.5,0) rectangle (axis cs:402.5,100);
\draw[draw=none,fill=darkorange25512714] (axis cs:402.5,0) rectangle (axis cs:403.5,100);
\draw[draw=none,fill=darkorange25512714] (axis cs:403.5,0) rectangle (axis cs:404.5,100);
\draw[draw=none,fill=darkorange25512714] (axis cs:404.5,0) rectangle (axis cs:405.5,100);
\draw[draw=none,fill=darkorange25512714] (axis cs:405.5,0) rectangle (axis cs:406.5,100);
\draw[draw=none,fill=darkorange25512714] (axis cs:406.5,0) rectangle (axis cs:407.5,100);
\draw[draw=none,fill=darkorange25512714] (axis cs:407.5,0) rectangle (axis cs:408.5,100);
\draw[draw=none,fill=darkorange25512714] (axis cs:408.5,0) rectangle (axis cs:409.5,100);
\draw[draw=none,fill=darkorange25512714] (axis cs:409.5,0) rectangle (axis cs:410.5,100);
\draw[draw=none,fill=darkorange25512714] (axis cs:410.5,0) rectangle (axis cs:411.5,100);
\draw[draw=none,fill=darkorange25512714] (axis cs:411.5,0) rectangle (axis cs:412.5,100);
\draw[draw=none,fill=darkorange25512714] (axis cs:412.5,0) rectangle (axis cs:413.5,100);
\draw[draw=none,fill=darkorange25512714] (axis cs:413.5,0) rectangle (axis cs:414.5,100);
\draw[draw=none,fill=darkorange25512714] (axis cs:414.5,0) rectangle (axis cs:415.5,100);
\draw[draw=none,fill=darkorange25512714] (axis cs:415.5,0) rectangle (axis cs:416.5,100);
\draw[draw=none,fill=darkorange25512714] (axis cs:416.5,0) rectangle (axis cs:417.5,100);
\draw[draw=none,fill=darkorange25512714] (axis cs:417.5,0) rectangle (axis cs:418.5,100);
\draw[draw=none,fill=darkorange25512714] (axis cs:418.5,0) rectangle (axis cs:419.5,100);
\draw[draw=none,fill=darkorange25512714] (axis cs:419.5,0) rectangle (axis cs:420.5,100);
\draw[draw=none,fill=darkorange25512714] (axis cs:420.5,0) rectangle (axis cs:421.5,100);
\draw[draw=none,fill=darkorange25512714] (axis cs:421.5,0) rectangle (axis cs:422.5,100);
\draw[draw=none,fill=darkorange25512714] (axis cs:422.5,0) rectangle (axis cs:423.5,100);
\draw[draw=none,fill=darkorange25512714] (axis cs:423.5,0) rectangle (axis cs:424.5,100);
\draw[draw=none,fill=darkorange25512714] (axis cs:424.5,0) rectangle (axis cs:425.5,100);
\draw[draw=none,fill=darkorange25512714] (axis cs:425.5,0) rectangle (axis cs:426.5,100);
\draw[draw=none,fill=darkorange25512714] (axis cs:426.5,0) rectangle (axis cs:427.5,100);
\draw[draw=none,fill=darkorange25512714] (axis cs:427.5,0) rectangle (axis cs:428.5,100);
\draw[draw=none,fill=darkorange25512714] (axis cs:428.5,0) rectangle (axis cs:429.5,100);
\draw[draw=none,fill=darkorange25512714] (axis cs:429.5,0) rectangle (axis cs:430.5,100);
\draw[draw=none,fill=darkorange25512714] (axis cs:430.5,0) rectangle (axis cs:431.5,100);
\draw[draw=none,fill=darkorange25512714] (axis cs:431.5,0) rectangle (axis cs:432.5,100);
\draw[draw=none,fill=darkorange25512714] (axis cs:432.5,0) rectangle (axis cs:433.5,100);
\draw[draw=none,fill=darkorange25512714] (axis cs:433.5,0) rectangle (axis cs:434.5,100);
\draw[draw=none,fill=darkorange25512714] (axis cs:434.5,0) rectangle (axis cs:435.5,100);
\draw[draw=none,fill=darkorange25512714] (axis cs:435.5,0) rectangle (axis cs:436.5,100);
\draw[draw=none,fill=darkorange25512714] (axis cs:436.5,0) rectangle (axis cs:437.5,100);
\draw[draw=none,fill=darkorange25512714] (axis cs:437.5,0) rectangle (axis cs:438.5,100);
\draw[draw=none,fill=darkorange25512714] (axis cs:438.5,0) rectangle (axis cs:439.5,100);
\draw[draw=none,fill=darkorange25512714] (axis cs:439.5,0) rectangle (axis cs:440.5,100);
\draw[draw=none,fill=darkorange25512714] (axis cs:440.5,0) rectangle (axis cs:441.5,100);
\draw[draw=none,fill=darkorange25512714] (axis cs:441.5,0) rectangle (axis cs:442.5,100);
\draw[draw=none,fill=darkorange25512714] (axis cs:442.5,0) rectangle (axis cs:443.5,100);
\draw[draw=none,fill=darkorange25512714] (axis cs:443.5,0) rectangle (axis cs:444.5,100);
\draw[draw=none,fill=darkorange25512714] (axis cs:444.5,0) rectangle (axis cs:445.5,100);
\draw[draw=none,fill=darkorange25512714] (axis cs:445.5,0) rectangle (axis cs:446.5,100);
\draw[draw=none,fill=darkorange25512714] (axis cs:446.5,0) rectangle (axis cs:447.5,100);
\draw[draw=none,fill=darkorange25512714] (axis cs:447.5,0) rectangle (axis cs:448.5,100);
\draw[draw=none,fill=darkorange25512714] (axis cs:448.5,0) rectangle (axis cs:449.5,100);
\draw[draw=none,fill=darkorange25512714] (axis cs:449.5,0) rectangle (axis cs:450.5,100);
\draw[draw=none,fill=darkorange25512714] (axis cs:450.5,0) rectangle (axis cs:451.5,100);
\draw[draw=none,fill=darkorange25512714] (axis cs:451.5,0) rectangle (axis cs:452.5,100);
\draw[draw=none,fill=darkorange25512714] (axis cs:452.5,0) rectangle (axis cs:453.5,100);
\draw[draw=none,fill=darkorange25512714] (axis cs:453.5,0) rectangle (axis cs:454.5,100);
\draw[draw=none,fill=darkorange25512714] (axis cs:454.5,0) rectangle (axis cs:455.5,100);
\draw[draw=none,fill=darkorange25512714] (axis cs:455.5,0) rectangle (axis cs:456.5,100);
\draw[draw=none,fill=darkorange25512714] (axis cs:456.5,0) rectangle (axis cs:457.5,100);
\draw[draw=none,fill=darkorange25512714] (axis cs:457.5,0) rectangle (axis cs:458.5,100);
\draw[draw=none,fill=darkorange25512714] (axis cs:458.5,0) rectangle (axis cs:459.5,100);
\draw[draw=none,fill=darkorange25512714] (axis cs:459.5,0) rectangle (axis cs:460.5,100);
\draw[draw=none,fill=darkorange25512714] (axis cs:460.5,0) rectangle (axis cs:461.5,100);
\draw[draw=none,fill=darkorange25512714] (axis cs:461.5,0) rectangle (axis cs:462.5,100);
\draw[draw=none,fill=darkorange25512714] (axis cs:462.5,0) rectangle (axis cs:463.5,100);
\draw[draw=none,fill=darkorange25512714] (axis cs:463.5,0) rectangle (axis cs:464.5,100);
\draw[draw=none,fill=darkorange25512714] (axis cs:464.5,0) rectangle (axis cs:465.5,100);
\draw[draw=none,fill=darkorange25512714] (axis cs:465.5,0) rectangle (axis cs:466.5,100);
\draw[draw=none,fill=darkorange25512714] (axis cs:466.5,0) rectangle (axis cs:467.5,100);
\draw[draw=none,fill=darkorange25512714] (axis cs:467.5,0) rectangle (axis cs:468.5,100);
\draw[draw=none,fill=darkorange25512714] (axis cs:468.5,0) rectangle (axis cs:469.5,100);
\draw[draw=none,fill=darkorange25512714] (axis cs:469.5,0) rectangle (axis cs:470.5,100);
\draw[draw=none,fill=darkorange25512714] (axis cs:470.5,0) rectangle (axis cs:471.5,100);
\draw[draw=none,fill=darkorange25512714] (axis cs:471.5,0) rectangle (axis cs:472.5,100);
\draw[draw=none,fill=darkorange25512714] (axis cs:472.5,0) rectangle (axis cs:473.5,100);
\draw[draw=none,fill=darkorange25512714] (axis cs:473.5,0) rectangle (axis cs:474.5,100);
\draw[draw=none,fill=darkorange25512714] (axis cs:474.5,0) rectangle (axis cs:475.5,100);
\draw[draw=none,fill=darkorange25512714] (axis cs:475.5,0) rectangle (axis cs:476.5,100);
\draw[draw=none,fill=darkorange25512714] (axis cs:476.5,0) rectangle (axis cs:477.5,100);
\draw[draw=none,fill=darkorange25512714] (axis cs:477.5,0) rectangle (axis cs:478.5,100);
\draw[draw=none,fill=darkorange25512714] (axis cs:478.5,0) rectangle (axis cs:479.5,100);
\draw[draw=none,fill=darkorange25512714] (axis cs:479.5,0) rectangle (axis cs:480.5,100);
\draw[draw=none,fill=darkorange25512714] (axis cs:480.5,0) rectangle (axis cs:481.5,100);
\draw[draw=none,fill=darkorange25512714] (axis cs:481.5,0) rectangle (axis cs:482.5,100);
\draw[draw=none,fill=darkorange25512714] (axis cs:482.5,0) rectangle (axis cs:483.5,100);
\draw[draw=none,fill=darkorange25512714] (axis cs:483.5,0) rectangle (axis cs:484.5,100);
\draw[draw=none,fill=darkorange25512714] (axis cs:484.5,0) rectangle (axis cs:485.5,100);
\draw[draw=none,fill=darkorange25512714] (axis cs:485.5,0) rectangle (axis cs:486.5,100);
\draw[draw=none,fill=darkorange25512714] (axis cs:486.5,0) rectangle (axis cs:487.5,100);
\draw[draw=none,fill=darkorange25512714] (axis cs:487.5,0) rectangle (axis cs:488.5,100);
\draw[draw=none,fill=darkorange25512714] (axis cs:488.5,0) rectangle (axis cs:489.5,100);
\draw[draw=none,fill=darkorange25512714] (axis cs:489.5,0) rectangle (axis cs:490.5,100);
\draw[draw=none,fill=darkorange25512714] (axis cs:490.5,0) rectangle (axis cs:491.5,100);
\draw[draw=none,fill=darkorange25512714] (axis cs:491.5,0) rectangle (axis cs:492.5,100);
\draw[draw=none,fill=darkorange25512714] (axis cs:492.5,0) rectangle (axis cs:493.5,100);
\draw[draw=none,fill=darkorange25512714] (axis cs:493.5,0) rectangle (axis cs:494.5,100);
\draw[draw=none,fill=darkorange25512714] (axis cs:494.5,0) rectangle (axis cs:495.5,100);
\draw[draw=none,fill=darkorange25512714] (axis cs:495.5,0) rectangle (axis cs:496.5,100);
\draw[draw=none,fill=darkorange25512714] (axis cs:496.5,0) rectangle (axis cs:497.5,100);
\draw[draw=none,fill=darkorange25512714] (axis cs:497.5,0) rectangle (axis cs:498.5,100);
\draw[draw=none,fill=darkorange25512714] (axis cs:498.5,0) rectangle (axis cs:499.5,100);
\draw[draw=none,fill=darkorange25512714] (axis cs:499.5,0) rectangle (axis cs:500.5,100);
\draw[draw=none,fill=darkorange25512714] (axis cs:500.5,0) rectangle (axis cs:501.5,100);
\draw[draw=none,fill=darkorange25512714] (axis cs:501.5,0) rectangle (axis cs:502.5,100);
\draw[draw=none,fill=darkorange25512714] (axis cs:502.5,0) rectangle (axis cs:503.5,100);
\draw[draw=none,fill=darkorange25512714] (axis cs:503.5,0) rectangle (axis cs:504.5,100);
\draw[draw=none,fill=darkorange25512714] (axis cs:504.5,0) rectangle (axis cs:505.5,100);
\draw[draw=none,fill=darkorange25512714] (axis cs:505.5,0) rectangle (axis cs:506.5,100);
\draw[draw=none,fill=darkorange25512714] (axis cs:506.5,0) rectangle (axis cs:507.5,100);
\draw[draw=none,fill=darkorange25512714] (axis cs:507.5,0) rectangle (axis cs:508.5,100);
\draw[draw=none,fill=darkorange25512714] (axis cs:508.5,0) rectangle (axis cs:509.5,100);
\draw[draw=none,fill=darkorange25512714] (axis cs:509.5,0) rectangle (axis cs:510.5,100);
\draw[draw=none,fill=darkorange25512714] (axis cs:510.5,0) rectangle (axis cs:511.5,100);
\draw[draw=none,fill=darkorange25512714] (axis cs:511.5,0) rectangle (axis cs:512.5,100);
\draw[draw=none,fill=darkorange25512714] (axis cs:512.5,0) rectangle (axis cs:513.5,100);
\draw[draw=none,fill=darkorange25512714] (axis cs:513.5,0) rectangle (axis cs:514.5,100);
\draw[draw=none,fill=darkorange25512714] (axis cs:514.5,0) rectangle (axis cs:515.5,100);
\draw[draw=none,fill=darkorange25512714] (axis cs:515.5,0) rectangle (axis cs:516.5,100);
\draw[draw=none,fill=darkorange25512714] (axis cs:516.5,0) rectangle (axis cs:517.5,100);
\draw[draw=none,fill=darkorange25512714] (axis cs:517.5,0) rectangle (axis cs:518.5,100);
\draw[draw=none,fill=darkorange25512714] (axis cs:518.5,0) rectangle (axis cs:519.5,100);
\draw[draw=none,fill=darkorange25512714] (axis cs:519.5,0) rectangle (axis cs:520.5,100);
\draw[draw=none,fill=darkorange25512714] (axis cs:520.5,0) rectangle (axis cs:521.5,100);
\draw[draw=none,fill=darkorange25512714] (axis cs:521.5,0) rectangle (axis cs:522.5,100);
\draw[draw=none,fill=darkorange25512714] (axis cs:522.5,0) rectangle (axis cs:523.5,100);
\draw[draw=none,fill=darkorange25512714] (axis cs:523.5,0) rectangle (axis cs:524.5,100);
\draw[draw=none,fill=darkorange25512714] (axis cs:524.5,0) rectangle (axis cs:525.5,100);
\draw[draw=none,fill=darkorange25512714] (axis cs:525.5,0) rectangle (axis cs:526.5,100);
\draw[draw=none,fill=darkorange25512714] (axis cs:526.5,0) rectangle (axis cs:527.5,100);
\draw[draw=none,fill=darkorange25512714] (axis cs:527.5,0) rectangle (axis cs:528.5,100);
\draw[draw=none,fill=darkorange25512714] (axis cs:528.5,0) rectangle (axis cs:529.5,100);
\draw[draw=none,fill=darkorange25512714] (axis cs:529.5,0) rectangle (axis cs:530.5,100);
\draw[draw=none,fill=darkorange25512714] (axis cs:530.5,0) rectangle (axis cs:531.5,100);
\draw[draw=none,fill=darkorange25512714] (axis cs:531.5,0) rectangle (axis cs:532.5,100);
\draw[draw=none,fill=darkorange25512714] (axis cs:532.5,0) rectangle (axis cs:533.5,100);
\draw[draw=none,fill=darkorange25512714] (axis cs:533.5,0) rectangle (axis cs:534.5,100);
\draw[draw=none,fill=darkorange25512714] (axis cs:534.5,0) rectangle (axis cs:535.5,100);
\draw[draw=none,fill=darkorange25512714] (axis cs:535.5,0) rectangle (axis cs:536.5,100);
\draw[draw=none,fill=darkorange25512714] (axis cs:536.5,0) rectangle (axis cs:537.5,100);
\draw[draw=none,fill=darkorange25512714] (axis cs:537.5,0) rectangle (axis cs:538.5,100);
\draw[draw=none,fill=darkorange25512714] (axis cs:538.5,0) rectangle (axis cs:539.5,100);
\draw[draw=none,fill=darkorange25512714] (axis cs:539.5,0) rectangle (axis cs:540.5,100);
\draw[draw=none,fill=darkorange25512714] (axis cs:540.5,0) rectangle (axis cs:541.5,100);
\draw[draw=none,fill=darkorange25512714] (axis cs:541.5,0) rectangle (axis cs:542.5,100);
\draw[draw=none,fill=darkorange25512714] (axis cs:542.5,0) rectangle (axis cs:543.5,100);
\draw[draw=none,fill=darkorange25512714] (axis cs:543.5,0) rectangle (axis cs:544.5,100);
\draw[draw=none,fill=darkorange25512714] (axis cs:544.5,0) rectangle (axis cs:545.5,100);
\draw[draw=none,fill=darkorange25512714] (axis cs:545.5,0) rectangle (axis cs:546.5,100);
\draw[draw=none,fill=darkorange25512714] (axis cs:546.5,0) rectangle (axis cs:547.5,100);
\draw[draw=none,fill=darkorange25512714] (axis cs:547.5,0) rectangle (axis cs:548.5,100);
\draw[draw=none,fill=darkorange25512714] (axis cs:548.5,0) rectangle (axis cs:549.5,100);
\draw[draw=none,fill=darkorange25512714] (axis cs:549.5,0) rectangle (axis cs:550.5,100);
\draw[draw=none,fill=darkorange25512714] (axis cs:550.5,0) rectangle (axis cs:551.5,100);
\draw[draw=none,fill=darkorange25512714] (axis cs:551.5,0) rectangle (axis cs:552.5,100);
\draw[draw=none,fill=darkorange25512714] (axis cs:552.5,0) rectangle (axis cs:553.5,100);
\draw[draw=none,fill=darkorange25512714] (axis cs:553.5,0) rectangle (axis cs:554.5,100);
\draw[draw=none,fill=darkorange25512714] (axis cs:554.5,0) rectangle (axis cs:555.5,100);
\draw[draw=none,fill=darkorange25512714] (axis cs:555.5,0) rectangle (axis cs:556.5,100);
\draw[draw=none,fill=darkorange25512714] (axis cs:556.5,0) rectangle (axis cs:557.5,100);
\draw[draw=none,fill=darkorange25512714] (axis cs:557.5,0) rectangle (axis cs:558.5,100);
\draw[draw=none,fill=darkorange25512714] (axis cs:558.5,0) rectangle (axis cs:559.5,100);
\draw[draw=none,fill=darkorange25512714] (axis cs:559.5,0) rectangle (axis cs:560.5,100);
\draw[draw=none,fill=darkorange25512714] (axis cs:560.5,0) rectangle (axis cs:561.5,100);
\draw[draw=none,fill=darkorange25512714] (axis cs:561.5,0) rectangle (axis cs:562.5,100);
\draw[draw=none,fill=darkorange25512714] (axis cs:562.5,0) rectangle (axis cs:563.5,100);
\draw[draw=none,fill=darkorange25512714] (axis cs:563.5,0) rectangle (axis cs:564.5,100);
\draw[draw=none,fill=darkorange25512714] (axis cs:564.5,0) rectangle (axis cs:565.5,100);
\draw[draw=none,fill=darkorange25512714] (axis cs:565.5,0) rectangle (axis cs:566.5,100);
\draw[draw=none,fill=darkorange25512714] (axis cs:566.5,0) rectangle (axis cs:567.5,100);
\draw[draw=none,fill=darkorange25512714] (axis cs:567.5,0) rectangle (axis cs:568.5,100);
\draw[draw=none,fill=darkorange25512714] (axis cs:568.5,0) rectangle (axis cs:569.5,100);
\draw[draw=none,fill=darkorange25512714] (axis cs:569.5,0) rectangle (axis cs:570.5,100);
\draw[draw=none,fill=darkorange25512714] (axis cs:570.5,0) rectangle (axis cs:571.5,100);
\draw[draw=none,fill=darkorange25512714] (axis cs:571.5,0) rectangle (axis cs:572.5,100);
\draw[draw=none,fill=darkorange25512714] (axis cs:572.5,0) rectangle (axis cs:573.5,100);
\draw[draw=none,fill=darkorange25512714] (axis cs:573.5,0) rectangle (axis cs:574.5,100);
\draw[draw=none,fill=darkorange25512714] (axis cs:574.5,0) rectangle (axis cs:575.5,100);
\draw[draw=none,fill=darkorange25512714] (axis cs:575.5,0) rectangle (axis cs:576.5,100);
\draw[draw=none,fill=darkorange25512714] (axis cs:576.5,0) rectangle (axis cs:577.5,100);
\draw[draw=none,fill=darkorange25512714] (axis cs:577.5,0) rectangle (axis cs:578.5,100);
\draw[draw=none,fill=darkorange25512714] (axis cs:578.5,0) rectangle (axis cs:579.5,100);
\draw[draw=none,fill=darkorange25512714] (axis cs:579.5,0) rectangle (axis cs:580.5,100);
\draw[draw=none,fill=darkorange25512714] (axis cs:580.5,0) rectangle (axis cs:581.5,100);
\draw[draw=none,fill=darkorange25512714] (axis cs:581.5,0) rectangle (axis cs:582.5,100);
\draw[draw=none,fill=darkorange25512714] (axis cs:582.5,0) rectangle (axis cs:583.5,100);
\draw[draw=none,fill=darkorange25512714] (axis cs:583.5,0) rectangle (axis cs:584.5,100);
\draw[draw=none,fill=darkorange25512714] (axis cs:584.5,0) rectangle (axis cs:585.5,100);
\draw[draw=none,fill=darkorange25512714] (axis cs:585.5,0) rectangle (axis cs:586.5,100);
\draw[draw=none,fill=darkorange25512714] (axis cs:586.5,0) rectangle (axis cs:587.5,100);
\draw[draw=none,fill=darkorange25512714] (axis cs:587.5,0) rectangle (axis cs:588.5,100);
\draw[draw=none,fill=darkorange25512714] (axis cs:588.5,0) rectangle (axis cs:589.5,100);
\draw[draw=none,fill=darkorange25512714] (axis cs:589.5,0) rectangle (axis cs:590.5,100);
\draw[draw=none,fill=darkorange25512714] (axis cs:590.5,0) rectangle (axis cs:591.5,100);
\draw[draw=none,fill=darkorange25512714] (axis cs:591.5,0) rectangle (axis cs:592.5,100);
\draw[draw=none,fill=darkorange25512714] (axis cs:592.5,0) rectangle (axis cs:593.5,100);
\draw[draw=none,fill=darkorange25512714] (axis cs:593.5,0) rectangle (axis cs:594.5,100);
\draw[draw=none,fill=darkorange25512714] (axis cs:594.5,0) rectangle (axis cs:595.5,100);
\draw[draw=none,fill=darkorange25512714] (axis cs:595.5,0) rectangle (axis cs:596.5,100);
\draw[draw=none,fill=darkorange25512714] (axis cs:596.5,0) rectangle (axis cs:597.5,100);
\draw[draw=none,fill=darkorange25512714] (axis cs:597.5,0) rectangle (axis cs:598.5,100);
\draw[draw=none,fill=darkorange25512714] (axis cs:598.5,0) rectangle (axis cs:599.5,100);
\draw[draw=none,fill=darkorange25512714] (axis cs:599.5,0) rectangle (axis cs:600.5,100);
\draw[draw=none,fill=darkorange25512714] (axis cs:600.5,0) rectangle (axis cs:601.5,100);
\draw[draw=none,fill=darkorange25512714] (axis cs:601.5,0) rectangle (axis cs:602.5,100);
\draw[draw=none,fill=darkorange25512714] (axis cs:602.5,0) rectangle (axis cs:603.5,100);
\draw[draw=none,fill=darkorange25512714] (axis cs:603.5,0) rectangle (axis cs:604.5,100);
\draw[draw=none,fill=darkorange25512714] (axis cs:604.5,0) rectangle (axis cs:605.5,100);
\draw[draw=none,fill=darkorange25512714] (axis cs:605.5,0) rectangle (axis cs:606.5,100);
\draw[draw=none,fill=darkorange25512714] (axis cs:606.5,0) rectangle (axis cs:607.5,100);
\draw[draw=none,fill=darkorange25512714] (axis cs:607.5,0) rectangle (axis cs:608.5,100);
\draw[draw=none,fill=darkorange25512714] (axis cs:608.5,0) rectangle (axis cs:609.5,100);
\draw[draw=none,fill=darkorange25512714] (axis cs:609.5,0) rectangle (axis cs:610.5,100);
\draw[draw=none,fill=darkorange25512714] (axis cs:610.5,0) rectangle (axis cs:611.5,100);
\draw[draw=none,fill=darkorange25512714] (axis cs:611.5,0) rectangle (axis cs:612.5,100);
\draw[draw=none,fill=darkorange25512714] (axis cs:612.5,0) rectangle (axis cs:613.5,100);
\draw[draw=none,fill=darkorange25512714] (axis cs:613.5,0) rectangle (axis cs:614.5,100);
\draw[draw=none,fill=darkorange25512714] (axis cs:614.5,0) rectangle (axis cs:615.5,100);
\draw[draw=none,fill=darkorange25512714] (axis cs:615.5,0) rectangle (axis cs:616.5,100);
\draw[draw=none,fill=darkorange25512714] (axis cs:616.5,0) rectangle (axis cs:617.5,100);
\draw[draw=none,fill=darkorange25512714] (axis cs:617.5,0) rectangle (axis cs:618.5,100);
\draw[draw=none,fill=darkorange25512714] (axis cs:618.5,0) rectangle (axis cs:619.5,100);
\draw[draw=none,fill=darkorange25512714] (axis cs:619.5,0) rectangle (axis cs:620.5,100);
\draw[draw=none,fill=darkorange25512714] (axis cs:620.5,0) rectangle (axis cs:621.5,100);
\draw[draw=none,fill=darkorange25512714] (axis cs:621.5,0) rectangle (axis cs:622.5,100);
\draw[draw=none,fill=darkorange25512714] (axis cs:622.5,0) rectangle (axis cs:623.5,100);
\draw[draw=none,fill=darkorange25512714] (axis cs:623.5,0) rectangle (axis cs:624.5,100);
\draw[draw=none,fill=darkorange25512714] (axis cs:624.5,0) rectangle (axis cs:625.5,100);
\draw[draw=none,fill=darkorange25512714] (axis cs:625.5,0) rectangle (axis cs:626.5,100);
\draw[draw=none,fill=darkorange25512714] (axis cs:626.5,0) rectangle (axis cs:627.5,100);
\draw[draw=none,fill=darkorange25512714] (axis cs:627.5,0) rectangle (axis cs:628.5,100);
\draw[draw=none,fill=darkorange25512714] (axis cs:628.5,0) rectangle (axis cs:629.5,100);
\draw[draw=none,fill=darkorange25512714] (axis cs:629.5,0) rectangle (axis cs:630.5,100);
\draw[draw=none,fill=darkorange25512714] (axis cs:630.5,0) rectangle (axis cs:631.5,100);
\draw[draw=none,fill=darkorange25512714] (axis cs:631.5,0) rectangle (axis cs:632.5,100);
\draw[draw=none,fill=darkorange25512714] (axis cs:632.5,0) rectangle (axis cs:633.5,100);
\draw[draw=none,fill=darkorange25512714] (axis cs:633.5,0) rectangle (axis cs:634.5,100);
\draw[draw=none,fill=darkorange25512714] (axis cs:634.5,0) rectangle (axis cs:635.5,100);
\draw[draw=none,fill=darkorange25512714] (axis cs:635.5,0) rectangle (axis cs:636.5,100);
\draw[draw=none,fill=darkorange25512714] (axis cs:636.5,0) rectangle (axis cs:637.5,100);
\draw[draw=none,fill=darkorange25512714] (axis cs:637.5,0) rectangle (axis cs:638.5,100);
\draw[draw=none,fill=darkorange25512714] (axis cs:638.5,0) rectangle (axis cs:639.5,100);
\draw[draw=none,fill=darkorange25512714] (axis cs:639.5,0) rectangle (axis cs:640.5,100);
\draw[draw=none,fill=darkorange25512714] (axis cs:640.5,0) rectangle (axis cs:641.5,100);
\draw[draw=none,fill=darkorange25512714] (axis cs:641.5,0) rectangle (axis cs:642.5,100);
\draw[draw=none,fill=darkorange25512714] (axis cs:642.5,0) rectangle (axis cs:643.5,100);
\draw[draw=none,fill=darkorange25512714] (axis cs:643.5,0) rectangle (axis cs:644.5,100);
\draw[draw=none,fill=darkorange25512714] (axis cs:644.5,0) rectangle (axis cs:645.5,100);
\draw[draw=none,fill=darkorange25512714] (axis cs:645.5,0) rectangle (axis cs:646.5,100);
\draw[draw=none,fill=darkorange25512714] (axis cs:646.5,0) rectangle (axis cs:647.5,100);
\draw[draw=none,fill=darkorange25512714] (axis cs:647.5,0) rectangle (axis cs:648.5,100);
\draw[draw=none,fill=darkorange25512714] (axis cs:648.5,0) rectangle (axis cs:649.5,100);
\draw[draw=none,fill=darkorange25512714] (axis cs:649.5,0) rectangle (axis cs:650.5,100);
\draw[draw=none,fill=darkorange25512714] (axis cs:650.5,0) rectangle (axis cs:651.5,100);
\draw[draw=none,fill=darkorange25512714] (axis cs:651.5,0) rectangle (axis cs:652.5,100);
\draw[draw=none,fill=darkorange25512714] (axis cs:652.5,0) rectangle (axis cs:653.5,100);
\draw[draw=none,fill=darkorange25512714] (axis cs:653.5,0) rectangle (axis cs:654.5,100);
\draw[draw=none,fill=darkorange25512714] (axis cs:654.5,0) rectangle (axis cs:655.5,100);
\draw[draw=none,fill=darkorange25512714] (axis cs:655.5,0) rectangle (axis cs:656.5,100);
\draw[draw=none,fill=darkorange25512714] (axis cs:656.5,0) rectangle (axis cs:657.5,100);
\draw[draw=none,fill=darkorange25512714] (axis cs:657.5,0) rectangle (axis cs:658.5,100);
\draw[draw=none,fill=darkorange25512714] (axis cs:658.5,0) rectangle (axis cs:659.5,100);
\draw[draw=none,fill=darkorange25512714] (axis cs:659.5,0) rectangle (axis cs:660.5,100);
\draw[draw=none,fill=darkorange25512714] (axis cs:660.5,0) rectangle (axis cs:661.5,100);
\draw[draw=none,fill=darkorange25512714] (axis cs:661.5,0) rectangle (axis cs:662.5,100);
\draw[draw=none,fill=darkorange25512714] (axis cs:662.5,0) rectangle (axis cs:663.5,100);
\draw[draw=none,fill=darkorange25512714] (axis cs:663.5,0) rectangle (axis cs:664.5,100);
\draw[draw=none,fill=darkorange25512714] (axis cs:664.5,0) rectangle (axis cs:665.5,100);
\draw[draw=none,fill=darkorange25512714] (axis cs:665.5,0) rectangle (axis cs:666.5,100);
\draw[draw=none,fill=darkorange25512714] (axis cs:666.5,0) rectangle (axis cs:667.5,100);
\draw[draw=none,fill=darkorange25512714] (axis cs:667.5,0) rectangle (axis cs:668.5,100);
\draw[draw=none,fill=darkorange25512714] (axis cs:668.5,0) rectangle (axis cs:669.5,100);
\draw[draw=none,fill=darkorange25512714] (axis cs:669.5,0) rectangle (axis cs:670.5,100);
\draw[draw=none,fill=darkorange25512714] (axis cs:670.5,0) rectangle (axis cs:671.5,100);
\draw[draw=none,fill=darkorange25512714] (axis cs:671.5,0) rectangle (axis cs:672.5,100);
\draw[draw=none,fill=darkorange25512714] (axis cs:672.5,0) rectangle (axis cs:673.5,100);
\draw[draw=none,fill=darkorange25512714] (axis cs:673.5,0) rectangle (axis cs:674.5,100);
\draw[draw=none,fill=darkorange25512714] (axis cs:674.5,0) rectangle (axis cs:675.5,100);
\draw[draw=none,fill=darkorange25512714] (axis cs:675.5,0) rectangle (axis cs:676.5,100);
\draw[draw=none,fill=darkorange25512714] (axis cs:676.5,0) rectangle (axis cs:677.5,100);
\draw[draw=none,fill=darkorange25512714] (axis cs:677.5,0) rectangle (axis cs:678.5,100);
\draw[draw=none,fill=darkorange25512714] (axis cs:678.5,0) rectangle (axis cs:679.5,100);
\draw[draw=none,fill=darkorange25512714] (axis cs:679.5,0) rectangle (axis cs:680.5,100);
\draw[draw=none,fill=darkorange25512714] (axis cs:680.5,0) rectangle (axis cs:681.5,100);
\draw[draw=none,fill=darkorange25512714] (axis cs:681.5,0) rectangle (axis cs:682.5,100);
\draw[draw=none,fill=darkorange25512714] (axis cs:682.5,0) rectangle (axis cs:683.5,100);
\draw[draw=none,fill=darkorange25512714] (axis cs:683.5,0) rectangle (axis cs:684.5,100);
\draw[draw=none,fill=darkorange25512714] (axis cs:684.5,0) rectangle (axis cs:685.5,100);
\draw[draw=none,fill=darkorange25512714] (axis cs:685.5,0) rectangle (axis cs:686.5,100);
\draw[draw=none,fill=darkorange25512714] (axis cs:686.5,0) rectangle (axis cs:687.5,100);
\draw[draw=none,fill=darkorange25512714] (axis cs:687.5,0) rectangle (axis cs:688.5,100);
\draw[draw=none,fill=darkorange25512714] (axis cs:688.5,0) rectangle (axis cs:689.5,100);
\draw[draw=none,fill=darkorange25512714] (axis cs:689.5,0) rectangle (axis cs:690.5,100);
\draw[draw=none,fill=darkorange25512714] (axis cs:690.5,0) rectangle (axis cs:691.5,100);
\draw[draw=none,fill=darkorange25512714] (axis cs:691.5,0) rectangle (axis cs:692.5,100);
\draw[draw=none,fill=darkorange25512714] (axis cs:692.5,0) rectangle (axis cs:693.5,100);
\draw[draw=none,fill=darkorange25512714] (axis cs:693.5,0) rectangle (axis cs:694.5,100);
\draw[draw=none,fill=darkorange25512714] (axis cs:694.5,0) rectangle (axis cs:695.5,100);
\draw[draw=none,fill=darkorange25512714] (axis cs:695.5,0) rectangle (axis cs:696.5,100);
\draw[draw=none,fill=darkorange25512714] (axis cs:696.5,0) rectangle (axis cs:697.5,100);
\draw[draw=none,fill=darkorange25512714] (axis cs:697.5,0) rectangle (axis cs:698.5,100);
\draw[draw=none,fill=darkorange25512714] (axis cs:698.5,0) rectangle (axis cs:699.5,100);
\draw[draw=none,fill=darkorange25512714] (axis cs:699.5,0) rectangle (axis cs:700.5,100);
\draw[draw=none,fill=darkorange25512714] (axis cs:700.5,0) rectangle (axis cs:701.5,100);
\draw[draw=none,fill=darkorange25512714] (axis cs:701.5,0) rectangle (axis cs:702.5,100);
\draw[draw=none,fill=darkorange25512714] (axis cs:702.5,0) rectangle (axis cs:703.5,100);
\draw[draw=none,fill=darkorange25512714] (axis cs:703.5,0) rectangle (axis cs:704.5,100);
\draw[draw=none,fill=darkorange25512714] (axis cs:704.5,0) rectangle (axis cs:705.5,100);
\draw[draw=none,fill=darkorange25512714] (axis cs:705.5,0) rectangle (axis cs:706.5,100);
\draw[draw=none,fill=darkorange25512714] (axis cs:706.5,0) rectangle (axis cs:707.5,100);
\draw[draw=none,fill=darkorange25512714] (axis cs:707.5,0) rectangle (axis cs:708.5,100);
\draw[draw=none,fill=darkorange25512714] (axis cs:708.5,0) rectangle (axis cs:709.5,100);
\draw[draw=none,fill=darkorange25512714] (axis cs:709.5,0) rectangle (axis cs:710.5,100);
\draw[draw=none,fill=darkorange25512714] (axis cs:710.5,0) rectangle (axis cs:711.5,100);
\draw[draw=none,fill=darkorange25512714] (axis cs:711.5,0) rectangle (axis cs:712.5,100);
\draw[draw=none,fill=darkorange25512714] (axis cs:712.5,0) rectangle (axis cs:713.5,100);
\draw[draw=none,fill=darkorange25512714] (axis cs:713.5,0) rectangle (axis cs:714.5,100);
\draw[draw=none,fill=darkorange25512714] (axis cs:714.5,0) rectangle (axis cs:715.5,100);
\draw[draw=none,fill=darkorange25512714] (axis cs:715.5,0) rectangle (axis cs:716.5,100);
\draw[draw=none,fill=darkorange25512714] (axis cs:716.5,0) rectangle (axis cs:717.5,100);
\draw[draw=none,fill=darkorange25512714] (axis cs:717.5,0) rectangle (axis cs:718.5,100);
\draw[draw=none,fill=darkorange25512714] (axis cs:718.5,0) rectangle (axis cs:719.5,100);
\draw[draw=none,fill=darkorange25512714] (axis cs:719.5,0) rectangle (axis cs:720.5,100);
\draw[draw=none,fill=darkorange25512714] (axis cs:720.5,0) rectangle (axis cs:721.5,100);
\draw[draw=none,fill=darkorange25512714] (axis cs:721.5,0) rectangle (axis cs:722.5,100);
\draw[draw=none,fill=darkorange25512714] (axis cs:722.5,0) rectangle (axis cs:723.5,100);
\draw[draw=none,fill=darkorange25512714] (axis cs:723.5,0) rectangle (axis cs:724.5,100);
\draw[draw=none,fill=darkorange25512714] (axis cs:724.5,0) rectangle (axis cs:725.5,100);
\draw[draw=none,fill=darkorange25512714] (axis cs:725.5,0) rectangle (axis cs:726.5,100);
\draw[draw=none,fill=darkorange25512714] (axis cs:726.5,0) rectangle (axis cs:727.5,100);
\draw[draw=none,fill=darkorange25512714] (axis cs:727.5,0) rectangle (axis cs:728.5,100);
\draw[draw=none,fill=darkorange25512714] (axis cs:728.5,0) rectangle (axis cs:729.5,100);
\draw[draw=none,fill=darkorange25512714] (axis cs:729.5,0) rectangle (axis cs:730.5,100);
\draw[draw=none,fill=darkorange25512714] (axis cs:730.5,0) rectangle (axis cs:731.5,100);
\draw[draw=none,fill=darkorange25512714] (axis cs:731.5,0) rectangle (axis cs:732.5,100);
\draw[draw=none,fill=darkorange25512714] (axis cs:732.5,0) rectangle (axis cs:733.5,100);
\draw[draw=none,fill=darkorange25512714] (axis cs:733.5,0) rectangle (axis cs:734.5,100);
\draw[draw=none,fill=darkorange25512714] (axis cs:734.5,0) rectangle (axis cs:735.5,100);
\draw[draw=none,fill=darkorange25512714] (axis cs:735.5,0) rectangle (axis cs:736.5,100);
\draw[draw=none,fill=darkorange25512714] (axis cs:736.5,0) rectangle (axis cs:737.5,100);
\draw[draw=none,fill=darkorange25512714] (axis cs:737.5,0) rectangle (axis cs:738.5,100);
\draw[draw=none,fill=darkorange25512714] (axis cs:738.5,0) rectangle (axis cs:739.5,100);
\draw[draw=none,fill=darkorange25512714] (axis cs:739.5,0) rectangle (axis cs:740.5,100);
\draw[draw=none,fill=darkorange25512714] (axis cs:740.5,0) rectangle (axis cs:741.5,100);
\draw[draw=none,fill=darkorange25512714] (axis cs:741.5,0) rectangle (axis cs:742.5,100);
\draw[draw=none,fill=darkorange25512714] (axis cs:742.5,0) rectangle (axis cs:743.5,100);
\draw[draw=none,fill=darkorange25512714] (axis cs:743.5,0) rectangle (axis cs:744.5,100);
\draw[draw=none,fill=darkorange25512714] (axis cs:744.5,0) rectangle (axis cs:745.5,100);
\draw[draw=none,fill=darkorange25512714] (axis cs:745.5,0) rectangle (axis cs:746.5,100);
\draw[draw=none,fill=darkorange25512714] (axis cs:746.5,0) rectangle (axis cs:747.5,100);
\draw[draw=none,fill=darkorange25512714] (axis cs:747.5,0) rectangle (axis cs:748.5,100);
\draw[draw=none,fill=darkorange25512714] (axis cs:748.5,0) rectangle (axis cs:749.5,100);
\draw[draw=none,fill=darkorange25512714] (axis cs:749.5,0) rectangle (axis cs:750.5,100);
\draw[draw=none,fill=darkorange25512714] (axis cs:750.5,0) rectangle (axis cs:751.5,100);
\draw[draw=none,fill=darkorange25512714] (axis cs:751.5,0) rectangle (axis cs:752.5,100);
\draw[draw=none,fill=darkorange25512714] (axis cs:752.5,0) rectangle (axis cs:753.5,100);
\draw[draw=none,fill=darkorange25512714] (axis cs:753.5,0) rectangle (axis cs:754.5,100);
\draw[draw=none,fill=darkorange25512714] (axis cs:754.5,0) rectangle (axis cs:755.5,100);
\draw[draw=none,fill=darkorange25512714] (axis cs:755.5,0) rectangle (axis cs:756.5,100);
\draw[draw=none,fill=darkorange25512714] (axis cs:756.5,0) rectangle (axis cs:757.5,100);
\draw[draw=none,fill=darkorange25512714] (axis cs:757.5,0) rectangle (axis cs:758.5,100);
\draw[draw=none,fill=darkorange25512714] (axis cs:758.5,0) rectangle (axis cs:759.5,100);
\draw[draw=none,fill=darkorange25512714] (axis cs:759.5,0) rectangle (axis cs:760.5,100);
\draw[draw=none,fill=darkorange25512714] (axis cs:760.5,0) rectangle (axis cs:761.5,100);
\draw[draw=none,fill=darkorange25512714] (axis cs:761.5,0) rectangle (axis cs:762.5,100);
\draw[draw=none,fill=darkorange25512714] (axis cs:762.5,0) rectangle (axis cs:763.5,100);
\draw[draw=none,fill=darkorange25512714] (axis cs:763.5,0) rectangle (axis cs:764.5,100);
\draw[draw=none,fill=darkorange25512714] (axis cs:764.5,0) rectangle (axis cs:765.5,100);
\draw[draw=none,fill=darkorange25512714] (axis cs:765.5,0) rectangle (axis cs:766.5,100);
\draw[draw=none,fill=darkorange25512714] (axis cs:766.5,0) rectangle (axis cs:767.5,100);
\draw[draw=none,fill=darkorange25512714] (axis cs:767.5,0) rectangle (axis cs:768.5,100);
\draw[draw=none,fill=darkorange25512714] (axis cs:768.5,0) rectangle (axis cs:769.5,100);
\draw[draw=none,fill=darkorange25512714] (axis cs:769.5,0) rectangle (axis cs:770.5,100);
\draw[draw=none,fill=darkorange25512714] (axis cs:770.5,0) rectangle (axis cs:771.5,100);
\draw[draw=none,fill=darkorange25512714] (axis cs:771.5,0) rectangle (axis cs:772.5,100);
\draw[draw=none,fill=darkorange25512714] (axis cs:772.5,0) rectangle (axis cs:773.5,100);
\draw[draw=none,fill=darkorange25512714] (axis cs:773.5,0) rectangle (axis cs:774.5,100);
\draw[draw=none,fill=darkorange25512714] (axis cs:774.5,0) rectangle (axis cs:775.5,100);
\draw[draw=none,fill=darkorange25512714] (axis cs:775.5,0) rectangle (axis cs:776.5,100);
\draw[draw=none,fill=darkorange25512714] (axis cs:776.5,0) rectangle (axis cs:777.5,100);
\draw[draw=none,fill=darkorange25512714] (axis cs:777.5,0) rectangle (axis cs:778.5,100);
\draw[draw=none,fill=darkorange25512714] (axis cs:778.5,0) rectangle (axis cs:779.5,100);
\draw[draw=none,fill=darkorange25512714] (axis cs:779.5,0) rectangle (axis cs:780.5,100);
\draw[draw=none,fill=darkorange25512714] (axis cs:780.5,0) rectangle (axis cs:781.5,100);
\draw[draw=none,fill=darkorange25512714] (axis cs:781.5,0) rectangle (axis cs:782.5,100);
\draw[draw=none,fill=darkorange25512714] (axis cs:782.5,0) rectangle (axis cs:783.5,100);
\draw[draw=none,fill=darkorange25512714] (axis cs:783.5,0) rectangle (axis cs:784.5,100);
\draw[draw=none,fill=darkorange25512714] (axis cs:784.5,0) rectangle (axis cs:785.5,100);
\draw[draw=none,fill=darkorange25512714] (axis cs:785.5,0) rectangle (axis cs:786.5,100);
\draw[draw=none,fill=forestgreen4416044] (axis cs:-0.5,0) rectangle (axis cs:0.5,0);
\addlegendimage{ybar,ybar legend,draw=none,fill=forestgreen4416044}
\addlegendentry{$\Delta_L=3$}

\draw[draw=none,fill=forestgreen4416044] (axis cs:0.5,0) rectangle (axis cs:1.5,0);
\draw[draw=none,fill=forestgreen4416044] (axis cs:1.5,0) rectangle (axis cs:2.5,0);
\draw[draw=none,fill=forestgreen4416044] (axis cs:2.5,0) rectangle (axis cs:3.5,0);
\draw[draw=none,fill=forestgreen4416044] (axis cs:3.5,0) rectangle (axis cs:4.5,0);
\draw[draw=none,fill=forestgreen4416044] (axis cs:4.5,0) rectangle (axis cs:5.5,0);
\draw[draw=none,fill=forestgreen4416044] (axis cs:5.5,0) rectangle (axis cs:6.5,0);
\draw[draw=none,fill=forestgreen4416044] (axis cs:6.5,0) rectangle (axis cs:7.5,0);
\draw[draw=none,fill=forestgreen4416044] (axis cs:7.5,0) rectangle (axis cs:8.5,0);
\draw[draw=none,fill=forestgreen4416044] (axis cs:8.5,0) rectangle (axis cs:9.5,0);
\draw[draw=none,fill=forestgreen4416044] (axis cs:9.5,0) rectangle (axis cs:10.5,0);
\draw[draw=none,fill=forestgreen4416044] (axis cs:10.5,0) rectangle (axis cs:11.5,0);
\draw[draw=none,fill=forestgreen4416044] (axis cs:11.5,0) rectangle (axis cs:12.5,0);
\draw[draw=none,fill=forestgreen4416044] (axis cs:12.5,0) rectangle (axis cs:13.5,2);
\draw[draw=none,fill=forestgreen4416044] (axis cs:13.5,0) rectangle (axis cs:14.5,2);
\draw[draw=none,fill=forestgreen4416044] (axis cs:14.5,0) rectangle (axis cs:15.5,4);
\draw[draw=none,fill=forestgreen4416044] (axis cs:15.5,0) rectangle (axis cs:16.5,4);
\draw[draw=none,fill=forestgreen4416044] (axis cs:16.5,0) rectangle (axis cs:17.5,6);
\draw[draw=none,fill=forestgreen4416044] (axis cs:17.5,0) rectangle (axis cs:18.5,6);
\draw[draw=none,fill=forestgreen4416044] (axis cs:18.5,0) rectangle (axis cs:19.5,6);
\draw[draw=none,fill=forestgreen4416044] (axis cs:19.5,0) rectangle (axis cs:20.5,6);
\draw[draw=none,fill=forestgreen4416044] (axis cs:20.5,0) rectangle (axis cs:21.5,8);
\draw[draw=none,fill=forestgreen4416044] (axis cs:21.5,0) rectangle (axis cs:22.5,10);
\draw[draw=none,fill=forestgreen4416044] (axis cs:22.5,0) rectangle (axis cs:23.5,10);
\draw[draw=none,fill=forestgreen4416044] (axis cs:23.5,0) rectangle (axis cs:24.5,12);
\draw[draw=none,fill=forestgreen4416044] (axis cs:24.5,0) rectangle (axis cs:25.5,12);
\draw[draw=none,fill=forestgreen4416044] (axis cs:25.5,0) rectangle (axis cs:26.5,12);
\draw[draw=none,fill=forestgreen4416044] (axis cs:26.5,0) rectangle (axis cs:27.5,12);
\draw[draw=none,fill=forestgreen4416044] (axis cs:27.5,0) rectangle (axis cs:28.5,14);
\draw[draw=none,fill=forestgreen4416044] (axis cs:28.5,0) rectangle (axis cs:29.5,14);
\draw[draw=none,fill=forestgreen4416044] (axis cs:29.5,0) rectangle (axis cs:30.5,14);
\draw[draw=none,fill=forestgreen4416044] (axis cs:30.5,0) rectangle (axis cs:31.5,16);
\draw[draw=none,fill=forestgreen4416044] (axis cs:31.5,0) rectangle (axis cs:32.5,16);
\draw[draw=none,fill=forestgreen4416044] (axis cs:32.5,0) rectangle (axis cs:33.5,16);
\draw[draw=none,fill=forestgreen4416044] (axis cs:33.5,0) rectangle (axis cs:34.5,16);
\draw[draw=none,fill=forestgreen4416044] (axis cs:34.5,0) rectangle (axis cs:35.5,18);
\draw[draw=none,fill=forestgreen4416044] (axis cs:35.5,0) rectangle (axis cs:36.5,18);
\draw[draw=none,fill=forestgreen4416044] (axis cs:36.5,0) rectangle (axis cs:37.5,18);
\draw[draw=none,fill=forestgreen4416044] (axis cs:37.5,0) rectangle (axis cs:38.5,18);
\draw[draw=none,fill=forestgreen4416044] (axis cs:38.5,0) rectangle (axis cs:39.5,18);
\draw[draw=none,fill=forestgreen4416044] (axis cs:39.5,0) rectangle (axis cs:40.5,18);
\draw[draw=none,fill=forestgreen4416044] (axis cs:40.5,0) rectangle (axis cs:41.5,18);
\draw[draw=none,fill=forestgreen4416044] (axis cs:41.5,0) rectangle (axis cs:42.5,18);
\draw[draw=none,fill=forestgreen4416044] (axis cs:42.5,0) rectangle (axis cs:43.5,18);
\draw[draw=none,fill=forestgreen4416044] (axis cs:43.5,0) rectangle (axis cs:44.5,20);
\draw[draw=none,fill=forestgreen4416044] (axis cs:44.5,0) rectangle (axis cs:45.5,20);
\draw[draw=none,fill=forestgreen4416044] (axis cs:45.5,0) rectangle (axis cs:46.5,20);
\draw[draw=none,fill=forestgreen4416044] (axis cs:46.5,0) rectangle (axis cs:47.5,22);
\draw[draw=none,fill=forestgreen4416044] (axis cs:47.5,0) rectangle (axis cs:48.5,22);
\draw[draw=none,fill=forestgreen4416044] (axis cs:48.5,0) rectangle (axis cs:49.5,22);
\draw[draw=none,fill=forestgreen4416044] (axis cs:49.5,0) rectangle (axis cs:50.5,22);
\draw[draw=none,fill=forestgreen4416044] (axis cs:50.5,0) rectangle (axis cs:51.5,22);
\draw[draw=none,fill=forestgreen4416044] (axis cs:51.5,0) rectangle (axis cs:52.5,24);
\draw[draw=none,fill=forestgreen4416044] (axis cs:52.5,0) rectangle (axis cs:53.5,26);
\draw[draw=none,fill=forestgreen4416044] (axis cs:53.5,0) rectangle (axis cs:54.5,26);
\draw[draw=none,fill=forestgreen4416044] (axis cs:54.5,0) rectangle (axis cs:55.5,26);
\draw[draw=none,fill=forestgreen4416044] (axis cs:55.5,0) rectangle (axis cs:56.5,26);
\draw[draw=none,fill=forestgreen4416044] (axis cs:56.5,0) rectangle (axis cs:57.5,26);
\draw[draw=none,fill=forestgreen4416044] (axis cs:57.5,0) rectangle (axis cs:58.5,26);
\draw[draw=none,fill=forestgreen4416044] (axis cs:58.5,0) rectangle (axis cs:59.5,26);
\draw[draw=none,fill=forestgreen4416044] (axis cs:59.5,0) rectangle (axis cs:60.5,26);
\draw[draw=none,fill=forestgreen4416044] (axis cs:60.5,0) rectangle (axis cs:61.5,26);
\draw[draw=none,fill=forestgreen4416044] (axis cs:61.5,0) rectangle (axis cs:62.5,26);
\draw[draw=none,fill=forestgreen4416044] (axis cs:62.5,0) rectangle (axis cs:63.5,26);
\draw[draw=none,fill=forestgreen4416044] (axis cs:63.5,0) rectangle (axis cs:64.5,26);
\draw[draw=none,fill=forestgreen4416044] (axis cs:64.5,0) rectangle (axis cs:65.5,26);
\draw[draw=none,fill=forestgreen4416044] (axis cs:65.5,0) rectangle (axis cs:66.5,26);
\draw[draw=none,fill=forestgreen4416044] (axis cs:66.5,0) rectangle (axis cs:67.5,26);
\draw[draw=none,fill=forestgreen4416044] (axis cs:67.5,0) rectangle (axis cs:68.5,26);
\draw[draw=none,fill=forestgreen4416044] (axis cs:68.5,0) rectangle (axis cs:69.5,26);
\draw[draw=none,fill=forestgreen4416044] (axis cs:69.5,0) rectangle (axis cs:70.5,26);
\draw[draw=none,fill=forestgreen4416044] (axis cs:70.5,0) rectangle (axis cs:71.5,26);
\draw[draw=none,fill=forestgreen4416044] (axis cs:71.5,0) rectangle (axis cs:72.5,26);
\draw[draw=none,fill=forestgreen4416044] (axis cs:72.5,0) rectangle (axis cs:73.5,26);
\draw[draw=none,fill=forestgreen4416044] (axis cs:73.5,0) rectangle (axis cs:74.5,26);
\draw[draw=none,fill=forestgreen4416044] (axis cs:74.5,0) rectangle (axis cs:75.5,26);
\draw[draw=none,fill=forestgreen4416044] (axis cs:75.5,0) rectangle (axis cs:76.5,28);
\draw[draw=none,fill=forestgreen4416044] (axis cs:76.5,0) rectangle (axis cs:77.5,28);
\draw[draw=none,fill=forestgreen4416044] (axis cs:77.5,0) rectangle (axis cs:78.5,28);
\draw[draw=none,fill=forestgreen4416044] (axis cs:78.5,0) rectangle (axis cs:79.5,30);
\draw[draw=none,fill=forestgreen4416044] (axis cs:79.5,0) rectangle (axis cs:80.5,30);
\draw[draw=none,fill=forestgreen4416044] (axis cs:80.5,0) rectangle (axis cs:81.5,30);
\draw[draw=none,fill=forestgreen4416044] (axis cs:81.5,0) rectangle (axis cs:82.5,30);
\draw[draw=none,fill=forestgreen4416044] (axis cs:82.5,0) rectangle (axis cs:83.5,30);
\draw[draw=none,fill=forestgreen4416044] (axis cs:83.5,0) rectangle (axis cs:84.5,30);
\draw[draw=none,fill=forestgreen4416044] (axis cs:84.5,0) rectangle (axis cs:85.5,30);
\draw[draw=none,fill=forestgreen4416044] (axis cs:85.5,0) rectangle (axis cs:86.5,32);
\draw[draw=none,fill=forestgreen4416044] (axis cs:86.5,0) rectangle (axis cs:87.5,32);
\draw[draw=none,fill=forestgreen4416044] (axis cs:87.5,0) rectangle (axis cs:88.5,32);
\draw[draw=none,fill=forestgreen4416044] (axis cs:88.5,0) rectangle (axis cs:89.5,32);
\draw[draw=none,fill=forestgreen4416044] (axis cs:89.5,0) rectangle (axis cs:90.5,32);
\draw[draw=none,fill=forestgreen4416044] (axis cs:90.5,0) rectangle (axis cs:91.5,32);
\draw[draw=none,fill=forestgreen4416044] (axis cs:91.5,0) rectangle (axis cs:92.5,36);
\draw[draw=none,fill=forestgreen4416044] (axis cs:92.5,0) rectangle (axis cs:93.5,36);
\draw[draw=none,fill=forestgreen4416044] (axis cs:93.5,0) rectangle (axis cs:94.5,36);
\draw[draw=none,fill=forestgreen4416044] (axis cs:94.5,0) rectangle (axis cs:95.5,36);
\draw[draw=none,fill=forestgreen4416044] (axis cs:95.5,0) rectangle (axis cs:96.5,36);
\draw[draw=none,fill=forestgreen4416044] (axis cs:96.5,0) rectangle (axis cs:97.5,36);
\draw[draw=none,fill=forestgreen4416044] (axis cs:97.5,0) rectangle (axis cs:98.5,36);
\draw[draw=none,fill=forestgreen4416044] (axis cs:98.5,0) rectangle (axis cs:99.5,36);
\draw[draw=none,fill=forestgreen4416044] (axis cs:99.5,0) rectangle (axis cs:100.5,36);
\draw[draw=none,fill=forestgreen4416044] (axis cs:100.5,0) rectangle (axis cs:101.5,36);
\draw[draw=none,fill=forestgreen4416044] (axis cs:101.5,0) rectangle (axis cs:102.5,36);
\draw[draw=none,fill=forestgreen4416044] (axis cs:102.5,0) rectangle (axis cs:103.5,36);
\draw[draw=none,fill=forestgreen4416044] (axis cs:103.5,0) rectangle (axis cs:104.5,36);
\draw[draw=none,fill=forestgreen4416044] (axis cs:104.5,0) rectangle (axis cs:105.5,36);
\draw[draw=none,fill=forestgreen4416044] (axis cs:105.5,0) rectangle (axis cs:106.5,36);
\draw[draw=none,fill=forestgreen4416044] (axis cs:106.5,0) rectangle (axis cs:107.5,36);
\draw[draw=none,fill=forestgreen4416044] (axis cs:107.5,0) rectangle (axis cs:108.5,36);
\draw[draw=none,fill=forestgreen4416044] (axis cs:108.5,0) rectangle (axis cs:109.5,36);
\draw[draw=none,fill=forestgreen4416044] (axis cs:109.5,0) rectangle (axis cs:110.5,36);
\draw[draw=none,fill=forestgreen4416044] (axis cs:110.5,0) rectangle (axis cs:111.5,36);
\draw[draw=none,fill=forestgreen4416044] (axis cs:111.5,0) rectangle (axis cs:112.5,36);
\draw[draw=none,fill=forestgreen4416044] (axis cs:112.5,0) rectangle (axis cs:113.5,36);
\draw[draw=none,fill=forestgreen4416044] (axis cs:113.5,0) rectangle (axis cs:114.5,36);
\draw[draw=none,fill=forestgreen4416044] (axis cs:114.5,0) rectangle (axis cs:115.5,36);
\draw[draw=none,fill=forestgreen4416044] (axis cs:115.5,0) rectangle (axis cs:116.5,36);
\draw[draw=none,fill=forestgreen4416044] (axis cs:116.5,0) rectangle (axis cs:117.5,38);
\draw[draw=none,fill=forestgreen4416044] (axis cs:117.5,0) rectangle (axis cs:118.5,38);
\draw[draw=none,fill=forestgreen4416044] (axis cs:118.5,0) rectangle (axis cs:119.5,38);
\draw[draw=none,fill=forestgreen4416044] (axis cs:119.5,0) rectangle (axis cs:120.5,38);
\draw[draw=none,fill=forestgreen4416044] (axis cs:120.5,0) rectangle (axis cs:121.5,38);
\draw[draw=none,fill=forestgreen4416044] (axis cs:121.5,0) rectangle (axis cs:122.5,38);
\draw[draw=none,fill=forestgreen4416044] (axis cs:122.5,0) rectangle (axis cs:123.5,38);
\draw[draw=none,fill=forestgreen4416044] (axis cs:123.5,0) rectangle (axis cs:124.5,38);
\draw[draw=none,fill=forestgreen4416044] (axis cs:124.5,0) rectangle (axis cs:125.5,38);
\draw[draw=none,fill=forestgreen4416044] (axis cs:125.5,0) rectangle (axis cs:126.5,38);
\draw[draw=none,fill=forestgreen4416044] (axis cs:126.5,0) rectangle (axis cs:127.5,38);
\draw[draw=none,fill=forestgreen4416044] (axis cs:127.5,0) rectangle (axis cs:128.5,38);
\draw[draw=none,fill=forestgreen4416044] (axis cs:128.5,0) rectangle (axis cs:129.5,38);
\draw[draw=none,fill=forestgreen4416044] (axis cs:129.5,0) rectangle (axis cs:130.5,38);
\draw[draw=none,fill=forestgreen4416044] (axis cs:130.5,0) rectangle (axis cs:131.5,40);
\draw[draw=none,fill=forestgreen4416044] (axis cs:131.5,0) rectangle (axis cs:132.5,40);
\draw[draw=none,fill=forestgreen4416044] (axis cs:132.5,0) rectangle (axis cs:133.5,40);
\draw[draw=none,fill=forestgreen4416044] (axis cs:133.5,0) rectangle (axis cs:134.5,40);
\draw[draw=none,fill=forestgreen4416044] (axis cs:134.5,0) rectangle (axis cs:135.5,40);
\draw[draw=none,fill=forestgreen4416044] (axis cs:135.5,0) rectangle (axis cs:136.5,40);
\draw[draw=none,fill=forestgreen4416044] (axis cs:136.5,0) rectangle (axis cs:137.5,40);
\draw[draw=none,fill=forestgreen4416044] (axis cs:137.5,0) rectangle (axis cs:138.5,40);
\draw[draw=none,fill=forestgreen4416044] (axis cs:138.5,0) rectangle (axis cs:139.5,42);
\draw[draw=none,fill=forestgreen4416044] (axis cs:139.5,0) rectangle (axis cs:140.5,42);
\draw[draw=none,fill=forestgreen4416044] (axis cs:140.5,0) rectangle (axis cs:141.5,42);
\draw[draw=none,fill=forestgreen4416044] (axis cs:141.5,0) rectangle (axis cs:142.5,42);
\draw[draw=none,fill=forestgreen4416044] (axis cs:142.5,0) rectangle (axis cs:143.5,42);
\draw[draw=none,fill=forestgreen4416044] (axis cs:143.5,0) rectangle (axis cs:144.5,42);
\draw[draw=none,fill=forestgreen4416044] (axis cs:144.5,0) rectangle (axis cs:145.5,42);
\draw[draw=none,fill=forestgreen4416044] (axis cs:145.5,0) rectangle (axis cs:146.5,44);
\draw[draw=none,fill=forestgreen4416044] (axis cs:146.5,0) rectangle (axis cs:147.5,44);
\draw[draw=none,fill=forestgreen4416044] (axis cs:147.5,0) rectangle (axis cs:148.5,44);
\draw[draw=none,fill=forestgreen4416044] (axis cs:148.5,0) rectangle (axis cs:149.5,44);
\draw[draw=none,fill=forestgreen4416044] (axis cs:149.5,0) rectangle (axis cs:150.5,44);
\draw[draw=none,fill=forestgreen4416044] (axis cs:150.5,0) rectangle (axis cs:151.5,44);
\draw[draw=none,fill=forestgreen4416044] (axis cs:151.5,0) rectangle (axis cs:152.5,44);
\draw[draw=none,fill=forestgreen4416044] (axis cs:152.5,0) rectangle (axis cs:153.5,44);
\draw[draw=none,fill=forestgreen4416044] (axis cs:153.5,0) rectangle (axis cs:154.5,44);
\draw[draw=none,fill=forestgreen4416044] (axis cs:154.5,0) rectangle (axis cs:155.5,44);
\draw[draw=none,fill=forestgreen4416044] (axis cs:155.5,0) rectangle (axis cs:156.5,44);
\draw[draw=none,fill=forestgreen4416044] (axis cs:156.5,0) rectangle (axis cs:157.5,44);
\draw[draw=none,fill=forestgreen4416044] (axis cs:157.5,0) rectangle (axis cs:158.5,44);
\draw[draw=none,fill=forestgreen4416044] (axis cs:158.5,0) rectangle (axis cs:159.5,44);
\draw[draw=none,fill=forestgreen4416044] (axis cs:159.5,0) rectangle (axis cs:160.5,44);
\draw[draw=none,fill=forestgreen4416044] (axis cs:160.5,0) rectangle (axis cs:161.5,44);
\draw[draw=none,fill=forestgreen4416044] (axis cs:161.5,0) rectangle (axis cs:162.5,44);
\draw[draw=none,fill=forestgreen4416044] (axis cs:162.5,0) rectangle (axis cs:163.5,44);
\draw[draw=none,fill=forestgreen4416044] (axis cs:163.5,0) rectangle (axis cs:164.5,44);
\draw[draw=none,fill=forestgreen4416044] (axis cs:164.5,0) rectangle (axis cs:165.5,44);
\draw[draw=none,fill=forestgreen4416044] (axis cs:165.5,0) rectangle (axis cs:166.5,44);
\draw[draw=none,fill=forestgreen4416044] (axis cs:166.5,0) rectangle (axis cs:167.5,44);
\draw[draw=none,fill=forestgreen4416044] (axis cs:167.5,0) rectangle (axis cs:168.5,44);
\draw[draw=none,fill=forestgreen4416044] (axis cs:168.5,0) rectangle (axis cs:169.5,44);
\draw[draw=none,fill=forestgreen4416044] (axis cs:169.5,0) rectangle (axis cs:170.5,44);
\draw[draw=none,fill=forestgreen4416044] (axis cs:170.5,0) rectangle (axis cs:171.5,46);
\draw[draw=none,fill=forestgreen4416044] (axis cs:171.5,0) rectangle (axis cs:172.5,46);
\draw[draw=none,fill=forestgreen4416044] (axis cs:172.5,0) rectangle (axis cs:173.5,46);
\draw[draw=none,fill=forestgreen4416044] (axis cs:173.5,0) rectangle (axis cs:174.5,46);
\draw[draw=none,fill=forestgreen4416044] (axis cs:174.5,0) rectangle (axis cs:175.5,46);
\draw[draw=none,fill=forestgreen4416044] (axis cs:175.5,0) rectangle (axis cs:176.5,46);
\draw[draw=none,fill=forestgreen4416044] (axis cs:176.5,0) rectangle (axis cs:177.5,46);
\draw[draw=none,fill=forestgreen4416044] (axis cs:177.5,0) rectangle (axis cs:178.5,46);
\draw[draw=none,fill=forestgreen4416044] (axis cs:178.5,0) rectangle (axis cs:179.5,46);
\draw[draw=none,fill=forestgreen4416044] (axis cs:179.5,0) rectangle (axis cs:180.5,46);
\draw[draw=none,fill=forestgreen4416044] (axis cs:180.5,0) rectangle (axis cs:181.5,46);
\draw[draw=none,fill=forestgreen4416044] (axis cs:181.5,0) rectangle (axis cs:182.5,46);
\draw[draw=none,fill=forestgreen4416044] (axis cs:182.5,0) rectangle (axis cs:183.5,46);
\draw[draw=none,fill=forestgreen4416044] (axis cs:183.5,0) rectangle (axis cs:184.5,46);
\draw[draw=none,fill=forestgreen4416044] (axis cs:184.5,0) rectangle (axis cs:185.5,46);
\draw[draw=none,fill=forestgreen4416044] (axis cs:185.5,0) rectangle (axis cs:186.5,46);
\draw[draw=none,fill=forestgreen4416044] (axis cs:186.5,0) rectangle (axis cs:187.5,46);
\draw[draw=none,fill=forestgreen4416044] (axis cs:187.5,0) rectangle (axis cs:188.5,46);
\draw[draw=none,fill=forestgreen4416044] (axis cs:188.5,0) rectangle (axis cs:189.5,46);
\draw[draw=none,fill=forestgreen4416044] (axis cs:189.5,0) rectangle (axis cs:190.5,46);
\draw[draw=none,fill=forestgreen4416044] (axis cs:190.5,0) rectangle (axis cs:191.5,46);
\draw[draw=none,fill=forestgreen4416044] (axis cs:191.5,0) rectangle (axis cs:192.5,48);
\draw[draw=none,fill=forestgreen4416044] (axis cs:192.5,0) rectangle (axis cs:193.5,48);
\draw[draw=none,fill=forestgreen4416044] (axis cs:193.5,0) rectangle (axis cs:194.5,48);
\draw[draw=none,fill=forestgreen4416044] (axis cs:194.5,0) rectangle (axis cs:195.5,48);
\draw[draw=none,fill=forestgreen4416044] (axis cs:195.5,0) rectangle (axis cs:196.5,48);
\draw[draw=none,fill=forestgreen4416044] (axis cs:196.5,0) rectangle (axis cs:197.5,48);
\draw[draw=none,fill=forestgreen4416044] (axis cs:197.5,0) rectangle (axis cs:198.5,50);
\draw[draw=none,fill=forestgreen4416044] (axis cs:198.5,0) rectangle (axis cs:199.5,50);
\draw[draw=none,fill=forestgreen4416044] (axis cs:199.5,0) rectangle (axis cs:200.5,50);
\draw[draw=none,fill=forestgreen4416044] (axis cs:200.5,0) rectangle (axis cs:201.5,50);
\draw[draw=none,fill=forestgreen4416044] (axis cs:201.5,0) rectangle (axis cs:202.5,50);
\draw[draw=none,fill=forestgreen4416044] (axis cs:202.5,0) rectangle (axis cs:203.5,50);
\draw[draw=none,fill=forestgreen4416044] (axis cs:203.5,0) rectangle (axis cs:204.5,50);
\draw[draw=none,fill=forestgreen4416044] (axis cs:204.5,0) rectangle (axis cs:205.5,50);
\draw[draw=none,fill=forestgreen4416044] (axis cs:205.5,0) rectangle (axis cs:206.5,50);
\draw[draw=none,fill=forestgreen4416044] (axis cs:206.5,0) rectangle (axis cs:207.5,50);
\draw[draw=none,fill=forestgreen4416044] (axis cs:207.5,0) rectangle (axis cs:208.5,50);
\draw[draw=none,fill=forestgreen4416044] (axis cs:208.5,0) rectangle (axis cs:209.5,50);
\draw[draw=none,fill=forestgreen4416044] (axis cs:209.5,0) rectangle (axis cs:210.5,50);
\draw[draw=none,fill=forestgreen4416044] (axis cs:210.5,0) rectangle (axis cs:211.5,50);
\draw[draw=none,fill=forestgreen4416044] (axis cs:211.5,0) rectangle (axis cs:212.5,52);
\draw[draw=none,fill=forestgreen4416044] (axis cs:212.5,0) rectangle (axis cs:213.5,52);
\draw[draw=none,fill=forestgreen4416044] (axis cs:213.5,0) rectangle (axis cs:214.5,52);
\draw[draw=none,fill=forestgreen4416044] (axis cs:214.5,0) rectangle (axis cs:215.5,56);
\draw[draw=none,fill=forestgreen4416044] (axis cs:215.5,0) rectangle (axis cs:216.5,56);
\draw[draw=none,fill=forestgreen4416044] (axis cs:216.5,0) rectangle (axis cs:217.5,56);
\draw[draw=none,fill=forestgreen4416044] (axis cs:217.5,0) rectangle (axis cs:218.5,56);
\draw[draw=none,fill=forestgreen4416044] (axis cs:218.5,0) rectangle (axis cs:219.5,56);
\draw[draw=none,fill=forestgreen4416044] (axis cs:219.5,0) rectangle (axis cs:220.5,56);
\draw[draw=none,fill=forestgreen4416044] (axis cs:220.5,0) rectangle (axis cs:221.5,56);
\draw[draw=none,fill=forestgreen4416044] (axis cs:221.5,0) rectangle (axis cs:222.5,56);
\draw[draw=none,fill=forestgreen4416044] (axis cs:222.5,0) rectangle (axis cs:223.5,56);
\draw[draw=none,fill=forestgreen4416044] (axis cs:223.5,0) rectangle (axis cs:224.5,56);
\draw[draw=none,fill=forestgreen4416044] (axis cs:224.5,0) rectangle (axis cs:225.5,58);
\draw[draw=none,fill=forestgreen4416044] (axis cs:225.5,0) rectangle (axis cs:226.5,58);
\draw[draw=none,fill=forestgreen4416044] (axis cs:226.5,0) rectangle (axis cs:227.5,58);
\draw[draw=none,fill=forestgreen4416044] (axis cs:227.5,0) rectangle (axis cs:228.5,58);
\draw[draw=none,fill=forestgreen4416044] (axis cs:228.5,0) rectangle (axis cs:229.5,58);
\draw[draw=none,fill=forestgreen4416044] (axis cs:229.5,0) rectangle (axis cs:230.5,58);
\draw[draw=none,fill=forestgreen4416044] (axis cs:230.5,0) rectangle (axis cs:231.5,58);
\draw[draw=none,fill=forestgreen4416044] (axis cs:231.5,0) rectangle (axis cs:232.5,58);
\draw[draw=none,fill=forestgreen4416044] (axis cs:232.5,0) rectangle (axis cs:233.5,58);
\draw[draw=none,fill=forestgreen4416044] (axis cs:233.5,0) rectangle (axis cs:234.5,58);
\draw[draw=none,fill=forestgreen4416044] (axis cs:234.5,0) rectangle (axis cs:235.5,58);
\draw[draw=none,fill=forestgreen4416044] (axis cs:235.5,0) rectangle (axis cs:236.5,58);
\draw[draw=none,fill=forestgreen4416044] (axis cs:236.5,0) rectangle (axis cs:237.5,60);
\draw[draw=none,fill=forestgreen4416044] (axis cs:237.5,0) rectangle (axis cs:238.5,60);
\draw[draw=none,fill=forestgreen4416044] (axis cs:238.5,0) rectangle (axis cs:239.5,60);
\draw[draw=none,fill=forestgreen4416044] (axis cs:239.5,0) rectangle (axis cs:240.5,60);
\draw[draw=none,fill=forestgreen4416044] (axis cs:240.5,0) rectangle (axis cs:241.5,60);
\draw[draw=none,fill=forestgreen4416044] (axis cs:241.5,0) rectangle (axis cs:242.5,60);
\draw[draw=none,fill=forestgreen4416044] (axis cs:242.5,0) rectangle (axis cs:243.5,60);
\draw[draw=none,fill=forestgreen4416044] (axis cs:243.5,0) rectangle (axis cs:244.5,60);
\draw[draw=none,fill=forestgreen4416044] (axis cs:244.5,0) rectangle (axis cs:245.5,60);
\draw[draw=none,fill=forestgreen4416044] (axis cs:245.5,0) rectangle (axis cs:246.5,60);
\draw[draw=none,fill=forestgreen4416044] (axis cs:246.5,0) rectangle (axis cs:247.5,60);
\draw[draw=none,fill=forestgreen4416044] (axis cs:247.5,0) rectangle (axis cs:248.5,64);
\draw[draw=none,fill=forestgreen4416044] (axis cs:248.5,0) rectangle (axis cs:249.5,64);
\draw[draw=none,fill=forestgreen4416044] (axis cs:249.5,0) rectangle (axis cs:250.5,64);
\draw[draw=none,fill=forestgreen4416044] (axis cs:250.5,0) rectangle (axis cs:251.5,64);
\draw[draw=none,fill=forestgreen4416044] (axis cs:251.5,0) rectangle (axis cs:252.5,64);
\draw[draw=none,fill=forestgreen4416044] (axis cs:252.5,0) rectangle (axis cs:253.5,64);
\draw[draw=none,fill=forestgreen4416044] (axis cs:253.5,0) rectangle (axis cs:254.5,64);
\draw[draw=none,fill=forestgreen4416044] (axis cs:254.5,0) rectangle (axis cs:255.5,64);
\draw[draw=none,fill=forestgreen4416044] (axis cs:255.5,0) rectangle (axis cs:256.5,64);
\draw[draw=none,fill=forestgreen4416044] (axis cs:256.5,0) rectangle (axis cs:257.5,64);
\draw[draw=none,fill=forestgreen4416044] (axis cs:257.5,0) rectangle (axis cs:258.5,64);
\draw[draw=none,fill=forestgreen4416044] (axis cs:258.5,0) rectangle (axis cs:259.5,64);
\draw[draw=none,fill=forestgreen4416044] (axis cs:259.5,0) rectangle (axis cs:260.5,64);
\draw[draw=none,fill=forestgreen4416044] (axis cs:260.5,0) rectangle (axis cs:261.5,64);
\draw[draw=none,fill=forestgreen4416044] (axis cs:261.5,0) rectangle (axis cs:262.5,64);
\draw[draw=none,fill=forestgreen4416044] (axis cs:262.5,0) rectangle (axis cs:263.5,64);
\draw[draw=none,fill=forestgreen4416044] (axis cs:263.5,0) rectangle (axis cs:264.5,64);
\draw[draw=none,fill=forestgreen4416044] (axis cs:264.5,0) rectangle (axis cs:265.5,64);
\draw[draw=none,fill=forestgreen4416044] (axis cs:265.5,0) rectangle (axis cs:266.5,64);
\draw[draw=none,fill=forestgreen4416044] (axis cs:266.5,0) rectangle (axis cs:267.5,64);
\draw[draw=none,fill=forestgreen4416044] (axis cs:267.5,0) rectangle (axis cs:268.5,64);
\draw[draw=none,fill=forestgreen4416044] (axis cs:268.5,0) rectangle (axis cs:269.5,64);
\draw[draw=none,fill=forestgreen4416044] (axis cs:269.5,0) rectangle (axis cs:270.5,64);
\draw[draw=none,fill=forestgreen4416044] (axis cs:270.5,0) rectangle (axis cs:271.5,64);
\draw[draw=none,fill=forestgreen4416044] (axis cs:271.5,0) rectangle (axis cs:272.5,64);
\draw[draw=none,fill=forestgreen4416044] (axis cs:272.5,0) rectangle (axis cs:273.5,64);
\draw[draw=none,fill=forestgreen4416044] (axis cs:273.5,0) rectangle (axis cs:274.5,64);
\draw[draw=none,fill=forestgreen4416044] (axis cs:274.5,0) rectangle (axis cs:275.5,64);
\draw[draw=none,fill=forestgreen4416044] (axis cs:275.5,0) rectangle (axis cs:276.5,64);
\draw[draw=none,fill=forestgreen4416044] (axis cs:276.5,0) rectangle (axis cs:277.5,64);
\draw[draw=none,fill=forestgreen4416044] (axis cs:277.5,0) rectangle (axis cs:278.5,64);
\draw[draw=none,fill=forestgreen4416044] (axis cs:278.5,0) rectangle (axis cs:279.5,64);
\draw[draw=none,fill=forestgreen4416044] (axis cs:279.5,0) rectangle (axis cs:280.5,64);
\draw[draw=none,fill=forestgreen4416044] (axis cs:280.5,0) rectangle (axis cs:281.5,64);
\draw[draw=none,fill=forestgreen4416044] (axis cs:281.5,0) rectangle (axis cs:282.5,64);
\draw[draw=none,fill=forestgreen4416044] (axis cs:282.5,0) rectangle (axis cs:283.5,64);
\draw[draw=none,fill=forestgreen4416044] (axis cs:283.5,0) rectangle (axis cs:284.5,64);
\draw[draw=none,fill=forestgreen4416044] (axis cs:284.5,0) rectangle (axis cs:285.5,64);
\draw[draw=none,fill=forestgreen4416044] (axis cs:285.5,0) rectangle (axis cs:286.5,64);
\draw[draw=none,fill=forestgreen4416044] (axis cs:286.5,0) rectangle (axis cs:287.5,64);
\draw[draw=none,fill=forestgreen4416044] (axis cs:287.5,0) rectangle (axis cs:288.5,64);
\draw[draw=none,fill=forestgreen4416044] (axis cs:288.5,0) rectangle (axis cs:289.5,68);
\draw[draw=none,fill=forestgreen4416044] (axis cs:289.5,0) rectangle (axis cs:290.5,68);
\draw[draw=none,fill=forestgreen4416044] (axis cs:290.5,0) rectangle (axis cs:291.5,70);
\draw[draw=none,fill=forestgreen4416044] (axis cs:291.5,0) rectangle (axis cs:292.5,70);
\draw[draw=none,fill=forestgreen4416044] (axis cs:292.5,0) rectangle (axis cs:293.5,70);
\draw[draw=none,fill=forestgreen4416044] (axis cs:293.5,0) rectangle (axis cs:294.5,70);
\draw[draw=none,fill=forestgreen4416044] (axis cs:294.5,0) rectangle (axis cs:295.5,70);
\draw[draw=none,fill=forestgreen4416044] (axis cs:295.5,0) rectangle (axis cs:296.5,70);
\draw[draw=none,fill=forestgreen4416044] (axis cs:296.5,0) rectangle (axis cs:297.5,70);
\draw[draw=none,fill=forestgreen4416044] (axis cs:297.5,0) rectangle (axis cs:298.5,70);
\draw[draw=none,fill=forestgreen4416044] (axis cs:298.5,0) rectangle (axis cs:299.5,70);
\draw[draw=none,fill=forestgreen4416044] (axis cs:299.5,0) rectangle (axis cs:300.5,70);
\draw[draw=none,fill=forestgreen4416044] (axis cs:300.5,0) rectangle (axis cs:301.5,70);
\draw[draw=none,fill=forestgreen4416044] (axis cs:301.5,0) rectangle (axis cs:302.5,70);
\draw[draw=none,fill=forestgreen4416044] (axis cs:302.5,0) rectangle (axis cs:303.5,70);
\draw[draw=none,fill=forestgreen4416044] (axis cs:303.5,0) rectangle (axis cs:304.5,70);
\draw[draw=none,fill=forestgreen4416044] (axis cs:304.5,0) rectangle (axis cs:305.5,72);
\draw[draw=none,fill=forestgreen4416044] (axis cs:305.5,0) rectangle (axis cs:306.5,72);
\draw[draw=none,fill=forestgreen4416044] (axis cs:306.5,0) rectangle (axis cs:307.5,72);
\draw[draw=none,fill=forestgreen4416044] (axis cs:307.5,0) rectangle (axis cs:308.5,72);
\draw[draw=none,fill=forestgreen4416044] (axis cs:308.5,0) rectangle (axis cs:309.5,74);
\draw[draw=none,fill=forestgreen4416044] (axis cs:309.5,0) rectangle (axis cs:310.5,74);
\draw[draw=none,fill=forestgreen4416044] (axis cs:310.5,0) rectangle (axis cs:311.5,74);
\draw[draw=none,fill=forestgreen4416044] (axis cs:311.5,0) rectangle (axis cs:312.5,76);
\draw[draw=none,fill=forestgreen4416044] (axis cs:312.5,0) rectangle (axis cs:313.5,76);
\draw[draw=none,fill=forestgreen4416044] (axis cs:313.5,0) rectangle (axis cs:314.5,76);
\draw[draw=none,fill=forestgreen4416044] (axis cs:314.5,0) rectangle (axis cs:315.5,76);
\draw[draw=none,fill=forestgreen4416044] (axis cs:315.5,0) rectangle (axis cs:316.5,76);
\draw[draw=none,fill=forestgreen4416044] (axis cs:316.5,0) rectangle (axis cs:317.5,76);
\draw[draw=none,fill=forestgreen4416044] (axis cs:317.5,0) rectangle (axis cs:318.5,76);
\draw[draw=none,fill=forestgreen4416044] (axis cs:318.5,0) rectangle (axis cs:319.5,76);
\draw[draw=none,fill=forestgreen4416044] (axis cs:319.5,0) rectangle (axis cs:320.5,76);
\draw[draw=none,fill=forestgreen4416044] (axis cs:320.5,0) rectangle (axis cs:321.5,76);
\draw[draw=none,fill=forestgreen4416044] (axis cs:321.5,0) rectangle (axis cs:322.5,76);
\draw[draw=none,fill=forestgreen4416044] (axis cs:322.5,0) rectangle (axis cs:323.5,76);
\draw[draw=none,fill=forestgreen4416044] (axis cs:323.5,0) rectangle (axis cs:324.5,76);
\draw[draw=none,fill=forestgreen4416044] (axis cs:324.5,0) rectangle (axis cs:325.5,76);
\draw[draw=none,fill=forestgreen4416044] (axis cs:325.5,0) rectangle (axis cs:326.5,76);
\draw[draw=none,fill=forestgreen4416044] (axis cs:326.5,0) rectangle (axis cs:327.5,76);
\draw[draw=none,fill=forestgreen4416044] (axis cs:327.5,0) rectangle (axis cs:328.5,76);
\draw[draw=none,fill=forestgreen4416044] (axis cs:328.5,0) rectangle (axis cs:329.5,76);
\draw[draw=none,fill=forestgreen4416044] (axis cs:329.5,0) rectangle (axis cs:330.5,76);
\draw[draw=none,fill=forestgreen4416044] (axis cs:330.5,0) rectangle (axis cs:331.5,76);
\draw[draw=none,fill=forestgreen4416044] (axis cs:331.5,0) rectangle (axis cs:332.5,76);
\draw[draw=none,fill=forestgreen4416044] (axis cs:332.5,0) rectangle (axis cs:333.5,76);
\draw[draw=none,fill=forestgreen4416044] (axis cs:333.5,0) rectangle (axis cs:334.5,76);
\draw[draw=none,fill=forestgreen4416044] (axis cs:334.5,0) rectangle (axis cs:335.5,76);
\draw[draw=none,fill=forestgreen4416044] (axis cs:335.5,0) rectangle (axis cs:336.5,76);
\draw[draw=none,fill=forestgreen4416044] (axis cs:336.5,0) rectangle (axis cs:337.5,76);
\draw[draw=none,fill=forestgreen4416044] (axis cs:337.5,0) rectangle (axis cs:338.5,76);
\draw[draw=none,fill=forestgreen4416044] (axis cs:338.5,0) rectangle (axis cs:339.5,76);
\draw[draw=none,fill=forestgreen4416044] (axis cs:339.5,0) rectangle (axis cs:340.5,76);
\draw[draw=none,fill=forestgreen4416044] (axis cs:340.5,0) rectangle (axis cs:341.5,76);
\draw[draw=none,fill=forestgreen4416044] (axis cs:341.5,0) rectangle (axis cs:342.5,76);
\draw[draw=none,fill=forestgreen4416044] (axis cs:342.5,0) rectangle (axis cs:343.5,76);
\draw[draw=none,fill=forestgreen4416044] (axis cs:343.5,0) rectangle (axis cs:344.5,76);
\draw[draw=none,fill=forestgreen4416044] (axis cs:344.5,0) rectangle (axis cs:345.5,76);
\draw[draw=none,fill=forestgreen4416044] (axis cs:345.5,0) rectangle (axis cs:346.5,76);
\draw[draw=none,fill=forestgreen4416044] (axis cs:346.5,0) rectangle (axis cs:347.5,76);
\draw[draw=none,fill=forestgreen4416044] (axis cs:347.5,0) rectangle (axis cs:348.5,76);
\draw[draw=none,fill=forestgreen4416044] (axis cs:348.5,0) rectangle (axis cs:349.5,76);
\draw[draw=none,fill=forestgreen4416044] (axis cs:349.5,0) rectangle (axis cs:350.5,76);
\draw[draw=none,fill=forestgreen4416044] (axis cs:350.5,0) rectangle (axis cs:351.5,76);
\draw[draw=none,fill=forestgreen4416044] (axis cs:351.5,0) rectangle (axis cs:352.5,76);
\draw[draw=none,fill=forestgreen4416044] (axis cs:352.5,0) rectangle (axis cs:353.5,76);
\draw[draw=none,fill=forestgreen4416044] (axis cs:353.5,0) rectangle (axis cs:354.5,76);
\draw[draw=none,fill=forestgreen4416044] (axis cs:354.5,0) rectangle (axis cs:355.5,76);
\draw[draw=none,fill=forestgreen4416044] (axis cs:355.5,0) rectangle (axis cs:356.5,76);
\draw[draw=none,fill=forestgreen4416044] (axis cs:356.5,0) rectangle (axis cs:357.5,76);
\draw[draw=none,fill=forestgreen4416044] (axis cs:357.5,0) rectangle (axis cs:358.5,76);
\draw[draw=none,fill=forestgreen4416044] (axis cs:358.5,0) rectangle (axis cs:359.5,76);
\draw[draw=none,fill=forestgreen4416044] (axis cs:359.5,0) rectangle (axis cs:360.5,76);
\draw[draw=none,fill=forestgreen4416044] (axis cs:360.5,0) rectangle (axis cs:361.5,76);
\draw[draw=none,fill=forestgreen4416044] (axis cs:361.5,0) rectangle (axis cs:362.5,76);
\draw[draw=none,fill=forestgreen4416044] (axis cs:362.5,0) rectangle (axis cs:363.5,76);
\draw[draw=none,fill=forestgreen4416044] (axis cs:363.5,0) rectangle (axis cs:364.5,76);
\draw[draw=none,fill=forestgreen4416044] (axis cs:364.5,0) rectangle (axis cs:365.5,76);
\draw[draw=none,fill=forestgreen4416044] (axis cs:365.5,0) rectangle (axis cs:366.5,76);
\draw[draw=none,fill=forestgreen4416044] (axis cs:366.5,0) rectangle (axis cs:367.5,78);
\draw[draw=none,fill=forestgreen4416044] (axis cs:367.5,0) rectangle (axis cs:368.5,78);
\draw[draw=none,fill=forestgreen4416044] (axis cs:368.5,0) rectangle (axis cs:369.5,78);
\draw[draw=none,fill=forestgreen4416044] (axis cs:369.5,0) rectangle (axis cs:370.5,78);
\draw[draw=none,fill=forestgreen4416044] (axis cs:370.5,0) rectangle (axis cs:371.5,78);
\draw[draw=none,fill=forestgreen4416044] (axis cs:371.5,0) rectangle (axis cs:372.5,78);
\draw[draw=none,fill=forestgreen4416044] (axis cs:372.5,0) rectangle (axis cs:373.5,78);
\draw[draw=none,fill=forestgreen4416044] (axis cs:373.5,0) rectangle (axis cs:374.5,78);
\draw[draw=none,fill=forestgreen4416044] (axis cs:374.5,0) rectangle (axis cs:375.5,78);
\draw[draw=none,fill=forestgreen4416044] (axis cs:375.5,0) rectangle (axis cs:376.5,78);
\draw[draw=none,fill=forestgreen4416044] (axis cs:376.5,0) rectangle (axis cs:377.5,78);
\draw[draw=none,fill=forestgreen4416044] (axis cs:377.5,0) rectangle (axis cs:378.5,78);
\draw[draw=none,fill=forestgreen4416044] (axis cs:378.5,0) rectangle (axis cs:379.5,78);
\draw[draw=none,fill=forestgreen4416044] (axis cs:379.5,0) rectangle (axis cs:380.5,78);
\draw[draw=none,fill=forestgreen4416044] (axis cs:380.5,0) rectangle (axis cs:381.5,78);
\draw[draw=none,fill=forestgreen4416044] (axis cs:381.5,0) rectangle (axis cs:382.5,78);
\draw[draw=none,fill=forestgreen4416044] (axis cs:382.5,0) rectangle (axis cs:383.5,78);
\draw[draw=none,fill=forestgreen4416044] (axis cs:383.5,0) rectangle (axis cs:384.5,78);
\draw[draw=none,fill=forestgreen4416044] (axis cs:384.5,0) rectangle (axis cs:385.5,78);
\draw[draw=none,fill=forestgreen4416044] (axis cs:385.5,0) rectangle (axis cs:386.5,78);
\draw[draw=none,fill=forestgreen4416044] (axis cs:386.5,0) rectangle (axis cs:387.5,78);
\draw[draw=none,fill=forestgreen4416044] (axis cs:387.5,0) rectangle (axis cs:388.5,78);
\draw[draw=none,fill=forestgreen4416044] (axis cs:388.5,0) rectangle (axis cs:389.5,78);
\draw[draw=none,fill=forestgreen4416044] (axis cs:389.5,0) rectangle (axis cs:390.5,78);
\draw[draw=none,fill=forestgreen4416044] (axis cs:390.5,0) rectangle (axis cs:391.5,78);
\draw[draw=none,fill=forestgreen4416044] (axis cs:391.5,0) rectangle (axis cs:392.5,78);
\draw[draw=none,fill=forestgreen4416044] (axis cs:392.5,0) rectangle (axis cs:393.5,78);
\draw[draw=none,fill=forestgreen4416044] (axis cs:393.5,0) rectangle (axis cs:394.5,78);
\draw[draw=none,fill=forestgreen4416044] (axis cs:394.5,0) rectangle (axis cs:395.5,80);
\draw[draw=none,fill=forestgreen4416044] (axis cs:395.5,0) rectangle (axis cs:396.5,80);
\draw[draw=none,fill=forestgreen4416044] (axis cs:396.5,0) rectangle (axis cs:397.5,80);
\draw[draw=none,fill=forestgreen4416044] (axis cs:397.5,0) rectangle (axis cs:398.5,80);
\draw[draw=none,fill=forestgreen4416044] (axis cs:398.5,0) rectangle (axis cs:399.5,80);
\draw[draw=none,fill=forestgreen4416044] (axis cs:399.5,0) rectangle (axis cs:400.5,80);
\draw[draw=none,fill=forestgreen4416044] (axis cs:400.5,0) rectangle (axis cs:401.5,80);
\draw[draw=none,fill=forestgreen4416044] (axis cs:401.5,0) rectangle (axis cs:402.5,80);
\draw[draw=none,fill=forestgreen4416044] (axis cs:402.5,0) rectangle (axis cs:403.5,80);
\draw[draw=none,fill=forestgreen4416044] (axis cs:403.5,0) rectangle (axis cs:404.5,80);
\draw[draw=none,fill=forestgreen4416044] (axis cs:404.5,0) rectangle (axis cs:405.5,80);
\draw[draw=none,fill=forestgreen4416044] (axis cs:405.5,0) rectangle (axis cs:406.5,80);
\draw[draw=none,fill=forestgreen4416044] (axis cs:406.5,0) rectangle (axis cs:407.5,80);
\draw[draw=none,fill=forestgreen4416044] (axis cs:407.5,0) rectangle (axis cs:408.5,80);
\draw[draw=none,fill=forestgreen4416044] (axis cs:408.5,0) rectangle (axis cs:409.5,80);
\draw[draw=none,fill=forestgreen4416044] (axis cs:409.5,0) rectangle (axis cs:410.5,80);
\draw[draw=none,fill=forestgreen4416044] (axis cs:410.5,0) rectangle (axis cs:411.5,80);
\draw[draw=none,fill=forestgreen4416044] (axis cs:411.5,0) rectangle (axis cs:412.5,80);
\draw[draw=none,fill=forestgreen4416044] (axis cs:412.5,0) rectangle (axis cs:413.5,80);
\draw[draw=none,fill=forestgreen4416044] (axis cs:413.5,0) rectangle (axis cs:414.5,80);
\draw[draw=none,fill=forestgreen4416044] (axis cs:414.5,0) rectangle (axis cs:415.5,80);
\draw[draw=none,fill=forestgreen4416044] (axis cs:415.5,0) rectangle (axis cs:416.5,82);
\draw[draw=none,fill=forestgreen4416044] (axis cs:416.5,0) rectangle (axis cs:417.5,84);
\draw[draw=none,fill=forestgreen4416044] (axis cs:417.5,0) rectangle (axis cs:418.5,84);
\draw[draw=none,fill=forestgreen4416044] (axis cs:418.5,0) rectangle (axis cs:419.5,84);
\draw[draw=none,fill=forestgreen4416044] (axis cs:419.5,0) rectangle (axis cs:420.5,84);
\draw[draw=none,fill=forestgreen4416044] (axis cs:420.5,0) rectangle (axis cs:421.5,84);
\draw[draw=none,fill=forestgreen4416044] (axis cs:421.5,0) rectangle (axis cs:422.5,84);
\draw[draw=none,fill=forestgreen4416044] (axis cs:422.5,0) rectangle (axis cs:423.5,84);
\draw[draw=none,fill=forestgreen4416044] (axis cs:423.5,0) rectangle (axis cs:424.5,84);
\draw[draw=none,fill=forestgreen4416044] (axis cs:424.5,0) rectangle (axis cs:425.5,84);
\draw[draw=none,fill=forestgreen4416044] (axis cs:425.5,0) rectangle (axis cs:426.5,84);
\draw[draw=none,fill=forestgreen4416044] (axis cs:426.5,0) rectangle (axis cs:427.5,84);
\draw[draw=none,fill=forestgreen4416044] (axis cs:427.5,0) rectangle (axis cs:428.5,84);
\draw[draw=none,fill=forestgreen4416044] (axis cs:428.5,0) rectangle (axis cs:429.5,84);
\draw[draw=none,fill=forestgreen4416044] (axis cs:429.5,0) rectangle (axis cs:430.5,84);
\draw[draw=none,fill=forestgreen4416044] (axis cs:430.5,0) rectangle (axis cs:431.5,84);
\draw[draw=none,fill=forestgreen4416044] (axis cs:431.5,0) rectangle (axis cs:432.5,84);
\draw[draw=none,fill=forestgreen4416044] (axis cs:432.5,0) rectangle (axis cs:433.5,84);
\draw[draw=none,fill=forestgreen4416044] (axis cs:433.5,0) rectangle (axis cs:434.5,84);
\draw[draw=none,fill=forestgreen4416044] (axis cs:434.5,0) rectangle (axis cs:435.5,84);
\draw[draw=none,fill=forestgreen4416044] (axis cs:435.5,0) rectangle (axis cs:436.5,84);
\draw[draw=none,fill=forestgreen4416044] (axis cs:436.5,0) rectangle (axis cs:437.5,84);
\draw[draw=none,fill=forestgreen4416044] (axis cs:437.5,0) rectangle (axis cs:438.5,84);
\draw[draw=none,fill=forestgreen4416044] (axis cs:438.5,0) rectangle (axis cs:439.5,84);
\draw[draw=none,fill=forestgreen4416044] (axis cs:439.5,0) rectangle (axis cs:440.5,84);
\draw[draw=none,fill=forestgreen4416044] (axis cs:440.5,0) rectangle (axis cs:441.5,86);
\draw[draw=none,fill=forestgreen4416044] (axis cs:441.5,0) rectangle (axis cs:442.5,86);
\draw[draw=none,fill=forestgreen4416044] (axis cs:442.5,0) rectangle (axis cs:443.5,86);
\draw[draw=none,fill=forestgreen4416044] (axis cs:443.5,0) rectangle (axis cs:444.5,86);
\draw[draw=none,fill=forestgreen4416044] (axis cs:444.5,0) rectangle (axis cs:445.5,86);
\draw[draw=none,fill=forestgreen4416044] (axis cs:445.5,0) rectangle (axis cs:446.5,86);
\draw[draw=none,fill=forestgreen4416044] (axis cs:446.5,0) rectangle (axis cs:447.5,86);
\draw[draw=none,fill=forestgreen4416044] (axis cs:447.5,0) rectangle (axis cs:448.5,86);
\draw[draw=none,fill=forestgreen4416044] (axis cs:448.5,0) rectangle (axis cs:449.5,86);
\draw[draw=none,fill=forestgreen4416044] (axis cs:449.5,0) rectangle (axis cs:450.5,86);
\draw[draw=none,fill=forestgreen4416044] (axis cs:450.5,0) rectangle (axis cs:451.5,86);
\draw[draw=none,fill=forestgreen4416044] (axis cs:451.5,0) rectangle (axis cs:452.5,86);
\draw[draw=none,fill=forestgreen4416044] (axis cs:452.5,0) rectangle (axis cs:453.5,86);
\draw[draw=none,fill=forestgreen4416044] (axis cs:453.5,0) rectangle (axis cs:454.5,86);
\draw[draw=none,fill=forestgreen4416044] (axis cs:454.5,0) rectangle (axis cs:455.5,86);
\draw[draw=none,fill=forestgreen4416044] (axis cs:455.5,0) rectangle (axis cs:456.5,86);
\draw[draw=none,fill=forestgreen4416044] (axis cs:456.5,0) rectangle (axis cs:457.5,86);
\draw[draw=none,fill=forestgreen4416044] (axis cs:457.5,0) rectangle (axis cs:458.5,86);
\draw[draw=none,fill=forestgreen4416044] (axis cs:458.5,0) rectangle (axis cs:459.5,86);
\draw[draw=none,fill=forestgreen4416044] (axis cs:459.5,0) rectangle (axis cs:460.5,86);
\draw[draw=none,fill=forestgreen4416044] (axis cs:460.5,0) rectangle (axis cs:461.5,86);
\draw[draw=none,fill=forestgreen4416044] (axis cs:461.5,0) rectangle (axis cs:462.5,86);
\draw[draw=none,fill=forestgreen4416044] (axis cs:462.5,0) rectangle (axis cs:463.5,86);
\draw[draw=none,fill=forestgreen4416044] (axis cs:463.5,0) rectangle (axis cs:464.5,86);
\draw[draw=none,fill=forestgreen4416044] (axis cs:464.5,0) rectangle (axis cs:465.5,86);
\draw[draw=none,fill=forestgreen4416044] (axis cs:465.5,0) rectangle (axis cs:466.5,86);
\draw[draw=none,fill=forestgreen4416044] (axis cs:466.5,0) rectangle (axis cs:467.5,86);
\draw[draw=none,fill=forestgreen4416044] (axis cs:467.5,0) rectangle (axis cs:468.5,86);
\draw[draw=none,fill=forestgreen4416044] (axis cs:468.5,0) rectangle (axis cs:469.5,86);
\draw[draw=none,fill=forestgreen4416044] (axis cs:469.5,0) rectangle (axis cs:470.5,86);
\draw[draw=none,fill=forestgreen4416044] (axis cs:470.5,0) rectangle (axis cs:471.5,86);
\draw[draw=none,fill=forestgreen4416044] (axis cs:471.5,0) rectangle (axis cs:472.5,86);
\draw[draw=none,fill=forestgreen4416044] (axis cs:472.5,0) rectangle (axis cs:473.5,86);
\draw[draw=none,fill=forestgreen4416044] (axis cs:473.5,0) rectangle (axis cs:474.5,86);
\draw[draw=none,fill=forestgreen4416044] (axis cs:474.5,0) rectangle (axis cs:475.5,86);
\draw[draw=none,fill=forestgreen4416044] (axis cs:475.5,0) rectangle (axis cs:476.5,86);
\draw[draw=none,fill=forestgreen4416044] (axis cs:476.5,0) rectangle (axis cs:477.5,86);
\draw[draw=none,fill=forestgreen4416044] (axis cs:477.5,0) rectangle (axis cs:478.5,86);
\draw[draw=none,fill=forestgreen4416044] (axis cs:478.5,0) rectangle (axis cs:479.5,86);
\draw[draw=none,fill=forestgreen4416044] (axis cs:479.5,0) rectangle (axis cs:480.5,86);
\draw[draw=none,fill=forestgreen4416044] (axis cs:480.5,0) rectangle (axis cs:481.5,86);
\draw[draw=none,fill=forestgreen4416044] (axis cs:481.5,0) rectangle (axis cs:482.5,86);
\draw[draw=none,fill=forestgreen4416044] (axis cs:482.5,0) rectangle (axis cs:483.5,86);
\draw[draw=none,fill=forestgreen4416044] (axis cs:483.5,0) rectangle (axis cs:484.5,86);
\draw[draw=none,fill=forestgreen4416044] (axis cs:484.5,0) rectangle (axis cs:485.5,86);
\draw[draw=none,fill=forestgreen4416044] (axis cs:485.5,0) rectangle (axis cs:486.5,86);
\draw[draw=none,fill=forestgreen4416044] (axis cs:486.5,0) rectangle (axis cs:487.5,86);
\draw[draw=none,fill=forestgreen4416044] (axis cs:487.5,0) rectangle (axis cs:488.5,86);
\draw[draw=none,fill=forestgreen4416044] (axis cs:488.5,0) rectangle (axis cs:489.5,86);
\draw[draw=none,fill=forestgreen4416044] (axis cs:489.5,0) rectangle (axis cs:490.5,86);
\draw[draw=none,fill=forestgreen4416044] (axis cs:490.5,0) rectangle (axis cs:491.5,88);
\draw[draw=none,fill=forestgreen4416044] (axis cs:491.5,0) rectangle (axis cs:492.5,88);
\draw[draw=none,fill=forestgreen4416044] (axis cs:492.5,0) rectangle (axis cs:493.5,88);
\draw[draw=none,fill=forestgreen4416044] (axis cs:493.5,0) rectangle (axis cs:494.5,88);
\draw[draw=none,fill=forestgreen4416044] (axis cs:494.5,0) rectangle (axis cs:495.5,88);
\draw[draw=none,fill=forestgreen4416044] (axis cs:495.5,0) rectangle (axis cs:496.5,88);
\draw[draw=none,fill=forestgreen4416044] (axis cs:496.5,0) rectangle (axis cs:497.5,88);
\draw[draw=none,fill=forestgreen4416044] (axis cs:497.5,0) rectangle (axis cs:498.5,88);
\draw[draw=none,fill=forestgreen4416044] (axis cs:498.5,0) rectangle (axis cs:499.5,88);
\draw[draw=none,fill=forestgreen4416044] (axis cs:499.5,0) rectangle (axis cs:500.5,88);
\draw[draw=none,fill=forestgreen4416044] (axis cs:500.5,0) rectangle (axis cs:501.5,88);
\draw[draw=none,fill=forestgreen4416044] (axis cs:501.5,0) rectangle (axis cs:502.5,88);
\draw[draw=none,fill=forestgreen4416044] (axis cs:502.5,0) rectangle (axis cs:503.5,88);
\draw[draw=none,fill=forestgreen4416044] (axis cs:503.5,0) rectangle (axis cs:504.5,88);
\draw[draw=none,fill=forestgreen4416044] (axis cs:504.5,0) rectangle (axis cs:505.5,88);
\draw[draw=none,fill=forestgreen4416044] (axis cs:505.5,0) rectangle (axis cs:506.5,88);
\draw[draw=none,fill=forestgreen4416044] (axis cs:506.5,0) rectangle (axis cs:507.5,88);
\draw[draw=none,fill=forestgreen4416044] (axis cs:507.5,0) rectangle (axis cs:508.5,88);
\draw[draw=none,fill=forestgreen4416044] (axis cs:508.5,0) rectangle (axis cs:509.5,88);
\draw[draw=none,fill=forestgreen4416044] (axis cs:509.5,0) rectangle (axis cs:510.5,88);
\draw[draw=none,fill=forestgreen4416044] (axis cs:510.5,0) rectangle (axis cs:511.5,88);
\draw[draw=none,fill=forestgreen4416044] (axis cs:511.5,0) rectangle (axis cs:512.5,88);
\draw[draw=none,fill=forestgreen4416044] (axis cs:512.5,0) rectangle (axis cs:513.5,88);
\draw[draw=none,fill=forestgreen4416044] (axis cs:513.5,0) rectangle (axis cs:514.5,90);
\draw[draw=none,fill=forestgreen4416044] (axis cs:514.5,0) rectangle (axis cs:515.5,90);
\draw[draw=none,fill=forestgreen4416044] (axis cs:515.5,0) rectangle (axis cs:516.5,90);
\draw[draw=none,fill=forestgreen4416044] (axis cs:516.5,0) rectangle (axis cs:517.5,90);
\draw[draw=none,fill=forestgreen4416044] (axis cs:517.5,0) rectangle (axis cs:518.5,90);
\draw[draw=none,fill=forestgreen4416044] (axis cs:518.5,0) rectangle (axis cs:519.5,90);
\draw[draw=none,fill=forestgreen4416044] (axis cs:519.5,0) rectangle (axis cs:520.5,90);
\draw[draw=none,fill=forestgreen4416044] (axis cs:520.5,0) rectangle (axis cs:521.5,90);
\draw[draw=none,fill=forestgreen4416044] (axis cs:521.5,0) rectangle (axis cs:522.5,90);
\draw[draw=none,fill=forestgreen4416044] (axis cs:522.5,0) rectangle (axis cs:523.5,92);
\draw[draw=none,fill=forestgreen4416044] (axis cs:523.5,0) rectangle (axis cs:524.5,92);
\draw[draw=none,fill=forestgreen4416044] (axis cs:524.5,0) rectangle (axis cs:525.5,92);
\draw[draw=none,fill=forestgreen4416044] (axis cs:525.5,0) rectangle (axis cs:526.5,92);
\draw[draw=none,fill=forestgreen4416044] (axis cs:526.5,0) rectangle (axis cs:527.5,92);
\draw[draw=none,fill=forestgreen4416044] (axis cs:527.5,0) rectangle (axis cs:528.5,92);
\draw[draw=none,fill=forestgreen4416044] (axis cs:528.5,0) rectangle (axis cs:529.5,92);
\draw[draw=none,fill=forestgreen4416044] (axis cs:529.5,0) rectangle (axis cs:530.5,92);
\draw[draw=none,fill=forestgreen4416044] (axis cs:530.5,0) rectangle (axis cs:531.5,92);
\draw[draw=none,fill=forestgreen4416044] (axis cs:531.5,0) rectangle (axis cs:532.5,92);
\draw[draw=none,fill=forestgreen4416044] (axis cs:532.5,0) rectangle (axis cs:533.5,92);
\draw[draw=none,fill=forestgreen4416044] (axis cs:533.5,0) rectangle (axis cs:534.5,92);
\draw[draw=none,fill=forestgreen4416044] (axis cs:534.5,0) rectangle (axis cs:535.5,92);
\draw[draw=none,fill=forestgreen4416044] (axis cs:535.5,0) rectangle (axis cs:536.5,92);
\draw[draw=none,fill=forestgreen4416044] (axis cs:536.5,0) rectangle (axis cs:537.5,92);
\draw[draw=none,fill=forestgreen4416044] (axis cs:537.5,0) rectangle (axis cs:538.5,92);
\draw[draw=none,fill=forestgreen4416044] (axis cs:538.5,0) rectangle (axis cs:539.5,92);
\draw[draw=none,fill=forestgreen4416044] (axis cs:539.5,0) rectangle (axis cs:540.5,92);
\draw[draw=none,fill=forestgreen4416044] (axis cs:540.5,0) rectangle (axis cs:541.5,92);
\draw[draw=none,fill=forestgreen4416044] (axis cs:541.5,0) rectangle (axis cs:542.5,92);
\draw[draw=none,fill=forestgreen4416044] (axis cs:542.5,0) rectangle (axis cs:543.5,92);
\draw[draw=none,fill=forestgreen4416044] (axis cs:543.5,0) rectangle (axis cs:544.5,92);
\draw[draw=none,fill=forestgreen4416044] (axis cs:544.5,0) rectangle (axis cs:545.5,92);
\draw[draw=none,fill=forestgreen4416044] (axis cs:545.5,0) rectangle (axis cs:546.5,92);
\draw[draw=none,fill=forestgreen4416044] (axis cs:546.5,0) rectangle (axis cs:547.5,92);
\draw[draw=none,fill=forestgreen4416044] (axis cs:547.5,0) rectangle (axis cs:548.5,92);
\draw[draw=none,fill=forestgreen4416044] (axis cs:548.5,0) rectangle (axis cs:549.5,92);
\draw[draw=none,fill=forestgreen4416044] (axis cs:549.5,0) rectangle (axis cs:550.5,92);
\draw[draw=none,fill=forestgreen4416044] (axis cs:550.5,0) rectangle (axis cs:551.5,92);
\draw[draw=none,fill=forestgreen4416044] (axis cs:551.5,0) rectangle (axis cs:552.5,92);
\draw[draw=none,fill=forestgreen4416044] (axis cs:552.5,0) rectangle (axis cs:553.5,92);
\draw[draw=none,fill=forestgreen4416044] (axis cs:553.5,0) rectangle (axis cs:554.5,92);
\draw[draw=none,fill=forestgreen4416044] (axis cs:554.5,0) rectangle (axis cs:555.5,94);
\draw[draw=none,fill=forestgreen4416044] (axis cs:555.5,0) rectangle (axis cs:556.5,94);
\draw[draw=none,fill=forestgreen4416044] (axis cs:556.5,0) rectangle (axis cs:557.5,94);
\draw[draw=none,fill=forestgreen4416044] (axis cs:557.5,0) rectangle (axis cs:558.5,94);
\draw[draw=none,fill=forestgreen4416044] (axis cs:558.5,0) rectangle (axis cs:559.5,94);
\draw[draw=none,fill=forestgreen4416044] (axis cs:559.5,0) rectangle (axis cs:560.5,94);
\draw[draw=none,fill=forestgreen4416044] (axis cs:560.5,0) rectangle (axis cs:561.5,94);
\draw[draw=none,fill=forestgreen4416044] (axis cs:561.5,0) rectangle (axis cs:562.5,94);
\draw[draw=none,fill=forestgreen4416044] (axis cs:562.5,0) rectangle (axis cs:563.5,94);
\draw[draw=none,fill=forestgreen4416044] (axis cs:563.5,0) rectangle (axis cs:564.5,94);
\draw[draw=none,fill=forestgreen4416044] (axis cs:564.5,0) rectangle (axis cs:565.5,94);
\draw[draw=none,fill=forestgreen4416044] (axis cs:565.5,0) rectangle (axis cs:566.5,94);
\draw[draw=none,fill=forestgreen4416044] (axis cs:566.5,0) rectangle (axis cs:567.5,94);
\draw[draw=none,fill=forestgreen4416044] (axis cs:567.5,0) rectangle (axis cs:568.5,94);
\draw[draw=none,fill=forestgreen4416044] (axis cs:568.5,0) rectangle (axis cs:569.5,94);
\draw[draw=none,fill=forestgreen4416044] (axis cs:569.5,0) rectangle (axis cs:570.5,94);
\draw[draw=none,fill=forestgreen4416044] (axis cs:570.5,0) rectangle (axis cs:571.5,94);
\draw[draw=none,fill=forestgreen4416044] (axis cs:571.5,0) rectangle (axis cs:572.5,94);
\draw[draw=none,fill=forestgreen4416044] (axis cs:572.5,0) rectangle (axis cs:573.5,94);
\draw[draw=none,fill=forestgreen4416044] (axis cs:573.5,0) rectangle (axis cs:574.5,94);
\draw[draw=none,fill=forestgreen4416044] (axis cs:574.5,0) rectangle (axis cs:575.5,94);
\draw[draw=none,fill=forestgreen4416044] (axis cs:575.5,0) rectangle (axis cs:576.5,94);
\draw[draw=none,fill=forestgreen4416044] (axis cs:576.5,0) rectangle (axis cs:577.5,94);
\draw[draw=none,fill=forestgreen4416044] (axis cs:577.5,0) rectangle (axis cs:578.5,94);
\draw[draw=none,fill=forestgreen4416044] (axis cs:578.5,0) rectangle (axis cs:579.5,94);
\draw[draw=none,fill=forestgreen4416044] (axis cs:579.5,0) rectangle (axis cs:580.5,94);
\draw[draw=none,fill=forestgreen4416044] (axis cs:580.5,0) rectangle (axis cs:581.5,94);
\draw[draw=none,fill=forestgreen4416044] (axis cs:581.5,0) rectangle (axis cs:582.5,94);
\draw[draw=none,fill=forestgreen4416044] (axis cs:582.5,0) rectangle (axis cs:583.5,94);
\draw[draw=none,fill=forestgreen4416044] (axis cs:583.5,0) rectangle (axis cs:584.5,96);
\draw[draw=none,fill=forestgreen4416044] (axis cs:584.5,0) rectangle (axis cs:585.5,96);
\draw[draw=none,fill=forestgreen4416044] (axis cs:585.5,0) rectangle (axis cs:586.5,96);
\draw[draw=none,fill=forestgreen4416044] (axis cs:586.5,0) rectangle (axis cs:587.5,96);
\draw[draw=none,fill=forestgreen4416044] (axis cs:587.5,0) rectangle (axis cs:588.5,96);
\draw[draw=none,fill=forestgreen4416044] (axis cs:588.5,0) rectangle (axis cs:589.5,96);
\draw[draw=none,fill=forestgreen4416044] (axis cs:589.5,0) rectangle (axis cs:590.5,96);
\draw[draw=none,fill=forestgreen4416044] (axis cs:590.5,0) rectangle (axis cs:591.5,96);
\draw[draw=none,fill=forestgreen4416044] (axis cs:591.5,0) rectangle (axis cs:592.5,96);
\draw[draw=none,fill=forestgreen4416044] (axis cs:592.5,0) rectangle (axis cs:593.5,96);
\draw[draw=none,fill=forestgreen4416044] (axis cs:593.5,0) rectangle (axis cs:594.5,96);
\draw[draw=none,fill=forestgreen4416044] (axis cs:594.5,0) rectangle (axis cs:595.5,96);
\draw[draw=none,fill=forestgreen4416044] (axis cs:595.5,0) rectangle (axis cs:596.5,96);
\draw[draw=none,fill=forestgreen4416044] (axis cs:596.5,0) rectangle (axis cs:597.5,96);
\draw[draw=none,fill=forestgreen4416044] (axis cs:597.5,0) rectangle (axis cs:598.5,98);
\draw[draw=none,fill=forestgreen4416044] (axis cs:598.5,0) rectangle (axis cs:599.5,98);
\draw[draw=none,fill=forestgreen4416044] (axis cs:599.5,0) rectangle (axis cs:600.5,98);
\draw[draw=none,fill=forestgreen4416044] (axis cs:600.5,0) rectangle (axis cs:601.5,98);
\draw[draw=none,fill=forestgreen4416044] (axis cs:601.5,0) rectangle (axis cs:602.5,98);
\draw[draw=none,fill=forestgreen4416044] (axis cs:602.5,0) rectangle (axis cs:603.5,98);
\draw[draw=none,fill=forestgreen4416044] (axis cs:603.5,0) rectangle (axis cs:604.5,98);
\draw[draw=none,fill=forestgreen4416044] (axis cs:604.5,0) rectangle (axis cs:605.5,98);
\draw[draw=none,fill=forestgreen4416044] (axis cs:605.5,0) rectangle (axis cs:606.5,98);
\draw[draw=none,fill=forestgreen4416044] (axis cs:606.5,0) rectangle (axis cs:607.5,98);
\draw[draw=none,fill=forestgreen4416044] (axis cs:607.5,0) rectangle (axis cs:608.5,98);
\draw[draw=none,fill=forestgreen4416044] (axis cs:608.5,0) rectangle (axis cs:609.5,98);
\draw[draw=none,fill=forestgreen4416044] (axis cs:609.5,0) rectangle (axis cs:610.5,98);
\draw[draw=none,fill=forestgreen4416044] (axis cs:610.5,0) rectangle (axis cs:611.5,98);
\draw[draw=none,fill=forestgreen4416044] (axis cs:611.5,0) rectangle (axis cs:612.5,98);
\draw[draw=none,fill=forestgreen4416044] (axis cs:612.5,0) rectangle (axis cs:613.5,98);
\draw[draw=none,fill=forestgreen4416044] (axis cs:613.5,0) rectangle (axis cs:614.5,98);
\draw[draw=none,fill=forestgreen4416044] (axis cs:614.5,0) rectangle (axis cs:615.5,98);
\draw[draw=none,fill=forestgreen4416044] (axis cs:615.5,0) rectangle (axis cs:616.5,98);
\draw[draw=none,fill=forestgreen4416044] (axis cs:616.5,0) rectangle (axis cs:617.5,98);
\draw[draw=none,fill=forestgreen4416044] (axis cs:617.5,0) rectangle (axis cs:618.5,98);
\draw[draw=none,fill=forestgreen4416044] (axis cs:618.5,0) rectangle (axis cs:619.5,98);
\draw[draw=none,fill=forestgreen4416044] (axis cs:619.5,0) rectangle (axis cs:620.5,98);
\draw[draw=none,fill=forestgreen4416044] (axis cs:620.5,0) rectangle (axis cs:621.5,98);
\draw[draw=none,fill=forestgreen4416044] (axis cs:621.5,0) rectangle (axis cs:622.5,98);
\draw[draw=none,fill=forestgreen4416044] (axis cs:622.5,0) rectangle (axis cs:623.5,98);
\draw[draw=none,fill=forestgreen4416044] (axis cs:623.5,0) rectangle (axis cs:624.5,98);
\draw[draw=none,fill=forestgreen4416044] (axis cs:624.5,0) rectangle (axis cs:625.5,98);
\draw[draw=none,fill=forestgreen4416044] (axis cs:625.5,0) rectangle (axis cs:626.5,98);
\draw[draw=none,fill=forestgreen4416044] (axis cs:626.5,0) rectangle (axis cs:627.5,98);
\draw[draw=none,fill=forestgreen4416044] (axis cs:627.5,0) rectangle (axis cs:628.5,98);
\draw[draw=none,fill=forestgreen4416044] (axis cs:628.5,0) rectangle (axis cs:629.5,98);
\draw[draw=none,fill=forestgreen4416044] (axis cs:629.5,0) rectangle (axis cs:630.5,98);
\draw[draw=none,fill=forestgreen4416044] (axis cs:630.5,0) rectangle (axis cs:631.5,98);
\draw[draw=none,fill=forestgreen4416044] (axis cs:631.5,0) rectangle (axis cs:632.5,98);
\draw[draw=none,fill=forestgreen4416044] (axis cs:632.5,0) rectangle (axis cs:633.5,98);
\draw[draw=none,fill=forestgreen4416044] (axis cs:633.5,0) rectangle (axis cs:634.5,98);
\draw[draw=none,fill=forestgreen4416044] (axis cs:634.5,0) rectangle (axis cs:635.5,98);
\draw[draw=none,fill=forestgreen4416044] (axis cs:635.5,0) rectangle (axis cs:636.5,98);
\draw[draw=none,fill=forestgreen4416044] (axis cs:636.5,0) rectangle (axis cs:637.5,98);
\draw[draw=none,fill=forestgreen4416044] (axis cs:637.5,0) rectangle (axis cs:638.5,98);
\draw[draw=none,fill=forestgreen4416044] (axis cs:638.5,0) rectangle (axis cs:639.5,98);
\draw[draw=none,fill=forestgreen4416044] (axis cs:639.5,0) rectangle (axis cs:640.5,98);
\draw[draw=none,fill=forestgreen4416044] (axis cs:640.5,0) rectangle (axis cs:641.5,98);
\draw[draw=none,fill=forestgreen4416044] (axis cs:641.5,0) rectangle (axis cs:642.5,98);
\draw[draw=none,fill=forestgreen4416044] (axis cs:642.5,0) rectangle (axis cs:643.5,98);
\draw[draw=none,fill=forestgreen4416044] (axis cs:643.5,0) rectangle (axis cs:644.5,98);
\draw[draw=none,fill=forestgreen4416044] (axis cs:644.5,0) rectangle (axis cs:645.5,98);
\draw[draw=none,fill=forestgreen4416044] (axis cs:645.5,0) rectangle (axis cs:646.5,98);
\draw[draw=none,fill=forestgreen4416044] (axis cs:646.5,0) rectangle (axis cs:647.5,98);
\draw[draw=none,fill=forestgreen4416044] (axis cs:647.5,0) rectangle (axis cs:648.5,98);
\draw[draw=none,fill=forestgreen4416044] (axis cs:648.5,0) rectangle (axis cs:649.5,98);
\draw[draw=none,fill=forestgreen4416044] (axis cs:649.5,0) rectangle (axis cs:650.5,98);
\draw[draw=none,fill=forestgreen4416044] (axis cs:650.5,0) rectangle (axis cs:651.5,98);
\draw[draw=none,fill=forestgreen4416044] (axis cs:651.5,0) rectangle (axis cs:652.5,98);
\draw[draw=none,fill=forestgreen4416044] (axis cs:652.5,0) rectangle (axis cs:653.5,98);
\draw[draw=none,fill=forestgreen4416044] (axis cs:653.5,0) rectangle (axis cs:654.5,98);
\draw[draw=none,fill=forestgreen4416044] (axis cs:654.5,0) rectangle (axis cs:655.5,98);
\draw[draw=none,fill=forestgreen4416044] (axis cs:655.5,0) rectangle (axis cs:656.5,98);
\draw[draw=none,fill=forestgreen4416044] (axis cs:656.5,0) rectangle (axis cs:657.5,98);
\draw[draw=none,fill=forestgreen4416044] (axis cs:657.5,0) rectangle (axis cs:658.5,98);
\draw[draw=none,fill=forestgreen4416044] (axis cs:658.5,0) rectangle (axis cs:659.5,98);
\draw[draw=none,fill=forestgreen4416044] (axis cs:659.5,0) rectangle (axis cs:660.5,98);
\draw[draw=none,fill=forestgreen4416044] (axis cs:660.5,0) rectangle (axis cs:661.5,98);
\draw[draw=none,fill=forestgreen4416044] (axis cs:661.5,0) rectangle (axis cs:662.5,98);
\draw[draw=none,fill=forestgreen4416044] (axis cs:662.5,0) rectangle (axis cs:663.5,98);
\draw[draw=none,fill=forestgreen4416044] (axis cs:663.5,0) rectangle (axis cs:664.5,98);
\draw[draw=none,fill=forestgreen4416044] (axis cs:664.5,0) rectangle (axis cs:665.5,98);
\draw[draw=none,fill=forestgreen4416044] (axis cs:665.5,0) rectangle (axis cs:666.5,98);
\draw[draw=none,fill=forestgreen4416044] (axis cs:666.5,0) rectangle (axis cs:667.5,98);
\draw[draw=none,fill=forestgreen4416044] (axis cs:667.5,0) rectangle (axis cs:668.5,98);
\draw[draw=none,fill=forestgreen4416044] (axis cs:668.5,0) rectangle (axis cs:669.5,98);
\draw[draw=none,fill=forestgreen4416044] (axis cs:669.5,0) rectangle (axis cs:670.5,98);
\draw[draw=none,fill=forestgreen4416044] (axis cs:670.5,0) rectangle (axis cs:671.5,98);
\draw[draw=none,fill=forestgreen4416044] (axis cs:671.5,0) rectangle (axis cs:672.5,98);
\draw[draw=none,fill=forestgreen4416044] (axis cs:672.5,0) rectangle (axis cs:673.5,98);
\draw[draw=none,fill=forestgreen4416044] (axis cs:673.5,0) rectangle (axis cs:674.5,98);
\draw[draw=none,fill=forestgreen4416044] (axis cs:674.5,0) rectangle (axis cs:675.5,98);
\draw[draw=none,fill=forestgreen4416044] (axis cs:675.5,0) rectangle (axis cs:676.5,98);
\draw[draw=none,fill=forestgreen4416044] (axis cs:676.5,0) rectangle (axis cs:677.5,98);
\draw[draw=none,fill=forestgreen4416044] (axis cs:677.5,0) rectangle (axis cs:678.5,98);
\draw[draw=none,fill=forestgreen4416044] (axis cs:678.5,0) rectangle (axis cs:679.5,98);
\draw[draw=none,fill=forestgreen4416044] (axis cs:679.5,0) rectangle (axis cs:680.5,98);
\draw[draw=none,fill=forestgreen4416044] (axis cs:680.5,0) rectangle (axis cs:681.5,98);
\draw[draw=none,fill=forestgreen4416044] (axis cs:681.5,0) rectangle (axis cs:682.5,98);
\draw[draw=none,fill=forestgreen4416044] (axis cs:682.5,0) rectangle (axis cs:683.5,98);
\draw[draw=none,fill=forestgreen4416044] (axis cs:683.5,0) rectangle (axis cs:684.5,98);
\draw[draw=none,fill=forestgreen4416044] (axis cs:684.5,0) rectangle (axis cs:685.5,98);
\draw[draw=none,fill=forestgreen4416044] (axis cs:685.5,0) rectangle (axis cs:686.5,98);
\draw[draw=none,fill=forestgreen4416044] (axis cs:686.5,0) rectangle (axis cs:687.5,98);
\draw[draw=none,fill=forestgreen4416044] (axis cs:687.5,0) rectangle (axis cs:688.5,98);
\draw[draw=none,fill=forestgreen4416044] (axis cs:688.5,0) rectangle (axis cs:689.5,98);
\draw[draw=none,fill=forestgreen4416044] (axis cs:689.5,0) rectangle (axis cs:690.5,98);
\draw[draw=none,fill=forestgreen4416044] (axis cs:690.5,0) rectangle (axis cs:691.5,98);
\draw[draw=none,fill=forestgreen4416044] (axis cs:691.5,0) rectangle (axis cs:692.5,98);
\draw[draw=none,fill=forestgreen4416044] (axis cs:692.5,0) rectangle (axis cs:693.5,98);
\draw[draw=none,fill=forestgreen4416044] (axis cs:693.5,0) rectangle (axis cs:694.5,98);
\draw[draw=none,fill=forestgreen4416044] (axis cs:694.5,0) rectangle (axis cs:695.5,98);
\draw[draw=none,fill=forestgreen4416044] (axis cs:695.5,0) rectangle (axis cs:696.5,98);
\draw[draw=none,fill=forestgreen4416044] (axis cs:696.5,0) rectangle (axis cs:697.5,98);
\draw[draw=none,fill=forestgreen4416044] (axis cs:697.5,0) rectangle (axis cs:698.5,98);
\draw[draw=none,fill=forestgreen4416044] (axis cs:698.5,0) rectangle (axis cs:699.5,98);
\draw[draw=none,fill=forestgreen4416044] (axis cs:699.5,0) rectangle (axis cs:700.5,98);
\draw[draw=none,fill=forestgreen4416044] (axis cs:700.5,0) rectangle (axis cs:701.5,98);
\draw[draw=none,fill=forestgreen4416044] (axis cs:701.5,0) rectangle (axis cs:702.5,98);
\draw[draw=none,fill=forestgreen4416044] (axis cs:702.5,0) rectangle (axis cs:703.5,98);
\draw[draw=none,fill=forestgreen4416044] (axis cs:703.5,0) rectangle (axis cs:704.5,98);
\draw[draw=none,fill=forestgreen4416044] (axis cs:704.5,0) rectangle (axis cs:705.5,98);
\draw[draw=none,fill=forestgreen4416044] (axis cs:705.5,0) rectangle (axis cs:706.5,98);
\draw[draw=none,fill=forestgreen4416044] (axis cs:706.5,0) rectangle (axis cs:707.5,98);
\draw[draw=none,fill=forestgreen4416044] (axis cs:707.5,0) rectangle (axis cs:708.5,98);
\draw[draw=none,fill=forestgreen4416044] (axis cs:708.5,0) rectangle (axis cs:709.5,98);
\draw[draw=none,fill=forestgreen4416044] (axis cs:709.5,0) rectangle (axis cs:710.5,98);
\draw[draw=none,fill=forestgreen4416044] (axis cs:710.5,0) rectangle (axis cs:711.5,98);
\draw[draw=none,fill=forestgreen4416044] (axis cs:711.5,0) rectangle (axis cs:712.5,98);
\draw[draw=none,fill=forestgreen4416044] (axis cs:712.5,0) rectangle (axis cs:713.5,98);
\draw[draw=none,fill=forestgreen4416044] (axis cs:713.5,0) rectangle (axis cs:714.5,98);
\draw[draw=none,fill=forestgreen4416044] (axis cs:714.5,0) rectangle (axis cs:715.5,98);
\draw[draw=none,fill=forestgreen4416044] (axis cs:715.5,0) rectangle (axis cs:716.5,98);
\draw[draw=none,fill=forestgreen4416044] (axis cs:716.5,0) rectangle (axis cs:717.5,98);
\draw[draw=none,fill=forestgreen4416044] (axis cs:717.5,0) rectangle (axis cs:718.5,98);
\draw[draw=none,fill=forestgreen4416044] (axis cs:718.5,0) rectangle (axis cs:719.5,98);
\draw[draw=none,fill=forestgreen4416044] (axis cs:719.5,0) rectangle (axis cs:720.5,98);
\draw[draw=none,fill=forestgreen4416044] (axis cs:720.5,0) rectangle (axis cs:721.5,98);
\draw[draw=none,fill=forestgreen4416044] (axis cs:721.5,0) rectangle (axis cs:722.5,98);
\draw[draw=none,fill=forestgreen4416044] (axis cs:722.5,0) rectangle (axis cs:723.5,98);
\draw[draw=none,fill=forestgreen4416044] (axis cs:723.5,0) rectangle (axis cs:724.5,98);
\draw[draw=none,fill=forestgreen4416044] (axis cs:724.5,0) rectangle (axis cs:725.5,98);
\draw[draw=none,fill=forestgreen4416044] (axis cs:725.5,0) rectangle (axis cs:726.5,98);
\draw[draw=none,fill=forestgreen4416044] (axis cs:726.5,0) rectangle (axis cs:727.5,98);
\draw[draw=none,fill=forestgreen4416044] (axis cs:727.5,0) rectangle (axis cs:728.5,98);
\draw[draw=none,fill=forestgreen4416044] (axis cs:728.5,0) rectangle (axis cs:729.5,98);
\draw[draw=none,fill=forestgreen4416044] (axis cs:729.5,0) rectangle (axis cs:730.5,98);
\draw[draw=none,fill=forestgreen4416044] (axis cs:730.5,0) rectangle (axis cs:731.5,98);
\draw[draw=none,fill=forestgreen4416044] (axis cs:731.5,0) rectangle (axis cs:732.5,98);
\draw[draw=none,fill=forestgreen4416044] (axis cs:732.5,0) rectangle (axis cs:733.5,98);
\draw[draw=none,fill=forestgreen4416044] (axis cs:733.5,0) rectangle (axis cs:734.5,98);
\draw[draw=none,fill=forestgreen4416044] (axis cs:734.5,0) rectangle (axis cs:735.5,98);
\draw[draw=none,fill=forestgreen4416044] (axis cs:735.5,0) rectangle (axis cs:736.5,98);
\draw[draw=none,fill=forestgreen4416044] (axis cs:736.5,0) rectangle (axis cs:737.5,98);
\draw[draw=none,fill=forestgreen4416044] (axis cs:737.5,0) rectangle (axis cs:738.5,98);
\draw[draw=none,fill=forestgreen4416044] (axis cs:738.5,0) rectangle (axis cs:739.5,98);
\draw[draw=none,fill=forestgreen4416044] (axis cs:739.5,0) rectangle (axis cs:740.5,98);
\draw[draw=none,fill=forestgreen4416044] (axis cs:740.5,0) rectangle (axis cs:741.5,98);
\draw[draw=none,fill=forestgreen4416044] (axis cs:741.5,0) rectangle (axis cs:742.5,98);
\draw[draw=none,fill=forestgreen4416044] (axis cs:742.5,0) rectangle (axis cs:743.5,98);
\draw[draw=none,fill=forestgreen4416044] (axis cs:743.5,0) rectangle (axis cs:744.5,98);
\draw[draw=none,fill=forestgreen4416044] (axis cs:744.5,0) rectangle (axis cs:745.5,98);
\draw[draw=none,fill=forestgreen4416044] (axis cs:745.5,0) rectangle (axis cs:746.5,98);
\draw[draw=none,fill=forestgreen4416044] (axis cs:746.5,0) rectangle (axis cs:747.5,98);
\draw[draw=none,fill=forestgreen4416044] (axis cs:747.5,0) rectangle (axis cs:748.5,98);
\draw[draw=none,fill=forestgreen4416044] (axis cs:748.5,0) rectangle (axis cs:749.5,98);
\draw[draw=none,fill=forestgreen4416044] (axis cs:749.5,0) rectangle (axis cs:750.5,98);
\draw[draw=none,fill=forestgreen4416044] (axis cs:750.5,0) rectangle (axis cs:751.5,98);
\draw[draw=none,fill=forestgreen4416044] (axis cs:751.5,0) rectangle (axis cs:752.5,98);
\draw[draw=none,fill=forestgreen4416044] (axis cs:752.5,0) rectangle (axis cs:753.5,98);
\draw[draw=none,fill=forestgreen4416044] (axis cs:753.5,0) rectangle (axis cs:754.5,98);
\draw[draw=none,fill=forestgreen4416044] (axis cs:754.5,0) rectangle (axis cs:755.5,98);
\draw[draw=none,fill=forestgreen4416044] (axis cs:755.5,0) rectangle (axis cs:756.5,98);
\draw[draw=none,fill=forestgreen4416044] (axis cs:756.5,0) rectangle (axis cs:757.5,98);
\draw[draw=none,fill=forestgreen4416044] (axis cs:757.5,0) rectangle (axis cs:758.5,98);
\draw[draw=none,fill=forestgreen4416044] (axis cs:758.5,0) rectangle (axis cs:759.5,98);
\draw[draw=none,fill=forestgreen4416044] (axis cs:759.5,0) rectangle (axis cs:760.5,98);
\draw[draw=none,fill=forestgreen4416044] (axis cs:760.5,0) rectangle (axis cs:761.5,98);
\draw[draw=none,fill=forestgreen4416044] (axis cs:761.5,0) rectangle (axis cs:762.5,98);
\draw[draw=none,fill=forestgreen4416044] (axis cs:762.5,0) rectangle (axis cs:763.5,98);
\draw[draw=none,fill=forestgreen4416044] (axis cs:763.5,0) rectangle (axis cs:764.5,98);
\draw[draw=none,fill=forestgreen4416044] (axis cs:764.5,0) rectangle (axis cs:765.5,98);
\draw[draw=none,fill=forestgreen4416044] (axis cs:765.5,0) rectangle (axis cs:766.5,98);
\draw[draw=none,fill=forestgreen4416044] (axis cs:766.5,0) rectangle (axis cs:767.5,98);
\draw[draw=none,fill=forestgreen4416044] (axis cs:767.5,0) rectangle (axis cs:768.5,98);
\draw[draw=none,fill=forestgreen4416044] (axis cs:768.5,0) rectangle (axis cs:769.5,98);
\draw[draw=none,fill=forestgreen4416044] (axis cs:769.5,0) rectangle (axis cs:770.5,98);
\draw[draw=none,fill=forestgreen4416044] (axis cs:770.5,0) rectangle (axis cs:771.5,98);
\draw[draw=none,fill=forestgreen4416044] (axis cs:771.5,0) rectangle (axis cs:772.5,98);
\draw[draw=none,fill=forestgreen4416044] (axis cs:772.5,0) rectangle (axis cs:773.5,98);
\draw[draw=none,fill=forestgreen4416044] (axis cs:773.5,0) rectangle (axis cs:774.5,98);
\draw[draw=none,fill=forestgreen4416044] (axis cs:774.5,0) rectangle (axis cs:775.5,98);
\draw[draw=none,fill=forestgreen4416044] (axis cs:775.5,0) rectangle (axis cs:776.5,98);
\draw[draw=none,fill=forestgreen4416044] (axis cs:776.5,0) rectangle (axis cs:777.5,98);
\draw[draw=none,fill=forestgreen4416044] (axis cs:777.5,0) rectangle (axis cs:778.5,98);
\draw[draw=none,fill=forestgreen4416044] (axis cs:778.5,0) rectangle (axis cs:779.5,98);
\draw[draw=none,fill=forestgreen4416044] (axis cs:779.5,0) rectangle (axis cs:780.5,98);
\draw[draw=none,fill=forestgreen4416044] (axis cs:780.5,0) rectangle (axis cs:781.5,98);
\draw[draw=none,fill=forestgreen4416044] (axis cs:781.5,0) rectangle (axis cs:782.5,98);
\draw[draw=none,fill=forestgreen4416044] (axis cs:782.5,0) rectangle (axis cs:783.5,98);
\draw[draw=none,fill=forestgreen4416044] (axis cs:783.5,0) rectangle (axis cs:784.5,98);
\draw[draw=none,fill=forestgreen4416044] (axis cs:784.5,0) rectangle (axis cs:785.5,98);
\draw[draw=none,fill=forestgreen4416044] (axis cs:785.5,0) rectangle (axis cs:786.5,98);
\end{axis}

\end{tikzpicture}
    \caption{Sample efficiency of LBH sampler at varying Levenshtein radii.}\label{fig:sample_efficiency}
  \end{figure}

  Next, measure the precision at various ranking cutoffs for varying wall-clock timeouts. Here, P@\{k=1, 5, 10, All\} indicates the percentage of syntax errors with a human repair of $\Delta=\{1, 2, 3, 4\}$ edits found in $\leq p$ seconds that were matched within the top-k results, using an n-gram likelihood model.

  \begin{figure}[h!]
%    \resizebox{.19\textwidth}{!}{% This file was created with tikzplotlib v0.10.1.
\begin{tikzpicture}

\definecolor{darkgray176}{RGB}{176,176,176}
\definecolor{darkviolet1270255}{RGB}{127,0,255}
\definecolor{deepskyblue3176236}{RGB}{3,176,236}
\definecolor{dodgerblue45123246}{RGB}{45,123,246}
\definecolor{royalblue8762253}{RGB}{87,62,253}

\begin{axis}[
tick align=outside,
tick pos=left,
title={\(\displaystyle \Delta\in[1,3]\) Repair Precision},
legend style={fill opacity=0.8, draw opacity=1, text opacity=1, legend columns=1, legend pos=north west},
x grid style={darkgray176},
xlabel={Seconds},
xmin=-0.5925, xmax=8.5925,
xtick style={color=black},
xtick={0,1,2,3,4,5,6,7,8},
xticklabels={1,2,3,4,5,6,7,8,9},
y grid style={darkgray176},
ylabel={Precision@k},
ymin=0, ymax=0.6993,
ytick style={color=black}
]
\draw[draw=none,fill=darkviolet1270255] (axis cs:-0.175,0) rectangle (axis cs:0.175,0.149);
\addlegendimage{ybar,ybar legend,draw=none,fill=darkviolet1270255}
\addlegendentry{P@All}

\draw[draw=none,fill=darkviolet1270255] (axis cs:0.825,0) rectangle (axis cs:1.175,0.303);
\draw[draw=none,fill=darkviolet1270255] (axis cs:1.825,0) rectangle (axis cs:2.175,0.412);
\draw[draw=none,fill=darkviolet1270255] (axis cs:2.825,0) rectangle (axis cs:3.175,0.501);
\draw[draw=none,fill=darkviolet1270255] (axis cs:3.825,0) rectangle (axis cs:4.175,0.569);
\draw[draw=none,fill=darkviolet1270255] (axis cs:4.825,0) rectangle (axis cs:5.175,0.605);
\draw[draw=none,fill=darkviolet1270255] (axis cs:5.825,0) rectangle (axis cs:6.175,0.63);
\draw[draw=none,fill=darkviolet1270255] (axis cs:6.825,0) rectangle (axis cs:7.175,0.645);
\draw[draw=none,fill=darkviolet1270255] (axis cs:7.825,0) rectangle (axis cs:8.175,0.666);
\draw[draw=none,fill=royalblue8762253] (axis cs:-0.175,0) rectangle (axis cs:0.175,0.128);
\addlegendimage{ybar,ybar legend,draw=none,fill=royalblue8762253}
\addlegendentry{P@10}

\draw[draw=none,fill=royalblue8762253] (axis cs:0.825,0) rectangle (axis cs:1.175,0.239);
\draw[draw=none,fill=royalblue8762253] (axis cs:1.825,0) rectangle (axis cs:2.175,0.323);
\draw[draw=none,fill=royalblue8762253] (axis cs:2.825,0) rectangle (axis cs:3.175,0.374);
\draw[draw=none,fill=royalblue8762253] (axis cs:3.825,0) rectangle (axis cs:4.175,0.421);
\draw[draw=none,fill=royalblue8762253] (axis cs:4.825,0) rectangle (axis cs:5.175,0.446);
\draw[draw=none,fill=royalblue8762253] (axis cs:5.825,0) rectangle (axis cs:6.175,0.459);
\draw[draw=none,fill=royalblue8762253] (axis cs:6.825,0) rectangle (axis cs:7.175,0.469);
\draw[draw=none,fill=royalblue8762253] (axis cs:7.825,0) rectangle (axis cs:8.175,0.477);
\draw[draw=none,fill=dodgerblue45123246] (axis cs:-0.175,0) rectangle (axis cs:0.175,0.11);
\addlegendimage{ybar,ybar legend,draw=none,fill=dodgerblue45123246}
\addlegendentry{P@5}

\draw[draw=none,fill=dodgerblue45123246] (axis cs:0.825,0) rectangle (axis cs:1.175,0.206);
\draw[draw=none,fill=dodgerblue45123246] (axis cs:1.825,0) rectangle (axis cs:2.175,0.282);
\draw[draw=none,fill=dodgerblue45123246] (axis cs:2.825,0) rectangle (axis cs:3.175,0.325);
\draw[draw=none,fill=dodgerblue45123246] (axis cs:3.825,0) rectangle (axis cs:4.175,0.369);
\draw[draw=none,fill=dodgerblue45123246] (axis cs:4.825,0) rectangle (axis cs:5.175,0.392);
\draw[draw=none,fill=dodgerblue45123246] (axis cs:5.825,0) rectangle (axis cs:6.175,0.405);
\draw[draw=none,fill=dodgerblue45123246] (axis cs:6.825,0) rectangle (axis cs:7.175,0.414);
\draw[draw=none,fill=dodgerblue45123246] (axis cs:7.825,0) rectangle (axis cs:8.175,0.422);
\draw[draw=none,fill=deepskyblue3176236] (axis cs:-0.175,0) rectangle (axis cs:0.175,0.092);
\addlegendimage{ybar,ybar legend,draw=none,fill=deepskyblue3176236}
\addlegendentry{P@1}

\draw[draw=none,fill=deepskyblue3176236] (axis cs:0.825,0) rectangle (axis cs:1.175,0.18);
\draw[draw=none,fill=deepskyblue3176236] (axis cs:1.825,0) rectangle (axis cs:2.175,0.244);
\draw[draw=none,fill=deepskyblue3176236] (axis cs:2.825,0) rectangle (axis cs:3.175,0.285);
\draw[draw=none,fill=deepskyblue3176236] (axis cs:3.825,0) rectangle (axis cs:4.175,0.324);
\draw[draw=none,fill=deepskyblue3176236] (axis cs:4.825,0) rectangle (axis cs:5.175,0.343);
\draw[draw=none,fill=deepskyblue3176236] (axis cs:5.825,0) rectangle (axis cs:6.175,0.356);
\draw[draw=none,fill=deepskyblue3176236] (axis cs:6.825,0) rectangle (axis cs:7.175,0.365);
\draw[draw=none,fill=deepskyblue3176236] (axis cs:7.825,0) rectangle (axis cs:8.175,0.37);
\addplot [red, thick] coordinates {(-0.8925,0.29) (14.8925,0.29)};
\end{axis}

\end{tikzpicture}
}
    \resizebox{.24\textwidth}{!}{% This file was created with tikzplotlib v0.10.1.
\begin{tikzpicture}
\begin{axis}[
legend cell align={left},
legend style={fill opacity=0.8, draw opacity=1, text opacity=1, draw=lightgray204, legend columns=-1, legend pos=north west},
tick align=outside,
tick pos=left,
axis lines*=left,
title={\(\displaystyle \Delta=1\) Repair Precision},
x grid style={darkgray176},
xlabel={Seconds},
xmin=-0.5925, xmax=8.5925,
xtick style={color=black},
xtick={0,1,2,3,4,5,6,7,8},
xticklabels={1,2,3,4,5,6,7,8,9},
y grid style={darkgray176},
ylabel={Precision@k},
ymin=0, ymax=1.0311,
ytick style={color=black}
]
\draw[draw=none,fill=darkviolet1270255] (axis cs:-0.175,0) rectangle (axis cs:0.175,0.546);
\addlegendimage{empty legend}
\addlegendentry{P@}
\addlegendimage{ybar,ybar legend,draw=none,fill=darkviolet1270255}
\addlegendentry{All}
\addlegendimage{ybar,ybar legend,draw=none,fill=royalblue8762253}
\addlegendentry{10}
\addlegendimage{ybar,ybar legend,draw=none,fill=dodgerblue45123246}
\addlegendentry{5}
\addlegendimage{ybar,ybar legend,draw=none,fill=deepskyblue3176236}
\addlegendentry{1}

\draw[draw=none,fill=darkviolet1270255] (axis cs:0.825,0) rectangle (axis cs:1.175,0.752);
\draw[draw=none,fill=darkviolet1270255] (axis cs:1.825,0) rectangle (axis cs:2.175,0.88);
\draw[draw=none,fill=darkviolet1270255] (axis cs:2.825,0) rectangle (axis cs:3.175,0.908);
\draw[draw=none,fill=darkviolet1270255] (axis cs:3.825,0) rectangle (axis cs:4.175,0.933);
\draw[draw=none,fill=darkviolet1270255] (axis cs:4.825,0) rectangle (axis cs:5.175,0.954);
\draw[draw=none,fill=darkviolet1270255] (axis cs:5.825,0) rectangle (axis cs:6.175,0.966);
\draw[draw=none,fill=darkviolet1270255] (axis cs:6.825,0) rectangle (axis cs:7.175,0.975);
\draw[draw=none,fill=darkviolet1270255] (axis cs:7.825,0) rectangle (axis cs:8.175,0.982);
\draw[draw=none,fill=royalblue8762253] (axis cs:-0.175,0) rectangle (axis cs:0.175,0.546);

\draw[draw=none,fill=royalblue8762253] (axis cs:0.825,0) rectangle (axis cs:1.175,0.752);
\draw[draw=none,fill=royalblue8762253] (axis cs:1.825,0) rectangle (axis cs:2.175,0.88);
\draw[draw=none,fill=royalblue8762253] (axis cs:2.825,0) rectangle (axis cs:3.175,0.908);
\draw[draw=none,fill=royalblue8762253] (axis cs:3.825,0) rectangle (axis cs:4.175,0.933);
\draw[draw=none,fill=royalblue8762253] (axis cs:4.825,0) rectangle (axis cs:5.175,0.954);
\draw[draw=none,fill=royalblue8762253] (axis cs:5.825,0) rectangle (axis cs:6.175,0.966);
\draw[draw=none,fill=royalblue8762253] (axis cs:6.825,0) rectangle (axis cs:7.175,0.975);
\draw[draw=none,fill=royalblue8762253] (axis cs:7.825,0) rectangle (axis cs:8.175,0.982);
\draw[draw=none,fill=dodgerblue45123246] (axis cs:-0.175,0) rectangle (axis cs:0.175,0.546);

\draw[draw=none,fill=dodgerblue45123246] (axis cs:0.825,0) rectangle (axis cs:1.175,0.752);
\draw[draw=none,fill=dodgerblue45123246] (axis cs:1.825,0) rectangle (axis cs:2.175,0.88);
\draw[draw=none,fill=dodgerblue45123246] (axis cs:2.825,0) rectangle (axis cs:3.175,0.908);
\draw[draw=none,fill=dodgerblue45123246] (axis cs:3.825,0) rectangle (axis cs:4.175,0.933);
\draw[draw=none,fill=dodgerblue45123246] (axis cs:4.825,0) rectangle (axis cs:5.175,0.954);
\draw[draw=none,fill=dodgerblue45123246] (axis cs:5.825,0) rectangle (axis cs:6.175,0.966);
\draw[draw=none,fill=dodgerblue45123246] (axis cs:6.825,0) rectangle (axis cs:7.175,0.975);
\draw[draw=none,fill=dodgerblue45123246] (axis cs:7.825,0) rectangle (axis cs:8.175,0.982);
\draw[draw=none,fill=deepskyblue3176236] (axis cs:-0.175,0) rectangle (axis cs:0.175,0.546);

\draw[draw=none,fill=deepskyblue3176236] (axis cs:0.825,0) rectangle (axis cs:1.175,0.752);
\draw[draw=none,fill=deepskyblue3176236] (axis cs:1.825,0) rectangle (axis cs:2.175,0.88);
\draw[draw=none,fill=deepskyblue3176236] (axis cs:2.825,0) rectangle (axis cs:3.175,0.908);
\draw[draw=none,fill=deepskyblue3176236] (axis cs:3.825,0) rectangle (axis cs:4.175,0.933);
\draw[draw=none,fill=deepskyblue3176236] (axis cs:4.825,0) rectangle (axis cs:5.175,0.954);
\draw[draw=none,fill=deepskyblue3176236] (axis cs:5.825,0) rectangle (axis cs:6.175,0.966);
\draw[draw=none,fill=deepskyblue3176236] (axis cs:6.825,0) rectangle (axis cs:7.175,0.975);
\draw[draw=none,fill=deepskyblue3176236] (axis cs:7.825,0) rectangle (axis cs:8.175,0.982);
\addplot [red, thick] coordinates {(-0.8925,0.39) (14.8925,0.39)};
\addplot [orange, thick] coordinates {(-0.8925,0.24) (14.8925,0.24)};
\end{axis}

\end{tikzpicture}
}
    \resizebox{.24\textwidth}{!}{% This file was created with tikzplotlib v0.10.1.
\begin{tikzpicture}
\begin{axis}[
tick align=outside,
tick pos=left,
axis lines=left,
title={\(\displaystyle \Delta=2\) Repair Precision},
x grid style={darkgray176},
xlabel={Seconds},
xmin=-0.5925, xmax=8.5925,
xtick style={color=black},
xtick={0,1,2,3,4,5,6,7,8},
xticklabels={10,20,30,40,50,60,70,80,90},
y grid style={darkgray176},
ylabel={Precision@k},
ymin=0, ymax=0.75285,
ytick style={color=black}
]

\draw[draw=none,fill=darkviolet1270255] (axis cs:-0.175,0) rectangle (axis cs:0.175,0.182);

\draw[draw=none,fill=darkviolet1270255] (axis cs:0.825,0) rectangle (axis cs:1.175,0.326);
\draw[draw=none,fill=darkviolet1270255] (axis cs:1.825,0) rectangle (axis cs:2.175,0.455);
\draw[draw=none,fill=darkviolet1270255] (axis cs:2.825,0) rectangle (axis cs:3.175,0.554);
\draw[draw=none,fill=darkviolet1270255] (axis cs:3.825,0) rectangle (axis cs:4.175,0.603);
\draw[draw=none,fill=darkviolet1270255] (axis cs:4.825,0) rectangle (axis cs:5.175,0.649);
\draw[draw=none,fill=darkviolet1270255] (axis cs:5.825,0) rectangle (axis cs:6.175,0.677);
\draw[draw=none,fill=darkviolet1270255] (axis cs:6.825,0) rectangle (axis cs:7.175,0.702);
\draw[draw=none,fill=darkviolet1270255] (axis cs:7.825,0) rectangle (axis cs:8.175,0.717);
\draw[draw=none,fill=royalblue8762253] (axis cs:-0.175,0) rectangle (axis cs:0.175,0.102);

\draw[draw=none,fill=royalblue8762253] (axis cs:0.825,0) rectangle (axis cs:1.175,0.178);
\draw[draw=none,fill=royalblue8762253] (axis cs:1.825,0) rectangle (axis cs:2.175,0.243);
\draw[draw=none,fill=royalblue8762253] (axis cs:2.825,0) rectangle (axis cs:3.175,0.289);
\draw[draw=none,fill=royalblue8762253] (axis cs:3.825,0) rectangle (axis cs:4.175,0.317);
\draw[draw=none,fill=royalblue8762253] (axis cs:4.825,0) rectangle (axis cs:5.175,0.335);
\draw[draw=none,fill=royalblue8762253] (axis cs:5.825,0) rectangle (axis cs:6.175,0.348);
\draw[draw=none,fill=royalblue8762253] (axis cs:6.825,0) rectangle (axis cs:7.175,0.357);
\draw[draw=none,fill=royalblue8762253] (axis cs:7.825,0) rectangle (axis cs:8.175,0.366);
\draw[draw=none,fill=dodgerblue45123246] (axis cs:-0.175,0) rectangle (axis cs:0.175,0.102);

\draw[draw=none,fill=dodgerblue45123246] (axis cs:0.825,0) rectangle (axis cs:1.175,0.178);
\draw[draw=none,fill=dodgerblue45123246] (axis cs:1.825,0) rectangle (axis cs:2.175,0.237);
\draw[draw=none,fill=dodgerblue45123246] (axis cs:2.825,0) rectangle (axis cs:3.175,0.283);
\draw[draw=none,fill=dodgerblue45123246] (axis cs:3.825,0) rectangle (axis cs:4.175,0.311);
\draw[draw=none,fill=dodgerblue45123246] (axis cs:4.825,0) rectangle (axis cs:5.175,0.329);
\draw[draw=none,fill=dodgerblue45123246] (axis cs:5.825,0) rectangle (axis cs:6.175,0.342);
\draw[draw=none,fill=dodgerblue45123246] (axis cs:6.825,0) rectangle (axis cs:7.175,0.351);
\draw[draw=none,fill=dodgerblue45123246] (axis cs:7.825,0) rectangle (axis cs:8.175,0.36);
\draw[draw=none,fill=deepskyblue3176236] (axis cs:-0.175,0) rectangle (axis cs:0.175,0.092);

\draw[draw=none,fill=deepskyblue3176236] (axis cs:0.825,0) rectangle (axis cs:1.175,0.16);
\draw[draw=none,fill=deepskyblue3176236] (axis cs:1.825,0) rectangle (axis cs:2.175,0.218);
\draw[draw=none,fill=deepskyblue3176236] (axis cs:2.825,0) rectangle (axis cs:3.175,0.258);
\draw[draw=none,fill=deepskyblue3176236] (axis cs:3.825,0) rectangle (axis cs:4.175,0.283);
\draw[draw=none,fill=deepskyblue3176236] (axis cs:4.825,0) rectangle (axis cs:5.175,0.302);
\draw[draw=none,fill=deepskyblue3176236] (axis cs:5.825,0) rectangle (axis cs:6.175,0.311);
\draw[draw=none,fill=deepskyblue3176236] (axis cs:6.825,0) rectangle (axis cs:7.175,0.32);
\draw[draw=none,fill=deepskyblue3176236] (axis cs:7.825,0) rectangle (axis cs:8.175,0.329);
\addplot [red, thick] coordinates {(-0.8925,0.15) (14.8925,0.15)};
\addplot [orange, thick] coordinates {(-0.8925,0.10) (14.8925,0.10)};
\end{axis}

\end{tikzpicture}
}
    \resizebox{.24\textwidth}{!}{% This file was created with tikzplotlib v0.10.1.
\begin{tikzpicture}

\definecolor{darkgray176}{RGB}{176,176,176}
\definecolor{darkviolet1270255}{RGB}{127,0,255}
\definecolor{deepskyblue3176236}{RGB}{3,176,236}
\definecolor{dodgerblue45123246}{RGB}{45,123,246}
\definecolor{royalblue8762253}{RGB}{87,62,253}

\begin{axis}[
tick align=outside,
tick pos=left,
axis lines=left,
title={\(\displaystyle \Delta=3\) Repair Precision},
legend style={fill opacity=0.8, draw opacity=1, text opacity=1, legend columns=1, legend pos=north west},
x grid style={darkgray176},
xlabel={Seconds},
xmin=-0.6425, xmax=9.6425,
xtick style={color=black},
xtick={0,1,2,3,4,5,6,7,8,9},
xticklabels={6,12,18,24,30,36,42,48,54,60},
y grid style={darkgray176},
ylabel={Precision@k},
ymin=0, ymax=0.5754,
ytick style={color=black}
]
\draw[draw=none,fill=darkviolet1270255] (axis cs:-0.175,0) rectangle (axis cs:0.175,0.014);
\addlegendimage{ybar,ybar legend,draw=none,fill=darkviolet1270255}
\addlegendentry{P@All}

\draw[draw=none,fill=darkviolet1270255] (axis cs:0.825,0) rectangle (axis cs:1.175,0.068);
\draw[draw=none,fill=darkviolet1270255] (axis cs:1.825,0) rectangle (axis cs:2.175,0.164);
\draw[draw=none,fill=darkviolet1270255] (axis cs:2.825,0) rectangle (axis cs:3.175,0.233);
\draw[draw=none,fill=darkviolet1270255] (axis cs:3.825,0) rectangle (axis cs:4.175,0.342);
\draw[draw=none,fill=darkviolet1270255] (axis cs:4.825,0) rectangle (axis cs:5.175,0.411);
\draw[draw=none,fill=darkviolet1270255] (axis cs:5.825,0) rectangle (axis cs:6.175,0.438);
\draw[draw=none,fill=darkviolet1270255] (axis cs:6.825,0) rectangle (axis cs:7.175,0.493);
\draw[draw=none,fill=darkviolet1270255] (axis cs:7.825,0) rectangle (axis cs:8.175,0.507);
\draw[draw=none,fill=darkviolet1270255] (axis cs:8.825,0) rectangle (axis cs:9.175,0.548);
\draw[draw=none,fill=royalblue8762253] (axis cs:-0.175,0) rectangle (axis cs:0.175,0);
\addlegendimage{ybar,ybar legend,draw=none,fill=royalblue8762253}
\addlegendentry{P@10}

\draw[draw=none,fill=royalblue8762253] (axis cs:0.825,0) rectangle (axis cs:1.175,0.014);
\draw[draw=none,fill=royalblue8762253] (axis cs:1.825,0) rectangle (axis cs:2.175,0.055);
\draw[draw=none,fill=royalblue8762253] (axis cs:2.825,0) rectangle (axis cs:3.175,0.068);
\draw[draw=none,fill=royalblue8762253] (axis cs:3.825,0) rectangle (axis cs:4.175,0.11);
\draw[draw=none,fill=royalblue8762253] (axis cs:4.825,0) rectangle (axis cs:5.175,0.137);
\draw[draw=none,fill=royalblue8762253] (axis cs:5.825,0) rectangle (axis cs:6.175,0.151);
\draw[draw=none,fill=royalblue8762253] (axis cs:6.825,0) rectangle (axis cs:7.175,0.164);
\draw[draw=none,fill=royalblue8762253] (axis cs:7.825,0) rectangle (axis cs:8.175,0.178);
\draw[draw=none,fill=royalblue8762253] (axis cs:8.825,0) rectangle (axis cs:9.175,0.178);
\draw[draw=none,fill=dodgerblue45123246] (axis cs:-0.175,0) rectangle (axis cs:0.175,0);
\addlegendimage{ybar,ybar legend,draw=none,fill=dodgerblue45123246}
\addlegendentry{P@5}

\draw[draw=none,fill=dodgerblue45123246] (axis cs:0.825,0) rectangle (axis cs:1.175,0.014);
\draw[draw=none,fill=dodgerblue45123246] (axis cs:1.825,0) rectangle (axis cs:2.175,0.055);
\draw[draw=none,fill=dodgerblue45123246] (axis cs:2.825,0) rectangle (axis cs:3.175,0.068);
\draw[draw=none,fill=dodgerblue45123246] (axis cs:3.825,0) rectangle (axis cs:4.175,0.11);
\draw[draw=none,fill=dodgerblue45123246] (axis cs:4.825,0) rectangle (axis cs:5.175,0.123);
\draw[draw=none,fill=dodgerblue45123246] (axis cs:5.825,0) rectangle (axis cs:6.175,0.137);
\draw[draw=none,fill=dodgerblue45123246] (axis cs:6.825,0) rectangle (axis cs:7.175,0.151);
\draw[draw=none,fill=dodgerblue45123246] (axis cs:7.825,0) rectangle (axis cs:8.175,0.164);
\draw[draw=none,fill=dodgerblue45123246] (axis cs:8.825,0) rectangle (axis cs:9.175,0.164);
\draw[draw=none,fill=deepskyblue3176236] (axis cs:-0.175,0) rectangle (axis cs:0.175,0);
\addlegendimage{ybar,ybar legend,draw=none,fill=deepskyblue3176236}
\addlegendentry{P@1}

\draw[draw=none,fill=deepskyblue3176236] (axis cs:0.825,0) rectangle (axis cs:1.175,0.014);
\draw[draw=none,fill=deepskyblue3176236] (axis cs:1.825,0) rectangle (axis cs:2.175,0.055);
\draw[draw=none,fill=deepskyblue3176236] (axis cs:2.825,0) rectangle (axis cs:3.175,0.068);
\draw[draw=none,fill=deepskyblue3176236] (axis cs:3.825,0) rectangle (axis cs:4.175,0.096);
\draw[draw=none,fill=deepskyblue3176236] (axis cs:4.825,0) rectangle (axis cs:5.175,0.11);
\draw[draw=none,fill=deepskyblue3176236] (axis cs:5.825,0) rectangle (axis cs:6.175,0.123);
\draw[draw=none,fill=deepskyblue3176236] (axis cs:6.825,0) rectangle (axis cs:7.175,0.137);
\draw[draw=none,fill=deepskyblue3176236] (axis cs:7.825,0) rectangle (axis cs:8.175,0.151);
\draw[draw=none,fill=deepskyblue3176236] (axis cs:8.825,0) rectangle (axis cs:9.175,0.151);
\addplot [red, thick] coordinates {(-0.8925,0.07) (14.8925,0.07)};
\end{axis}

\end{tikzpicture}
}
    \resizebox{.24\textwidth}{!}{% This file was created with tikzplotlib v0.10.1.
\begin{tikzpicture}

\definecolor{darkgray176}{RGB}{176,176,176}
\definecolor{darkviolet1270255}{RGB}{127,0,255}
\definecolor{deepskyblue3176236}{RGB}{3,176,236}
\definecolor{dodgerblue45123246}{RGB}{45,123,246}
\definecolor{royalblue8762253}{RGB}{87,62,253}

\begin{axis}[
tick align=outside,
tick pos=left,
axis lines=left,
title={\(\displaystyle \Delta=4\) Repair Precision},
legend style={fill opacity=0.8, draw opacity=1, text opacity=1, legend columns=1, legend pos=north west},
x grid style={darkgray176},
xlabel={Seconds},
xmin=-0.5925, xmax=8.5925,
xtick style={color=black},
xtick={0,1,2,3,4,5,6,7,8},
xticklabels={10,20,30,40,50,60,70,80,90},
y grid style={darkgray176},
ylabel={Precision@k},
ymin=0, ymax=0.54915,
ytick style={color=black}
]
\draw[draw=none,fill=darkviolet1270255] (axis cs:-0.175,0) rectangle (axis cs:0.175,0.123);
\addlegendimage{ybar,ybar legend,draw=none,fill=darkviolet1270255}
\addlegendentry{P@All}

\draw[draw=none,fill=darkviolet1270255] (axis cs:0.825,0) rectangle (axis cs:1.175,0.323);
\draw[draw=none,fill=darkviolet1270255] (axis cs:1.825,0) rectangle (axis cs:2.175,0.415);
\draw[draw=none,fill=darkviolet1270255] (axis cs:2.825,0) rectangle (axis cs:3.175,0.431);
\draw[draw=none,fill=darkviolet1270255] (axis cs:3.825,0) rectangle (axis cs:4.175,0.477);
\draw[draw=none,fill=darkviolet1270255] (axis cs:4.825,0) rectangle (axis cs:5.175,0.508);
\draw[draw=none,fill=darkviolet1270255] (axis cs:5.825,0) rectangle (axis cs:6.175,0.523);
\draw[draw=none,fill=darkviolet1270255] (axis cs:6.825,0) rectangle (axis cs:7.175,0.523);
\draw[draw=none,fill=darkviolet1270255] (axis cs:7.825,0) rectangle (axis cs:8.175,0.523);
\draw[draw=none,fill=royalblue8762253] (axis cs:-0.175,0) rectangle (axis cs:0.175,0.092);
\addlegendimage{ybar,ybar legend,draw=none,fill=royalblue8762253}
\addlegendentry{P@10}

\draw[draw=none,fill=royalblue8762253] (axis cs:0.825,0) rectangle (axis cs:1.175,0.169);
\draw[draw=none,fill=royalblue8762253] (axis cs:1.825,0) rectangle (axis cs:2.175,0.215);
\draw[draw=none,fill=royalblue8762253] (axis cs:2.825,0) rectangle (axis cs:3.175,0.231);
\draw[draw=none,fill=royalblue8762253] (axis cs:3.825,0) rectangle (axis cs:4.175,0.277);
\draw[draw=none,fill=royalblue8762253] (axis cs:4.825,0) rectangle (axis cs:5.175,0.308);
\draw[draw=none,fill=royalblue8762253] (axis cs:5.825,0) rectangle (axis cs:6.175,0.323);
\draw[draw=none,fill=royalblue8762253] (axis cs:6.825,0) rectangle (axis cs:7.175,0.323);
\draw[draw=none,fill=royalblue8762253] (axis cs:7.825,0) rectangle (axis cs:8.175,0.323);
\draw[draw=none,fill=dodgerblue45123246] (axis cs:-0.175,0) rectangle (axis cs:0.175,0.092);
\addlegendimage{ybar,ybar legend,draw=none,fill=dodgerblue45123246}
\addlegendentry{P@5}

\draw[draw=none,fill=dodgerblue45123246] (axis cs:0.825,0) rectangle (axis cs:1.175,0.169);
\draw[draw=none,fill=dodgerblue45123246] (axis cs:1.825,0) rectangle (axis cs:2.175,0.215);
\draw[draw=none,fill=dodgerblue45123246] (axis cs:2.825,0) rectangle (axis cs:3.175,0.231);
\draw[draw=none,fill=dodgerblue45123246] (axis cs:3.825,0) rectangle (axis cs:4.175,0.277);
\draw[draw=none,fill=dodgerblue45123246] (axis cs:4.825,0) rectangle (axis cs:5.175,0.308);
\draw[draw=none,fill=dodgerblue45123246] (axis cs:5.825,0) rectangle (axis cs:6.175,0.323);
\draw[draw=none,fill=dodgerblue45123246] (axis cs:6.825,0) rectangle (axis cs:7.175,0.323);
\draw[draw=none,fill=dodgerblue45123246] (axis cs:7.825,0) rectangle (axis cs:8.175,0.323);
\draw[draw=none,fill=deepskyblue3176236] (axis cs:-0.175,0) rectangle (axis cs:0.175,0.077);
\addlegendimage{ybar,ybar legend,draw=none,fill=deepskyblue3176236}
\addlegendentry{P@1}

\draw[draw=none,fill=deepskyblue3176236] (axis cs:0.825,0) rectangle (axis cs:1.175,0.138);
\draw[draw=none,fill=deepskyblue3176236] (axis cs:1.825,0) rectangle (axis cs:2.175,0.185);
\draw[draw=none,fill=deepskyblue3176236] (axis cs:2.825,0) rectangle (axis cs:3.175,0.2);
\draw[draw=none,fill=deepskyblue3176236] (axis cs:3.825,0) rectangle (axis cs:4.175,0.246);
\draw[draw=none,fill=deepskyblue3176236] (axis cs:4.825,0) rectangle (axis cs:5.175,0.277);
\draw[draw=none,fill=deepskyblue3176236] (axis cs:5.825,0) rectangle (axis cs:6.175,0.292);
\draw[draw=none,fill=deepskyblue3176236] (axis cs:6.825,0) rectangle (axis cs:7.175,0.292);
\draw[draw=none,fill=deepskyblue3176236] (axis cs:7.825,0) rectangle (axis cs:8.175,0.292);
\addplot [red, thick] coordinates {(-0.8925,0.10) (14.8925,0.10)};
\end{axis}

\end{tikzpicture}
}
%    \resizebox{.24\textwidth}{!}{% This file was created with tikzplotlib v0.10.1.
\begin{tikzpicture}

\definecolor{darkgray176}{RGB}{176,176,176}
\definecolor{darkviolet1270255}{RGB}{127,0,255}
\definecolor{deepskyblue3176236}{RGB}{3,176,236}
\definecolor{dodgerblue45123246}{RGB}{45,123,246}
\definecolor{royalblue8762253}{RGB}{87,62,253}

\begin{axis}[
tick align=outside,
tick pos=left,
axis lines=left,
title={\(\displaystyle \Delta=4\) Repair Precision},
legend style={fill opacity=0.8, draw opacity=1, text opacity=1, legend columns=1, legend pos=north west},
x grid style={darkgray176},
xlabel={Seconds},
xmin=-0.5925, xmax=8.5925,
xtick={0,1,2,3,4,5,6,7,8},
xticklabels={10,20,30,40,50,60,70,80,90},
y grid style={darkgray176},
ylabel={Precision@k},
ymin=0, ymax=0.54915,
]
\draw[draw=none,fill=darkviolet1270255] (axis cs:-0.175,0) rectangle (axis cs:0.175,0.123);
\addlegendimage{ybar,ybar legend,draw=none,fill=darkviolet1270255}
\addlegendentry{P@All}

\draw[draw=none,fill=darkviolet1270255] (axis cs:0.825,0) rectangle (axis cs:1.175,0.323);
\draw[draw=none,fill=darkviolet1270255] (axis cs:1.825,0) rectangle (axis cs:2.175,0.415);
\draw[draw=none,fill=darkviolet1270255] (axis cs:2.825,0) rectangle (axis cs:3.175,0.431);
\draw[draw=none,fill=darkviolet1270255] (axis cs:3.825,0) rectangle (axis cs:4.175,0.477);
\draw[draw=none,fill=darkviolet1270255] (axis cs:4.825,0) rectangle (axis cs:5.175,0.508);
\draw[draw=none,fill=darkviolet1270255] (axis cs:5.825,0) rectangle (axis cs:6.175,0.523);
\draw[draw=none,fill=darkviolet1270255] (axis cs:6.825,0) rectangle (axis cs:7.175,0.523);
\draw[draw=none,fill=darkviolet1270255] (axis cs:7.825,0) rectangle (axis cs:8.175,0.523);
\draw[draw=none,fill=royalblue8762253] (axis cs:-0.175,0) rectangle (axis cs:0.175,0.092);
\addlegendimage{ybar,ybar legend,draw=none,fill=royalblue8762253}
\addlegendentry{P@10}

\draw[draw=none,fill=royalblue8762253] (axis cs:0.825,0) rectangle (axis cs:1.175,0.169);
\draw[draw=none,fill=royalblue8762253] (axis cs:1.825,0) rectangle (axis cs:2.175,0.215);
\draw[draw=none,fill=royalblue8762253] (axis cs:2.825,0) rectangle (axis cs:3.175,0.231);
\draw[draw=none,fill=royalblue8762253] (axis cs:3.825,0) rectangle (axis cs:4.175,0.277);
\draw[draw=none,fill=royalblue8762253] (axis cs:4.825,0) rectangle (axis cs:5.175,0.308);
\draw[draw=none,fill=royalblue8762253] (axis cs:5.825,0) rectangle (axis cs:6.175,0.323);
\draw[draw=none,fill=royalblue8762253] (axis cs:6.825,0) rectangle (axis cs:7.175,0.323);
\draw[draw=none,fill=royalblue8762253] (axis cs:7.825,0) rectangle (axis cs:8.175,0.323);
\draw[draw=none,fill=dodgerblue45123246] (axis cs:-0.175,0) rectangle (axis cs:0.175,0.092);
\addlegendimage{ybar,ybar legend,draw=none,fill=dodgerblue45123246}
\addlegendentry{P@5}

\draw[draw=none,fill=dodgerblue45123246] (axis cs:0.825,0) rectangle (axis cs:1.175,0.169);
\draw[draw=none,fill=dodgerblue45123246] (axis cs:1.825,0) rectangle (axis cs:2.175,0.215);
\draw[draw=none,fill=dodgerblue45123246] (axis cs:2.825,0) rectangle (axis cs:3.175,0.231);
\draw[draw=none,fill=dodgerblue45123246] (axis cs:3.825,0) rectangle (axis cs:4.175,0.277);
\draw[draw=none,fill=dodgerblue45123246] (axis cs:4.825,0) rectangle (axis cs:5.175,0.308);
\draw[draw=none,fill=dodgerblue45123246] (axis cs:5.825,0) rectangle (axis cs:6.175,0.323);
\draw[draw=none,fill=dodgerblue45123246] (axis cs:6.825,0) rectangle (axis cs:7.175,0.323);
\draw[draw=none,fill=dodgerblue45123246] (axis cs:7.825,0) rectangle (axis cs:8.175,0.323);
\draw[draw=none,fill=deepskyblue3176236] (axis cs:-0.175,0) rectangle (axis cs:0.175,0.077);
\addlegendimage{ybar,ybar legend,draw=none,fill=deepskyblue3176236}
\addlegendentry{P@1}

\draw[draw=none,fill=deepskyblue3176236] (axis cs:0.825,0) rectangle (axis cs:1.175,0.138);
\draw[draw=none,fill=deepskyblue3176236] (axis cs:1.825,0) rectangle (axis cs:2.175,0.185);
\draw[draw=none,fill=deepskyblue3176236] (axis cs:2.825,0) rectangle (axis cs:3.175,0.2);
\draw[draw=none,fill=deepskyblue3176236] (axis cs:3.825,0) rectangle (axis cs:4.175,0.246);
\draw[draw=none,fill=deepskyblue3176236] (axis cs:4.825,0) rectangle (axis cs:5.175,0.277);
\draw[draw=none,fill=deepskyblue3176236] (axis cs:5.825,0) rectangle (axis cs:6.175,0.292);
\draw[draw=none,fill=deepskyblue3176236] (axis cs:6.825,0) rectangle (axis cs:7.175,0.292);
\draw[draw=none,fill=deepskyblue3176236] (axis cs:7.825,0) rectangle (axis cs:8.175,0.292);
\addplot [red, thick] coordinates {(-0.8925,0.10) (14.8925,0.10)};
\end{axis}

\end{tikzpicture}
}
%\resizebox{.3\textwidth}{!}{% This file was created with tikzplotlib v0.10.1.
\begin{tikzpicture}

\definecolor{darkgray176}{RGB}{176,176,176}
\definecolor{darkviolet1270255}{RGB}{127,0,255}
\definecolor{deepskyblue3176236}{RGB}{3,176,236}
\definecolor{dodgerblue45123246}{RGB}{45,123,246}
\definecolor{lightgray204}{RGB}{204,204,204}
\definecolor{royalblue8762253}{RGB}{87,62,253}

\begin{axis}[
legend cell align={left},
legend style={fill opacity=0.8, draw opacity=1, text opacity=1, draw=lightgray204, legend columns=-1, legend pos=north west},
tick align=outside,
tick pos=left,
title={$\Delta=1$ Repair Precision},
x grid style={darkgray176},
xlabel={Seconds},
xmin=-0.4925, xmax=6.4925,
xtick style={color=black},
xtick={0,1,2,3,4,5,6},
xticklabels={20,60,100,140,180,220,260},
y grid style={darkgray176},
ylabel={\phantom{Precision@k}},
ymin=0, ymax=0.77595,
ytick style={color=black}
]
\addlegendimage{empty legend}
\addlegendentry{P@}
\draw[draw=none,fill=darkviolet1270255] (axis cs:-0.175,0) rectangle (axis cs:0.175,0.145);
\addlegendimage{ybar,ybar legend,draw=none,fill=darkviolet1270255}
\addlegendentry{All}

\draw[draw=none,fill=darkviolet1270255] (axis cs:0.825,0) rectangle (axis cs:1.175,0.304);
\draw[draw=none,fill=darkviolet1270255] (axis cs:1.825,0) rectangle (axis cs:2.175,0.42);
\draw[draw=none,fill=darkviolet1270255] (axis cs:2.825,0) rectangle (axis cs:3.175,0.507);
\draw[draw=none,fill=darkviolet1270255] (axis cs:3.825,0) rectangle (axis cs:4.175,0.594);
\draw[draw=none,fill=darkviolet1270255] (axis cs:4.825,0) rectangle (axis cs:5.175,0.71);
\draw[draw=none,fill=darkviolet1270255] (axis cs:5.825,0) rectangle (axis cs:6.175,0.739);
\draw[draw=none,fill=royalblue8762253] (axis cs:-0.175,0) rectangle (axis cs:0.175,0.145);
\addlegendimage{ybar,ybar legend,draw=none,fill=royalblue8762253}
\addlegendentry{10}

\draw[draw=none,fill=royalblue8762253] (axis cs:0.825,0) rectangle (axis cs:1.175,0.304);
\draw[draw=none,fill=royalblue8762253] (axis cs:1.825,0) rectangle (axis cs:2.175,0.42);
\draw[draw=none,fill=royalblue8762253] (axis cs:2.825,0) rectangle (axis cs:3.175,0.507);
\draw[draw=none,fill=royalblue8762253] (axis cs:3.825,0) rectangle (axis cs:4.175,0.594);
\draw[draw=none,fill=royalblue8762253] (axis cs:4.825,0) rectangle (axis cs:5.175,0.703);
\draw[draw=none,fill=royalblue8762253] (axis cs:5.825,0) rectangle (axis cs:6.175,0.732);
\draw[draw=none,fill=dodgerblue45123246] (axis cs:-0.175,0) rectangle (axis cs:0.175,0.145);
\addlegendimage{ybar,ybar legend,draw=none,fill=dodgerblue45123246}
\addlegendentry{5}

\draw[draw=none,fill=dodgerblue45123246] (axis cs:0.825,0) rectangle (axis cs:1.175,0.304);
\draw[draw=none,fill=dodgerblue45123246] (axis cs:1.825,0) rectangle (axis cs:2.175,0.42);
\draw[draw=none,fill=dodgerblue45123246] (axis cs:2.825,0) rectangle (axis cs:3.175,0.493);
\draw[draw=none,fill=dodgerblue45123246] (axis cs:3.825,0) rectangle (axis cs:4.175,0.565);
\draw[draw=none,fill=dodgerblue45123246] (axis cs:4.825,0) rectangle (axis cs:5.175,0.638);
\draw[draw=none,fill=dodgerblue45123246] (axis cs:5.825,0) rectangle (axis cs:6.175,0.667);
\draw[draw=none,fill=deepskyblue3176236] (axis cs:-0.175,0) rectangle (axis cs:0.175,0.116);
\addlegendimage{ybar,ybar legend,draw=none,fill=deepskyblue3176236}
\addlegendentry{1}

\draw[draw=none,fill=deepskyblue3176236] (axis cs:0.825,0) rectangle (axis cs:1.175,0.203);
\draw[draw=none,fill=deepskyblue3176236] (axis cs:1.825,0) rectangle (axis cs:2.175,0.268);
\draw[draw=none,fill=deepskyblue3176236] (axis cs:2.825,0) rectangle (axis cs:3.175,0.312);
\draw[draw=none,fill=deepskyblue3176236] (axis cs:3.825,0) rectangle (axis cs:4.175,0.355);
\draw[draw=none,fill=deepskyblue3176236] (axis cs:4.825,0) rectangle (axis cs:5.175,0.428);
\draw[draw=none,fill=deepskyblue3176236] (axis cs:5.825,0) rectangle (axis cs:6.175,0.442);
\end{axis}

\end{tikzpicture}
}
%\resizebox{.307\textwidth}{!}{% This file was created with tikzplotlib v0.10.1.
\begin{tikzpicture}

\definecolor{darkgray176}{RGB}{176,176,176}
\definecolor{darkviolet1270255}{RGB}{127,0,255}
\definecolor{deepskyblue3176236}{RGB}{3,176,236}
\definecolor{dodgerblue45123246}{RGB}{45,123,246}
\definecolor{lightgray204}{RGB}{204,204,204}
\definecolor{royalblue8762253}{RGB}{87,62,253}

\begin{axis}[
legend cell align={left},
legend style={fill opacity=0.8, draw opacity=1, text opacity=1, draw=lightgray204, legend columns=-1, legend pos=north west},
tick align=outside,
tick pos=left,
title={$\Delta=2$ Repair Precision},
x grid style={darkgray176},
xlabel={Seconds},
xmin=-0.4925, xmax=6.4925,
xtick style={color=black},
xtick={0,1,2,3,4,5,6},
xticklabels={20,60,100,140,180,220,260},
y grid style={darkgray176},
ylabel={\phantom{Precision@k}},
ymin=0, ymax=0.40635,
ytick style={color=black}
]
\addlegendimage{empty legend}
\addlegendentry{P@}
\draw[draw=none,fill=darkviolet1270255] (axis cs:-0.175,0) rectangle (axis cs:0.175,0.048);
\addlegendimage{ybar,ybar legend,draw=none,fill=darkviolet1270255}
\addlegendentry{All}

\draw[draw=none,fill=darkviolet1270255] (axis cs:0.825,0) rectangle (axis cs:1.175,0.081);
\draw[draw=none,fill=darkviolet1270255] (axis cs:1.825,0) rectangle (axis cs:2.175,0.177);
\draw[draw=none,fill=darkviolet1270255] (axis cs:2.825,0) rectangle (axis cs:3.175,0.274);
\draw[draw=none,fill=darkviolet1270255] (axis cs:3.825,0) rectangle (axis cs:4.175,0.306);
\draw[draw=none,fill=darkviolet1270255] (axis cs:4.825,0) rectangle (axis cs:5.175,0.355);
\draw[draw=none,fill=darkviolet1270255] (axis cs:5.825,0) rectangle (axis cs:6.175,0.387);
\draw[draw=none,fill=royalblue8762253] (axis cs:-0.175,0) rectangle (axis cs:0.175,0.048);
\addlegendimage{ybar,ybar legend,draw=none,fill=royalblue8762253}
\addlegendentry{10}

\draw[draw=none,fill=royalblue8762253] (axis cs:0.825,0) rectangle (axis cs:1.175,0.081);
\draw[draw=none,fill=royalblue8762253] (axis cs:1.825,0) rectangle (axis cs:2.175,0.113);
\draw[draw=none,fill=royalblue8762253] (axis cs:2.825,0) rectangle (axis cs:3.175,0.21);
\draw[draw=none,fill=royalblue8762253] (axis cs:3.825,0) rectangle (axis cs:4.175,0.226);
\draw[draw=none,fill=royalblue8762253] (axis cs:4.825,0) rectangle (axis cs:5.175,0.258);
\draw[draw=none,fill=royalblue8762253] (axis cs:5.825,0) rectangle (axis cs:6.175,0.242);
\draw[draw=none,fill=dodgerblue45123246] (axis cs:-0.175,0) rectangle (axis cs:0.175,0.048);
\addlegendimage{ybar,ybar legend,draw=none,fill=dodgerblue45123246}
\addlegendentry{5}

\draw[draw=none,fill=dodgerblue45123246] (axis cs:0.825,0) rectangle (axis cs:1.175,0.032);
\draw[draw=none,fill=dodgerblue45123246] (axis cs:1.825,0) rectangle (axis cs:2.175,0.065);
\draw[draw=none,fill=dodgerblue45123246] (axis cs:2.825,0) rectangle (axis cs:3.175,0.129);
\draw[draw=none,fill=dodgerblue45123246] (axis cs:3.825,0) rectangle (axis cs:4.175,0.113);
\draw[draw=none,fill=dodgerblue45123246] (axis cs:4.825,0) rectangle (axis cs:5.175,0.129);
\draw[draw=none,fill=dodgerblue45123246] (axis cs:5.825,0) rectangle (axis cs:6.175,0.113);
\draw[draw=none,fill=deepskyblue3176236] (axis cs:-0.175,0) rectangle (axis cs:0.175,0.016);
\addlegendimage{ybar,ybar legend,draw=none,fill=deepskyblue3176236}
\addlegendentry{1}

\draw[draw=none,fill=deepskyblue3176236] (axis cs:0.825,0) rectangle (axis cs:1.175,0);
\draw[draw=none,fill=deepskyblue3176236] (axis cs:1.825,0) rectangle (axis cs:2.175,0.016);
\draw[draw=none,fill=deepskyblue3176236] (axis cs:2.825,0) rectangle (axis cs:3.175,0.065);
\draw[draw=none,fill=deepskyblue3176236] (axis cs:3.825,0) rectangle (axis cs:4.175,0.048);
\draw[draw=none,fill=deepskyblue3176236] (axis cs:4.825,0) rectangle (axis cs:5.175,0.032);
\draw[draw=none,fill=deepskyblue3176236] (axis cs:5.825,0) rectangle (axis cs:6.175,0.032);
\end{axis}

\end{tikzpicture}
}
    \caption{Human repair benchmark. Note the y-axis across different edit distance plots has varying ranges.}\label{fig:human}
  \end{figure}

%  In the following benchmark, we measure the precision@k of our repair procedure against human repairs of varying edit distances and latency cutoffs, using an adaptive resampling procedure described in \S\ref{sec:adaptive}. This sampler maintains a buffer of successful repairs ranked by perplexity and uses stochastic local search to resample edits within a neighborhood. Initially, edits are sampled uniformly at random. Over time and as the admissible set grows, it prioritizes edits nearby low-perplexity repairs. This technique offers a significant advantage in the low-latency setting.

  Previously, we used an adaptive sampler, which was much worse, because it did not sample directly from the admissible set, but the entire Levenshtein ball. We present the results from the earlier evaluation for comparison below.

  \begin{figure}[H]
    \resizebox{.24\textwidth}{!}{% This file was created with tikzplotlib v0.10.1.
\begin{tikzpicture}

\definecolor{darkgray176}{RGB}{176,176,176}
\definecolor{darkviolet1270255}{RGB}{127,0,255}
\definecolor{deepskyblue3176236}{RGB}{3,176,236}
\definecolor{dodgerblue45123246}{RGB}{45,123,246}
\definecolor{lightgray204}{RGB}{204,204,204}
\definecolor{royalblue8762253}{RGB}{87,62,253}

\begin{axis}[
legend cell align={left},
legend style={fill opacity=0.8, draw opacity=1, text opacity=1, draw=lightgray204, legend columns=-1, legend pos=north west},
tick align=outside,
tick pos=left,
title={\(\displaystyle \Delta\in[1,3]\) Repair Precision},
x grid style={darkgray176},
xlabel={Seconds},
xmin=-0.5925, xmax=8.5925,
ymin=0.0, ymax=0.7,
xtick style={color=black},
xtick={0,1,2,3,4,5,6,7,8},
xticklabels={0.1,0.3,0.5,0.7,0.9,1.1,1.3,1.5,1.7},
y grid style={darkgray176},
ylabel={Precision@k},
ytick style={color=black},
width=0.5\textwidth,
height=0.48\textwidth,
]
\addlegendimage{empty legend}
\addlegendentry{P@}
\draw[draw=none,fill=darkviolet1270255] (axis cs:-0.175,0) rectangle (axis cs:0.175,0.245);
\addlegendimage{ybar,ybar legend,draw=none,fill=darkviolet1270255}
\addlegendentry{All}

\draw[draw=none,fill=darkviolet1270255] (axis cs:0.825,0) rectangle (axis cs:1.175,0.417);
\draw[draw=none,fill=darkviolet1270255] (axis cs:1.825,0) rectangle (axis cs:2.175,0.496);
\draw[draw=none,fill=darkviolet1270255] (axis cs:2.825,0) rectangle (axis cs:3.175,0.535);
\draw[draw=none,fill=darkviolet1270255] (axis cs:3.825,0) rectangle (axis cs:4.175,0.562);
\draw[draw=none,fill=darkviolet1270255] (axis cs:4.825,0) rectangle (axis cs:5.175,0.584);
\draw[draw=none,fill=darkviolet1270255] (axis cs:5.825,0) rectangle (axis cs:6.175,0.603);
\draw[draw=none,fill=darkviolet1270255] (axis cs:6.825,0) rectangle (axis cs:7.175,0.619);
\draw[draw=none,fill=darkviolet1270255] (axis cs:7.825,0) rectangle (axis cs:8.175,0.629);
\draw[draw=none,fill=royalblue8762253] (axis cs:-0.175,0) rectangle (axis cs:0.175,0.234);
\addlegendimage{ybar,ybar legend,draw=none,fill=royalblue8762253}
\addlegendentry{10}

\draw[draw=none,fill=royalblue8762253] (axis cs:0.825,0) rectangle (axis cs:1.175,0.387);
\draw[draw=none,fill=royalblue8762253] (axis cs:1.825,0) rectangle (axis cs:2.175,0.451);
\draw[draw=none,fill=royalblue8762253] (axis cs:2.825,0) rectangle (axis cs:3.175,0.481);
\draw[draw=none,fill=royalblue8762253] (axis cs:3.825,0) rectangle (axis cs:4.175,0.502);
\draw[draw=none,fill=royalblue8762253] (axis cs:4.825,0) rectangle (axis cs:5.175,0.517);
\draw[draw=none,fill=royalblue8762253] (axis cs:5.825,0) rectangle (axis cs:6.175,0.53);
\draw[draw=none,fill=royalblue8762253] (axis cs:6.825,0) rectangle (axis cs:7.175,0.54);
\draw[draw=none,fill=royalblue8762253] (axis cs:7.825,0) rectangle (axis cs:8.175,0.546);
\draw[draw=none,fill=dodgerblue45123246] (axis cs:-0.175,0) rectangle (axis cs:0.175,0.218);
\addlegendimage{ybar,ybar legend,draw=none,fill=dodgerblue45123246}
\addlegendentry{5}

\draw[draw=none,fill=dodgerblue45123246] (axis cs:0.825,0) rectangle (axis cs:1.175,0.353);
\draw[draw=none,fill=dodgerblue45123246] (axis cs:1.825,0) rectangle (axis cs:2.175,0.406);
\draw[draw=none,fill=dodgerblue45123246] (axis cs:2.825,0) rectangle (axis cs:3.175,0.43);
\draw[draw=none,fill=dodgerblue45123246] (axis cs:3.825,0) rectangle (axis cs:4.175,0.444);
\draw[draw=none,fill=dodgerblue45123246] (axis cs:4.825,0) rectangle (axis cs:5.175,0.455);
\draw[draw=none,fill=dodgerblue45123246] (axis cs:5.825,0) rectangle (axis cs:6.175,0.467);
\draw[draw=none,fill=dodgerblue45123246] (axis cs:6.825,0) rectangle (axis cs:7.175,0.474);
\draw[draw=none,fill=dodgerblue45123246] (axis cs:7.825,0) rectangle (axis cs:8.175,0.479);
\draw[draw=none,fill=deepskyblue3176236] (axis cs:-0.175,0) rectangle (axis cs:0.175,0.117);
\addlegendimage{ybar,ybar legend,draw=none,fill=deepskyblue3176236}
\addlegendentry{1}

\draw[draw=none,fill=deepskyblue3176236] (axis cs:0.825,0) rectangle (axis cs:1.175,0.187);
\draw[draw=none,fill=deepskyblue3176236] (axis cs:1.825,0) rectangle (axis cs:2.175,0.219);
\draw[draw=none,fill=deepskyblue3176236] (axis cs:2.825,0) rectangle (axis cs:3.175,0.233);
\draw[draw=none,fill=deepskyblue3176236] (axis cs:3.825,0) rectangle (axis cs:4.175,0.244);
\draw[draw=none,fill=deepskyblue3176236] (axis cs:4.825,0) rectangle (axis cs:5.175,0.25);
\draw[draw=none,fill=deepskyblue3176236] (axis cs:5.825,0) rectangle (axis cs:6.175,0.258);
\draw[draw=none,fill=deepskyblue3176236] (axis cs:6.825,0) rectangle (axis cs:7.175,0.262);
\draw[draw=none,fill=deepskyblue3176236] (axis cs:7.825,0) rectangle (axis cs:8.175,0.265);
\addplot [red, thick] coordinates {(-0.8925,0.35) (14.8925,0.35)};
\end{axis}

\end{tikzpicture}
}
    \resizebox{.25\textwidth}{!}{% This file was created with tikzplotlib v0.10.1.
\begin{tikzpicture}

\definecolor{darkgray176}{RGB}{176,176,176}
\definecolor{darkviolet1270255}{RGB}{127,0,255}
\definecolor{deepskyblue3176236}{RGB}{3,176,236}
\definecolor{dodgerblue45123246}{RGB}{45,123,246}
\definecolor{lightgray204}{RGB}{204,204,204}
\definecolor{royalblue8762253}{RGB}{87,62,253}

\begin{axis}[
legend cell align={left},
legend style={fill opacity=0.8, draw opacity=1, text opacity=1, draw=lightgray204, legend columns=-1, legend pos=north west},
tick align=outside,
tick pos=left,
title={$\Delta=1$ Repair Precision},
x grid style={darkgray176},
xlabel={Seconds},
xmin=-0.5925, xmax=8.5925,
xtick style={color=black},
xtick={0,1,2,3,4,5,6,7,8},
xticklabels={0.1,0.3,0.5,0.7,0.9,1.1,1.3,1.5,1.7},
y grid style={darkgray176},
ymin=0, ymax=1.0,
ytick style={color=black}
width=0.43\textwidth,
height=0.49\textwidth,
]
\addlegendimage{empty legend}
\addlegendentry{P@}
\draw[draw=none,fill=darkviolet1270255] (axis cs:-0.175,0) rectangle (axis cs:0.175,0.416);
\addlegendimage{ybar,ybar legend,draw=none,fill=darkviolet1270255}
\addlegendentry{All}

\draw[draw=none,fill=darkviolet1270255] (axis cs:0.825,0) rectangle (axis cs:1.175,0.672);
\draw[draw=none,fill=darkviolet1270255] (axis cs:1.825,0) rectangle (axis cs:2.175,0.773);
\draw[draw=none,fill=darkviolet1270255] (axis cs:2.825,0) rectangle (axis cs:3.175,0.822);
\draw[draw=none,fill=darkviolet1270255] (axis cs:3.825,0) rectangle (axis cs:4.175,0.849);
\draw[draw=none,fill=darkviolet1270255] (axis cs:4.825,0) rectangle (axis cs:5.175,0.873);
\draw[draw=none,fill=darkviolet1270255] (axis cs:5.825,0) rectangle (axis cs:6.175,0.891);
\draw[draw=none,fill=darkviolet1270255] (axis cs:6.825,0) rectangle (axis cs:7.175,0.904);
\draw[draw=none,fill=darkviolet1270255] (axis cs:7.825,0) rectangle (axis cs:8.175,0.912);
\draw[draw=none,fill=royalblue8762253] (axis cs:-0.175,0) rectangle (axis cs:0.175,0.402);
\addlegendimage{ybar,ybar legend,draw=none,fill=royalblue8762253}
\addlegendentry{10}

\draw[draw=none,fill=royalblue8762253] (axis cs:0.825,0) rectangle (axis cs:1.175,0.645);
\draw[draw=none,fill=royalblue8762253] (axis cs:1.825,0) rectangle (axis cs:2.175,0.741);
\draw[draw=none,fill=royalblue8762253] (axis cs:2.825,0) rectangle (axis cs:3.175,0.786);
\draw[draw=none,fill=royalblue8762253] (axis cs:3.825,0) rectangle (axis cs:4.175,0.812);
\draw[draw=none,fill=royalblue8762253] (axis cs:4.825,0) rectangle (axis cs:5.175,0.834);
\draw[draw=none,fill=royalblue8762253] (axis cs:5.825,0) rectangle (axis cs:6.175,0.851);
\draw[draw=none,fill=royalblue8762253] (axis cs:6.825,0) rectangle (axis cs:7.175,0.862);
\draw[draw=none,fill=royalblue8762253] (axis cs:7.825,0) rectangle (axis cs:8.175,0.87);
\draw[draw=none,fill=dodgerblue45123246] (axis cs:-0.175,0) rectangle (axis cs:0.175,0.379);
\addlegendimage{ybar,ybar legend,draw=none,fill=dodgerblue45123246}
\addlegendentry{5}

\draw[draw=none,fill=dodgerblue45123246] (axis cs:0.825,0) rectangle (axis cs:1.175,0.601);
\draw[draw=none,fill=dodgerblue45123246] (axis cs:1.825,0) rectangle (axis cs:2.175,0.686);
\draw[draw=none,fill=dodgerblue45123246] (axis cs:2.825,0) rectangle (axis cs:3.175,0.724);
\draw[draw=none,fill=dodgerblue45123246] (axis cs:3.825,0) rectangle (axis cs:4.175,0.743);
\draw[draw=none,fill=dodgerblue45123246] (axis cs:4.825,0) rectangle (axis cs:5.175,0.762);
\draw[draw=none,fill=dodgerblue45123246] (axis cs:5.825,0) rectangle (axis cs:6.175,0.777);
\draw[draw=none,fill=dodgerblue45123246] (axis cs:6.825,0) rectangle (axis cs:7.175,0.785);
\draw[draw=none,fill=dodgerblue45123246] (axis cs:7.825,0) rectangle (axis cs:8.175,0.791);
\draw[draw=none,fill=deepskyblue3176236] (axis cs:-0.175,0) rectangle (axis cs:0.175,0.211);
\addlegendimage{ybar,ybar legend,draw=none,fill=deepskyblue3176236}
\addlegendentry{1}

\draw[draw=none,fill=deepskyblue3176236] (axis cs:0.825,0) rectangle (axis cs:1.175,0.327);
\draw[draw=none,fill=deepskyblue3176236] (axis cs:1.825,0) rectangle (axis cs:2.175,0.378);
\draw[draw=none,fill=deepskyblue3176236] (axis cs:2.825,0) rectangle (axis cs:3.175,0.401);
\draw[draw=none,fill=deepskyblue3176236] (axis cs:3.825,0) rectangle (axis cs:4.175,0.418);
\draw[draw=none,fill=deepskyblue3176236] (axis cs:4.825,0) rectangle (axis cs:5.175,0.428);
\draw[draw=none,fill=deepskyblue3176236] (axis cs:5.825,0) rectangle (axis cs:6.175,0.437);
\draw[draw=none,fill=deepskyblue3176236] (axis cs:6.825,0) rectangle (axis cs:7.175,0.442);
\draw[draw=none,fill=deepskyblue3176236] (axis cs:7.825,0) rectangle (axis cs:8.175,0.446);
\addplot [red, thick] coordinates {(-0.8925,0.45) (14.8925,0.45)};
\end{axis}

\end{tikzpicture}
}
    \resizebox{.24\textwidth}{!}{% This file was created with tikzplotlib v0.10.1.
\begin{tikzpicture}

\definecolor{darkgray176}{RGB}{176,176,176}
\definecolor{darkviolet1270255}{RGB}{127,0,255}
\definecolor{deepskyblue3176236}{RGB}{3,176,236}
\definecolor{dodgerblue45123246}{RGB}{45,123,246}
\definecolor{lightgray204}{RGB}{204,204,204}
\definecolor{royalblue8762253}{RGB}{87,62,253}

\begin{axis}[
legend cell align={left},
legend style={fill opacity=0.8, draw opacity=1, text opacity=1, draw=lightgray204, legend columns=-1, legend pos=north west},
tick align=outside,
tick pos=left,
axis lines*=left,
title={$\Delta=2$ Repair Precision},
x grid style={darkgray176},
xlabel={Seconds},
xmin=-0.5925, xmax=8.5925,
xtick style={color=black},
xtick={0,1,2,3,4,5,6,7,8},
xticklabels={0.1,0.3,0.5,0.7,0.9,1.1,1.3,1.5,1.7},
y grid style={darkgray176},
ymin=0, ymax=0.3843,
ytick style={color=black},
width=0.55\textwidth,
height=0.49\textwidth,
]
\addlegendimage{empty legend}
\addlegendentry{P@}
\draw[draw=none,fill=darkviolet1270255] (axis cs:-0.175,0) rectangle (axis cs:0.175,0.067);
\addlegendimage{ybar,ybar legend,draw=none,fill=darkviolet1270255}
\addlegendentry{All}

\draw[draw=none,fill=darkviolet1270255] (axis cs:0.825,0) rectangle (axis cs:1.175,0.165);
\draw[draw=none,fill=darkviolet1270255] (axis cs:1.825,0) rectangle (axis cs:2.175,0.225);
\draw[draw=none,fill=darkviolet1270255] (axis cs:2.825,0) rectangle (axis cs:3.175,0.261);
\draw[draw=none,fill=darkviolet1270255] (axis cs:3.825,0) rectangle (axis cs:4.175,0.292);
\draw[draw=none,fill=darkviolet1270255] (axis cs:4.825,0) rectangle (axis cs:5.175,0.309);
\draw[draw=none,fill=darkviolet1270255] (axis cs:5.825,0) rectangle (axis cs:6.175,0.335);
\draw[draw=none,fill=darkviolet1270255] (axis cs:6.825,0) rectangle (axis cs:7.175,0.353);
\draw[draw=none,fill=darkviolet1270255] (axis cs:7.825,0) rectangle (axis cs:8.175,0.366);
\draw[draw=none,fill=royalblue8762253] (axis cs:-0.175,0) rectangle (axis cs:0.175,0.06);
\addlegendimage{ybar,ybar legend,draw=none,fill=royalblue8762253}
\addlegendentry{10}

\draw[draw=none,fill=royalblue8762253] (axis cs:0.825,0) rectangle (axis cs:1.175,0.128);
\draw[draw=none,fill=royalblue8762253] (axis cs:1.825,0) rectangle (axis cs:2.175,0.164);
\draw[draw=none,fill=royalblue8762253] (axis cs:2.825,0) rectangle (axis cs:3.175,0.182);
\draw[draw=none,fill=royalblue8762253] (axis cs:3.825,0) rectangle (axis cs:4.175,0.2);
\draw[draw=none,fill=royalblue8762253] (axis cs:4.825,0) rectangle (axis cs:5.175,0.208);
\draw[draw=none,fill=royalblue8762253] (axis cs:5.825,0) rectangle (axis cs:6.175,0.221);
\draw[draw=none,fill=royalblue8762253] (axis cs:6.825,0) rectangle (axis cs:7.175,0.23);
\draw[draw=none,fill=royalblue8762253] (axis cs:7.825,0) rectangle (axis cs:8.175,0.238);
\draw[draw=none,fill=dodgerblue45123246] (axis cs:-0.175,0) rectangle (axis cs:0.175,0.05);
\addlegendimage{ybar,ybar legend,draw=none,fill=dodgerblue45123246}
\addlegendentry{5}

\draw[draw=none,fill=dodgerblue45123246] (axis cs:0.825,0) rectangle (axis cs:1.175,0.1);
\draw[draw=none,fill=dodgerblue45123246] (axis cs:1.825,0) rectangle (axis cs:2.175,0.123);
\draw[draw=none,fill=dodgerblue45123246] (axis cs:2.825,0) rectangle (axis cs:3.175,0.132);
\draw[draw=none,fill=dodgerblue45123246] (axis cs:3.825,0) rectangle (axis cs:4.175,0.142);
\draw[draw=none,fill=dodgerblue45123246] (axis cs:4.825,0) rectangle (axis cs:5.175,0.144);
\draw[draw=none,fill=dodgerblue45123246] (axis cs:5.825,0) rectangle (axis cs:6.175,0.155);
\draw[draw=none,fill=dodgerblue45123246] (axis cs:6.825,0) rectangle (axis cs:7.175,0.163);
\draw[draw=none,fill=dodgerblue45123246] (axis cs:7.825,0) rectangle (axis cs:8.175,0.167);
\draw[draw=none,fill=deepskyblue3176236] (axis cs:-0.175,0) rectangle (axis cs:0.175,0.016);
\addlegendimage{ybar,ybar legend,draw=none,fill=deepskyblue3176236}
\addlegendentry{1}

\draw[draw=none,fill=deepskyblue3176236] (axis cs:0.825,0) rectangle (axis cs:1.175,0.038);
\draw[draw=none,fill=deepskyblue3176236] (axis cs:1.825,0) rectangle (axis cs:2.175,0.052);
\draw[draw=none,fill=deepskyblue3176236] (axis cs:2.825,0) rectangle (axis cs:3.175,0.057);
\draw[draw=none,fill=deepskyblue3176236] (axis cs:3.825,0) rectangle (axis cs:4.175,0.061);
\draw[draw=none,fill=deepskyblue3176236] (axis cs:4.825,0) rectangle (axis cs:5.175,0.064);
\draw[draw=none,fill=deepskyblue3176236] (axis cs:5.825,0) rectangle (axis cs:6.175,0.072);
\draw[draw=none,fill=deepskyblue3176236] (axis cs:6.825,0) rectangle (axis cs:7.175,0.076);
\draw[draw=none,fill=deepskyblue3176236] (axis cs:7.825,0) rectangle (axis cs:8.175,0.078);
\addplot [red, thick] coordinates {(-0.8925,0.16) (14.8925,0.16)};
\end{axis}

\end{tikzpicture}
}
    \resizebox{.24\textwidth}{!}{% This file was created with tikzplotlib v0.10.1.
\begin{tikzpicture}

\definecolor{darkgray176}{RGB}{176,176,176}
\definecolor{darkviolet1270255}{RGB}{127,0,255}
\definecolor{deepskyblue3176236}{RGB}{3,176,236}
\definecolor{dodgerblue45123246}{RGB}{45,123,246}
\definecolor{lightgray204}{RGB}{204,204,204}
\definecolor{royalblue8762253}{RGB}{87,62,253}

\begin{axis}[
legend cell align={left},
legend style={fill opacity=0.8, draw opacity=1, text opacity=1, draw=lightgray204, legend columns=-1, legend pos=north west},
tick align=outside,
tick pos=left,
axis lines*=left,
title={$\Delta=3$ Repair Precision},
x grid style={darkgray176},
xlabel={Seconds},
xmin=-0.5925, xmax=8.5925,
xtick style={color=black},
xtick={0,1,2,3,4,5,6,7,8},
xticklabels={0.1,0.3,0.5,0.7,0.9,1.1,1.3,1.5,1.7},
yticklabels={0, 0.05, 0.10, 0.15, 0.20},
y grid style={darkgray176},
ymin=0, ymax=0.15225,
ytick style={color=black},
width=0.53\textwidth,
height=0.48\textwidth,
]
\addlegendimage{empty legend}
\addlegendentry{P@}
\draw[draw=none,fill=darkviolet1270255] (axis cs:-0.175,0) rectangle (axis cs:0.175,0.01);
\addlegendimage{ybar,ybar legend,draw=none,fill=darkviolet1270255}
\addlegendentry{All}

\draw[draw=none,fill=darkviolet1270255] (axis cs:0.825,0) rectangle (axis cs:1.175,0.035);
\draw[draw=none,fill=darkviolet1270255] (axis cs:1.825,0) rectangle (axis cs:2.175,0.064);
\draw[draw=none,fill=darkviolet1270255] (axis cs:2.825,0) rectangle (axis cs:3.175,0.072);
\draw[draw=none,fill=darkviolet1270255] (axis cs:3.825,0) rectangle (axis cs:4.175,0.09);
\draw[draw=none,fill=darkviolet1270255] (axis cs:4.825,0) rectangle (axis cs:5.175,0.109);
\draw[draw=none,fill=darkviolet1270255] (axis cs:5.825,0) rectangle (axis cs:6.175,0.116);
\draw[draw=none,fill=darkviolet1270255] (axis cs:6.825,0) rectangle (axis cs:7.175,0.133);
\draw[draw=none,fill=darkviolet1270255] (axis cs:7.825,0) rectangle (axis cs:8.175,0.145);
\draw[draw=none,fill=royalblue8762253] (axis cs:-0.175,0) rectangle (axis cs:0.175,0.003);
\addlegendimage{ybar,ybar legend,draw=none,fill=royalblue8762253}
\addlegendentry{10}

\draw[draw=none,fill=royalblue8762253] (axis cs:0.825,0) rectangle (axis cs:1.175,0.01);
\draw[draw=none,fill=royalblue8762253] (axis cs:1.825,0) rectangle (axis cs:2.175,0.01);
\draw[draw=none,fill=royalblue8762253] (axis cs:2.825,0) rectangle (axis cs:3.175,0.014);
\draw[draw=none,fill=royalblue8762253] (axis cs:3.825,0) rectangle (axis cs:4.175,0.018);
\draw[draw=none,fill=royalblue8762253] (axis cs:4.825,0) rectangle (axis cs:5.175,0.019);
\draw[draw=none,fill=royalblue8762253] (axis cs:5.825,0) rectangle (axis cs:6.175,0.018);
\draw[draw=none,fill=royalblue8762253] (axis cs:6.825,0) rectangle (axis cs:7.175,0.023);
\draw[draw=none,fill=royalblue8762253] (axis cs:7.825,0) rectangle (axis cs:8.175,0.023);
\draw[draw=none,fill=dodgerblue45123246] (axis cs:-0.175,0) rectangle (axis cs:0.175,0.002);
\addlegendimage{ybar,ybar legend,draw=none,fill=dodgerblue45123246}
\addlegendentry{5}

\draw[draw=none,fill=dodgerblue45123246] (axis cs:0.825,0) rectangle (axis cs:1.175,0.005);
\draw[draw=none,fill=dodgerblue45123246] (axis cs:1.825,0) rectangle (axis cs:2.175,0.008);
\draw[draw=none,fill=dodgerblue45123246] (axis cs:2.825,0) rectangle (axis cs:3.175,0.01);
\draw[draw=none,fill=dodgerblue45123246] (axis cs:3.825,0) rectangle (axis cs:4.175,0.011);
\draw[draw=none,fill=dodgerblue45123246] (axis cs:4.825,0) rectangle (axis cs:5.175,0.011);
\draw[draw=none,fill=dodgerblue45123246] (axis cs:5.825,0) rectangle (axis cs:6.175,0.011);
\draw[draw=none,fill=dodgerblue45123246] (axis cs:6.825,0) rectangle (axis cs:7.175,0.014);
\draw[draw=none,fill=dodgerblue45123246] (axis cs:7.825,0) rectangle (axis cs:8.175,0.018);
\draw[draw=none,fill=deepskyblue3176236] (axis cs:-0.175,0) rectangle (axis cs:0.175,0);
\addlegendimage{ybar,ybar legend,draw=none,fill=deepskyblue3176236}
\addlegendentry{1}

\draw[draw=none,fill=deepskyblue3176236] (axis cs:0.825,0) rectangle (axis cs:1.175,0.003);
\draw[draw=none,fill=deepskyblue3176236] (axis cs:1.825,0) rectangle (axis cs:2.175,0.008);
\draw[draw=none,fill=deepskyblue3176236] (axis cs:2.825,0) rectangle (axis cs:3.175,0.008);
\draw[draw=none,fill=deepskyblue3176236] (axis cs:3.825,0) rectangle (axis cs:4.175,0.011);
\draw[draw=none,fill=deepskyblue3176236] (axis cs:4.825,0) rectangle (axis cs:5.175,0.011);
\draw[draw=none,fill=deepskyblue3176236] (axis cs:5.825,0) rectangle (axis cs:6.175,0.011);
\draw[draw=none,fill=deepskyblue3176236] (axis cs:6.825,0) rectangle (axis cs:7.175,0.013);
\draw[draw=none,fill=deepskyblue3176236] (axis cs:7.825,0) rectangle (axis cs:8.175,0.013);
\addplot [red, thick] coordinates {(-0.8925,0.096) (14.8925,0.096)};
\addplot [orange, thick] coordinates {(-0.8925,0.02) (14.8925,0.02)};
\end{axis}

\end{tikzpicture}
}
    \caption{Adaptive sampling repairs. The red line indicates Seq2Parse precision@1 on the same dataset. Since it only supports generating one repair, we do not report precision@k or the intermediate latency cutoffs.}\label{fig:adaptive}
  \end{figure}

  We also evaluate Seq2Parse on the same dataset. Seq2Parse only supports precision@1 repairs, and so we only report Seq2parse precision@1 from the StackOverflow benchmark for comparison. Unlike our approach which only produces syntactically correct repairs, Seq2Parse also produces syntactically incorrect repairs and so we report the percentage of repairs matching the human repair for both our method and Seq2Parse. Seq2Parse latency varies depending on the length of the repair, averaging 1.5s for $\Delta=1$ to 2.7s for $\Delta=3$, across the entire StackOverflow dataset.

%  While adaptive sampling is able to saturate the admissible set for 1- and 2-edit repairs before the timeout elapses, 3-edit throughput is heavily constrained by compute around 16 lexical tokens, when Python's Levenshtein ball has a volume of roughly $6\times 10^8$ edits. This bottleneck can be relaxed with a longer timeout or additional CPU cores. Despite the high computational cost of sampling multi-edit repairs, our precision@all remains competitive with the Seq2Parse neurosymbolic baseline at the same latency. We provide some qualitative examples of repairs in Table~\ref{sec:appendix}.

  \begin{wrapfigure}{r}{0.35\textwidth}
    \vspace{-10pt}
    \resizebox{.35\textwidth}{!}{% This file was created with tikzplotlib v0.10.1.
\begin{tikzpicture}
\begin{axis}[
  legend cell align={left},
  legend style={fill opacity=0.8, draw opacity=1, text opacity=1, legend columns=-1, legend pos=north east},
  tick align=outside,
  tick pos=left,
  axis lines*=left,
  title={\textbf{Average throughput}},
  x grid style={darkgray176},
  xlabel={Lexical tokens},
  xmin=-0.95, xmax=19.95,
  xtick style={color=black},
  xtick={0,1,2,3,4,5,6,7,8,9,10,11,12,13,14,15,16,17,18,19},
  xticklabels={4, , 8, , 12, , 16, , 20, , 24, , 28, , 32, , 36, , 40,},
  y grid style={darkgray176},
  ylabel={Total repairs discovered},
  ymin=0, ymax=200000,
  ymode=log,
  legend pos=north east,
  ytick style={color=black}
]
\addlegendimage{empty legend}
\addlegendentry{$\Delta(\err\sigma, \sigma')=$}
\addlegendentry{\footnotesize{$1,$}}
\addlegendentry{\footnotesize{$2,$}}
\addlegendentry{\footnotesize{$3,$}}
\addlegendentry{\footnotesize{$4$}}
\addplot [semithick, green, mark=*, mark size=3, mark options={solid}]
table {%
  0 7.14285714285714
  1 12.1666666666667
  2 9.27272727272727
  3 10.6041666666667
  4 7.22857142857143
  5 6.4
  6 10.3956043956044
  7 7.86075949367089
  8 9.81818181818182
  9 8.28089887640449
  10 8.08695652173913
  11 10.0430107526882
  12 6.90476190476191
  13 6.94029850746269
  14 8.74698795180723
  15 7.70422535211268
  16 12.1414141414141
  17 4.1
};
\addplot [semithick, blue, mark=*, mark size=3, mark options={solid}]
table {%
0 46.625
1 135.653846153846
2 175.676923076923
3 79.0090909090909
4 87.6589147286822
5 69.9683544303797
6 83.9896373056995
7 85.9777777777778
8 73.3684210526316
9 67.7715517241379
10 92.1878172588832
11 186.568527918782
12 116.195945945946
13 87.8901734104046
14 99.6514285714286
15 154.44
16 123.658064516129
17 36.6111111111111
};
\addplot [semithick, orange, mark=*, mark size=3, mark options={solid}]
table {%
0 1506.8
1 2076.57142857143
2 5194.91304347826
3 413.939393939394
4 3749.34693877551
5 1130.83076923077
6 431.141025641026
7 601.172839506173
8 2135.26666666667
9 1233.85106382979
10 960.263888888889
11 932.225352112676
12 612.833333333333
13 880.2625
14 965.863636363636
15 725.139784946237
16 1199.05882352941
17 114.5
};
\addplot [semithick, red, mark=*, mark size=3, mark options={solid}]
table {%
0 nan
1 73461
2 24965
3 4220.4
4 35145.8181818182
5 2613.28571428571
6 68402.25
7 3300.5
8 10010
9 24536.5217391304
10 17497.6315789474
11 8883.58333333333
12 1948.23076923077
13 4175.27272727273
14 5145.45454545455
15 27785.8095238095
16 2468
17 4027.25
};
\end{axis}

\end{tikzpicture}}
    \caption{Total repairs found in 30s.}
    \label{fig:throughput}
    \vspace{-10pt}
  \end{wrapfigure}

  End-to-end throughput varies significantly with the edit distance of the repair. Some errors are trivial to fix, while others require a large number of edits to be sampled before a syntactically valid edit is discovered. We evaluate throughput by sampling edits across invalid strings $|\err\sigma| \leq 40$ from the StackOverflow dataset of varying length, and measure the total number of syntactically valid edits discovered, as a function of string length and language edit distance $\Delta\in[1, 3]$. Each trial is terminated after 10 seconds, and the experiment is repeated across 7.3k total repairs. Note the y-axis is log-scaled, as the number of admissible repairs increases sharply with language edit distance. Our approach discovers a large number of syntactically valid repairs in a relatively short amount of time, and is able to quickly saturate the admissible set for 1- and 2-edit repairs before timeout. As the Seq2Parse baseline is unable to generate more than one syntactically valid repair per string, we do not report its throughput.

  \begin{wrapfigure}{l}{0.35\textwidth}
    \resizebox{.35\textwidth}{!}{% This file was created with tikzplotlib v0.10.1.
\begin{tikzpicture}

\definecolor{green}{RGB}{0,128,0}

\begin{axis}[
log basis y={10},
tick align=outside,
tick pos=left,
x grid style={darkgray176},
xlabel={Snippet length},
xmin=9.5, xmax=42.5,
title={Latency across repair size},
xtick style={color=black},
y grid style={darkgray176},
log basis y={10},
axis lines*=left,
ylabel={Repair Latency (ms)},
ymin=65.2014650905023, ymax=509998.877384793,
ymode=log,
mark size=1.2pt,
ytick style={color=black},
legend style={fill opacity=0.8, draw opacity=1, text opacity=1, legend columns=2, legend pos=south east},
yticklabels={
  \(\displaystyle {10^{0}}\),
  \(\displaystyle {10^{1}}\),
  \(\displaystyle {10^{2}}\),
  \(\displaystyle {10^{3}}\),
  \(\displaystyle {10^{4}}\),
  \(\displaystyle {10^{5}}\),
  \(\displaystyle {10^{6}}\),
  \(\displaystyle {10^{7}}\)
}
]
\addlegendentry{\footnotesize{$\Delta = 1$}}
\addlegendentry{\footnotesize{$\Delta = 4$}}
\addlegendentry{\footnotesize{$\Delta = 2$}}
\addlegendentry{\footnotesize{$\Delta = 3$}}
\addplot [draw=green, fill=green, mark=*, only marks]
table{%
x  y
32 5148
29 1085
38 5383
19 1011
18 763
21 334
39 14789
29 3437
22 2883
17 764
28 2089
22 4131
20 960
21 1591
30 3769
35 4770
15 429
28 1943
25 2659
22 1590
27 851
14 173
36 2007
16 473
18 779
33 1866
39 2838
28 815
31 4644
40 4755
16 338
11 158
29 4790
23 3520
18 750
21 3491
27 2735
40 12157
22 4114
25 988
26 2727
13 193
11 151
25 3897
11 153
15 409
17 523
16 376
33 5443
16 494
34 8625
26 933
20 331
24 1043
12 192
21 3604
36 2593
15 2718
15 252
13 377
28 3815
13 245
12 146
29 2346
22 958
18 641
26 1208
26 3619
39 4494
16 674
30 4110
34 4124
28 3496
21 1387
27 3300
16 398
17 2755
13 282
36 4758
39 7288
35 1486
25 3472
34 1606
15 453
24 1210
19 782
12 173
35 4494
37 6857
12 127
21 1142
16 3631
18 574
39 3727
26 3361
36 1863
39 3152
39 4837
14 416
14 455
12 252
27 2046
11 114
27 5193
26 2286
25 1292
22 607
35 3958
16 263
26 1651
34 7934
31 1462
12 125
20 330
26 1767
23 479
17 539
37 3038
31 1084
18 880
32 2256
31 4098
37 1607
16 227
17 1050
24 3343
28 5643
40 2478
28 3858
27 1994
22 1267
24 1994
24 617
27 2141
19 414
11 2761
13 169
39 9225
28 3785
21 1130
34 1299
17 1016
20 3013
34 2118
21 4000
34 6822
21 303
21 400
20 857
37 1936
20 1063
18 2657
16 357
31 1229
25 1999
26 748
15 404
18 438
26 583
27 1157
31 6206
26 2629
16 229
17 612
33 10128
37 5953
11 196
14 304
12 303
41 4299
25 1908
20 3028
33 4834
40 2165
27 4406
18 396
38 2757
25 569
30 2030
29 3317
15 544
21 731
15 3022
31 1594
31 1047
27 2908
25 5094
12 173
11 150
30 3979
16 339
33 1593
31 3451
36 4199
28 6294
35 1726
16 637
11 165
15 750
25 2259
17 329
15 242
12 202
22 384
29 4869
39 5517
23 765
33 5001
32 1396
34 8403
20 1191
27 2758
33 1190
16 3353
13 186
26 5396
34 2910
33 2022
13 2877
28 1016
18 499
25 1246
33 3780
13 163
23 3873
19 806
39 8507
19 3459
34 3809
25 1802
20 585
26 2232
23 3450
32 4602
26 5929
18 347
19 3242
37 2207
17 746
22 6797
19 704
15 353
12 2762
23 580
22 1373
22 3589
25 1160
31 1695
20 1950
29 3053
39 2198
30 3956
32 2285
13 294
32 3911
13 237
27 4485
24 2690
19 934
29 4071
27 2122
24 995
24 507
18 625
24 3284
37 8257
12 3040
35 9141
14 2982
28 734
28 1306
17 2695
24 930
30 950
19 2784
32 6357
18 659
20 348
15 308
14 244
16 408
13 2797
27 3366
19 960
31 3625
11 193
30 3366
30 4727
33 5697
33 6494
29 1069
36 5357
37 3787
27 643
41 6979
32 1640
17 474
23 402
19 398
20 1209
25 1904
12 206
12 350
36 2470
35 1505
14 478
12 221
15 2811
28 3637
19 1122
35 10465
14 316
11 98
15 2795
13 3284
26 1998
12 177
36 7469
20 1307
36 3463
19 643
37 2210
36 2022
33 5903
35 4781
13 260
20 290
40 12774
35 4327
23 1640
14 318
13 163
14 276
39 7636
16 176
18 3151
32 3490
20 773
39 2417
11 143
14 326
11 135
37 9387
23 618
15 532
};
\addplot [draw=red, fill=red, mark=*, only marks]
table{%
x  y
20 58108
12 51666
15 52872
27 107841
22 45114
31 195597
29 107688
28 173325
39 307961
20 19806
41 191265
28 119322
38 116176
34 142491
23 67212
25 72841
23 72985
22 62959
36 136040
36 178035
14 48610
28 170679
36 111571
33 123007
22 61457
26 96451
41 175655
26 77194
29 35609
28 169819
14 46105
39 195626
24 50751
21 48576
37 231823
11 47882
21 84741
22 53728
33 90374
35 208976
33 162398
21 90295
39 322474
29 202614
35 168255
31 140663
28 148913
37 269034
26 115719
33 216191
25 71448
38 173519
13 43689
34 102807
13 57551
31 112047
21 61169
32 163450
22 64456
18 60343
33 109713
36 225779
19 57508
33 166085
12 51689
16 50373
19 59679
25 264645
22 78703
14 78482
33 87321
31 110919
20 82528
26 100040
37 339313
40 102371
32 145489
24 99148
18 55190
25 110437
24 80558
36 184598
18 53084
20 71368
33 313482
12 40595
13 45971
24 81069
39 137454
34 171818
20 111020
32 111393
15 51481
22 64615
20 63536
30 179273
16 15228
19 61004
31 125079
36 93285
34 142710
20 64707
19 46655
24 157899
34 112240
16 52176
34 172346
14 18390
13 19427
13 46764
31 134228
36 240347
27 120224
20 98838
21 34748
15 69259
13 50058
31 100419
20 47676
14 30778
27 69207
37 156018
40 172073
18 41322
36 84168
19 171901
17 71764
33 186628
26 127454
24 97074
23 68230
18 55273
32 154745
24 72679
16 60711
22 67811
35 112980
35 109433
14 49207
23 85428
21 67406
38 190668
34 311037
39 124408
};
\addplot [draw=blue, fill=blue, mark=*, only marks]
table{%
x  y
15 1546
36 31862
38 31248
26 16119
33 29070
20 7825
19 12970
23 4255
32 11597
35 21227
15 5376
13 2141
30 39009
19 15202
38 15296
25 13377
32 9555
24 5876
20 15019
38 22100
33 35495
11 719
28 20635
27 9206
31 6117
23 20579
36 42335
28 15167
25 4452
20 9944
11 934
19 8080
12 1903
24 5029
18 2931
37 44732
34 17401
11 1400
30 17034
33 18639
34 34009
22 19118
12 5107
18 12158
19 6940
30 15976
12 5113
18 4296
21 6390
28 21456
34 13034
25 5045
22 4333
34 22718
23 2505
30 39479
16 4844
33 35755
15 2477
22 8883
22 13650
23 5820
22 4991
14 3418
35 7532
31 11363
25 29496
20 5739
12 3618
25 18211
17 16099
18 9777
28 8551
18 8232
17 1239
15 2497
39 9304
32 25118
14 1484
26 10561
21 11388
29 7176
30 17223
21 8634
37 7977
17 5433
21 8490
19 17680
18 4292
25 5967
34 13023
28 14019
24 8835
16 3006
24 25420
25 12155
13 1941
29 11108
31 18410
26 8491
19 7119
17 3180
17 3805
28 15120
32 32599
12 5333
14 2206
13 8839
23 8906
33 30905
38 14948
31 22032
38 9392
27 9183
36 13507
37 18594
26 22120
17 4754
34 35382
21 10508
24 22208
36 40758
36 31016
37 28948
14 4118
14 5684
20 3662
29 14922
12 1274
13 4091
22 7839
20 15066
23 12974
14 973
21 8622
37 46911
14 3107
32 32976
37 9228
38 26957
22 3555
37 27104
34 9138
13 3353
15 8690
26 16640
34 23305
23 4885
22 10846
28 7864
15 7957
15 8747
30 6519
12 4219
20 4469
35 47446
19 4889
27 8550
39 57886
23 14245
23 21774
20 2934
26 14781
32 37106
35 11013
20 8515
15 2539
29 24498
31 25400
22 13575
28 32976
33 37348
30 17818
39 11397
29 20198
12 4619
29 17768
28 18812
14 9168
26 21072
29 28933
16 11018
12 2620
17 8626
23 19402
41 61109
33 25388
16 7368
33 6903
20 7721
29 9658
15 1256
11 4263
34 9959
18 2161
37 47273
30 21254
31 10892
26 11905
11 1763
26 13949
19 7328
29 28416
16 12111
20 6234
17 4271
27 11313
29 12137
21 7872
35 21597
27 21491
39 10804
19 15818
19 3881
38 47207
19 3066
28 25870
30 12676
17 5394
23 20199
28 15303
22 29638
39 56108
38 7822
30 13293
31 9281
26 37151
30 19684
17 3761
39 20861
32 14503
22 13901
16 7406
17 13836
28 11129
19 2229
28 33492
38 16182
15 4529
33 31173
23 12037
22 9964
20 5659
32 24562
12 680
24 14414
30 15023
23 12023
38 17828
18 2136
34 13575
28 22060
21 16814
25 13082
38 47841
35 7347
26 18942
32 14484
34 10717
25 15950
34 32785
27 5235
23 19239
12 1326
30 20466
27 22080
26 16954
23 13937
30 18121
32 18321
15 5283
15 10042
29 19455
32 22029
12 3696
25 14987
38 15618
15 2085
38 13985
37 23038
18 4363
17 3734
21 6877
36 30820
36 40020
23 11704
36 28883
33 47148
41 13381
34 11999
19 4825
29 35034
15 5870
20 8343
23 9744
36 8368
31 37536
37 67404
38 14815
25 17157
31 8009
22 11478
23 17884
27 16520
11 1158
16 7454
28 9036
26 7471
14 1098
15 6672
14 4076
};
\addplot [draw=orange, fill=orange, mark=*, only marks]
table{%
x  y
40 42016
20 24189
33 33848
21 10526
19 29639
30 36650
20 40015
23 8429
38 32819
27 69823
29 43136
26 63140
17 42342
32 17391
12 40540
28 54691
26 43793
19 9754
34 173208
36 80096
12 16827
33 33888
16 39156
29 64485
23 14735
37 126065
36 118322
39 101483
25 57385
16 36604
31 66360
11 10512
38 126586
38 137665
41 32922
25 32976
37 65246
40 65888
22 66806
27 51204
29 24094
32 39285
37 72613
26 65788
28 44251
16 5413
35 63518
32 110601
31 36090
29 47350
19 12132
31 52986
36 149418
11 37909
41 32579
26 72372
29 87247
22 50063
29 56161
23 9463
11 39435
32 48779
35 73904
28 67568
29 71226
29 30317
25 26198
35 47708
37 132547
39 174702
17 13890
30 38364
15 41787
28 83550
29 55774
14 15458
30 46008
39 60620
22 45732
37 102710
14 42288
30 87541
38 102615
30 319242
23 56070
33 58231
20 31861
15 29942
14 22163
15 15176
19 48163
19 16105
29 53572
37 159305
39 62454
28 18109
20 30293
15 11821
33 89003
27 30070
31 63481
30 55899
29 61285
19 50656
32 16841
19 16676
21 42955
15 36952
19 11367
15 14400
29 64709
30 48081
29 53699
21 57294
12 38724
31 63965
23 49107
12 37097
12 37806
34 19660
17 17929
31 64083
33 88215
21 58401
15 39393
12 36692
17 39872
32 77042
35 127043
35 126005
12 18445
20 22042
26 44232
12 17708
11 38234
36 94269
35 120430
33 75175
21 39828
31 81582
36 96767
11 5778
19 43704
13 8415
23 44485
32 79172
29 60849
12 3793
27 41624
20 26682
17 14129
33 126568
24 46496
29 48634
16 42893
27 77985
24 65710
27 85586
31 62940
41 68300
33 102270
11 5412
41 191452
31 59553
17 52058
15 38902
32 28359
17 45407
34 135818
21 5610
31 112023
13 37517
19 14117
18 33770
29 58688
38 110486
21 44776
23 8391
19 23909
25 60406
30 28910
36 138052
36 135891
39 80921
30 108163
21 52436
20 21960
37 106582
35 78295
31 69280
22 27743
37 79928
12 15280
30 116392
14 7227
37 35707
26 49205
35 58959
35 149408
23 12075
37 101615
33 60528
18 59862
37 177210
33 77206
28 57257
13 7439
12 17730
24 68492
39 112291
22 47825
24 76251
13 34728
23 42931
39 145463
37 98906
33 72393
23 44026
28 23606
16 27186
22 40794
18 46860
26 34273
21 9723
26 74168
24 44922
38 118479
30 70306
27 67938
22 61961
35 95075
26 111028
17 42197
13 41792
13 10854
34 61113
15 39520
34 62196
29 50751
41 133897
18 22246
27 21455
16 9835
14 11967
37 35868
21 45493
37 69815
35 27949
28 44323
32 47729
32 83133
22 47076
};
\end{axis}

\end{tikzpicture}
}
    \caption{End-to-end repair timings across snippet length.}
    \label{fig:timings}
    \vspace{-25pt}
  \end{wrapfigure}

  In Fig.~\ref{fig:timings}, we plot the end-to-end repair timings across snippet length. To measure the timings, we collect 1000 repair samples of varying lengths and edit distances, then measure the time until the sampler locates the exact human repair. For each code snippet, we measure the total time taken, for repairs between 1-4 Levenshtein edits. Note the logarithmic scale on the y-axis.

  \subsection{Synthetic repair benchmark}\label{sec:latency}

  Pairwise naturally-occurring errors and human fixes are the most authentic dataset, but can be difficult to obtain in as great a quantity as we would like. An alternative nearly as good is to collect a dataset of clean code, and synthetically corrupt it. The original source code then becomes the ground truth repair for the synthetically introduced typo, and the target for evaluating the precision of our repair procedure.

  In our synthetic repair bechmark, we synthetically corrupt BIFI code snippets using a learned corruption model. To train the model computing the Levenshtein alignment on a pairwise dataset, then approximating the conditional probability of each edit given the local context. We then sample a corruption of a good string matching the distribution of typos in the dataset, and measure the precision at recovering the originally valid string.

  Suppose we have a dataset of single token edits and their local context. For simplicity, we shall assume a trigram language model, i.e., $P(\sigma_i' \mid \sigma_{i-1}, \sigma_i, \sigma_{i+1})$, however the approach can be generalized to higher-order Markov models. Given a string $\sigma$, we can sample error trajectories $q^1(\sigma), q^2(\sigma), \ldots, q^n(\sigma)$ by defining $q(\sigma)$ to sample a single edit from the set of all relevant edit actions $Q(\sigma)$, then recursively applying $q$ to the resulting string. More formally,

  \begin{enumerate}
    \item Given a string $\sigma$, compute $Q(\sigma)$, the set of all relevant edit actions for all possible edit locations by unioning the set of all possible edits at each location, i.e., $Q(\sigma) \coloneqq \bigcup_{i=1}^{|\sigma| - 1} \big\{\sigma_i' \mid  0 < P(\sigma_i \mid \sigma_{i-1}, \sigma_i, \sigma_{i+1})\big\}$.
    \item Renormalize the probabilities of each edit $P(q \mid \sigma)$ by $\sum_{q \in Q(\sigma)} P(q)$. This ensures the probability of sampling a particular edit is proportional to its relative probability under the language model and sums to 1.
    \item Sample an edit $q(\sigma) \sim Q(\sigma)$, then repeat for $n$ steps where $n$ is sampled from a geometric distribution with mean $\mu$ matching the average edit distance of the dataset (this assumes the edit distance is independent of the edits).
  \end{enumerate}

  This allows us to sample a corruption of a good string matching the distribution of typos in the dataset. We then measure the precision at recovering the originally valid string.

  In addition to the StackOverflow dataset, we also evaluate our approach on two datasets containing synthetic strings generated by a Dyck language, and bracketing errors of synthetic and organic provenance in organic source code. The first dataset contains length-50 strings sampled from various Dyck languages, i.e., the Dyck language containing n different types of balanced parentheses. The second contains abstracted Java and Python source code mined from GitHub repositories. The Dyck languages used in the remaining experiments are defined by the following context-free grammar(s):

  \begin{wholetidyinput}
    Dyck-1 -> ( ) | ( Dyck-1 ) | Dyck-1 Dyck-1
Dyck-2 -> Dyck-1 | [ ] | ( Dyck-2 ) | [ Dyck-2 ] | Dyck-2 Dyck-2
Dyck-3 -> Dyck-2 | { } | ( Dyck-3 ) | [ Dyck-3 ] | { Dyck-3 } | Dyck-3 Dyck-3
  \end{wholetidyinput}

  \noindent In experiment (1a), we sample a random valid string $\sigma \sim \Sigma^{50} \cap \mathcal{L}_{\text{Dyck-n}}$, then replace a fixed number of indices in $[0, |\sigma|)$ with holes and measure the average time required to decode ten syntactically-admissible repairs across 100 trial runs. In experiment (1b), we sample a random valid string as before, but delete p tokens at random and rather than provide their location(s), ask our model to solve for both the location(s) and repair by sampling uniformly from all n-token HCs, then measure the total time required to decode the first admissible repair. Note the logarithmic scale on the y-axis.

  \begin{figure}[H]
    \begin{minipage}{.48\textwidth}
      \begin{center}\footnotesize\textbf{Synthetic bracket language}\end{center}
    \end{minipage}
    \begin{minipage}{.48\textwidth}
      \begin{center}\footnotesize\textbf{Organic bracket language}\end{center}
    \end{minipage}\\
    \vspace{10pt}
    \hspace{-0.25cm}\begin{tikzpicture}[scale=0.35]
    \begin{axis}[
      width=8.3cm,
      height=7cm,
      title={\hspace{-1cm}\textbf{(1a) Latency with known locations}},
      ybar,
      bar width=6pt,
      xlabel={Number of holes},
      ylabel={ms to synthesize 10 repairs},
      xtick=data,
      axis x line*=bottom,
      axis y line*=left,
      ytick pos=left,
      xticklabels from table={\loctimings}{holes},
      ymajorgrids,
      legend pos=north west,
      legend columns=2,
      error bars/y dir=both,
      error bars/y explicit
    ]
      \addplot table [x expr=\coordindex, y=d1, y error=d1err]{\loctimings};
      \addplot table [x expr=\coordindex, y=d2, y error=d2err]{\loctimings};
      \addplot table [x expr=\coordindex, y=d3, y error=d3err]{\loctimings};
      \addplot table [x expr=\coordindex, y=d4, y error=d4err]{\loctimings};
      \legend{Dyck-1, Dyck-2, Dyck-3, Dyck-4}
    \end{axis}
    \end{tikzpicture}
    \begin{tikzpicture}[scale=0.35]
      \begin{axis}[
        width=8.3cm,
        height=7cm,
        title={\hspace{-1cm}\textbf{(1b) Latency with unknown locations}},
        ybar,
        bar width=20pt,
        xlabel={Number of errors},
%ylabel={ms to synthesize 1 repair},
        xtick=data,
        axis x line*=bottom,
        axis y line*=left,
        enlarge x limits={abs=0.5},
        ymode=log,
        ytick pos=left,
        xticklabels from table={\unloctimings}{errors},
        ymajorgrids,
        legend pos=north west,
        error bars/y dir=both,
        error bars/y explicit
      ]
        \addplot table [x expr=\coordindex, y=d1]{\unloctimings};
        \addplot table [x expr=\coordindex, y=d2]{\unloctimings};
        \addplot table [x expr=\coordindex, y=d3]{\unloctimings};
        \legend{Dyck-1, Dyck-2, Dyck-3}
      \end{axis}
    \end{tikzpicture}
    \hspace{20pt}
    \begin{tikzpicture}[scale=0.35]
      \begin{axis}[
      width=8.3cm,
      height=7cm,
%    title={\hspace{-1cm}\textbf{Java Brackets}},
      ybar,
      bar width=10pt,
      xlabel={$|\Sigma^*|$},
      ylabel={Top-1 parser acceptance},
      title={\textbf{(2a) Synthetic Java bracket error correction}},
      xtick=data,
      axis x line*=bottom,
      axis y line*=left,
      enlarge x limits={abs=0.5},
      ytick pos=left,
      xticklabels from table={\syntheticerrors}{len},
      ymajorgrids,
      legend pos=north east,
      legend columns=3,
      error bars/y dir=both,
      error bars/y explicit
      ]
      \addplot table [x expr=\coordindex, y=10s]{\syntheticerrors};
      \addplot table [x expr=\coordindex, y=30s]{\syntheticerrors};
      \addplot table [x expr=\coordindex, y=60s]{\syntheticerrors};
      \legend{10s, 30s, 60s}
      \end{axis}
    \end{tikzpicture}
    \begin{tikzpicture}[scale=0.35]
      \begin{axis}[
      width=8.3cm,
      height=7cm,
      title={\textbf{(2b) Organic Python bracket error correction}},
      ybar,
      bar width=10pt,
      xlabel={$|\Sigma^*|$},
%ylabel={Parser acceptance},
      xtick=data,
      axis x line*=bottom,
      axis y line*=left,
      enlarge x limits={abs=0.5},
      ytick pos=left,
      xticklabels from table={\naturalerrors}{len},
      ymajorgrids,
      y tick label style={/pgf/number format/.cd,%
      scaled y ticks = false,
      set thousands separator={},
      fixed},
      legend pos=north east,
      legend columns=3,
      error bars/y dir=both,
      error bars/y explicit
      ]
      \addplot table [x expr=\coordindex, y=10s]{\naturalerrors};
      \addplot table [x expr=\coordindex, y=30s]{\naturalerrors};
      \addplot table [x expr=\coordindex, y=60s]{\naturalerrors};
      \legend{10s, 30s, 60s}
      \end{axis}
    \end{tikzpicture}
    \caption{Benchmarking bracket correction latency and accuracy across two bracketing languages, one generated from Dyck-n, and the second uses an abstracted source code snippet with imbalanced parentheses.}
  \end{figure}

  In the second set of experiments, we analyze bracketing errors in a dataset of Java and Python code snippets mined from open-source repositories on GitHub using the Dyck-nw\footnote{Using the Dyck-n grammar augmented with a single additional production, \texttt{Dyck-1} {\color{blue}\texttt{->}} \texttt{w} {\color{blue}\texttt{|}} \texttt{Dyck-1}. Contiguous non-bracket characters are substituted with a single placeholder token, \texttt{w}, and restored verbatim after bracket repair.}, in which all source code tokens except brackets are replaced with a \texttt{w} token. For Java (2a), we sample valid single-line statements with bracket nesting more than two levels deep, synthetically delete one bracket uniformly at random, and repair using Tidyparse, then take the top-1 repair after $t$ seconds, and validate using ANTLR's Java 8 parser. For Python (2b), we sample invalid code fragments uniformly from the imbalanced bracket category of the Break-It-Fix-It (BIFI) dataset~\cite{yasunaga2021break}, a dataset of organic Python errors, which we repair using Tidyparse, take the top-1 repair after $t$ seconds, and validate repairs using Python's \texttt{ast.parse()} method. Since the Java and Python datasets do not have a ground-truth human fix, we report the percentage of repairs that are accepted by the language's official parser for repairs generated under a fixed time cutoff. Although the Java and Python datasets are not directly comparable, we observe that Tidyparse can detect and repair a significant fraction of bracket errors in both languages with a relatively unsophisticated grammar.

  \section{Related Work}\label{sec:related}

  Three important questions arise when repairing syntax errors: (1) is the program broken in the first place? (2) if so, where are the errors located? (3) how should those locations then be altered? In the case of syntax correction, those questions are addressed by three related research areas, (1) parsing, (2) language equations and (3) repair. We survey each of those areas in turn.

  \subsection{Parsing}

  Context-free language (CFL) parsing is the well-studied problem of how to turn a string into a unique tree, with many different algorithms and implementations (e.g., shift-reduce, recursive-descent, LR). Many of those algorithms expect grammars to be expressed in a certain form (e.g., left- or right- recursive) or are optimized for a narrow class of grammars (e.g., regular, linear).

  General CFL parsing allows ambiguity (non-unique trees) and can be formulated as a dynamic programming problem, as shown by Cocke-Younger-Kasami (CYK)~\cite{sakai1961syntax}, Earley~\cite{earley1970efficient} and others. These parsers have roughly cubic complexity with respect to the length of the input string.

  As shown by Valiant~\cite{valiant1975general}, Lee~\cite{lee2002fast} and others, general CFL recognition is in some sense equivalent to binary matrix multiplication, another well-studied combinatorial problem with broad applications, known to be at worst subcubic. This reduction unlocks the door to a wide range of complexity-theoretic and practical speedups to CFL recognition and fast general parsing algorithms.

  Okhotin (2001)~\cite{okhotin2001conjunctive} extends CFGs with language conjunction in \textit{conjunctive grammars}, followed by Zhang \& Su (2017)~\cite{zhang2017context} who apply conjunctive language reachability to dataflow analysis.

  \subsection{Language equations}

  Language equations are a powerful tool for reasoning about formal languages and their inhabitants. First proposed by Ginsburg et al.~\cite{ginsburg1962two} for the ALGOL language, language equations are essentially systems of inequalities with variables representing \textit{holes}, i.e., unknown values, in the language or grammar. Solutions to these equations can be obtained using various fixpoint techniques, yielding members of the language. This insight reveals the true algebraic nature of CFLs and their cousins.

  Being an algebraic formalism, language equations naturally give rise to a kind of calculus, vaguely reminiscent of Leibniz' and Newton's. First studied by Brzozowski~\cite{brzozowski1964derivatives, brzozowski1980equations} and Antimirov~\cite{antimirov1996partial}, one can take the derivative of a language equation, yielding another equation. This can be interpreted as a kind of continuation or language quotient, revealing the suffixes that complete a given prefix. This technique leads to an elegant family of algorithms for incremental parsing~\cite{might2011parsing, adams2016complexity} and automata minimization~\cite{brzozowski1962canonical}. In our setting, differentiation corresponds to code completion.

  In this paper, we restrict our attention to language equations over context-free and weakly context-sensitive languages, whose variables coincide with edit locations in the source code of a computer program, and solutions correspond to syntax repairs. Although prior work has studied the use of language equations for parsing~\cite{might2011parsing}, to our knowledge they have never previously been considered for the purpose of code completion or syntax error correction.

%Rosenkrantz~\cite{rosenkrantz1967matrix} introduces , and

%When a sentence contains so-called \textit{holes}, parsing becomes slightly more challenging. Holes are special tokens that can be filled by words in the grammar. A sentence might have multiple holes, representing a set of contextually valid words. For example, ``\textit{\_ \_ to eat \_}'' which could be filled by, e.g., ``\textit{I like\ldots sushi}'', ``\textit{They want\ldots pizza}'', or ``\textit{We need\ldots something}''.

  \subsection{Syntax repair}

  In finite languages, syntax repair corresponds to spelling correction, a more restrictive and largely solved problem. Schulz and Stoyan~\cite{schulz2002fast} construct a finite automaton that returns the nearest dictionary entry by Levenshtein edit distance. Though considerably simpler than syntax correction, their work shares similar challenges and offers insights for handling more general repair scenarios.

  When a sentence is grammatically invalid, parsing grows more challenging. Like spelling, the problem is to find the minimum number of edits required to transform an arbitrary string into a syntactically valid one, where validity is defined as containment in a (typically) context-free language. Early work, including Irons~\cite{irons1963error} and Aho~\cite{aho1972minimum} propose a dynamic programming algorithm to compute the minimum number of edits required to fix an invalid string. Prior work on error correcting parsing only considers the shortest edit(s), and does not study multiple edits over the Levenshtein ball. Furthermore, the problem of actually generating the repairs is not well-posed, as there are usually many valid strings that can be obtained within a given number of edits. We instead focus on bounded Levenshtein reachability, which is the problem of finding useful repairs within a fixed Levenshtein distance of the broken string, which requires language intersection.

%Bar-Hillel~\cite{bar1961formal} establishes the closure of CFLs under intersection with regular languages, and can be used to construct the corresponding minimum-cost CFG in a straightforward manner. However while generally quite expressive, CFLs are themselves not closed under intersection and have other practical limitations, i.e., are unable to express indentation or variable binding. These limitations motivate us towards more expressive yet still efficiently parsable formalisms.

%Conveniently, Okhotin~\cite{okhotin2001conjunctive} offers exactly the formalism needed: language equations augmented with logical operators like conjunction or disjunction. These operators afford us the flexibility to encode both language union and intersection, which are difficult or impossible to express using a single grammar, as well as incorporate extra-syntactic side-constraints during solving.

%Naively, one might think to enumerate all edits within a certain Levenshtein radius and reject invalid strings. However, this approach is intractable for long strings. Another approach is to compute the \textit{edit distance} between the string and the language as proposed by Aho. This approach is also intractable depending on the size of the grammar. A third approach is to use a language equation to compute the intersection between the Levenshtein hypersphere and the language.

%The literature on parsers is vast and deep, covering far more ground than we could possibly hope to survey. We take inspiration from their findings, but restrict ourselves to a dozen most closely related papers, in four major research areas: (1) formal language theory, (2) constraint satisfaction (3) program synthesis, and (4) error correction. We survey each of these areas in turn.
%
%It was Noam Chomsky who first developed the algebraic theory of \textit{context-free grammars} (CFGs) in 1959~\cite{chomsky1959algebraic}, and since then, CFGs have been the subject of extensive research in formal language theory and program analysis. In particular, we take inspiration from Leslie Valiant~\cite{valiant1975general} who first discovered the connection to matrix multiplication in 1975 and Alexander Okhotin~\cite{okhotin2001conjunctive} who later introduced the idea of \textit{conjunctive grammars} in 2001. More recently, Azimov \& Grigorev~\cite{azimov2018cfpq} and Qirun Zhang shed light into the practical applicability of these ideas and made the theory of CFG reachability more accessible to the general public. In our work, we show how to compile their ideas onto a SAT solver to fix syntax errors, which seems like a perfectly natural extension, but was heretofore previously never considered to the best of our knowledge.

  \subsection{Classical program synthesis}

  There is related work on string constraint solving in the constraint programing literature, featuring solvers like CFGAnalyzer and HAMPI~\cite{kiezun2009hampi}, which consider bounded context free grammars and intersections thereof. Axelson et al. (2008)~\cite{axelsson2008analyzing} has some work on incremental SAT encoding but does not exploit the linear-algebraic structure of parsing, conjunctive reachability nor provide real-time guarantees. D'Antoni et al. (2014) introduces \textit{symbolic automata}~\cite{dantoni2014minimization}, a generalization of finite automata which allow infinite alphabets and symbolic expressions over them. In none of the constraint programming literature we surveyed do any of the approaches employ matrix-based parsing, and therefor do not enjoy the optimality guarantees of Valiant's parser. Our solver can handle context-free and conjunctive grammars with finite alphabets and does not require any special grammar encoding. The matrix encoding makes it particularly amenable to parallelization.

  \subsection{Error correcting codes}

  Our work focuses on errors arising from human factors in computer programming, in particular \textit{syntax error correction}, which is the problem of fixing partially corrupted programs. Modern research on error correction, however, can be traced back to the early days of coding theory when researchers designed \textit{error-correcting codes} (ECCs) to denoise transmission errors induced by external interference, e.g., collision with a high-energy proton, manipulation by an adversary or even typographical mistake. In this context, \textit{code} can be any logical representation for communicating information between two parties (such as a human and a computer), and an ECC is a carefully-designed scheme which ensures that even if some portion of the message should become corrupted, one can still recover the original message by solving a linear system of equations. When designing ECCs, one typically assumes a noise model over a certain event space, such as the Hamming~\cite{titsias2017hamming, dong2023number} or Levenshtein~\cite{becerra2008learning, barlev2021levenshtein} balls, from which we draw inspiration for our work.

  \subsection{Neural program repair}

  The recent success of deep learning has lead to a variety of work on neural program repair~\cite{allamanis2021self, chirkova2021empirical, drain2021generating}. These approaches typically employ Transformer-based neural language models (NLMs) and model the problem as a sequence-to-sequence transformation. Although recent work on circuit lower bounds have cast doubt on the ability of Transformers to truly learn formal languages~\cite{merrill2022saturated, chiang2023tighter}, expressivity aside, these models have been widely adopted for practical program repair tasks. In particular, two papers stand out being most closely related to our own: Break-It-Fix-It (BIFI)~\cite{yasunaga2021break} and Seq2Parse~\cite{sakkas2022seq2parse}. BIFI adapts techniques from semi-supervised machine translation to generate synthetic errors in clean code and fixes them. This reduces the amount of pairwise training data, but may generalize poorly to natural errors. Seq2Parse combines a transformer-based model with an augmented version of the Early parser to suggest error rules, but only suggests a single repair. Our work differs from both in that we suggest multiple repairs with much lower latency, do not require a pairwise repair dataset, and can fix syntax errors in any language with a well-defined grammar. We note our approach is complementary to existing work in neural program repair, and may be used to generate synthetic repairs for training or employ a NLM model to rank its repairs.

  \section{Discussion}\label{sec:discussion}

  The main lesson we draw from our experiments is that it is possible to leverage compute to compete with large language models on practical program repair tasks. Though extremely sample-efficient, their size comes at the cost of higher latency, and domain adaptation requires fine-tuning or retraining on pairwise repairs. Our approach uses a tiny grammar and a relatively cheap ranking metric to achieve comparable accuracy at the same latency. This allows us to repair errors in languages with little to no training data and provides far more flexibility and controllability.

  Our primary insight leading to SoTA precision@\{5,10\} is that repairs are typically concentrated near a small number of edit locations, and by biasing the search toward previously successful edit locations, we can achieve a significant speedup over a na\"ive search. We note this heuristic may not be applicable to all grammars, and it may be possible to construct less natural counterexamples where this heuristic fails, although we consider these unlikely in practice.

  Latency can vary depending on many factors including string length and grammar size, and critically the Levenshtein edit distance. This can be an advantage because in the absence of any contextual or statistical information, syntax and Levenshtein edits are often sufficiently constrained to determine a small number of valid repairs. It is also a limitation because as the number of edits grows, the admissible set grows rapidly and the number of valid repairs may become too large to be useful without a good metric, depending on the language and source code snippet under repair.

  Tidyparse in its current form has several other technical shortcomings: firstly, it does not incorporate any neural language modeling technology at present, an omission we hope to address. Training a language model to predict likely repair locations and rank admissible results could lead to lower overall latency and more natural repairs. We also hope to explore the use of Metropolis-Hastings and determinantal point processes to encourage sampling diversity.

  Secondly, our current method does not specialize language intersection to the grammar family, nor employ Bar-Hillel's~\cite{bar1961formal} construction for REG-CFL intersection, which would lead to a more efficient encoding of Levenshtein-CFL reachability. Furthermore, considering recent extensions of Boolean matrix-based parsing to linear context-free rewriting systems (LCFRS)~\cite{cohen2016parsing}, it may be feasible to search through richer language families within the SAT solver without employing an external stochastic search to generate and validate candidate repairs.

  Lastly and perhaps most significantly, Tidyparse does not incorporate any semantic constraints, so its repairs whilst syntactically admissible, are not guaranteed to be semantically valid. This can be partly alleviated by filtering the results through an incremental compiler or linter, however, the latency introduced may be non-negligible. It is also possible to encode type-based semantic constraints into the solver and we intend to explore this direction more fully in future work.

  We envision a few primary use cases for our tool: (1) helping novice programmers become more quickly familiar with a new programming language, (2) autocorrecting common typos among proficient but forgetful programmers, (3) as a prototyping tool for PL designers and educators, and (4) as a pluggable library or service for parser-generators and language servers. Featuring a grammar editor and built-in SAT solver, Tidyparse helps developers navigate the language design space, visualize syntax trees, debug parsing errors and quickly generate simple examples and counterexamples for benchmarking and testing.

  \section{Conclusion}\label{sec:conclusion}


  The great compromise in program synthesis is one of efficiency versus expressiveness. The more expressive a language, the more concise and varied the programs it can represent, but the harder those programs are to synthesize without resorting to domain-specific heuristics. Likewise, the simpler a language is to synthesize, the weaker its concision and expressive power.

  Most existing work on program synthesis has focused on general $\lambda$-calculi, or narrow languages such as finite sets or regular expressions. The former are too expressive to be efficiently synthesized or verified, whilst the latter are too restrictive to be useful. In our work, we focus on context-free and mildly context-sensitive grammars, which are expressive enough to capture a variety of useful programming language features, but not so expressive as to be unsynthesizable.

  The second great compromise in program synthesis is that of reusability versus specialization. In programming, as in human communications, there is a vast constellation of languages, each requiring specialized generators and interpreters. Are these languages truly irreconcilable? Or, as Noam Chomsky argues, are these merely dialects of a universal language? \textit{Synthesis} then, might be a misnomer, and more aptly called \textit{recognition}, in the analytic tradition.

  In our work, we argue these two compromises are not mutually exclusive, but complementary and reciprocal. Programs and the languages they inhabit are indeed synthetic, but can be analyzed and reused in the metalanguage of context-free grammars closed under conjunction. Not only does this admit an efficient synthesis algorithm, but allows users to introduce additional constraints without breaking compositionality, one of the most sacred tenets in programming language design.

  Over the last few years, there has been a surge of progress in applying language models to write programs. That work is primarily based on methods from differential calculus and continuous optimization, leading to the so-called \textit{naturalness hypothesis}, which suggests programming languages are not so different from natural ones. In contrast, programming language theory takes the view that languages are essentially discrete and finitely-generated sets, although perhaps uncomputably large ones, governed by logical calculi. Programming, thus viewed, is more like a mathematical exercise in constraint satisfaction, an oft-overlooked but unavoidable aspect of translating abstract ideas into computing machinery. These constraints naturally arise at various stages of syntax validation, type-checking and runtime verification, and help to ensure programs fulfill their intended purpose.

  As our work shows, not only is linear algebra over finite fields an expressive language for probabilistic inference, but also an efficient framework for inference on languages themselves. Borrowing analysis techniques from multilinear algebra and tensor completion in the machine learning setting, we develop an equational theory that allows us to translate various decision problems on formal languages into a system of inequalities over finite fields. As a practical consequence, this means we can efficiently encode a number of problems in parsing, code completion and program repair using SAT solvers, which are known to be highly efficient, scalable and flexible to domain-specific constraints. We demonstrate the effectiveness of our approach in a variety of practical scenarios, including code completion and program repair for linear conjunctive languages, and show that our approach is competitive with state-of-the-art methods in terms of both accuracy and efficiency. In contrast with LL and LR-style parsers, our technique can recover partial forests from invalid strings, and handle arbitrary conjunctive languages. In future work, we hope to extend our method to more natural grammars like PCFG, LCFRS and other mildly context-sensitive languages.

  Syntax correction tools should be as user-friendly and widely-accessible as autocorrection tools in word processors. From a practical standpoint, we argue it is possible to reduce disruption from manual syntax repair and improve the efficiency of working programmers by driving down the latency needed to synthesize an acceptable repair. In contrast with program synthesizers that require intermediate editor states to be well-formed, our synthesizer does not impose any constraints on the code itself being written and can be used in a live programming environment.

  Despite its computational complexity, the design of the tool itself is relatively simple. Tidyparse accepts a linear conjunctive language and a string. If the string is valid, it returns the parse forest, otherwise, it returns a set of repairs, ordered by their perplexity. Tidyparse compiles the grammar and candidate string onto a discrete dynamical system using an extended version of Valiant's algorithm and solves for its fixedpoints using an incremental SAT solver. By allowing the string to contain holes, repairs may contain either concrete tokens or nonterminals, which can be expanded by the user or a neural-guided search procedure to produce a valid program. This approach to parsing has many advantages, enabling us to repair syntax errors, correct typos and recover from errors in real time, as well as being provably sound and complete with respect to the grammatical specification. It is also compatible with neural program synthesis and repair techniques and shares the same masked language modeling (MLM) target used by Transformer-based LLMs.

  We have implemented our approach as an IDE plugin and demonstrated its viability as a practical tool for realtime programming. A considerable amount of effort was devoted to supporting fast error correction functionality. Tidyparse is capable of generating repairs for invalid code in a range of practical languages with very little downstream language integration required. We plan to continue expanding its grammar and autocorrection functionality to cover a broader range of languages and hope to conduct a more thorough user study to validate its effectiveness.

%\subsection{Ranking}
%
%Since the number of solutions can be very large, we can use a language model to rank the results maximizing likelihood, or minimizing perplexity, subject to the constraints. This ranking can be used to guide the propagation, sample the choice function, sample hole locations or as a post-processing step after a fixed timeout has expired.

%Alternatively, this expression can be rewritten as a polynomial over GF(2):
%
%\[
%  (v_1 \times w_2 + y_3 + 1) \Leftrightarrow [S \in Y] \Leftrightarrow [Q R \in L(G)]
%\]

%
% ---- Bibliography ----
%
% BibTeX users should specify bibliography style 'splncs04'.
% References will then be sorted and formatted in the correct style.
%
%  \bibliographystyle{splncs04}
  \bibliography{../bib/acmart}
\end{document}